% 自 Notizen/W06_1019/ 各表由 Skripte/gen_tex.py 產生；改字改表後重產，不手改本檔。
\chapter{Waldron 論違憲審查的核心論據（2026-10-19）}

2026 年 10 月 19 日上午課堂錄音（約 145 分鐘）之逐字稿。由同學亦凡導讀 Waldron 2006 年之文，湯德宗老師穿插評論。左邊界為段首時間碼；未標說話人者為湯德宗老師。同週的導讀問答見課程問答集第 6 部。


\section{本文七部分概述}

\tc{0:00:00}
% cue 1 0.580 2.760
\hypertarget{w06:cue:1}{}但是他現在要講的是，
% cue 2 2.760 4.880
\hypertarget{w06:cue:2}{}反對裡面的核心問題\footnote{Jeremy Waldron, The Core of the Case Against Judicial Review, 115 Yale Law Journal 1346–1406（2006），課綱第 6 週指定文獻。Waldron 其時任 Columbia University 法學院 University Professor，2006 年 7 月起轉任 New York University 法學院教授。全文除導論與結論外分七部分：第壹部分界定違憲審查，第貳部分設四項假設，第參部分談論證形式，第肆、伍部分分論結果面與程序面的理由，第陸部分討論多數暴政，第柒部分處理非核心情形。}，
% cue 3 4.880 6.910
\hypertarget{w06:cue:3}{}那他對這個核心，
% cue 4 6.910 9.100
\hypertarget{w06:cue:4}{}他有一個定義，
% cue 5 9.100 10.100
\hypertarget{w06:cue:5}{}他要使他存化，
% cue 6 10.100 11.100
\hypertarget{w06:cue:6}{}要使他怎樣怎樣，
% cue 7 11.100 12.060
\hypertarget{w06:cue:7}{}好啦，
% cue 8 12.060 14.180
\hypertarget{w06:cue:8}{}這個是一部分，
% cue 9 14.180 15.180
\hypertarget{w06:cue:9}{}然後我們看一下這個，
% cue 10 15.180 17.770
\hypertarget{w06:cue:10}{}你們現在講一這個，
% cue 11 17.770 18.770
\hypertarget{w06:cue:11}{}這個鄭偉，
% cue 12 18.770 19.770
\hypertarget{w06:cue:12}{}我甚至想罵你，
% cue 13 19.770 21.830
\hypertarget{w06:cue:13}{}你這個每一個都整個不一樣，
% cue 14 21.830 22.830
\hypertarget{w06:cue:14}{}有大有小，
% cue 15 22.830 23.830
\hypertarget{w06:cue:15}{}還有橫的，
% cue 16 23.830 24.830
\hypertarget{w06:cue:16}{}還有直的，
% cue 17 24.830 25.840
\hypertarget{w06:cue:17}{}還有罵最發糟的，
% cue 18 25.840 28.110
\hypertarget{w06:cue:18}{}大家你們都沒有罵他，
% cue 19 28.110 29.110
\hypertarget{w06:cue:19}{}這個人，
% cue 20 29.110 30.110
\hypertarget{w06:cue:20}{}你看這個人，
% cue 21 30.110 32.110
\hypertarget{w06:cue:21}{}如果彈簧請教 一定會老師罵的，
% cue 22 32.110 37.110
\hypertarget{w06:cue:22}{}你可以有忽強忽弱 忽大忽小 忽高忽低，
% cue 23 37.110 40.270
\hypertarget{w06:cue:23}{}從事要把調子定這一樣。

\tc{0:00:40}
% cue 24 40.270 41.270
\hypertarget{w06:cue:24}{}那我們看一下這個，
% cue 25 41.270 44.230
\hypertarget{w06:cue:25}{}你們現在講義的翻過來的那頁有沒有，
% cue 26 44.230 47.640
\hypertarget{w06:cue:26}{}我們有些什麼樣的目次，
% cue 27 47.640 49.640
\hypertarget{w06:cue:27}{}這個是亦凡你要講的 我方你講一下而已，
% cue 28 49.640 51.840
\hypertarget{w06:cue:28}{}第一塊是 Introduction，
% cue 29 51.840 55.050
\hypertarget{w06:cue:29}{}Introduction 大概是問題的說明，
% cue 30 55.050 58.580
\hypertarget{w06:cue:30}{}我覺得應該等一下你大概兩分鐘就講完。

\tc{0:00:58}
% cue 31 58.580 60.760
\hypertarget{w06:cue:31}{}下面就是 Definition of Judicial Review，
% cue 32 60.760 61.640
\hypertarget{w06:cue:32}{}這個比較活躍，
% cue 33 61.640 63.960
\hypertarget{w06:cue:33}{}因為他的 Definition，
% cue 34 63.960 65.510
\hypertarget{w06:cue:34}{}跟我的 Definition 不一樣，
% cue 35 65.510 67.620
\hypertarget{w06:cue:35}{}我實在太受大陸法系的影響，
% cue 36 67.620 69.880
\hypertarget{w06:cue:36}{}他是非常英美法系的意義，
% cue 37 69.880 72.050
\hypertarget{w06:cue:37}{}所以他對這個 Definition 是不一樣的，
% cue 38 72.050 73.970
\hypertarget{w06:cue:38}{}我不能說完全不一樣，
% cue 39 73.970 75.450
\hypertarget{w06:cue:39}{}就是我那個重點不一樣，
% cue 40 75.450 77.660
\hypertarget{w06:cue:40}{}他說什麼叫做司法審查，
% cue 41 77.660 79.240
\hypertarget{w06:cue:41}{}他的司法審查，
% cue 42 79.240 81.300
\hypertarget{w06:cue:42}{}他甚至沒有完全使用，
% cue 43 81.300 82.790
\hypertarget{w06:cue:43}{}Constitutional Review
% cue 44 82.790 84.210
\hypertarget{w06:cue:44}{}因為他有些可能不施定，
% cue 45 84.210 85.840
\hypertarget{w06:cue:45}{}不一定是 Constitutional Review。

\tc{0:01:25}
% cue 46 85.840 87.500
\hypertarget{w06:cue:46}{}然後第二大塊呢，
% cue 47 87.500 89.340
\hypertarget{w06:cue:47}{}他有四個 Assumptions，
% cue 48 89.340 94.560
\hypertarget{w06:cue:48}{}有四個假設，他是在這四個假設之下來討論，
% cue 49 94.560 100.750
\hypertarget{w06:cue:49}{}所以他是有條件的，他自己說他是有條件的反違憲審查論，
% cue 50 100.750 106.650
\hypertarget{w06:cue:50}{}不是全部條件，他如果這個條件這四個假設不存在的時候，
% cue 51 106.650 110.130
\hypertarget{w06:cue:51}{}他覺得有違憲審查並不見得是錯誤，
% cue 52 110.130 114.190
\hypertarget{w06:cue:52}{}第三塊，the form of the argument，
% cue 53 114.190 117.920
\hypertarget{w06:cue:53}{}就是論證的這個形式，
% cue 54 117.920 121.110
\hypertarget{w06:cue:54}{}這個地方等一下，剛才亦凡要怎麼講，
% cue 55 121.110 124.660
\hypertarget{w06:cue:55}{}這裡我覺得是比較費解的，
% cue 56 124.660 125.860
\hypertarget{w06:cue:56}{}我不懂為什麼要講。

\tc{0:02:05}
% cue 57 125.860 128.050
\hypertarget{w06:cue:57}{}第四大塊，
% cue 58 128.050 130.460
\hypertarget{w06:cue:58}{}outcome related reasons
% cue 59 130.460 133.020
\hypertarget{w06:cue:59}{}是跟結果有關的，
% cue 60 133.020 138.240
\hypertarget{w06:cue:60}{}跟結果有關的理論，
% cue 61 138.240 140.010
\hypertarget{w06:cue:61}{}它有兩塊，
% cue 62 140.010 142.650
\hypertarget{w06:cue:62}{}有一塊是跟程序有關的理論，
% cue 63 142.650 144.210
\hypertarget{w06:cue:63}{}有一塊是跟結果，
% cue 64 144.210 145.610
\hypertarget{w06:cue:64}{}所以分別出現在，
% cue 65 145.610 147.170
\hypertarget{w06:cue:65}{}第四跟第五，
% cue 66 147.170 148.120
\hypertarget{w06:cue:66}{}換句話說，
% cue 67 148.120 150.290
\hypertarget{w06:cue:67}{}你支持違憲審查，
% cue 68 150.290 152.790
\hypertarget{w06:cue:68}{}那有一些是因為，
% cue 69 152.790 154.630
\hypertarget{w06:cue:69}{}考慮到違憲審查，
% cue 70 154.630 155.970
\hypertarget{w06:cue:70}{}它會產生，
% cue 71 155.970 160.250
\hypertarget{w06:cue:71}{}好多結果 正確的結果 公平的結果 這是一個，
% cue 72 160.250 164.410
\hypertarget{w06:cue:72}{}另外有一種第五小節 他叫做 Chapter 5，
% cue 73 164.410 168.920
\hypertarget{w06:cue:73}{}他認為是張了 沒關係 總之 Section Chapter 一樣，
% cue 74 168.920 174.610
\hypertarget{w06:cue:74}{}那第五節呢 這一塊呢 他認為第二種支持違憲審查，
% cue 75 174.610 177.900
\hypertarget{w06:cue:75}{}不論呢 是跟程序相關的理由，
% cue 76 177.900 183.760
\hypertarget{w06:cue:76}{}支持違憲審查有結果相關的理由 也有程序相關的理由。

\tc{0:03:03}
% cue 77 183.760 192.190
\hypertarget{w06:cue:77}{}好，然後另外還有一個很常見的理由呢，是在第六塊，所謂的 tyranny of the majority\footnote{英文，多數暴政。本文第陸部分（刊本 1395 頁起）以此為題；該部分註 127 指出此語經 Tocqueville《Democracy in America》推廣，Mill 於 1840 年 10 月刊於 Edinburgh Review 的書評則認為，此語推廣的政治效果會是替保守勢力多添一套反對進步立法的說辭。}，
% cue 78 192.190 200.590
\hypertarget{w06:cue:78}{}就是因為我們擔心會有多數暴力，這個可能因為我們受到了美國所謂的 Federalist Papers\footnote{Alexander Hamilton、James Madison、John Jay 以 Publius 為筆名所寫的 85 篇文章，1787 年 10 月至 1788 年 5 月間主要刊於紐約州兩家報紙，為促成聯邦憲法的批准而作。Madison 執筆的第 10 篇（1787 年 11 月 23 日刊於 New York Packet）把 faction（黨派）界定為一群公民，不論占全體的多數或少數，受共同的激情或利益驅使，與其他公民的權利或社會整體的長遠利益相違，並討論多數結成黨派時民選政府所受的威脅。\url{https://guides.loc.gov/federalist-papers/full-text}。} 的影響，
% cue 79 200.590 210.260
\hypertarget{w06:cue:79}{}大家就會朗朗上口說，因為我們多數呢會欺負少數啊，多數欺負少數，少數就沒有辦法反抗啊，
% cue 80 210.260 213.470
\hypertarget{w06:cue:80}{}所以我們要有個別的機構來反應，
% cue 81 213.470 215.530
\hypertarget{w06:cue:81}{}來壓抑多數的暴力，
% cue 82 215.530 217.960
\hypertarget{w06:cue:82}{}制裁控制多數的暴力，
% cue 83 217.960 219.250
\hypertarget{w06:cue:83}{}那這是第六。

\tc{0:03:39}
% cue 84 219.250 220.770
\hypertarget{w06:cue:84}{}然後第七塊，
% cue 85 220.770 222.440
\hypertarget{w06:cue:85}{}Non-core case
% cue 86 222.440 225.370
\hypertarget{w06:cue:86}{}我們注意它的標題開始是 core case，
% cue 87 225.370 230.050
\hypertarget{w06:cue:87}{}我現在講反對違憲審查的核心，
% cue 88 230.050 232.710
\hypertarget{w06:cue:88}{}那麼如果是非核心的情境，
% cue 89 232.710 233.880
\hypertarget{w06:cue:89}{}那又不一樣，
% cue 90 233.880 235.640
\hypertarget{w06:cue:90}{}它不會完全反對，
% cue 91 235.640 236.720
\hypertarget{w06:cue:91}{}在非核心，
% cue 92 236.720 238.480
\hypertarget{w06:cue:92}{}就是剛才前面的第二節，
% cue 93 238.480 239.760
\hypertarget{w06:cue:93}{}有四個假設，
% cue 94 239.760 241.600
\hypertarget{w06:cue:94}{}如果有一個或者幾個，
% cue 95 241.600 242.560
\hypertarget{w06:cue:95}{}不存在的時候，
% cue 96 242.560 244.660
\hypertarget{w06:cue:96}{}那他不會真正的反對，
% cue 97 244.660 245.780
\hypertarget{w06:cue:97}{}或者他有些，
% cue 98 245.780 247.540
\hypertarget{w06:cue:98}{}他認為不適合心的思想，
% cue 99 247.540 248.940
\hypertarget{w06:cue:99}{}最後是他的結論，
% cue 100 248.940 249.220
\hypertarget{w06:cue:100}{}好了，
% cue 101 249.220 250.440
\hypertarget{w06:cue:101}{}我能講呢，
% cue 102 250.440 250.920
\hypertarget{w06:cue:102}{}我講完了。

\section{摘要與導論}

\tc{0:04:10}
% cue 103 250.920 251.160
\hypertarget{w06:cue:103}{}來，
% cue 104 251.160 252.440
\hypertarget{w06:cue:104}{}你開始吧，
% cue 105 252.440 253.670
\hypertarget{w06:cue:105}{}前面 introduction，
% cue 106 253.670 254.590
\hypertarget{w06:cue:106}{}你跟我們介紹一下，
% cue 107 254.590 255.150
\hypertarget{w06:cue:107}{}他說什麼？

\tc{0:04:15}
\speaker{導讀人：亦凡}
% cue 108 255.150 256.830
\hypertarget{w06:cue:108}{}前面 introduction 就是，
% cue 109 256.830 257.690
\hypertarget{w06:cue:109}{}他，

\tc{0:04:17}
\speaker{老師}
% cue 110 257.690 258.830
\hypertarget{w06:cue:110}{}你告訴我們在哪裡？
% cue 111 258.830 261.380
\hypertarget{w06:cue:111}{}還是你不要玩，
% cue 112 261.380 262.050
\hypertarget{w06:cue:112}{}在哪裡？

\tc{0:04:22}
\speaker{導讀人：亦凡}
% cue 113 262.050 267.350
\hypertarget{w06:cue:113}{}大家看的是哪一份，

\tc{0:04:27}
\speaker{老師}
% cue 114 267.350 267.830
\hypertarget{w06:cue:114}{}哪一段，
% cue 115 267.830 268.150
\hypertarget{w06:cue:115}{}對，
% cue 116 268.150 269.550
\hypertarget{w06:cue:116}{}你要我們眼睛，
% cue 117 269.550 270.610
\hypertarget{w06:cue:117}{}跟著你到哪裡去，
% cue 118 270.610 271.110
\hypertarget{w06:cue:118}{}OK

\tc{0:04:31}
\speaker{導讀人：亦凡}
% cue 119 271.110 272.410
\hypertarget{w06:cue:119}{}就是，
% cue 120 272.410 274.290
\hypertarget{w06:cue:120}{}1346
% cue 121 274.290 275.690
\hypertarget{w06:cue:121}{}就是最開始的，
% cue 122 275.690 278.450
\hypertarget{w06:cue:122}{}1346 的 abstract 的部分，

\tc{0:04:38}
\speaker{老師}
% cue 123 278.450 280.650
\hypertarget{w06:cue:123}{}前面的 1346 考，
% cue 124 280.650 282.330
\hypertarget{w06:cue:124}{}你要先講 abstract，
% cue 125 282.330 282.770
\hypertarget{w06:cue:125}{}對，

\tc{0:04:42}
\speaker{導讀人：亦凡}
% cue 126 282.770 286.120
\hypertarget{w06:cue:126}{}他就是先講，
% cue 127 286.120 288.060
\hypertarget{w06:cue:127}{}他這篇文章要寫什麼，
% cue 128 288.060 291.770
\hypertarget{w06:cue:128}{}要批評的是什麼，
% cue 129 291.770 293.920
\hypertarget{w06:cue:129}{}標的，
% cue 130 293.920 297.830
\hypertarget{w06:cue:130}{}然後通過什麼方式去執行，

\tc{0:04:57}
\speaker{老師}
% cue 131 297.830 301.930
\hypertarget{w06:cue:131}{}講完了嗎？
% cue 132 301.930 302.490
\hypertarget{w06:cue:132}{}很快啊，
% cue 133 302.490 303.110
\hypertarget{w06:cue:133}{}跳過來了，
% cue 134 303.110 304.780
\hypertarget{w06:cue:134}{}你覺得他講什麼？

\tc{0:05:04}
\speaker{導讀人：亦凡}
% cue 135 304.780 306.840
\hypertarget{w06:cue:135}{}Introduction 的部分就是，
% cue 136 306.840 309.550
\hypertarget{w06:cue:136}{}他先舉了，
% cue 137 309.550 315.390
\hypertarget{w06:cue:137}{}一些例子去顯示，
% cue 138 315.390 318.550
\hypertarget{w06:cue:138}{}如果當法官確信某項法律，
% cue 139 318.550 321.620
\hypertarget{w06:cue:139}{}有侵害個人權利自由的時候，
% cue 140 321.620 324.240
\hypertarget{w06:cue:140}{}他們有怎麼樣的一些案件，
% cue 141 324.240 331.270
\hypertarget{w06:cue:141}{}能夠做出這種跟民主程序的互動，
% cue 142 331.270 334.250
\hypertarget{w06:cue:142}{}然後他就說，
% cue 143 334.250 337.330
\hypertarget{w06:cue:143}{}後面 1348 頁的最後，
% cue 144 337.330 337.790
\hypertarget{w06:cue:144}{}他就說，
% cue 145 337.790 348.160
\hypertarget{w06:cue:145}{}就是他先承認了這些上述的一些判決，
% cue 146 348.160 355.300
\hypertarget{w06:cue:146}{}對於個人權利的保護的有一些起到了一些正面的作用，
% cue 147 355.300 357.800
\hypertarget{w06:cue:147}{}然後他就引出了後面的一些論證，
% cue 148 357.800 363.940
\hypertarget{w06:cue:148}{}就是他想要說明在他的文章後面的一些假設之下，
% cue 149 363.940 365.620
\hypertarget{w06:cue:149}{}這並不是最好的方法。

\tc{0:06:05}
\speaker{老師}
% cue 150 365.620 368.260
\hypertarget{w06:cue:150}{}我們看一下 P348 頁最後一句，
% cue 151 368.260 369.540
\hypertarget{w06:cue:151}{}這本文，
% cue 152 369.540 375.270
\hypertarget{w06:cue:152}{}我們讓司法來審查立法，
% cue 153 375.270 383.390
\hypertarget{w06:cue:153}{}作為一種最終決策的形式，
% cue 154 383.390 388.070
\hypertarget{w06:cue:154}{}它本文中指是，
% cue 155 388.070 390.390
\hypertarget{w06:cue:155}{}它認為在一個自由民主社會，
% cue 156 390.390 394.600
\hypertarget{w06:cue:156}{}我們讓司法來審查立法，
% cue 157 394.600 396.700
\hypertarget{w06:cue:157}{}然後作為一個最終的決策，
% cue 158 396.700 398.600
\hypertarget{w06:cue:158}{}這是一個不適當的，
% cue 159 398.600 399.520
\hypertarget{w06:cue:159}{}inappropriate\footnote{刊本 1348 頁末句：「This Essay will argue that judicial review of legislation is inappropriate as a mode of final decisionmaking in a free and democratic society.」（本文主張，在自由民主的社會，以法院審查立法作為終局決定的方式並不適當。）}
% cue 160 399.520 400.590
\hypertarget{w06:cue:160}{}tappropriate
% cue 161 400.590 401.410
\hypertarget{w06:cue:161}{}所以它反對，
% cue 162 401.410 402.090
\hypertarget{w06:cue:162}{}OK
% cue 163 402.090 402.850
\hypertarget{w06:cue:163}{}好 下面，
% cue 164 402.850 404.890
\hypertarget{w06:cue:164}{}第 3、4、9 還有比較，
% cue 165 404.890 405.890
\hypertarget{w06:cue:165}{}還有什麼要講嗎？

\tc{0:06:45}
\speaker{導讀人：亦凡}
% cue 166 405.890 407.390
\hypertarget{w06:cue:166}{}第 3、4、9，
% cue 167 407.390 408.390
\hypertarget{w06:cue:167}{}嗯，

\tc{0:06:48}
\speaker{老師}
% cue 168 408.390 409.390
\hypertarget{w06:cue:168}{}有嗎？
% cue 169 409.390 410.390
\hypertarget{w06:cue:169}{}有你就講，
% cue 170 410.390 411.390
\hypertarget{w06:cue:170}{}沒有就要跳，
% cue 171 411.390 412.390
\hypertarget{w06:cue:171}{}我不知道你要講什麼，
% cue 172 412.390 413.390
\hypertarget{w06:cue:172}{}你動作要快啊，
% cue 173 413.390 414.690
\hypertarget{w06:cue:173}{}不然我今天就講過來了。

\tc{0:06:54}
\speaker{導讀人：亦凡}
% cue 174 414.690 417.690
\hypertarget{w06:cue:174}{}好 那我跳到後面好了，

\tc{0:06:57}
\speaker{老師}
% cue 175 417.690 418.190
\hypertarget{w06:cue:175}{}多少，
% cue 176 418.190 420.770
\hypertarget{w06:cue:176}{}你多少年還在第一大段嗎，
% cue 177 420.770 421.750
\hypertarget{w06:cue:177}{}還在一嗎？
% cue 178 421.750 423.750
\hypertarget{w06:cue:178}{}還在 introduction 嗎？

\tc{0:07:03}
\speaker{導讀人：亦凡}
% cue 179 423.750 426.680
\hypertarget{w06:cue:179}{}可以看 introduction，
% cue 180 426.680 429.430
\hypertarget{w06:cue:180}{}但是就是他舉了一些例子，

\tc{0:07:09}
\speaker{老師}
% cue 181 429.430 431.530
\hypertarget{w06:cue:181}{}好那你講在哪裡啊？
% cue 182 431.530 432.500
\hypertarget{w06:cue:182}{}那什麼例子？
% cue 183 432.500 434.240
\hypertarget{w06:cue:183}{}他下面還會說到，
% cue 184 434.240 436.360
\hypertarget{w06:cue:184}{}他不斷的都是用這幾個例子，
% cue 185 436.360 437.720
\hypertarget{w06:cue:185}{}好你講在哪裡？

\tc{0:07:17}
\speaker{導讀人：亦凡}
% cue 186 437.720 440.550
\hypertarget{w06:cue:186}{}我覺得比較有意思的，
% cue 187 440.550 442.630
\hypertarget{w06:cue:187}{}就是他寫的方式比較有意思的是，
% cue 188 442.630 444.110
\hypertarget{w06:cue:188}{}一三五一頁後面，
% cue 189 444.110 446.750
\hypertarget{w06:cue:189}{}就是他就講，
% cue 190 446.750 451.780
\hypertarget{w06:cue:190}{}in recent years a number of books have appeared
% cue 191 451.780 454.180
\hypertarget{w06:cue:191}{}attacking judicial review in American
% cue 192 454.180 457.450
\hypertarget{w06:cue:192}{}就是他就是先講了一下，
% cue 193 457.450 460.690
\hypertarget{w06:cue:193}{}那個時間當時的，
% cue 194 460.690 463.810
\hypertarget{w06:cue:194}{}對於違憲審查的攻擊的狀，
% cue 195 463.810 466.930
\hypertarget{w06:cue:195}{}違憲審查的運行的狀態，
% cue 196 466.930 468.670
\hypertarget{w06:cue:196}{}跟被攻擊的狀態，
% cue 197 468.670 471.470
\hypertarget{w06:cue:197}{}就是他講隨著這個。

\tc{0:07:51}
\speaker{老師}
% cue 198 471.470 472.530
\hypertarget{w06:cue:198}{}在什麼地方？

\tc{0:07:52}
\speaker{學生}
% cue 199 472.530 477.710
\hypertarget{w06:cue:199}{}1350 的最下，

\tc{0:07:57}
\speaker{老師}
% cue 200 477.710 482.790
\hypertarget{w06:cue:200}{}1350 in recent years
% cue 201 482.790 484.500
\hypertarget{w06:cue:201}{}OK 好 趕快，
% cue 202 484.500 486.560
\hypertarget{w06:cue:202}{}The number of books have appeared
% cue 203 486.560 488.840
\hypertarget{w06:cue:203}{}Attacking judicial review in America
% cue 204 488.840 491.020
\hypertarget{w06:cue:204}{}我們這個小註十一\footnote{刊本 1350 頁註 11 所舉兩部攻擊違憲審查的著作：Larry D. Kramer, The People Themselves: Popular Constitutionalism and Judicial Review（2004）；Mark Tushnet, Taking the Constitution Away from the Courts（1999）。}，
% cue 205 491.020 492.890
\hypertarget{w06:cue:205}{}我們這兩個人都會讀到，
% cue 206 492.890 494.440
\hypertarget{w06:cue:206}{}It's a primary institution
% cue 207 494.440 495.890
\hypertarget{w06:cue:207}{}Question 然後下禮拜，
% cue 208 495.890 498.660
\hypertarget{w06:cue:208}{}For years, support of the practice
% cue 209 498.660 501.330
\hypertarget{w06:cue:209}{}Has come from liberals
% cue 210 501.330 505.270
\hypertarget{w06:cue:210}{}And opposition from conservative opponents of the right
% cue 211 505.270 508.070
\hypertarget{w06:cue:211}{}好 然後呢 要講什麼，
% cue 212 508.070 508.870
\hypertarget{w06:cue:212}{}你要講什麼？

\tc{0:08:28}
\speaker{導讀人：亦凡}
% cue 213 508.870 513.090
\hypertarget{w06:cue:213}{}就是他的這個寫法，
% cue 214 513.090 514.970
\hypertarget{w06:cue:214}{}我覺得蠻有意思的就是。

\tc{0:09:01}
\speaker{老師}
% cue 215 541.580 544.920
\hypertarget{w06:cue:215}{}當家座會有民主原則或國民主權原則，
% cue 216 544.920 547.420
\hypertarget{w06:cue:216}{}所以就擰不了讀問一個問題，
% cue 217 547.420 551.120
\hypertarget{w06:cue:217}{}那我們人民或者我們人民的代表所做的決定，
% cue 218 551.120 555.770
\hypertarget{w06:cue:218}{}那為什麼可以由沒有民意基礎的機關把它推翻掉呢？
% cue 219 555.770 559.060
\hypertarget{w06:cue:219}{}那這種反多數決的困難\footnote{語出 Alexander M. Bickel, The Least Dangerous Branch（1962）第 16 頁起題為 The Counter-Majoritarian Difficulty 的一節，首句為「The root difficulty is that judicial review is a counter-majoritarian force in our system.」（根本的難題在於，違憲審查在我們的體制裡是一股反多數的力量）。該節為課綱第 4 週指定文獻。\cite{bickel2004}。}究竟要怎麼理解，
% cue 220 559.060 560.270
\hypertarget{w06:cue:220}{}這是典型的，
% cue 221 560.270 563.920
\hypertarget{w06:cue:221}{}典型的討論贊成或者是反對的，
% cue 222 563.920 566.440
\hypertarget{w06:cue:222}{}所以他也承認了小註 8 那裡，
% cue 223 566.440 572.320
\hypertarget{w06:cue:223}{}違憲審查基本上是美式民主下的一種脫軌的制度\footnote{刊本 1349 頁引 Bickel 之語：「judicial review is a deviant institution in the American democracy」（違憲審查是美國民主之中一項偏離常軌的制度），註 8 註明出處為 The Least Dangerous Branch 第 17 至 18 頁。}，
% cue 224 572.660 575.380
\hypertarget{w06:cue:224}{}正常情形你就是在民主程序裡去決定，
% cue 225 575.380 577.910
\hypertarget{w06:cue:225}{}你為什麼就要給一個非民主的程序，
% cue 226 577.910 580.090
\hypertarget{w06:cue:226}{}來審查民主決定做的決定，
% cue 227 580.090 581.340
\hypertarget{w06:cue:227}{}能不能接受啊？

\tc{0:09:41}
% cue 228 581.340 584.370
\hypertarget{w06:cue:228}{}那這裡面有一些我們以前有講過，
% cue 229 584.370 587.790
\hypertarget{w06:cue:229}{}比方說我們認為這是因為憲法是高性，
% cue 230 587.790 590.100
\hypertarget{w06:cue:230}{}我們不是法官高於立法者，
% cue 231 590.100 592.940
\hypertarget{w06:cue:231}{}而是立法者通過的法律，
% cue 232 592.940 595.780
\hypertarget{w06:cue:232}{}按照憲法規定說他不能抵觸憲法，
% cue 233 595.780 597.460
\hypertarget{w06:cue:233}{}那麼什麼叫做法，
% cue 234 597.460 599.070
\hypertarget{w06:cue:234}{}以及什麼叫做憲法，
% cue 235 599.070 600.770
\hypertarget{w06:cue:235}{}就是我們法官來解釋，
% cue 236 600.770 602.350
\hypertarget{w06:cue:236}{}這不是我比他高，
% cue 237 602.350 604.930
\hypertarget{w06:cue:237}{}而是因為我正好有這個解釋的權，
% cue 238 604.930 605.660
\hypertarget{w06:cue:238}{}就這樣子，
% cue 239 605.660 607.000
\hypertarget{w06:cue:239}{}所以我就可以審查。

\tc{0:10:07}
% cue 240 607.000 608.380
\hypertarget{w06:cue:240}{}臺灣講太多了，
% cue 241 608.380 609.900
\hypertarget{w06:cue:241}{}那剩下的問題它沒有解決，
% cue 242 609.900 611.060
\hypertarget{w06:cue:242}{}這是刑事的論證，
% cue 243 611.060 613.500
\hypertarget{w06:cue:243}{}那為什麼你的決定就是最終的決定？
% cue 244 613.500 615.360
\hypertarget{w06:cue:244}{}為什麼你的決定有拘束力？
% cue 245 615.360 616.360
\hypertarget{w06:cue:245}{}你可以審查，
% cue 246 616.360 617.520
\hypertarget{w06:cue:246}{}你可以說不合，
% cue 247 617.520 619.960
\hypertarget{w06:cue:247}{}但是你沒有理由說我不適用，
% cue 248 619.960 622.260
\hypertarget{w06:cue:248}{}甚至於把整個的法律宣告違憲，
% cue 249 622.260 623.060
\hypertarget{w06:cue:249}{}失效，
% cue 250 623.060 624.410
\hypertarget{w06:cue:250}{}臺灣現在是失效\footnote{憲法訴訟法第 52 條第 1 項規定：「判決宣告法規範違憲且應失效者，該法規範自判決生效日起失效。但主文另有諭知溯及失效或定期失效者，依其諭知。」第 2 項規定定期失效的期間，法律位階法規範不得逾二年，命令位階法規範不得逾一年。\url{https://law.moj.gov.tw/LawClass/LawSingle.aspx?pcode=A0030159&flno=52}。}，
% cue 251 624.410 626.170
\hypertarget{w06:cue:251}{}這是他講的 strong case，
% cue 252 626.170 628.190
\hypertarget{w06:cue:252}{}strong judicial review\footnote{刊本 1354 頁的定義：在強式違憲審查之下，法院有權在個案中拒絕適用法律，或修正法律的效果（modify the effect of a statute），使其適用合於個人權利；法院並得以法律見解確立某一法律不予適用，經由 stare decisis（遵循先例）與爭點排除，使該法律實際上成為具文。更強的形式則授權法院把法律自法典中整個刪除。}
% cue 253 628.190 629.140
\hypertarget{w06:cue:253}{}不是 weak，
% cue 254 629.140 632.520
\hypertarget{w06:cue:254}{}那他認為 weak 的是另外一件事，
% cue 255 632.520 635.160
\hypertarget{w06:cue:255}{}他講的那個 core 主要是針對 strong，
% cue 256 635.160 635.980
\hypertarget{w06:cue:256}{}好了，
% cue 257 635.980 638.600
\hypertarget{w06:cue:257}{}這 1351 有什麼要講的嗎？

\tc{0:10:38}
\speaker{導讀人：亦凡}
% cue 258 638.600 640.220
\hypertarget{w06:cue:258}{}好，
% cue 259 640.220 641.420
\hypertarget{w06:cue:259}{}那沒，

\tc{0:10:41}
\speaker{老師}
% cue 260 641.420 643.180
\hypertarget{w06:cue:260}{}1351 有還是沒有？

\tc{0:10:43}
\speaker{導讀人：亦凡}
% cue 261 643.180 644.540
\hypertarget{w06:cue:261}{}有但是。

\section{核心論證的意義}

\tc{0:10:44}
\speaker{老師}
% cue 262 644.540 645.880
\hypertarget{w06:cue:262}{}好那就講吧，
% cue 263 645.880 648.770
\hypertarget{w06:cue:263}{}第一大段一開始的地方是不是，
% cue 264 648.770 653.750
\hypertarget{w06:cue:264}{}What I want to do is identify a core argument
% cue 265 653.750 655.090
\hypertarget{w06:cue:265}{}我剛開始講了，
% cue 266 655.090 656.950
\hypertarget{w06:cue:266}{}你是一個核心的論證，
% cue 267 656.950 660.600
\hypertarget{w06:cue:267}{}against that is independent of both
% cue 268 660.600 661.660
\hypertarget{w06:cue:268}{}and 的什麼？
% cue 269 661.660 663.040
\hypertarget{w06:cue:269}{}both A and B
% cue 270 663.040 667.750
\hypertarget{w06:cue:270}{}跟它的歷史上的 manifestation\footnote{刊本 1351 頁：「What I want to do is identify a core argument against judicial review that is independent of both its historical manifestations and questions about its particular effects」，historical manifestations 指違憲審查在各國歷史上呈現的具體樣貌，particular effects 指它作成的個別裁判及其好壞。} 無關，
% cue 271 667.750 671.860
\hypertarget{w06:cue:271}{}跟它的歷史的這個表現無關，
% cue 272 671.860 675.310
\hypertarget{w06:cue:272}{}也跟下面這個 question 無關，
% cue 273 675.310 677.850
\hypertarget{w06:cue:273}{}about its particular effects
% cue 274 677.850 680.710
\hypertarget{w06:cue:274}{}跟它的特別的效力，
% cue 275 680.710 682.270
\hypertarget{w06:cue:275}{}特定的效力無關，
% cue 276 682.270 683.970
\hypertarget{w06:cue:276}{}有些決定的，
% cue 277 683.970 684.670
\hypertarget{w06:cue:277}{}你可能喜歡，
% cue 278 684.670 685.810
\hypertarget{w06:cue:278}{}有些決定你不喜歡，
% cue 279 685.810 687.120
\hypertarget{w06:cue:279}{}這我不管，
% cue 280 687.120 688.500
\hypertarget{w06:cue:280}{}我現在要講的是，
% cue 281 688.500 692.390
\hypertarget{w06:cue:281}{}我要建構一個反違憲審查的核心論。

\tc{0:11:32}
% cue 282 692.390 694.400
\hypertarget{w06:cue:282}{}那怎麼樣叫做核心？
% cue 283 694.400 696.350
\hypertarget{w06:cue:283}{}他到這裡解釋了嗎？

\tc{0:11:36}
\speaker{導讀人：亦凡}
% cue 284 696.350 698.350
\hypertarget{w06:cue:284}{}馬上要解釋了，

\tc{0:11:38}
\speaker{老師}
% cue 285 698.350 700.510
\hypertarget{w06:cue:285}{}好 馬上看哪邊解釋了，

\tc{0:11:40}
\speaker{導讀人：亦凡}
% cue 286 700.510 700.870
\hypertarget{w06:cue:286}{}就是，

\tc{0:11:40}
\speaker{老師}
% cue 287 700.870 702.290
\hypertarget{w06:cue:287}{}叫到哪裡去，

\tc{0:11:42}
\speaker{導讀人：亦凡}
% cue 288 702.290 704.790
\hypertarget{w06:cue:288}{}1351 到 1353 的這，

\tc{0:11:44}
\speaker{老師}
% cue 289 704.790 706.650
\hypertarget{w06:cue:289}{}52 到 53，
% cue 290 706.650 707.490
\hypertarget{w06:cue:290}{}對，
% cue 291 707.490 708.710
\hypertarget{w06:cue:291}{}找一下吧，
% cue 292 708.710 710.170
\hypertarget{w06:cue:292}{}1353 嗎？

\tc{0:11:50}
\speaker{導讀人：亦凡}
% cue 293 710.170 711.710
\hypertarget{w06:cue:293}{}1355
% cue 294 711.710 714.570
\hypertarget{w06:cue:294}{}如果具體點的話就是 1352 吧，

\tc{0:11:54}
\speaker{老師}
% cue 295 714.570 715.510
\hypertarget{w06:cue:295}{}好 1352，
% cue 296 715.510 715.990
\hypertarget{w06:cue:296}{}來，

\tc{0:11:55}
\speaker{導讀人：亦凡}
% cue 297 715.990 718.900
\hypertarget{w06:cue:297}{}就是核心，
% cue 298 718.900 720.480
\hypertarget{w06:cue:298}{}他的想法核心就是，
% cue 299 720.480 723.080
\hypertarget{w06:cue:299}{}他要把那些跟美國，
% cue 300 723.080 726.830
\hypertarget{w06:cue:300}{}跟美國憲政歷史的，
% cue 301 726.830 728.530
\hypertarget{w06:cue:301}{}這種文化的，
% cue 302 728.530 731.010
\hypertarget{w06:cue:302}{}或者是個案性具體的這種，
% cue 303 731.010 733.450
\hypertarget{w06:cue:303}{}血肉\footnote{刊本 1351 頁，Waldron 引 Frank Michelman 為 Kramer《The People Themselves》所寫的封底推薦語，稱該書的歷史敘述「puts flesh on the bones of debates over judicial review and popular constitutionalism」；Waldron 自述要「take off some of the flesh and boil down the normative argument to its bare bones」（剔去部分血肉，把規範論證熬到只剩骨架）。}，
% cue 304 733.450 736.170
\hypertarget{w06:cue:304}{}從它的核心的骨架當中，
% cue 305 736.170 737.630
\hypertarget{w06:cue:305}{}剔除掉，

\tc{0:12:17}
\speaker{老師}
% cue 306 737.630 738.770
\hypertarget{w06:cue:306}{}那它怎麼還能，
% cue 307 738.770 740.910
\hypertarget{w06:cue:307}{}剔除掉它怎麼還能做骨架呢？
% cue 308 740.910 744.110
\hypertarget{w06:cue:308}{}再看一看你的理解，
% cue 309 744.110 745.910
\hypertarget{w06:cue:309}{}那下面還有嗎？
% cue 310 745.910 746.860
\hypertarget{w06:cue:310}{}1353

\tc{0:12:26}
\speaker{導讀人：亦凡}
% cue 311 746.860 752.980
\hypertarget{w06:cue:311}{}1353 到最後，
% cue 312 752.980 753.780
\hypertarget{w06:cue:312}{}最後一點吧，
% cue 313 753.780 755.540
\hypertarget{w06:cue:313}{}我覺得最後一段。

\section{本文的兩項攻擊理由}

\tc{0:12:35}
\speaker{老師}
% cue 314 755.540 758.080
\hypertarget{w06:cue:314}{}我們看一下 1353 的，
% cue 315 758.080 760.580
\hypertarget{w06:cue:315}{}In this essay 本文裡面，
% cue 316 760.580 767.010
\hypertarget{w06:cue:316}{}I should argue that a judicial review is vulnerable to attack on two fronts.
% cue 317 767.010 770.850
\hypertarget{w06:cue:317}{}違憲審查的兩個點子上是脆弱的。
% cue 318 770.850 783.320
\hypertarget{w06:cue:318}{}It does not provide a way for a society to focus clearly on the real issues at stake when citizens disagree about the right.
% cue 319 783.320 789.190
\hypertarget{w06:cue:319}{}當公民對於權利有不同見解的時候，
% cue 320 789.950 798.690
\hypertarget{w06:cue:320}{}違憲審查並沒有辦法讓這個社會聚焦在這個權利和真實的爭點上。

\tc{0:13:18}
% cue 321 798.690 800.770
\hypertarget{w06:cue:321}{}比方說墮胎跟反墮胎，
% cue 322 800.770 801.950
\hypertarget{w06:cue:322}{}你到底贊成什麼，
% cue 323 801.950 802.690
\hypertarget{w06:cue:323}{}你反對什麼，
% cue 324 802.690 804.570
\hypertarget{w06:cue:324}{}透過違憲審查，
% cue 325 804.570 809.190
\hypertarget{w06:cue:325}{}無助於讓公民聚焦在問題的本質上，
% cue 326 809.190 812.290
\hypertarget{w06:cue:326}{}你可能會講大法官的解釋方法是什麼，
% cue 327 812.290 813.530
\hypertarget{w06:cue:327}{}你憑什麼這樣子講，
% cue 328 813.530 814.640
\hypertarget{w06:cue:328}{}它會歪曲掉，
% cue 329 814.640 816.360
\hypertarget{w06:cue:329}{}會扭到別的地方去，
% cue 330 816.360 818.130
\hypertarget{w06:cue:330}{}對於解決問題，
% cue 331 818.130 819.070
\hypertarget{w06:cue:331}{}真實的問題，
% cue 332 819.070 820.990
\hypertarget{w06:cue:332}{}real issue 沒有幫助。

\tc{0:13:40}
% cue 333 820.990 828.930
\hypertarget{w06:cue:333}{}另外 on the contrary, it is distracting with a side issues about the precedent of text and interpretation\footnote{刊本 1353 頁原文為「on the contrary, it distracts them with side-issues about precedent, texts, and interpretation」，所列為先例、文本與解釋三項旁枝問題。}，
% cue 334 828.930 833.780
\hypertarget{w06:cue:334}{}違憲審查做的決定，大法官做的決定會使人們導向，
% cue 335 833.780 838.650
\hypertarget{w06:cue:335}{}有沒有符合解釋先例啊，這樣子有沒有脫離文義啊，
% cue 336 838.650 842.320
\hypertarget{w06:cue:336}{}這樣的解釋方法是不正確啊，不變得到那裡去，
% cue 337 842.320 846.990
\hypertarget{w06:cue:337}{}但是對於墮胎與否，墮胎要不要承認它是一個權利，
% cue 338 846.990 850.480
\hypertarget{w06:cue:338}{}這裡面牽涉的真實爭點，人會忽略掉。

\tc{0:14:10}
% cue 339 850.480 857.500
\hypertarget{w06:cue:339}{}And it is politically legitimate so far as democratic values are concerned
% cue 340 857.500 860.570
\hypertarget{w06:cue:340}{}那講到政治的正當性，
% cue 341 860.570 863.390
\hypertarget{w06:cue:341}{}如果牽涉到這個民主的價值的話，
% cue 342 863.390 865.820
\hypertarget{w06:cue:342}{}By privileging majority voting
% cue 343 865.820 870.550
\hypertarget{w06:cue:343}{}among a small number of unelected, unaccountable judges
% cue 344 870.550 875.330
\hypertarget{w06:cue:344}{}你為什麼讓那些沒有選舉的極少數的大法官，
% cue 345 875.330 877.480
\hypertarget{w06:cue:345}{}他們有權利來做這個決定。

\tc{0:14:37}
% cue 346 877.480 881.820
\hypertarget{w06:cue:346}{}It's it disenfranchises ordinary citizens
% cue 347 881.820 885.690
\hypertarget{w06:cue:347}{}你這樣做等於是剝奪了一般公民的投票權，
% cue 348 885.690 899.750
\hypertarget{w06:cue:348}{}當我們最終要決定要不要在這個社會承認某種權利的時候，
% cue 349 899.750 906.790
\hypertarget{w06:cue:349}{}那我們一向引以一向珍惜一向遵守的所謂代議的原則，
% cue 350 906.790 910.370
\hypertarget{w06:cue:350}{}還有政治平等每人都可以參與的這個原則，
% cue 351 910.370 914.700
\hypertarget{w06:cue:351}{}實際上就被你們當法官給掃除了，
% cue 352 914.700 917.050
\hypertarget{w06:cue:352}{}你們當法官也不是我們選的，
% cue 353 917.050 917.910
\hypertarget{w06:cue:353}{}你們有什麼代意，
% cue 354 917.910 918.850
\hypertarget{w06:cue:354}{}你們有什麼，
% cue 355 918.850 922.890
\hypertarget{w06:cue:355}{}我們怎麼在你們身上體現出政治參與的平等，
% cue 356 922.890 924.060
\hypertarget{w06:cue:356}{}沒有啊，

\tc{0:15:24}
\speaker{導讀人：亦凡}
% cue 357 924.060 925.690
\hypertarget{w06:cue:357}{}好了。

\section{本文架構}

\tc{0:15:25}
\speaker{老師}
% cue 358 925.690 931.050
\hypertarget{w06:cue:358}{}那麼下面看一下他怎麼做，
% cue 359 931.050 933.000
\hypertarget{w06:cue:359}{}本文如何建構，
% cue 360 933.000 935.270
\hypertarget{w06:cue:360}{}第三五三頁來你講，
% cue 361 935.270 940.850
\hypertarget{w06:cue:361}{}我將進步以下列的方法來進行，

\tc{0:15:40}
\speaker{導讀人：亦凡}
% cue 362 940.850 947.360
\hypertarget{w06:cue:362}{}In part one, I will define the target of my argument

\tc{0:15:47}
\speaker{老師}
% cue 363 947.360 950.910
\hypertarget{w06:cue:363}{}OK，它叫做 strong judicial review of legislation，
% cue 364 950.910 954.640
\hypertarget{w06:cue:364}{}它講的是強烈的違憲審查，
% cue 365 954.640 956.310
\hypertarget{w06:cue:365}{}不是 weak judicial review\footnote{刊本 1355 頁的定義：「courts may scrutinize legislation for its conformity to individual rights but they may not decline to apply it (or moderate its application) simply because rights would otherwise be violated.」（法院得審查立法是否合於個人權利，但不得僅因權利將受侵害而拒絕適用或緩和其適用。）Waldron 以英國、紐西蘭為例，並引 Stephen Gardbaum, The New Commonwealth Model of Constitutionalism, 49 American Journal of Comparative Law 707（2001）。}，
% cue 366 956.310 961.660
\hypertarget{w06:cue:366}{}OK，然後呢，weak 不是它要討論，
% cue 367 961.660 964.780
\hypertarget{w06:cue:367}{}它要講的是具有拘束力，
% cue 368 964.780 967.220
\hypertarget{w06:cue:368}{}甚至於可以讓它失效。

\tc{0:16:07}
% cue 369 967.220 968.160
\hypertarget{w06:cue:369}{}像我們這種，
% cue 370 968.160 970.430
\hypertarget{w06:cue:370}{}有個案拒絕適用，
% cue 371 970.430 972.340
\hypertarget{w06:cue:371}{}像我們的行政法院的法官，
% cue 372 972.340 973.890
\hypertarget{w06:cue:372}{}當它發現呢，
% cue 373 973.890 976.050
\hypertarget{w06:cue:373}{}它裁判所需要依據的命令，
% cue 374 976.050 979.910
\hypertarget{w06:cue:374}{}他沒有權利宣告法律失效，
% cue 375 979.910 984.820
\hypertarget{w06:cue:375}{}但是他可以拒絕適用他認為違法的命令\footnote{司法院釋字第 216 號解釋（1987 年 6 月 19 日）：「各機關依其職掌就有關法規為釋示之行政命令，法官於審判案件時，固可予以引用，但仍得依據法律，表示適當之不同見解，並不受其拘束」，並稱釋字第 137 號解釋即本此意旨。中華民國憲法第 80 條規定法官依據法律獨立審判，不受任何干涉。}在本案，
% cue 376 984.820 986.150
\hypertarget{w06:cue:376}{}這是一種，
% cue 377 986.150 987.860
\hypertarget{w06:cue:377}{}這也是他所謂的 strong，
% cue 378 987.860 989.770
\hypertarget{w06:cue:378}{}另外一個 strong 是更 strong，
% cue 379 989.770 993.830
\hypertarget{w06:cue:379}{}就像我們大法官直接宣告法律違憲而失效，
% cue 380 993.830 995.390
\hypertarget{w06:cue:380}{}從此失去效力，
% cue 381 995.390 998.090
\hypertarget{w06:cue:381}{}那失效都是向後失效，
% cue 382 998.090 1000.050
\hypertarget{w06:cue:382}{}但是如果跟上次上禮拜講過，
% cue 383 1000.050 1002.170
\hypertarget{w06:cue:383}{}如果更嚴重也有一種可能，
% cue 384 1002.170 1003.750
\hypertarget{w06:cue:384}{}他也許會溯及失效，
% cue 385 1003.750 1006.030
\hypertarget{w06:cue:385}{}那因為這個牽涉到法的安定性，
% cue 386 1006.030 1009.020
\hypertarget{w06:cue:386}{}所以通常大法官，違憲審查不會允許這樣。

\tc{0:16:49}
% cue 387 1009.020 1035.480
\hypertarget{w06:cue:387}{}但不管怎樣，他要講的本文要針對的對象，違憲審查，是指 strong judicial review，不是 weak，好吧，這是第一件事情，以上是 part one，然後他先，他先界定他要打擊的對象，本文所謂的違憲審查是指強的違憲審查，OK，下一遍，第二，part two,part two will set up some assumptions on which my argument is predicated。

\tc{0:17:15}
% cue 388 1035.480 1050.630
\hypertarget{w06:cue:388}{}我們講一解釋一下這個 part two，
% cue 389 1050.630 1053.460
\hypertarget{w06:cue:389}{}為什麼剛才講 part two 有四個 assumption，
% cue 390 1053.460 1053.900
\hypertarget{w06:cue:390}{}因為，

\tc{0:17:33}
\speaker{導讀人：亦凡}
% cue 391 1053.900 1056.560
\hypertarget{w06:cue:391}{}因為它就是理想，
% cue 392 1056.560 1058.520
\hypertarget{w06:cue:392}{}我認為它就是一個理想當中的，
% cue 393 1058.520 1062.030
\hypertarget{w06:cue:393}{}民主社會運行的一個狀態。

\tc{0:17:42}
\speaker{老師}
% cue 394 1062.030 1062.930
\hypertarget{w06:cue:394}{}你可以這樣子講，
% cue 395 1062.930 1064.850
\hypertarget{w06:cue:395}{}但是它為什麼要有四個假設，
% cue 396 1064.850 1068.160
\hypertarget{w06:cue:396}{}因為它是在這四個假設之下，
% cue 397 1068.160 1070.590
\hypertarget{w06:cue:397}{}它來攻擊違憲審查，
% cue 398 1070.590 1072.740
\hypertarget{w06:cue:398}{}那四個假設如果不在的時候，
% cue 399 1072.740 1074.620
\hypertarget{w06:cue:399}{}還要不要批評違憲審查，
% cue 400 1074.620 1075.550
\hypertarget{w06:cue:400}{}是另一個方面，
% cue 401 1075.550 1079.390
\hypertarget{w06:cue:401}{}也就是說它不是一個頑固分子，
% cue 402 1079.390 1080.750
\hypertarget{w06:cue:402}{}它有通融的一面，
% cue 403 1080.750 1082.800
\hypertarget{w06:cue:403}{}我們一般講，
% cue 404 1082.800 1085.060
\hypertarget{w06:cue:404}{}我是泛泛地講我反對，
% cue 405 1085.060 1086.520
\hypertarget{w06:cue:405}{}但是我之所以反對，
% cue 406 1086.520 1089.600
\hypertarget{w06:cue:406}{}是基於下面有四個假設的情況下，
% cue 407 1089.600 1090.160
\hypertarget{w06:cue:407}{}我反對，
% cue 408 1090.160 1091.490
\hypertarget{w06:cue:408}{}四個假設，
% cue 409 1091.490 1093.710
\hypertarget{w06:cue:409}{}那我覺得是沒有理由，
% cue 410 1093.710 1097.240
\hypertarget{w06:cue:410}{}讓違憲審查來取代一般的立法者，
% cue 411 1097.240 1103.070
\hypertarget{w06:cue:411}{}可是，如果這四個假設不存在，那我就不一定統統反對。

\tc{0:18:23}
% cue 412 1103.070 1105.920
\hypertarget{w06:cue:412}{}繼續，In Part 3,

\tc{0:18:25}
\speaker{導讀人：亦凡}
% cue 413 1105.920 1111.980
\hypertarget{w06:cue:413}{}In Part 3, I will reveal the general character of the arguments I proposed to make.

\tc{0:18:31}
\speaker{老師}
% cue 414 1111.980 1120.020
\hypertarget{w06:cue:414}{}好，General Character，剛才我不是跟你講說，我們如果從名稱上是看不懂嗎？
% cue 415 1120.020 1124.140
\hypertarget{w06:cue:415}{}論證，什麼叫做 General Character of Arguing?
% cue 416 1124.140 1128.090
\hypertarget{w06:cue:416}{}它打算要提出來的論證的一般性質。
% cue 417 1128.090 1132.890
\hypertarget{w06:cue:417}{}那這個一般性質他剛剛他用的那個用語各位還有印象嗎，
% cue 418 1132.890 1136.220
\hypertarget{w06:cue:418}{}他的用語叫做什麼 form of argument，
% cue 419 1136.220 1138.050
\hypertarget{w06:cue:419}{}論證的形式吧，
% cue 420 1138.050 1139.260
\hypertarget{w06:cue:420}{}各位有沒有記住啊，
% cue 421 1139.260 1140.430
\hypertarget{w06:cue:421}{}一般，
% cue 422 1140.430 1141.850
\hypertarget{w06:cue:422}{}沒有啊，
% cue 423 1141.850 1143.580
\hypertarget{w06:cue:423}{}the form of the argument
% cue 424 1143.580 1145.010
\hypertarget{w06:cue:424}{}論證的形式，
% cue 425 1145.010 1147.680
\hypertarget{w06:cue:425}{}他論證的形式就是在講，
% cue 426 1147.680 1151.080
\hypertarget{w06:cue:426}{}我批評違憲審查，我提出來的論證，
% cue 427 1151.080 1153.120
\hypertarget{w06:cue:427}{}他的一般性質是怎樣，
% cue 428 1153.120 1154.770
\hypertarget{w06:cue:428}{}等下我們看看他怎麼一般性質。

\tc{0:19:14}
% cue 429 1154.770 1162.990
\hypertarget{w06:cue:429}{}然後，that argument will attend to both outcomes and process related reasons，
% cue 430 1162.990 1165.080
\hypertarget{w06:cue:430}{}你繼續講這句什麼意思？

\tc{0:19:25}
\speaker{導讀人：亦凡}
% cue 431 1165.080 1170.080
\hypertarget{w06:cue:431}{}它就是分為從結果面上來看一個違憲審查的案，

\tc{0:19:30}
\speaker{老師}
% cue 432 1170.080 1173.840
\hypertarget{w06:cue:432}{}你覺得它是分嗎，還是 both,and，
% cue 433 1173.840 1177.630
\hypertarget{w06:cue:433}{}它說我的論證會兼顧到，

\tc{0:19:37}
\speaker{導讀人：亦凡}
% cue 434 1177.630 1179.300
\hypertarget{w06:cue:434}{}對，兼顧到，

\tc{0:19:39}
\speaker{老師}
% cue 435 1179.300 1183.020
\hypertarget{w06:cue:435}{}and these will be discussed in part 4 and 5
% cue 436 1183.020 1187.440
\hypertarget{w06:cue:436}{}它的論證形式，就是它的論證的性質，
% cue 437 1187.440 1192.010
\hypertarget{w06:cue:437}{}它會同時兼顧結果相關的理由，
% cue 438 1192.010 1193.950
\hypertarget{w06:cue:438}{}和程序相關的理由，
% cue 439 1193.950 1196.210
\hypertarget{w06:cue:439}{}有些是結果面的理由，
% cue 440 1196.210 1197.430
\hypertarget{w06:cue:440}{}有些是程序面的理由，
% cue 441 1197.430 1198.070
\hypertarget{w06:cue:441}{}它會兼顧，
% cue 442 1198.070 1200.330
\hypertarget{w06:cue:442}{}然後呢，
% cue 443 1200.330 1202.190
\hypertarget{w06:cue:443}{}In part 4 來繼續，
% cue 444 1202.190 1203.740
\hypertarget{w06:cue:444}{}In part 4 and 5
% cue 445 1203.740 1204.960
\hypertarget{w06:cue:445}{}它分別講的就是，
% cue 446 1204.960 1207.110
\hypertarget{w06:cue:446}{}結果相關的理由，
% cue 447 1207.110 1211.090
\hypertarget{w06:cue:447}{}第 5 個部分就是程序相關的理由，
% cue 448 1211.090 1211.750
\hypertarget{w06:cue:448}{}好 講完了。

\tc{0:20:11}
% cue 449 1211.750 1213.820
\hypertarget{w06:cue:449}{}In part 6 第 6 部分，

\tc{0:20:13}
\speaker{導讀人：亦凡}
% cue 450 1213.820 1215.620
\hypertarget{w06:cue:450}{}In part 6
% cue 451 1215.620 1219.440
\hypertarget{w06:cue:451}{}I will expose the fallacy of the most common argument
% cue 452 1219.440 1224.070
\hypertarget{w06:cue:452}{}against allowing representative institutions to prevail
% cue 453 1224.070 1233.210
\hypertarget{w06:cue:453}{}that such a system inevitably leads to the tyranny of the majority。

\tc{0:20:33}
\speaker{老師}
% cue 454 1233.210 1234.450
\hypertarget{w06:cue:454}{}好 這是什麼意思？
% cue 455 1234.450 1237.210
\hypertarget{w06:cue:455}{}這句話這一段有點翻過來，
% cue 456 1237.210 1239.070
\hypertarget{w06:cue:456}{}它要怎麼樣，
% cue 457 1239.070 1242.940
\hypertarget{w06:cue:457}{}它要去 expose the fallacy，
% cue 458 1242.940 1247.060
\hypertarget{w06:cue:458}{}要去闡述這樣的論證，
% cue 459 1247.060 1250.400
\hypertarget{w06:cue:459}{}一般的人反對，
% cue 460 1250.400 1254.990
\hypertarget{w06:cue:460}{}讓這個代議機關勝出，
% cue 461 1254.990 1258.430
\hypertarget{w06:cue:461}{}就是大家認為應該有違憲審查，
% cue 462 1258.430 1260.740
\hypertarget{w06:cue:462}{}不應該讓代議機關勝出，
% cue 463 1260.740 1261.910
\hypertarget{w06:cue:463}{}主要的理由就是，
% cue 464 1261.910 1265.570
\hypertarget{w06:cue:464}{}因為代議機關會有多數欺負少數的多數暴力問題。

\tc{0:21:05}
% cue 465 1265.570 1270.830
\hypertarget{w06:cue:465}{}但是他要 expose the fallacy of this argument，
% cue 466 1270.830 1275.550
\hypertarget{w06:cue:466}{}他要破解這樣的論證的錯誤，
% cue 467 1275.550 1278.690
\hypertarget{w06:cue:467}{}你們都認為有多數暴力，
% cue 468 1278.690 1280.990
\hypertarget{w06:cue:468}{}所以要有違憲審查，
% cue 469 1280.990 1282.490
\hypertarget{w06:cue:469}{}跟你們講錯了，
% cue 470 1282.490 1284.430
\hypertarget{w06:cue:470}{}這一節第六節，
% cue 471 1284.430 1286.100
\hypertarget{w06:cue:471}{}就專門跟你講是錯戰，
% cue 472 1286.100 1288.300
\hypertarget{w06:cue:472}{}所以它是反過來，
% cue 473 1288.300 1289.820
\hypertarget{w06:cue:473}{}它平常就是我反對，
% cue 474 1289.820 1290.700
\hypertarget{w06:cue:474}{}我反對我反對，
% cue 475 1290.700 1293.020
\hypertarget{w06:cue:475}{}它現在是反對你們反對，
% cue 476 1293.020 1296.120
\hypertarget{w06:cue:476}{}所以它也是反對，
% cue 477 1296.120 1298.400
\hypertarget{w06:cue:477}{}因為你們搞錯了，
% cue 478 1298.400 1299.140
\hypertarget{w06:cue:478}{}所以你們反對。

\tc{0:21:39}
% cue 479 1299.140 1300.150
\hypertarget{w06:cue:479}{}finally

\tc{0:21:40}
\speaker{導讀人：亦凡}
% cue 480 1300.150 1303.940
\hypertarget{w06:cue:480}{}finally in part 7
% cue 481 1303.940 1307.440
\hypertarget{w06:cue:481}{}I shall say a little bit about non-core cases
% cue 482 1307.440 1310.590
\hypertarget{w06:cue:482}{}that is cases in which there is reason
% cue 483 1310.590 1312.070
\hypertarget{w06:cue:483}{}to depart from the assumptions
% cue 484 1312.070 1314.910
\hypertarget{w06:cue:484}{}on which the core argument depends

\tc{0:21:54}
\speaker{老師}
% cue 485 1314.910 1320.160
\hypertarget{w06:cue:485}{}跟我剛才前面講的那四個假設不一樣的情況，
% cue 486 1320.160 1323.760
\hypertarget{w06:cue:486}{}那我批評它們是基於這四個假設，
% cue 487 1323.760 1325.680
\hypertarget{w06:cue:487}{}那這個四個假設如果變了呢，
% cue 488 1325.680 1327.020
\hypertarget{w06:cue:488}{}那我可能會有點改變，
% cue 489 1327.020 1329.780
\hypertarget{w06:cue:489}{}我講一下如果非典型的情況是怎麼樣。

\section{強式與弱式違憲審查}

\tc{0:22:09}
% cue 490 1329.780 1330.980
\hypertarget{w06:cue:490}{}好，我們往下來，
% cue 491 1330.980 1332.840
\hypertarget{w06:cue:491}{}一、二、一間三百五十三頁，
% cue 492 1332.840 1334.500
\hypertarget{w06:cue:492}{}代分層選後的絕對修理，
% cue 493 1334.500 1336.500
\hypertarget{w06:cue:493}{}應該是中間才有代分層，
% cue 494 1336.500 1338.840
\hypertarget{w06:cue:494}{}來，違憲審查的定義，
% cue 495 1338.840 1340.710
\hypertarget{w06:cue:495}{}或者司法審查的定義，
% cue 496 1340.710 1342.850
\hypertarget{w06:cue:496}{}從哪裡看，試看？

\tc{0:22:22}
\speaker{導讀人：亦凡}
% cue 497 1342.850 1345.530
\hypertarget{w06:cue:497}{}從一、三、五、四中間，
% cue 498 1345.530 1354.340
\hypertarget{w06:cue:498}{}如果核心句子的話，
% cue 499 1354.340 1355.080
\hypertarget{w06:cue:499}{}就是它說，
% cue 500 1355.080 1357.400
\hypertarget{w06:cue:500}{}My target is strong judicial review
% cue 501 1357.400 1358.590
\hypertarget{w06:cue:501}{}然後它。

\tc{0:22:38}
\speaker{老師}
% cue 502 1358.590 1360.310
\hypertarget{w06:cue:502}{}下面嘛，

\tc{0:22:40}
\speaker{導讀人：亦凡}
% cue 503 1360.310 1363.700
\hypertarget{w06:cue:503}{}In a system of strong judicial review

\tc{0:22:43}
\speaker{老師}
% cue 504 1363.700 1365.060
\hypertarget{w06:cue:504}{}從這一段開始來，

\tc{0:22:45}
\speaker{導讀人：亦凡}
% cue 505 1365.060 1377.150
\hypertarget{w06:cue:505}{}法院有權利在個案當中，

\tc{0:22:57}
\speaker{老師}
% cue 506 1377.150 1379.760
\hypertarget{w06:cue:506}{}拒絕適用某個法律，
% cue 507 1379.760 1382.010
\hypertarget{w06:cue:507}{}or

\tc{0:23:02}
\speaker{導讀人：亦凡}
% cue 508 1382.010 1389.370
\hypertarget{w06:cue:508}{}甚至於去修改修正那個法律的效果，

\tc{0:23:09}
\speaker{老師}
% cue 509 1389.370 1394.890
\hypertarget{w06:cue:509}{}讓它的適用，
% cue 510 1394.890 1396.090
\hypertarget{w06:cue:510}{}讓那個法律的適用，
% cue 511 1396.090 1397.890
\hypertarget{w06:cue:511}{}能夠符合個案的權利，
% cue 512 1397.890 1401.210
\hypertarget{w06:cue:512}{}人民的個人的權利，
% cue 513 1401.210 1404.230
\hypertarget{w06:cue:513}{}就是你認為它是違憲，
% cue 514 1404.230 1405.810
\hypertarget{w06:cue:514}{}甚至也還有權利，
% cue 515 1405.810 1408.230
\hypertarget{w06:cue:515}{}可以去修改那個法律的效力。

\tc{0:23:28}
% cue 516 1408.230 1409.570
\hypertarget{w06:cue:516}{}這句話講得很抽象，
% cue 517 1409.570 1410.730
\hypertarget{w06:cue:517}{}我不知道它怎麼修改，
% cue 518 1410.730 1412.600
\hypertarget{w06:cue:518}{}比方說我舉個例子，
% cue 519 1412.600 1414.760
\hypertarget{w06:cue:519}{}大法官解釋的方法當中，
% cue 520 1414.760 1417.930
\hypertarget{w06:cue:520}{}是不是有用到所謂的合憲性解釋\footnote{法律有數種解釋可能時，採取與憲法相符的解釋，排除牴觸憲法的解釋。司法院釋字第 572 號解釋理由書謂法官聲請解釋時，「如僅對法律是否違憲發生疑義，或系爭法律有合憲解釋之可能者，尚難謂已提出客觀上形成確信法律為違憲之具體理由」。}，
% cue 521 1417.930 1419.320
\hypertarget{w06:cue:521}{}它有很多種可能，
% cue 522 1419.320 1422.870
\hypertarget{w06:cue:522}{}但是我把它限縮成這一種，
% cue 523 1422.870 1424.220
\hypertarget{w06:cue:523}{}跟憲法意旨，
% cue 524 1424.220 1426.440
\hypertarget{w06:cue:524}{}或者跟人權保障相符的，
% cue 525 1426.440 1428.040
\hypertarget{w06:cue:525}{}其他的我就不用了，
% cue 526 1428.040 1429.870
\hypertarget{w06:cue:526}{}這個算不算是他講的這個。

\tc{0:23:49}
% cue 527 1429.870 1432.280
\hypertarget{w06:cue:527}{}我們現在是用大陸法系的話，
% cue 528 1432.280 1433.540
\hypertarget{w06:cue:528}{}用我們這邊理解的話，
% cue 529 1433.540 1435.450
\hypertarget{w06:cue:529}{}來講他這句很抽象的話，
% cue 530 1435.450 1437.000
\hypertarget{w06:cue:530}{}他這句話講得太抽象，
% cue 531 1437.000 1439.810
\hypertarget{w06:cue:531}{}因為 Modify 有幾千幾百種方法，
% cue 532 1439.810 1443.810
\hypertarget{w06:cue:532}{}然後它 Modify 的是 Effect 效果，
% cue 533 1443.810 1446.100
\hypertarget{w06:cue:533}{}一個違憲的法律，
% cue 534 1446.100 1447.840
\hypertarget{w06:cue:534}{}一個法律的效果，
% cue 535 1447.840 1449.340
\hypertarget{w06:cue:535}{}我們指它修改成，
% cue 536 1449.340 1452.410
\hypertarget{w06:cue:536}{}它符合憲法有關於人權規定，
% cue 537 1452.410 1454.110
\hypertarget{w06:cue:537}{}跟人權規定相符，
% cue 538 1454.110 1456.370
\hypertarget{w06:cue:538}{}那最我能最能想象到，
% cue 539 1456.370 1458.830
\hypertarget{w06:cue:539}{}是不是就是合憲性解釋方法，
% cue 540 1458.830 1462.420
\hypertarget{w06:cue:540}{}為什麼我只採用其中一部分，
% cue 541 1462.420 1465.870
\hypertarget{w06:cue:541}{}因為其他都跟憲法意旨相反，
% cue 542 1465.870 1467.570
\hypertarget{w06:cue:542}{}都跟人權保障相反，
% cue 543 1467.570 1469.120
\hypertarget{w06:cue:543}{}所以我只取這一塊。

\tc{0:24:29}
% cue 544 1469.120 1471.190
\hypertarget{w06:cue:544}{}你們有興趣看一下，
% cue 545 1471.190 1473.830
\hypertarget{w06:cue:545}{}釋字 272 號幾次，
% cue 546 1473.830 1475.550
\hypertarget{w06:cue:546}{}你們就聽懂我在講什麼，
% cue 547 1475.550 1477.740
\hypertarget{w06:cue:547}{}我不知道，
% cue 548 1477.740 1480.440
\hypertarget{w06:cue:548}{}反正現在老師錄的視頻有太多，
% cue 549 1480.440 1481.420
\hypertarget{w06:cue:549}{}我也搞不清楚，
% cue 550 1481.420 1482.520
\hypertarget{w06:cue:550}{}老師我講到這個，
% cue 551 1482.520 1484.300
\hypertarget{w06:cue:551}{}但是總之你們可以看得到，
% cue 552 1484.300 1485.680
\hypertarget{w06:cue:552}{}我們有這樣的課，
% cue 553 1485.680 1487.070
\hypertarget{w06:cue:553}{}正在講這些東西，
% cue 554 1487.070 1487.640
\hypertarget{w06:cue:554}{}好來，
% cue 555 1487.640 1488.320
\hypertarget{w06:cue:555}{}Moreover。

\tc{0:24:48}
\speaker{導讀人：亦凡}
% cue 556 1488.320 1491.360
\hypertarget{w06:cue:556}{}因為我覺得如果他上面這句話，
% cue 557 1491.360 1493.620
\hypertarget{w06:cue:557}{}在下面一段有一個對應，
% cue 558 1493.620 1495.470
\hypertarget{w06:cue:558}{}就是在那個，
% cue 559 1495.470 1496.290
\hypertarget{w06:cue:559}{}另外一個，
% cue 560 1496.290 1498.250
\hypertarget{w06:cue:560}{}in a system of weak judicial
% cue 561 1498.250 1499.910
\hypertarget{w06:cue:561}{}就是他那個後面，
% cue 562 1499.910 1501.090
\hypertarget{w06:cue:562}{}就是裡面他就講，
% cue 563 1501.090 1503.630
\hypertarget{w06:cue:563}{}一三，
% cue 564 1503.630 1505.270
\hypertarget{w06:cue:564}{}一三五五頁，
% cue 565 1505.270 1508.630
\hypertarget{w06:cue:565}{}一三五五頁的那一段裡面，
% cue 566 1508.630 1509.850
\hypertarget{w06:cue:566}{}他就是講，
% cue 567 1509.850 1513.120
\hypertarget{w06:cue:567}{}他就是在相反的情況下講，
% cue 568 1513.120 1515.300
\hypertarget{w06:cue:568}{}若是 Weak Judicial Review 的時候，
% cue 569 1515.300 1515.900
\hypertarget{w06:cue:569}{}他就講，
% cue 570 1515.900 1518.160
\hypertarget{w06:cue:570}{}在那個括號裡面就講說，
% cue 571 1518.160 1520.260
\hypertarget{w06:cue:571}{}or moderate its application。

\tc{0:25:20}
\speaker{老師}
% cue 572 1520.260 1524.540
\hypertarget{w06:cue:572}{}他是適用或者修正其適用，
% cue 573 1524.540 1525.070
\hypertarget{w06:cue:573}{}是，
% cue 574 1525.070 1526.230
\hypertarget{w06:cue:574}{}然後你要說什麼？
% cue 575 1526.230 1529.250
\hypertarget{w06:cue:575}{}你覺得修正其適用，

\tc{0:25:29}
\speaker{導讀人：亦凡}
% cue 576 1529.250 1532.030
\hypertarget{w06:cue:576}{}這一段他是講，
% cue 577 1532.030 1533.370
\hypertarget{w06:cue:577}{}Weak Judicial Review
% cue 578 1533.370 1535.110
\hypertarget{w06:cue:578}{}不能做的事情，
% cue 579 1535.110 1536.890
\hypertarget{w06:cue:579}{}所以我在想他這個是不是就是，
% cue 580 1536.890 1540.380
\hypertarget{w06:cue:580}{}前面 strong judicial review 可以做的事情，
% cue 581 1540.380 1541.260
\hypertarget{w06:cue:581}{}是啊。

\tc{0:25:41}
\speaker{老師}
% cue 582 1541.260 1545.730
\hypertarget{w06:cue:582}{}那現在我說 strong judicial review 可以做的事情，
% cue 583 1545.730 1548.330
\hypertarget{w06:cue:583}{}那我沒有挑戰，
% cue 584 1548.330 1550.550
\hypertarget{w06:cue:584}{}也沒有質疑他的分類，
% cue 585 1550.550 1554.810
\hypertarget{w06:cue:585}{}我只是說他現在講的這個話不容易懂，
% cue 586 1554.810 1559.170
\hypertarget{w06:cue:586}{}什麼叫做 modify the effect of a statute？
% cue 587 1559.170 1563.250
\hypertarget{w06:cue:587}{}你如何修改一個綠的效率，
% cue 588 1563.250 1564.350
\hypertarget{w06:cue:588}{}你怎麼修改？
% cue 589 1564.350 1566.560
\hypertarget{w06:cue:589}{}你怎麼把它調整，
% cue 590 1566.560 1568.460
\hypertarget{w06:cue:590}{}那我的意思就是，
% cue 591 1568.460 1569.700
\hypertarget{w06:cue:591}{}我隨便指一個例子，
% cue 592 1569.700 1571.120
\hypertarget{w06:cue:592}{}我馬上想到的就是，
% cue 593 1571.120 1572.400
\hypertarget{w06:cue:593}{}合憲性解釋方法，
% cue 594 1572.400 1574.090
\hypertarget{w06:cue:594}{}合憲性解釋方法。

\tc{0:26:14}
% cue 595 1574.090 1575.170
\hypertarget{w06:cue:595}{}我需要再多講嗎？
% cue 596 1575.170 1576.610
\hypertarget{w06:cue:596}{}有沒有畫個圈圈講？
% cue 597 1576.610 1577.830
\hypertarget{w06:cue:597}{}現在是個大圈圈，
% cue 598 1577.830 1580.270
\hypertarget{w06:cue:598}{}這個法條或者這個法律，
% cue 599 1580.270 1582.530
\hypertarget{w06:cue:599}{}它可能解釋的範圍這麼大，
% cue 600 1582.530 1584.080
\hypertarget{w06:cue:600}{}但是大法官說，
% cue 601 1584.080 1585.380
\hypertarget{w06:cue:601}{}這個法條啊，
% cue 602 1585.380 1585.880
\hypertarget{w06:cue:602}{}注意啊，
% cue 603 1585.880 1586.500
\hypertarget{w06:cue:603}{}這個框框，
% cue 604 1586.500 1588.330
\hypertarget{w06:cue:604}{}它的意思就是這裡，
% cue 605 1588.330 1589.990
\hypertarget{w06:cue:605}{}這個方法我們的一小半，
% cue 606 1589.990 1591.110
\hypertarget{w06:cue:606}{}這個才對，
% cue 607 1591.110 1592.730
\hypertarget{w06:cue:607}{}其他的我們不認為，
% cue 608 1592.730 1594.870
\hypertarget{w06:cue:608}{}我簡單的講是客氣，
% cue 609 1594.870 1597.250
\hypertarget{w06:cue:609}{}我沒有把你直接宣告違憲。

\tc{0:26:37}
% cue 610 1597.250 1601.060
\hypertarget{w06:cue:610}{}因為你這個解釋的可能性當中，
% cue 611 1601.060 1602.610
\hypertarget{w06:cue:611}{}有一種是合憲的，
% cue 612 1602.610 1604.120
\hypertarget{w06:cue:612}{}那我就乾脆告訴你，
% cue 613 1604.120 1605.680
\hypertarget{w06:cue:613}{}這個解釋只能做這個，
% cue 614 1605.680 1607.320
\hypertarget{w06:cue:614}{}這條文只能做這個解釋，
% cue 615 1607.320 1608.990
\hypertarget{w06:cue:615}{}我只限在這個方法，
% cue 616 1608.990 1612.220
\hypertarget{w06:cue:616}{}這樣子就是我已經修改了，
% cue 617 1612.220 1614.480
\hypertarget{w06:cue:617}{}這個法律適用的效果，
% cue 618 1614.480 1616.600
\hypertarget{w06:cue:618}{}它本來可以適用這麼大一個方法，
% cue 619 1616.600 1617.700
\hypertarget{w06:cue:619}{}我跟你講不對，
% cue 620 1617.700 1619.620
\hypertarget{w06:cue:620}{}合憲性解釋中間的這一塊，
% cue 621 1619.620 1620.140
\hypertarget{w06:cue:621}{}就這裡，
% cue 622 1620.140 1621.200
\hypertarget{w06:cue:622}{}這裡才對，
% cue 623 1621.200 1623.970
\hypertarget{w06:cue:623}{}我這樣講有點抽象，
% cue 624 1623.970 1624.830
\hypertarget{w06:cue:624}{}但是沒有問題，
% cue 625 1624.830 1626.710
\hypertarget{w06:cue:625}{}你們自己去看 272 號解釋，
% cue 626 1626.710 1628.830
\hypertarget{w06:cue:626}{}還有其他的解釋也還有，
% cue 627 1628.830 1630.920
\hypertarget{w06:cue:627}{}我現在想不起來的解釋，
% cue 628 1630.920 1632.620
\hypertarget{w06:cue:628}{}509 嘛\footnote{司法院釋字第 509 號解釋（2000 年 7 月 7 日）認刑法第 310 條第 1 項及第 2 項誹謗罪符合憲法第 23 條之意旨；就同條第 3 項前段「能證明其為真實者不罰」，解釋文謂「行為人雖不能證明言論內容為真實，但依其所提證據資料，認為行為人有相當理由確信其為真實者，即不能以誹謗罪之刑責相繩」，並據此認該項與憲法保障言論自由之旨趣並無牴觸。}，
% cue 629 1632.620 1636.190
\hypertarget{w06:cue:629}{}去看看就明白了。

\tc{0:27:16}
% cue 630 1636.190 1636.430
\hypertarget{w06:cue:630}{}好了，
% cue 631 1636.430 1637.960
\hypertarget{w06:cue:631}{}往下跳，
% cue 632 1637.960 1638.840
\hypertarget{w06:cue:632}{}不要講這個，
% cue 633 1638.840 1639.620
\hypertarget{w06:cue:633}{}I'm Moreover
% cue 634 1639.620 1640.480
\hypertarget{w06:cue:634}{}此外，

\tc{0:27:20}
\speaker{導讀人：亦凡}
% cue 635 1640.480 1642.780
\hypertarget{w06:cue:635}{}Moreover courts in this system
% cue 636 1642.780 1645.560
\hypertarget{w06:cue:636}{}have the authority to establish as a matter of law
% cue 637 1645.560 1649.440
\hypertarget{w06:cue:637}{}that a given statute or legislative provision
% cue 638 1649.440 1650.880
\hypertarget{w06:cue:638}{}will not be applied

\tc{0:27:30}
\speaker{老師}
% cue 639 1650.880 1652.640
\hypertarget{w06:cue:639}{}不會被適用，
% cue 640 1652.640 1654.080
\hypertarget{w06:cue:640}{}不被適用。

\tc{0:27:34}
% cue 641 1654.080 1655.680
\hypertarget{w06:cue:641}{}然後往下跳，
% cue 642 1655.680 1656.300
\hypertarget{w06:cue:642}{}A form

\tc{0:27:36}
\speaker{導讀人：亦凡}
% cue 643 1656.300 1659.680
\hypertarget{w06:cue:643}{}A form of even stronger judicial review
% cue 644 1659.680 1662.360
\hypertarget{w06:cue:644}{}would empower the courts
% cue 645 1662.360 1665.600
\hypertarget{w06:cue:645}{}to actually strike a piece of legislation
% cue 646 1665.600 1668.330
\hypertarget{w06:cue:646}{}out of the statute book altogether。

\tc{0:28:10}
\speaker{老師}
% cue 647 1690.070 1693.010
\hypertarget{w06:cue:647}{}或者他應調整為如下的方法適用，
% cue 648 1693.010 1695.510
\hypertarget{w06:cue:648}{}但是甚至於他也有可能，
% cue 649 1695.510 1698.130
\hypertarget{w06:cue:649}{}乾脆就說這個法律整個無效，
% cue 650 1698.130 1700.690
\hypertarget{w06:cue:650}{}這都是講的 strong，
% cue 651 1700.690 1703.100
\hypertarget{w06:cue:651}{}strong judicial review
% cue 652 1703.100 1705.550
\hypertarget{w06:cue:652}{}他要講的是 strong judicial review，
% cue 653 1705.550 1708.180
\hypertarget{w06:cue:653}{}他要反違憲審查，
% cue 654 1708.180 1709.740
\hypertarget{w06:cue:654}{}是指反這種審查，
% cue 655 1709.740 1711.260
\hypertarget{w06:cue:655}{}比之相對的，
% cue 656 1711.260 1712.740
\hypertarget{w06:cue:656}{}是 weak judicial review。

\tc{0:28:32}
% cue 657 1712.740 1714.160
\hypertarget{w06:cue:657}{}其實你們在這裡，
% cue 658 1714.160 1715.020
\hypertarget{w06:cue:658}{}你們看到呢，
% cue 659 1715.020 1717.150
\hypertarget{w06:cue:659}{}這個論文大陸法系的話來講，
% cue 660 1717.150 1718.110
\hypertarget{w06:cue:660}{}他就不會這樣講，
% cue 661 1718.110 1720.300
\hypertarget{w06:cue:661}{}他說的違憲審查的效力有不同，
% cue 662 1720.300 1721.720
\hypertarget{w06:cue:662}{}you're strong and weak
% cue 663 1723.230 1726.210
\hypertarget{w06:cue:663}{}你不會覺得這個是兩種不同的違憲審查，
% cue 664 1726.210 1728.220
\hypertarget{w06:cue:664}{}但他認為是不同的違憲審查，
% cue 665 1728.220 1730.740
\hypertarget{w06:cue:665}{}他反的是前者不反對後者。

\tc{0:28:50}
% cue 666 1730.740 1731.960
\hypertarget{w06:cue:666}{}來 看一下 weak，
% cue 667 1731.960 1734.010
\hypertarget{w06:cue:667}{}1355 頁，

\tc{0:28:54}
\speaker{導讀人：亦凡}
% cue 668 1734.830 1753.090
\hypertarget{w06:cue:668}{}"In a system of weak judicial review, by contrast, courts may scrutinize legislation for its conformity to individual rights, but they may not decline to apply it."

\tc{0:29:13}
\speaker{老師}
% cue 669 1753.090 1755.230
\hypertarget{w06:cue:669}{}他可以審查是不是相符，
% cue 670 1755.230 1757.760
\hypertarget{w06:cue:670}{}有沒有跟人權的規定相符，
% cue 671 1757.760 1761.050
\hypertarget{w06:cue:671}{}但是他沒有權利拒絕適用，
% cue 672 1761.050 1765.080
\hypertarget{w06:cue:672}{}或者是修改他的適用的範圍，
% cue 673 1765.080 1769.250
\hypertarget{w06:cue:673}{}Simply because rights would otherwise be violated

\tc{0:29:29}
\speaker{導讀人：亦凡}
% cue 674 1769.250 1776.490
\hypertarget{w06:cue:674}{}Nevertheless the scrutiny may have some effect

\tc{0:29:36}
\speaker{老師}
% cue 675 1776.490 1778.670
\hypertarget{w06:cue:675}{}你雖然這樣子講，
% cue 676 1778.670 1780.850
\hypertarget{w06:cue:676}{}你說我沒有權利拒絕適用，
% cue 677 1780.850 1783.970
\hypertarget{w06:cue:677}{}並不表示我所做的違宣告，
% cue 678 1783.970 1785.210
\hypertarget{w06:cue:678}{}一點用都沒有，
% cue 679 1785.210 1786.740
\hypertarget{w06:cue:679}{}沒有任何銷利品，
% cue 680 1786.740 1787.840
\hypertarget{w06:cue:680}{}不是這樣的意思。

\tc{0:29:47}
% cue 681 1787.840 1789.880
\hypertarget{w06:cue:681}{}像這裡就是我們，
% cue 682 1789.880 1791.650
\hypertarget{w06:cue:682}{}我今天舉這個例子，
% cue 683 1791.650 1792.470
\hypertarget{w06:cue:683}{}臺灣，
% cue 684 1792.470 1793.670
\hypertarget{w06:cue:684}{}除了我，
% cue 685 1793.670 1794.650
\hypertarget{w06:cue:685}{}我不知道還有誰，
% cue 686 1794.650 1796.770
\hypertarget{w06:cue:686}{}我們現在講比較憲法，
% cue 687 1796.770 1798.730
\hypertarget{w06:cue:687}{}我們太受到德國的影響，
% cue 688 1798.730 1800.800
\hypertarget{w06:cue:688}{}其他的都看不見，
% cue 689 1800.800 1802.930
\hypertarget{w06:cue:689}{}但是事實上不真的這樣，
% cue 690 1802.930 1804.350
\hypertarget{w06:cue:690}{}好 你看下面這個句子，
% cue 691 1804.350 1805.460
\hypertarget{w06:cue:691}{}In the United Kingdom
% cue 692 1805.460 1806.300
\hypertarget{w06:cue:692}{}在英國，
% cue 693 1806.300 1809.440
\hypertarget{w06:cue:693}{}繼續啊 繼續啊，
% cue 694 1809.440 1813.810
\hypertarget{w06:cue:694}{}The courts may review a statute

\tc{0:30:13}
\speaker{導讀人：亦凡}
% cue 695 1813.810 1821.370
\hypertarget{w06:cue:695}{}A statute with a view to issuing a declaration of incompatibility\footnote{英國 Human Rights Act 1998 第 4 條第 2 項規定，法院認定主要立法的條文與公約權利不相容者，得宣告其不相容；第 6 項規定此項宣告「does not affect the validity, continuing operation or enforcement of the provision in respect of which it is given」，對該訴訟的當事人亦無拘束力。依第 5 項，得為宣告的法院包括最高法院、樞密院司法委員會，以及英格蘭與威爾斯的高等法院與上訴法院等。\url{https://www.legislation.gov.uk/ukpga/1998/42/section/4}。}。

\tc{0:30:21}
\speaker{老師}
% cue 696 1821.370 1823.580
\hypertarget{w06:cue:696}{}它可以發布什麼，
% cue 697 1823.580 1825.310
\hypertarget{w06:cue:697}{}我宣告不合，
% cue 698 1825.310 1828.560
\hypertarget{w06:cue:698}{}用我們臺灣的話講就是宣告違憲，
% cue 699 1828.560 1831.100
\hypertarget{w06:cue:699}{}違憲就是違反憲法，
% cue 700 1831.100 1834.130
\hypertarget{w06:cue:700}{}宣告與憲法意旨不相符合，
% cue 701 1834.130 1834.900
\hypertarget{w06:cue:701}{}宣告，
% cue 702 1834.900 1837.760
\hypertarget{w06:cue:702}{}A declaration of incompatibility
% cue 703 1837.760 1839.270
\hypertarget{w06:cue:703}{}兩者不相容。

\tc{0:30:39}
% cue 704 1839.270 1844.880
\hypertarget{w06:cue:704}{}當法院發現它是跟 convention right，
% cue 705 1844.880 1847.680
\hypertarget{w06:cue:705}{}你可以解釋一下什麼叫 convention right 嗎？
% cue 706 1847.680 1850.170
\hypertarget{w06:cue:706}{}convention 是什麼 convention，
% cue 707 1850.170 1851.010
\hypertarget{w06:cue:707}{}你們聽得懂嗎？
% cue 708 1851.010 1852.490
\hypertarget{w06:cue:708}{}下面那個字，
% cue 709 1852.490 1855.650
\hypertarget{w06:cue:709}{}European Convention of Human Rights\footnote{正式名稱為 Convention for the Protection of Human Rights and Fundamental Freedoms（保障人權及基本自由公約），1950 年 11 月 4 日由歐洲理事會會員國於羅馬簽署。Human Rights Act 1998 第 1 條所稱 Convention rights，指公約第 2 條至第 12 條、第 14 條，第一議定書第 1 條至第 3 條，及第十三議定書第 1 條所列的權利與基本自由。\url{https://www.echr.coe.int/documents/d/echr/convention_ENG}。}。

\tc{0:30:55}
% cue 710 1855.890 1857.220
\hypertarget{w06:cue:710}{}英國本來沒有，
% cue 711 1857.220 1859.300
\hypertarget{w06:cue:711}{}老師在視頻裡面有講過了，
% cue 712 1859.300 1863.140
\hypertarget{w06:cue:712}{}英國本來沒有成文憲法，
% cue 713 1863.140 1864.920
\hypertarget{w06:cue:713}{}英國是不成文憲法的國家，
% cue 714 1864.920 1866.670
\hypertarget{w06:cue:714}{}到現在為止還是，
% cue 715 1866.670 1869.340
\hypertarget{w06:cue:715}{}它可能是世界上唯一的一個，
% cue 716 1869.340 1871.300
\hypertarget{w06:cue:716}{}不成文憲法的國家，
% cue 717 1871.300 1873.440
\hypertarget{w06:cue:717}{}它就沒有一部法典叫憲法，
% cue 718 1873.440 1875.970
\hypertarget{w06:cue:718}{}它不是沒有憲法，
% cue 719 1875.970 1878.110
\hypertarget{w06:cue:719}{}它的憲法散落在不同的地方。

\tc{0:31:18}
% cue 720 1878.110 1879.990
\hypertarget{w06:cue:720}{}好那現在一個根本的問題，
% cue 721 1879.990 1882.310
\hypertarget{w06:cue:721}{}請問那麼英國的人權規定，
% cue 722 1882.310 1883.050
\hypertarget{w06:cue:722}{}藏在哪裡？
% cue 723 1883.050 1884.300
\hypertarget{w06:cue:723}{}那就是藏在這裡，
% cue 724 1884.300 1886.340
\hypertarget{w06:cue:724}{}它通過了一個法律，
% cue 725 1886.340 1888.100
\hypertarget{w06:cue:725}{}叫做 Human Rights Act\footnote{Human Rights Act 1998（1998 c. 42），1998 年 11 月 9 日制定，立法目的為使歐洲人權公約保障的權利與自由在英國法上進一步生效。依 Human Rights Act 1998 (Commencement No. 2) Order 2000，該法尚未施行的條文自 2000 年 10 月 2 日施行。\url{https://www.legislation.gov.uk/ukpga/1998/42/section/1}。}，
% cue 726 1888.100 1889.400
\hypertarget{w06:cue:726}{}下一方再下一方，
% cue 727 1889.400 1890.160
\hypertarget{w06:cue:727}{}看到嗎？
% cue 728 1890.160 1893.680
\hypertarget{w06:cue:728}{}這是 1998 年通過這個 Human Rights Act，
% cue 729 1893.680 1896.080
\hypertarget{w06:cue:729}{}然後規定了 2000 年開始施行，
% cue 730 1896.080 1898.660
\hypertarget{w06:cue:730}{}所以他現在零施行二十幾年。

\tc{0:31:38}
% cue 731 1898.660 1902.940
\hypertarget{w06:cue:731}{}那他就把歐洲人權公約裡面的所有的那些人權，
% cue 732 1902.940 1905.490
\hypertarget{w06:cue:732}{}就直接搬進英國內國法來，
% cue 733 1905.490 1910.580
\hypertarget{w06:cue:733}{}然後說這個政府的行為不可以違反，
% cue 734 1910.580 1913.530
\hypertarget{w06:cue:734}{}歐洲人權公約裡面所規定，
% cue 735 1913.530 1917.990
\hypertarget{w06:cue:735}{}然後搬進了我們國家人權法裡面的人權，
% cue 736 1917.990 1919.890
\hypertarget{w06:cue:736}{}所以他叫做 convention right，
% cue 737 1919.890 1924.100
\hypertarget{w06:cue:737}{}Convention 是指歐洲人權公約的公約權利，
% cue 738 1924.100 1925.600
\hypertarget{w06:cue:738}{}公約權利就是人權，
% cue 739 1925.600 1928.400
\hypertarget{w06:cue:739}{}我這樣講有沒有聽懂？
% cue 740 1928.400 1929.820
\hypertarget{w06:cue:740}{}它有點像金蟬脫殼，
% cue 741 1929.820 1932.210
\hypertarget{w06:cue:741}{}這個事情是有趣的事情，
% cue 742 1932.210 1933.250
\hypertarget{w06:cue:742}{}這是立法技術。

\tc{0:32:13}
% cue 743 1933.250 1936.730
\hypertarget{w06:cue:743}{}你們會不會覺得英國人做事情，
% cue 744 1936.730 1938.490
\hypertarget{w06:cue:744}{}常常超過你的想象，
% cue 745 1938.490 1940.890
\hypertarget{w06:cue:745}{}你會這樣子嗎？
% cue 746 1940.890 1942.030
\hypertarget{w06:cue:746}{}你會覺得不行，
% cue 747 1942.030 1944.270
\hypertarget{w06:cue:747}{}我們一定要在憲法裡面，
% cue 748 1944.270 1945.490
\hypertarget{w06:cue:748}{}然後 闢一個專章，
% cue 749 1945.490 1946.970
\hypertarget{w06:cue:749}{}像我們第二章一樣，
% cue 750 1946.970 1948.630
\hypertarget{w06:cue:750}{}像德國人第一章一樣，
% cue 751 1948.630 1949.510
\hypertarget{w06:cue:751}{}搶基本權。

\tc{0:32:29}
% cue 752 1949.510 1952.250
\hypertarget{w06:cue:752}{}這叫做我們國家引進了人權，
% cue 753 1952.250 1954.460
\hypertarget{w06:cue:753}{}他說不要那麼麻煩，
% cue 754 1954.460 1956.480
\hypertarget{w06:cue:754}{}我們國家反正就沒有憲法，
% cue 755 1956.480 1958.490
\hypertarget{w06:cue:755}{}我們就制定一個法律，
% cue 756 1958.490 1961.200
\hypertarget{w06:cue:756}{}他說這個人權法不能違反，
% cue 757 1961.200 1962.700
\hypertarget{w06:cue:757}{}那這個裡面規定的東西，
% cue 758 1962.700 1964.780
\hypertarget{w06:cue:758}{}就說他跟歐洲人權公約長的一樣，
% cue 759 1964.780 1965.440
\hypertarget{w06:cue:759}{}就搬進來了，
% cue 760 1965.440 1967.810
\hypertarget{w06:cue:760}{}臺灣想要做，
% cue 761 1967.810 1970.430
\hypertarget{w06:cue:761}{}我們想要把兩公約內國法化\footnote{公民與政治權利國際公約及經濟社會文化權利國際公約施行法，2009 年 4 月 22 日制定公布，行政院以同年 10 月 29 日院臺外字第 0980067638 號令定自同年 12 月 10 日施行。第 2 條規定：「兩公約所揭示保障人權之規定，具有國內法律之效力。」\url{https://law.moj.gov.tw/LawClass/LawAll.aspx?pcode=I0020028}。}，
% cue 762 1970.430 1972.850
\hypertarget{w06:cue:762}{}但是我們說兩公約的那東西，
% cue 763 1972.850 1974.490
\hypertarget{w06:cue:763}{}在臺灣只具有法的效力，
% cue 764 1974.490 1975.890
\hypertarget{w06:cue:764}{}不具有憲法的效力，
% cue 765 1975.890 1977.510
\hypertarget{w06:cue:765}{}這是一個有趣的問題，
% cue 766 1977.510 1978.740
\hypertarget{w06:cue:766}{}你覺得立法者，
% cue 767 1978.740 1979.880
\hypertarget{w06:cue:767}{}可不可以搬一個法界法，
% cue 768 1979.880 1981.280
\hypertarget{w06:cue:768}{}還是說他具有憲法效力，
% cue 769 1981.280 1983.060
\hypertarget{w06:cue:769}{}那為什麼英國可以，
% cue 770 1983.060 1984.540
\hypertarget{w06:cue:770}{}有沒有人可以回答我這個問題？

\tc{0:33:04}
% cue 771 1984.540 1986.320
\hypertarget{w06:cue:771}{}我今天這個一定講不完，
% cue 772 1986.320 1988.450
\hypertarget{w06:cue:772}{}我也不相信，
% cue 773 1988.450 1990.260
\hypertarget{w06:cue:773}{}這個亦凡已經全部開完了，
% cue 774 1990.260 1990.960
\hypertarget{w06:cue:774}{}沒關係，
% cue 775 1990.960 1992.400
\hypertarget{w06:cue:775}{}反正您講究竟然算多少，
% cue 776 1992.400 1994.110
\hypertarget{w06:cue:776}{}現在的重點是，
% cue 777 1994.110 1995.690
\hypertarget{w06:cue:777}{}我們講到的東西，
% cue 778 1995.690 1997.540
\hypertarget{w06:cue:778}{}跟我國有關的部分，
% cue 779 1997.540 1999.380
\hypertarget{w06:cue:779}{}我需要利用這個機會，
% cue 780 1999.380 2000.660
\hypertarget{w06:cue:780}{}讓你們去做一個對比，
% cue 781 2000.660 2003.370
\hypertarget{w06:cue:781}{}因為這是你們一定可以做到，
% cue 782 2003.370 2005.490
\hypertarget{w06:cue:782}{}我這樣講話聽說話，
% cue 783 2005.490 2007.320
\hypertarget{w06:cue:783}{}我剛剛講的問題怎麼回答我，
% cue 784 2007.320 2008.840
\hypertarget{w06:cue:784}{}在座這麼多人，
% cue 785 2008.840 2011.300
\hypertarget{w06:cue:785}{}我們臺灣能不能透過法律，
% cue 786 2011.300 2012.860
\hypertarget{w06:cue:786}{}比方說我想把兩公約，
% cue 787 2012.860 2014.490
\hypertarget{w06:cue:787}{}我們也把它內國法化，
% cue 788 2014.490 2015.650
\hypertarget{w06:cue:788}{}因為定義個兩公約，
% cue 789 2015.650 2016.410
\hypertarget{w06:cue:789}{}施行法，
% cue 790 2016.410 2019.160
\hypertarget{w06:cue:790}{}把兩公約搬進來，
% cue 791 2019.160 2021.210
\hypertarget{w06:cue:791}{}但是為什麼我說，
% cue 792 2021.210 2022.650
\hypertarget{w06:cue:792}{}沒有像英國一樣說，
% cue 793 2022.650 2025.000
\hypertarget{w06:cue:793}{}他就給予他憲法的效力，
% cue 794 2025.000 2027.340
\hypertarget{w06:cue:794}{}我們還是說他只具有法律的效力，
% cue 795 2027.340 2029.170
\hypertarget{w06:cue:795}{}為什麼我們不敢這樣做，
% cue 796 2029.170 2032.100
\hypertarget{w06:cue:796}{}跟我剛才講的，
% cue 797 2032.100 2034.460
\hypertarget{w06:cue:797}{}英國沒有成文憲法是有關的，
% cue 798 2034.460 2037.660
\hypertarget{w06:cue:798}{}英國所有的東西到目前為止都叫法律，
% cue 799 2037.660 2038.760
\hypertarget{w06:cue:799}{}有的不叫法律，
% cue 800 2038.760 2039.680
\hypertarget{w06:cue:800}{}有的叫慣例，
% cue 801 2039.680 2041.220
\hypertarget{w06:cue:801}{}有的藏在什麼地方，
% cue 802 2041.220 2042.770
\hypertarget{w06:cue:802}{}那不重要，
% cue 803 2042.770 2045.320
\hypertarget{w06:cue:803}{}大家通通都是用法律的形式。

\tc{0:34:05}
% cue 804 2045.320 2047.880
\hypertarget{w06:cue:804}{}記不記得你們以前學憲法的時候說，
% cue 805 2047.880 2050.050
\hypertarget{w06:cue:805}{}有成文憲法有不成文憲法，
% cue 806 2050.050 2053.320
\hypertarget{w06:cue:806}{}那他的最大的差別不是他的形式，
% cue 807 2053.320 2057.080
\hypertarget{w06:cue:807}{}而是他修改的機關，
% cue 808 2057.080 2058.550
\hypertarget{w06:cue:808}{}還有修改的程序，
% cue 809 2058.550 2060.190
\hypertarget{w06:cue:809}{}不成文憲法，
% cue 810 2060.190 2063.940
\hypertarget{w06:cue:810}{}他的修改機關跟普通法律，
% cue 811 2063.940 2064.640
\hypertarget{w06:cue:811}{}是一樣的，
% cue 812 2064.640 2065.940
\hypertarget{w06:cue:812}{}都是立法制，
% cue 813 2065.940 2068.330
\hypertarget{w06:cue:813}{}因為英國的國會是至上的，
% cue 814 2068.330 2070.850
\hypertarget{w06:cue:814}{}他沒有不能做的事情，
% cue 815 2070.850 2072.510
\hypertarget{w06:cue:815}{}OK 第一件事，
% cue 816 2072.510 2073.230
\hypertarget{w06:cue:816}{}第二件事情，
% cue 817 2073.230 2074.420
\hypertarget{w06:cue:817}{}修改的程序，
% cue 818 2074.420 2075.640
\hypertarget{w06:cue:818}{}我們是很複雜，
% cue 819 2075.640 2076.400
\hypertarget{w06:cue:819}{}要四分之三\footnote{中華民國憲法增修條文第 12 條規定：「憲法之修改，須經立法院立法委員四分之一之提議，四分之三之出席，及出席委員四分之三之決議，提出憲法修正案，並於公告半年後，經中華民國自由地區選舉人投票複決，有效同意票過選舉人總額之半數，即通過之」。\url{https://law.moj.gov.tw/LawClass/LawSingle.aspx?pcode=A0000002&flno=12}。}，
% cue 820 2076.400 2077.360
\hypertarget{w06:cue:820}{}美國要四分之三\footnote{美國聯邦憲法第 5 條規定，修正案由國會兩院各三分之二認為必要時提出，或依三分之二州議會的請求召開制憲會議提出；經四分之三州議會或四分之三州的制憲會議批准，即成為憲法的一部分。\url{https://www.archives.gov/founding-docs/constitution-transcript}。}，
% cue 821 2077.360 2078.560
\hypertarget{w06:cue:821}{}還有很多道手續，
% cue 822 2078.560 2079.360
\hypertarget{w06:cue:822}{}他沒有，
% cue 823 2079.360 2082.060
\hypertarget{w06:cue:823}{}他通過跟普通法律，
% cue 824 2082.060 2083.440
\hypertarget{w06:cue:824}{}一樣的方法，
% cue 825 2083.440 2084.800
\hypertarget{w06:cue:825}{}但他就在裡面寫說，
% cue 826 2084.800 2085.820
\hypertarget{w06:cue:826}{}你不能違反他，
% cue 827 2085.820 2087.450
\hypertarget{w06:cue:827}{}這下去對我調好，
% cue 828 2087.550 2089.020
\hypertarget{w06:cue:828}{}那有違反怎麼辦呢，
% cue 829 2089.020 2091.550
\hypertarget{w06:cue:829}{}那法院現在是英國最高法院\footnote{Supreme Court of the United Kingdom，依 Constitutional Reform Act 2005 第 23 條設立；依 Constitutional Reform Act 2005 (Commencement No. 11) Order 2009，第 23 條至第 44 條自 2009 年 10 月 1 日施行。Human Rights Act 1998 第 4 條第 5 項第 a 款所列得為不相容宣告的法院，同日改為最高法院。\url{https://www.legislation.gov.uk/ukpga/2005/4/section/23}。}，
% cue 830 2091.550 2093.440
\hypertarget{w06:cue:830}{}你就直接宣告不符合，
% cue 831 2093.440 2095.230
\hypertarget{w06:cue:831}{}不符合就是違憲了。

\tc{0:34:55}
% cue 832 2095.230 2096.030
\hypertarget{w06:cue:832}{}繼續，
% cue 833 2096.030 2097.860
\hypertarget{w06:cue:833}{}它不是說都沒有效果，
% cue 834 2097.860 2098.560
\hypertarget{w06:cue:834}{}有什麼效果？

\tc{0:34:58}
\speaker{導讀人：亦凡}
% cue 835 2098.560 2109.120
\hypertarget{w06:cue:835}{}The act provides that such declaration。

\tc{0:35:09}
\speaker{老師}
% cue 836 2109.120 2110.180
\hypertarget{w06:cue:836}{}這樣的一個宣告，
% cue 837 2110.180 2111.220
\hypertarget{w06:cue:837}{}違憲的宣告，
% cue 838 2111.220 2113.020
\hypertarget{w06:cue:838}{}does not affect the validity
% cue 839 2113.020 2115.210
\hypertarget{w06:cue:839}{}continuing operation
% cue 840 2115.210 2117.500
\hypertarget{w06:cue:840}{}or enforcement of the provision
% cue 841 2117.500 2119.680
\hypertarget{w06:cue:841}{}in respect of which it is given
% cue 842 2119.680 2121.520
\hypertarget{w06:cue:842}{}這句話講得太複雜了，
% cue 843 2121.520 2122.420
\hypertarget{w06:cue:843}{}意思就是說，
% cue 844 2122.420 2124.900
\hypertarget{w06:cue:844}{}不影響到系爭規定的效力，
% cue 845 2124.900 2125.800
\hypertarget{w06:cue:845}{}繼續執行，
% cue 846 2125.800 2127.580
\hypertarget{w06:cue:846}{}它不影響，
% cue 847 2127.580 2130.640
\hypertarget{w06:cue:847}{}and it's not a binding
% cue 848 2130.640 2132.220
\hypertarget{w06:cue:848}{}對於當事人，
% cue 849 2132.220 2135.930
\hypertarget{w06:cue:849}{}提起訴訟的雙方或多方當事人來講，
% cue 850 2135.930 2136.930
\hypertarget{w06:cue:850}{}也沒有拘束力，
% cue 851 2136.930 2138.930
\hypertarget{w06:cue:851}{}所以照這樣看起來，
% cue 852 2138.930 2140.060
\hypertarget{w06:cue:852}{}它不是廢話一句嗎，
% cue 853 2140.060 2141.720
\hypertarget{w06:cue:853}{}你宣告不符合，
% cue 854 2141.720 2142.540
\hypertarget{w06:cue:854}{}然後呢，
% cue 855 2142.540 2144.250
\hypertarget{w06:cue:855}{}原法繼續執行，
% cue 856 2144.250 2146.650
\hypertarget{w06:cue:856}{}對我們雙方當事人也沒有拘束力。

\tc{0:35:46}
% cue 857 2146.650 2148.450
\hypertarget{w06:cue:857}{}but still。

\tc{0:35:48}
\speaker{導讀人：亦凡}
% cue 858 2148.450 2153.500
\hypertarget{w06:cue:858}{}but still it has an effect。

\tc{0:35:53}
\speaker{老師}
% cue 859 2153.500 2154.380
\hypertarget{w06:cue:859}{}有效果，
% cue 860 2154.380 2155.120
\hypertarget{w06:cue:860}{}什麼效果？

\tc{0:35:55}
\speaker{導讀人：亦凡}
% cue 861 2155.120 2170.660
\hypertarget{w06:cue:861}{}A minister may use such a declaration as authorization to initiate a fast-track legislative procedure\footnote{Waldron 所指為 Human Rights Act 1998 第 10 條（本文註 28）：法律條文經依第 4 條宣告不相容，且當事人已表示不上訴、上訴期間已過或上訴已終結者，部長認有不得已的理由（compelling reasons）時，得以命令修正該法律至消除不相容所必要的程度。此種命令稱為 remedial order。\url{https://www.legislation.gov.uk/ukpga/1998/42/section/10}。} to remedy the incompatibility。

\tc{0:36:10}
\speaker{老師}
% cue 862 2170.660 2172.930
\hypertarget{w06:cue:862}{}好，這句話最重要，這句話什麼意思？

\tc{0:36:12}
\speaker{導讀人：亦凡}
% cue 863 2172.930 2179.320
\hypertarget{w06:cue:863}{}就是在依據這個指引嘛，依據這個審查的。

\tc{0:36:19}
\speaker{老師}
% cue 864 2179.320 2182.740
\hypertarget{w06:cue:864}{}我說宣告他違憲嘛，我剛剛說宣告不符合，
% cue 865 2182.740 2186.000
\hypertarget{w06:cue:865}{}依據這個法院的宣告，不符合的宣告，
% cue 866 2186.000 2187.800
\hypertarget{w06:cue:866}{}然後呢？

\tc{0:36:27}
\speaker{導讀人：亦凡}
% cue 867 2187.800 2189.430
\hypertarget{w06:cue:867}{}可以發動一個。

\tc{0:36:29}
\speaker{老師}
% cue 868 2189.430 2189.950
\hypertarget{w06:cue:868}{}誰誰，
% cue 869 2189.950 2190.490
\hypertarget{w06:cue:869}{}主持。

\tc{0:36:30}
\speaker{導讀人：亦凡}
% cue 870 2190.490 2191.790
\hypertarget{w06:cue:870}{}minister。

\tc{0:36:31}
\speaker{老師}
% cue 871 2191.790 2192.750
\hypertarget{w06:cue:871}{}部長，
% cue 872 2192.750 2195.760
\hypertarget{w06:cue:872}{}部長就獲得了一個授權，
% cue 873 2195.760 2198.060
\hypertarget{w06:cue:873}{}他可以啟動一個什麼？

\tc{0:36:38}
\speaker{導讀人：亦凡}
% cue 874 2198.060 2200.510
\hypertarget{w06:cue:874}{}啟動一個快速修法程序。

\tc{0:36:40}
\speaker{老師}
% cue 875 2200.510 2201.810
\hypertarget{w06:cue:875}{}對這樣子我聽懂嗎，
% cue 876 2201.810 2203.690
\hypertarget{w06:cue:876}{}相當於臺灣的逕付二讀\footnote{立法院職權行使法第 8 條第 2 項規定，政府機關提出之議案或立法委員提出之法律案，應先送程序委員會，提報院會朗讀標題後，即應交付有關委員會審查；「但有出席委員提議，二十人以上連署或附議，經表決通過，得逕付二讀。」\url{https://law.moj.gov.tw/LawClass/LawSingle.aspx?pcode=A0020058&flno=8}。}，
% cue 877 2203.690 2205.030
\hypertarget{w06:cue:877}{}不要一讀了，
% cue 878 2205.030 2206.510
\hypertarget{w06:cue:878}{}直接跳進去，
% cue 879 2206.510 2207.210
\hypertarget{w06:cue:879}{}然後呢，
% cue 880 2207.210 2210.880
\hypertarget{w06:cue:880}{}如果你想修，
% cue 881 2210.880 2213.950
\hypertarget{w06:cue:881}{}你想要去掉這個不違線的部分，
% cue 882 2213.950 2215.300
\hypertarget{w06:cue:882}{}那也是修法，
% cue 883 2215.300 2216.980
\hypertarget{w06:cue:883}{}但是你可以快速修法，
% cue 884 2216.980 2219.420
\hypertarget{w06:cue:884}{}所以它不是全然沒有效果，
% cue 885 2219.420 2221.730
\hypertarget{w06:cue:885}{}它的效果是程序效果，
% cue 886 2221.730 2223.260
\hypertarget{w06:cue:886}{}而且是得不是應，
% cue 887 2223.260 2224.000
\hypertarget{w06:cue:887}{}May
% cue 888 2224.000 2225.060
\hypertarget{w06:cue:888}{}注意這個字，
% cue 889 2225.060 2227.280
\hypertarget{w06:cue:889}{}部長不受拘束，
% cue 890 2227.280 2229.330
\hypertarget{w06:cue:890}{}你不認為應該修改，
% cue 891 2229.330 2231.130
\hypertarget{w06:cue:891}{}你就把它當作廢紙放在旁邊，
% cue 892 2231.130 2232.680
\hypertarget{w06:cue:892}{}但如果你認為，
% cue 893 2232.680 2233.660
\hypertarget{w06:cue:893}{}講得很有道理，
% cue 894 2233.660 2234.700
\hypertarget{w06:cue:894}{}我想要接近修改，
% cue 895 2234.700 2234.920
\hypertarget{w06:cue:895}{}好，
% cue 896 2234.920 2236.390
\hypertarget{w06:cue:896}{}它就給你的尚方寶劍，
% cue 897 2236.390 2237.510
\hypertarget{w06:cue:897}{}我們就直接進入 20，
% cue 898 2237.510 2239.590
\hypertarget{w06:cue:898}{}Fast Track Legislative Procedure
% cue 899 2239.590 2241.660
\hypertarget{w06:cue:899}{}我這樣講有聽懂嗎？

\tc{0:37:21}
% cue 900 2241.660 2243.320
\hypertarget{w06:cue:900}{}你們會不會覺得，
% cue 901 2243.320 2244.500
\hypertarget{w06:cue:901}{}英國人做的事情，
% cue 902 2244.500 2246.100
\hypertarget{w06:cue:902}{}超乎你的想象，
% cue 903 2246.100 2248.250
\hypertarget{w06:cue:903}{}這就是我們比較憲法的，
% cue 904 2248.250 2250.470
\hypertarget{w06:cue:904}{}我們的眼界可以大開，
% cue 905 2250.470 2252.390
\hypertarget{w06:cue:905}{}因為世界上有很多事情，
% cue 906 2252.390 2254.240
\hypertarget{w06:cue:906}{}我們現在是被我們的腦袋，
% cue 907 2254.240 2256.920
\hypertarget{w06:cue:907}{}被這個框框條條給束縛住了，
% cue 908 2256.920 2259.230
\hypertarget{w06:cue:908}{}我們怎麼能想象這樣的事情呢，
% cue 909 2259.230 2261.170
\hypertarget{w06:cue:909}{}有那麼嚴重嗎？

\tc{0:37:41}
% cue 910 2261.170 2261.930
\hypertarget{w06:cue:910}{}好了講完了，
% cue 911 2261.930 2265.900
\hypertarget{w06:cue:911}{}然後 2356 頁，
% cue 912 2265.900 2267.720
\hypertarget{w06:cue:912}{}講到紐西蘭\footnote{New Zealand Bill of Rights Act 1990 第 4 條規定，法院不得僅因法律條文與該法所列權利不一致，即認其已被默示廢止、無效，或拒絕適用；第 6 條規定法律可解釋為與該法權利相符者，應優先採取該解釋。刊本 1356 頁兩條均有引錄，Waldron 並指出紐西蘭法院即使自行作出不一致的宣告，該宣告對立法程序亦不生英國法那樣的效果。}它是怎麼樣，
% cue 913 2267.720 2269.280
\hypertarget{w06:cue:913}{}大概講一下就好了，
% cue 914 2269.280 2270.400
\hypertarget{w06:cue:914}{}它也會宣告不符，
% cue 915 2270.400 2272.780
\hypertarget{w06:cue:915}{}然後呢？

\tc{0:37:52}
\speaker{導讀人：亦凡}
% cue 916 2272.780 2277.460
\hypertarget{w06:cue:916}{}對那個立法程序也不產生任何法律上的效果。

\tc{0:37:57}
\speaker{老師}
% cue 917 2277.460 2279.440
\hypertarget{w06:cue:917}{}所以紐西蘭的，
% cue 918 2279.440 2280.440
\hypertarget{w06:cue:918}{}這是更弱，
% cue 919 2280.440 2283.000
\hypertarget{w06:cue:919}{}紐西蘭是 weak and weak，
% cue 920 2283.000 2284.220
\hypertarget{w06:cue:920}{}weaker and weaker。

\tc{0:38:04}
% cue 921 2284.220 2286.190
\hypertarget{w06:cue:921}{}然後呢，
% cue 922 2286.190 2286.950
\hypertarget{w06:cue:922}{}還有沒有，
% cue 923 2286.950 2287.790
\hypertarget{w06:cue:923}{}Canada 呢，
% cue 924 2287.790 2289.440
\hypertarget{w06:cue:924}{}Canada 是強還是弱？

\tc{0:38:09}
\speaker{導讀人：亦凡}
% cue 925 2289.440 2291.550
\hypertarget{w06:cue:925}{}中間。

\tc{0:38:11}
\speaker{老師}
% cue 926 2291.550 2292.430
\hypertarget{w06:cue:926}{}中間，
% cue 927 2292.430 2293.490
\hypertarget{w06:cue:927}{}為什麼是中間，
% cue 928 2293.490 2294.250
\hypertarget{w06:cue:928}{}他認為是。

\tc{0:38:14}
\speaker{導讀人：亦凡}
% cue 929 2294.250 2295.110
\hypertarget{w06:cue:929}{}對。

\tc{0:38:15}
\speaker{老師}
% cue 930 2295.110 2296.290
\hypertarget{w06:cue:930}{}他有說是中間，
% cue 931 2296.290 2298.290
\hypertarget{w06:cue:931}{}為什麼他認為是中間，
% cue 932 2298.290 2300.060
\hypertarget{w06:cue:932}{}我也不懂，
% cue 933 2300.060 2301.920
\hypertarget{w06:cue:933}{}我覺得他很強啊。

\tc{0:38:21}
\speaker{導讀人：亦凡}
% cue 934 2301.920 2303.980
\hypertarget{w06:cue:934}{}就是很少的使用。

\tc{0:38:23}
\speaker{老師}
% cue 935 2303.980 2304.380
\hypertarget{w06:cue:935}{}它是叫，
% cue 936 2304.380 2305.960
\hypertarget{w06:cue:936}{}intermediate cases
% cue 937 2305.960 2307.960
\hypertarget{w06:cue:937}{}介於中間的案子，
% cue 938 2307.960 2309.290
\hypertarget{w06:cue:938}{}那你先念這句。

\tc{0:38:29}
% cue 939 2309.290 2309.850
\hypertarget{w06:cue:939}{}In Canada。

\tc{0:38:29}
\speaker{導讀人：亦凡}
% cue 940 2309.850 2311.670
\hypertarget{w06:cue:940}{}In Canada
% cue 941 2311.670 2316.110
\hypertarget{w06:cue:941}{}There is a provision for the review of legislation by courts。

\tc{0:38:36}
\speaker{老師}
% cue 942 2316.110 2317.730
\hypertarget{w06:cue:942}{}有沒有憲法明文依據，
% cue 943 2317.730 2318.130
\hypertarget{w06:cue:943}{}有，
% cue 944 2318.130 2320.940
\hypertarget{w06:cue:944}{}其次，美國沒有，它有，
% cue 945 2320.940 2324.580
\hypertarget{w06:cue:945}{}然後，the courts there，中間跳過，
% cue 946 2324.580 2326.600
\hypertarget{w06:cue:946}{}may decline to apply
% cue 947 2326.600 2329.180
\hypertarget{w06:cue:947}{}它可以拒絕適用，
% cue 948 2329.180 2335.760
\hypertarget{w06:cue:948}{}如果被宣告違反加拿大憲法裡面的人權規定\footnote{指 Canadian Charter of Rights and Freedoms（加拿大權利與自由憲章），為 Constitution Act, 1982 第一部分。同法第 52 條第 1 項規定，加拿大憲法為加拿大的最高法，任何法律與憲法規定牴觸者，於牴觸的範圍內不生效力。\url{https://laws-lois.justice.gc.ca/eng/Const/FullText.html}。}的話，
% cue 949 2335.760 2336.980
\hypertarget{w06:cue:949}{}它有人權清單，
% cue 950 2336.980 2340.150
\hypertarget{w06:cue:950}{}一樣的，名字不一樣，意思是一樣的。

\tc{0:39:00}
% cue 951 2340.150 2341.630
\hypertarget{w06:cue:951}{}所以到這裡，一樣嘛，
% cue 952 2341.630 2344.980
\hypertarget{w06:cue:952}{}這是 strong 呢？我不懂它為什麼是 intermediate，
% cue 953 2344.980 2346.930
\hypertarget{w06:cue:953}{}把它講嗎？but。

\tc{0:39:06}
\speaker{導讀人：亦凡}
% cue 954 2346.930 2363.260
\hypertarget{w06:cue:954}{}But Canadian legislation may be couched in a form that insulates it from this scrutiny。

\tc{0:39:23}
\speaker{老師}
% cue 955 2363.260 2366.960
\hypertarget{w06:cue:955}{}可以避免掉這樣的嚴格性。
% cue 956 2366.960 2374.340
\hypertarget{w06:cue:956}{}Canadian assemblies legislate notwithstanding the right in the charter。

\tc{0:39:34}
% cue 957 2374.340 2377.530
\hypertarget{w06:cue:957}{}這句話什麼意思？我看不懂。

\tc{0:39:37}
\speaker{導讀人：亦凡}
% cue 958 2377.530 2379.950
\hypertarget{w06:cue:958}{}就是有一個制度，
% cue 959 2379.950 2383.050
\hypertarget{w06:cue:959}{}或者是有一種可能性。

\tc{0:39:43}
\speaker{老師}
% cue 960 2383.050 2384.070
\hypertarget{w06:cue:960}{}重點是這句，
% cue 961 2384.070 2386.190
\hypertarget{w06:cue:961}{}Canadian Assemblies
% cue 962 2386.190 2389.010
\hypertarget{w06:cue:962}{}他們的議會還是可以立法，
% cue 963 2389.010 2391.680
\hypertarget{w06:cue:963}{}儘管，
% cue 964 2391.680 2399.160
\hypertarget{w06:cue:964}{}The notwithstanding clause\footnote{指加拿大權利與自由憲章第 33 條，刊本 1356 頁註 32 錄其第 1、2 項：國會或省議會得於法律中明示宣告，該法律或其條文排除憲章第 2 條或第 7 條至第 15 條的適用而施行。依第 3 項，此項宣告自生效起五年失效，第 4 項許其重新制定。Waldron 在 1356 至 1357 頁稱此條款實際上很少援用，故把加拿大的制度計為強式違憲審查。\url{https://laws-lois.justice.gc.ca/eng/Const/FullText.html}。} is rarely invoked
% cue 965 2399.160 2400.500
\hypertarget{w06:cue:965}{}很少被使用，
% cue 966 2400.500 2403.070
\hypertarget{w06:cue:966}{}但是因為這個樣子，
% cue 967 2403.070 2404.690
\hypertarget{w06:cue:967}{}他還把它算成實，
% cue 968 2404.690 2406.370
\hypertarget{w06:cue:968}{}介在兩者之中，
% cue 969 2406.370 2408.050
\hypertarget{w06:cue:969}{}但是這句話，
% cue 970 2408.050 2409.230
\hypertarget{w06:cue:970}{}我其實說真的，
% cue 971 2409.230 2410.740
\hypertarget{w06:cue:971}{}我沒有完全看懂，
% cue 972 2410.740 2412.490
\hypertarget{w06:cue:972}{}這個可能還要再去查，
% cue 973 2412.490 2413.910
\hypertarget{w06:cue:973}{}最主要應該是，
% cue 974 2413.910 2415.590
\hypertarget{w06:cue:974}{}小註 32 這裡，
% cue 975 2415.590 2417.390
\hypertarget{w06:cue:975}{}有沒有，
% cue 976 2417.390 2419.000
\hypertarget{w06:cue:976}{}所以你看不懂啊，
% cue 977 2419.000 2421.880
\hypertarget{w06:cue:977}{}那你就要去看看這個小註 32 在講什麼？

\tc{0:40:21}
% cue 978 2421.880 2423.440
\hypertarget{w06:cue:978}{}換句話說，
% cue 979 2423.440 2425.090
\hypertarget{w06:cue:979}{}從前面看起來，
% cue 980 2425.090 2429.230
\hypertarget{w06:cue:980}{}這個法院是可以審查立法，
% cue 981 2429.230 2430.020
\hypertarget{w06:cue:981}{}沒問題，
% cue 982 2430.020 2431.470
\hypertarget{w06:cue:982}{}因為寫在憲法裡，
% cue 983 2431.470 2433.460
\hypertarget{w06:cue:983}{}而且它也可以拒絕適用，
% cue 984 2433.460 2434.360
\hypertarget{w06:cue:984}{}也沒問題，
% cue 985 2434.360 2437.410
\hypertarget{w06:cue:985}{}然後它現在唯一的差別，
% cue 986 2437.410 2440.670
\hypertarget{w06:cue:986}{}就是即便是被宣告為違憲，
% cue 987 2440.670 2442.170
\hypertarget{w06:cue:987}{}這個時候呢，
% cue 988 2442.170 2444.850
\hypertarget{w06:cue:988}{}還是可以，
% cue 989 2444.850 2447.340
\hypertarget{w06:cue:989}{}他們的立法者，
% cue 990 2447.340 2451.090
\hypertarget{w06:cue:990}{}還是可以 legislate the right in the charter，
% cue 991 2451.090 2455.020
\hypertarget{w06:cue:991}{}他還是可以就這個權利來加以立法，
% cue 992 2455.020 2456.550
\hypertarget{w06:cue:992}{}這句話意思不清楚，
% cue 993 2456.550 2458.560
\hypertarget{w06:cue:993}{}所以不清楚是要開槍度 32，
% cue 994 2458.560 2462.430
\hypertarget{w06:cue:994}{}那這裡一定是顯然對於剛才的效果有了打折。

\tc{0:41:02}
% cue 995 2462.430 2464.020
\hypertarget{w06:cue:995}{}好了，然後往下講，
% cue 996 2464.020 2467.820
\hypertarget{w06:cue:996}{}美國呢，第 3、第 7 頁都不快，
% cue 997 2467.820 2475.780
\hypertarget{w06:cue:997}{}這個在老師的文章裡面有講到，
% cue 998 2475.780 2480.690
\hypertarget{w06:cue:998}{}我是認為我們跟大陸法系站在一起看這個文章，
% cue 999 2480.690 2482.550
\hypertarget{w06:cue:999}{}主要的是應該要強調，
% cue 1000 2482.550 2484.200
\hypertarget{w06:cue:1000}{}是抽象的違憲審查，
% cue 1001 2484.200 2487.570
\hypertarget{w06:cue:1001}{}還是個案當中所做的具體違憲審查，
% cue 1002 2487.570 2489.170
\hypertarget{w06:cue:1002}{}或者是兩個都有，
% cue 1003 2489.170 2490.990
\hypertarget{w06:cue:1003}{}像德國人是兩個都有\footnote{德國基本法第 94 條第 1 項第 2 款規定，聯邦憲法法院經聯邦政府、邦政府或聯邦議會議員四分之一聲請，就聯邦法或邦法與基本法是否相符的爭議或疑義為裁判，即抽象規範審查（abstrakte Normenkontrolle）；第 100 條第 1 項規定，法院認為裁判所據的法律違憲者，應停止訴訟程序，聲請聯邦憲法法院裁判，即具體規範審查（konkrete Normenkontrolle）。\url{https://www.gesetze-im-internet.de/gg/art_94.html}。}，
% cue 1004 2490.990 2493.410
\hypertarget{w06:cue:1004}{}像法國人只有抽象沒有具體\footnote{法國第五共和憲法第 61 條規定，法律於公布前，得由總統、總理、國民議會議長、參議院議長或六十名國民議會議員、六十名參議員提請憲法委員會審查。2008 年修憲增訂第 61-1 條：訴訟繫屬中主張法律條文侵害憲法保障的權利與自由者，經中央行政法院（Conseil d’État）或最高法院（Cour de cassation）移送，憲法委員會得受理之，稱合憲性優先問題（question prioritaire de constitutionnalité）。此程序自 2010 年 3 月 1 日起實施。\cite{const1958}。}，
% cue 1005 2493.410 2495.770
\hypertarget{w06:cue:1005}{}像美國人只有具體沒有抽象\footnote{美國聯邦憲法第 3 條第 2 項第 1 款規定，司法權及於依憲法、合眾國法律與條約所生的一切案件（Cases），以及該款所列各類爭議（Controversies），聯邦法院的違憲審查在具體案件中為之。\url{https://www.archives.gov/founding-docs/constitution-transcript}。}，
% cue 1006 2495.770 2498.710
\hypertarget{w06:cue:1006}{}我會從比較憲法的觀點看重的是這個，
% cue 1007 2498.710 2500.570
\hypertarget{w06:cue:1007}{}但是這個畢竟是個美國人，
% cue 1008 2500.570 2503.300
\hypertarget{w06:cue:1008}{}他不會去想這樣的事情，
% cue 1009 2503.300 2505.290
\hypertarget{w06:cue:1009}{}所以他現在直接講到美國，
% cue 1010 2505.290 2506.730
\hypertarget{w06:cue:1010}{}就是直接講下去，
% cue 1011 2506.730 2510.490
\hypertarget{w06:cue:1011}{}他告訴你們下面有些什麼樣的問題。

\section{第二種區分：權利在憲法體系中的位置}

\tc{0:41:50}
% cue 1012 2510.490 2513.130
\hypertarget{w06:cue:1012}{}然後第三個，
% cue 1013 2513.130 2514.420
\hypertarget{w06:cue:1013}{}1358 頁，
% cue 1014 2514.420 2521.180
\hypertarget{w06:cue:1014}{}剛才講的是第一個，
% cue 1015 2521.180 2521.940
\hypertarget{w06:cue:1015}{}Distinction
% cue 1016 2521.940 2524.340
\hypertarget{w06:cue:1016}{}第一個 Distinction，
% cue 1017 2524.340 2528.110
\hypertarget{w06:cue:1017}{}到 1357 頁之前，
% cue 1018 2528.110 2529.770
\hypertarget{w06:cue:1018}{}是第一個 distinction，
% cue 1019 2529.770 2531.280
\hypertarget{w06:cue:1019}{}簡單的講就是，
% cue 1020 2531.280 2533.350
\hypertarget{w06:cue:1020}{}strong versus weak
% cue 1021 2533.350 2535.470
\hypertarget{w06:cue:1021}{}聽懂嗎，
% cue 1022 2535.470 2541.310
\hypertarget{w06:cue:1022}{}強的跟弱的，
% cue 1023 2541.310 2544.730
\hypertarget{w06:cue:1023}{}第二個 distinction，
% cue 1024 2544.730 2547.220
\hypertarget{w06:cue:1024}{}從 1357 頁開始，
% cue 1025 2547.220 2547.720
\hypertarget{w06:cue:1025}{}有沒有，
% cue 1026 2547.720 2549.270
\hypertarget{w06:cue:1026}{}a second distinction
% cue 1027 2549.270 2552.840
\hypertarget{w06:cue:1027}{}那你繼續講。

\tc{0:42:32}
\speaker{導讀人：亦凡}
% cue 1028 2552.840 2555.020
\hypertarget{w06:cue:1028}{}我先想一下，
% cue 1029 2555.020 2556.660
\hypertarget{w06:cue:1029}{}大家可以看我的，
% cue 1030 2556.660 2558.740
\hypertarget{w06:cue:1030}{}那個檔案的第 80，
% cue 1031 2558.740 2560.970
\hypertarget{w06:cue:1031}{}PDF 的第 81 頁，
% cue 1032 2560.970 2562.590
\hypertarget{w06:cue:1032}{}就是我有整理，
% cue 1033 2562.590 2565.370
\hypertarget{w06:cue:1033}{}那個表格跟圖形，
% cue 1034 2565.370 2566.130
\hypertarget{w06:cue:1034}{}表格，
% cue 1035 2566.130 2567.130
\hypertarget{w06:cue:1035}{}就是把剛剛講的。

\tc{0:42:47}
\speaker{老師}
% cue 1036 2567.130 2567.910
\hypertarget{w06:cue:1036}{}沒關系，
% cue 1037 2567.910 2568.670
\hypertarget{w06:cue:1037}{}你先講，
% cue 1038 2568.670 2570.210
\hypertarget{w06:cue:1038}{}你現在看到的這個是什麼，
% cue 1039 2570.210 2570.510
\hypertarget{w06:cue:1039}{}好，
% cue 1040 2570.510 2572.010
\hypertarget{w06:cue:1040}{}它是第二種分類，
% cue 1041 2572.010 2574.540
\hypertarget{w06:cue:1041}{}現在的第二種分類，
% cue 1042 2574.540 2578.340
\hypertarget{w06:cue:1042}{}違憲審查的第二種分類，
% cue 1043 2578.340 2579.280
\hypertarget{w06:cue:1043}{}是什麼跟什麼？

\tc{0:42:59}
\speaker{導讀人：亦凡}
% cue 1044 2579.280 2580.140
\hypertarget{w06:cue:1044}{}是。

\tc{0:43:00}
\speaker{老師}
% cue 1045 2580.140 2582.200
\hypertarget{w06:cue:1045}{}the place of individual right\footnote{刊本 1357 頁：「A second distinction among types of judicial review pays attention to the place of individual rights in the constitutional system of a society.」在美國，權利審查屬一般憲法審查的一部分，法院基於個人權利宣告法律違憲，與基於其他憲法限制宣告違憲同其精神；1358 頁另舉英國依 Human Rights Act 進行的弱式審查，其所據的權利法案不屬於（結構性的）憲法。}
% cue 1046 2582.200 2585.320
\hypertarget{w06:cue:1046}{}對啊，
% cue 1047 2585.320 2590.440
\hypertarget{w06:cue:1047}{}在第一個分類是按照效力，
% cue 1048 2590.440 2592.120
\hypertarget{w06:cue:1048}{}現在第二個分類，
% cue 1049 2592.120 2595.020
\hypertarget{w06:cue:1049}{}按照這個權利所在的位置，
% cue 1050 2595.020 2596.960
\hypertarget{w06:cue:1050}{}你是放在憲法裡的，
% cue 1051 2596.960 2598.810
\hypertarget{w06:cue:1051}{}還是不放在憲法裡，
% cue 1052 2598.810 2600.060
\hypertarget{w06:cue:1052}{}對嗎，
% cue 1053 2600.060 2601.280
\hypertarget{w06:cue:1053}{}是這個意思嗎？

\tc{0:43:21}
\speaker{導讀人：亦凡}
% cue 1054 2601.280 2601.470
\hypertarget{w06:cue:1054}{}嗯。

\tc{0:43:21}
\speaker{老師}
% cue 1055 2601.470 2607.400
\hypertarget{w06:cue:1055}{}這種分類我覺得是很少看到的，
% cue 1056 2607.400 2609.220
\hypertarget{w06:cue:1056}{}因為沒有這樣分了，
% cue 1057 2609.220 2610.480
\hypertarget{w06:cue:1057}{}這是作者，
% cue 1058 2610.480 2636.130
\hypertarget{w06:cue:1058}{}當你把人權放進憲法裡，
% cue 1059 2636.130 2638.140
\hypertarget{w06:cue:1059}{}這個時候所做的違憲審查，
% cue 1060 2638.140 2639.770
\hypertarget{w06:cue:1060}{}就如同大法官說，
% cue 1061 2639.770 2642.470
\hypertarget{w06:cue:1061}{}我們現在違反了中央地方分權，
% cue 1062 2642.470 2644.570
\hypertarget{w06:cue:1062}{}或者違反權力分立原則一樣，
% cue 1063 2644.570 2647.110
\hypertarget{w06:cue:1063}{}你違反了憲法這一章的根據保障，
% cue 1064 2647.110 2649.050
\hypertarget{w06:cue:1064}{}他認為這個意思是相同的，
% cue 1065 2649.050 2653.120
\hypertarget{w06:cue:1065}{}有沒有清楚最以上是 2，
% cue 1066 2653.120 2656.810
\hypertarget{w06:cue:1066}{}但是這個 2 有沒有對照組，
% cue 1067 2656.810 2658.490
\hypertarget{w06:cue:1067}{}有沒有不是的，
% cue 1068 2658.490 2676.980
\hypertarget{w06:cue:1068}{}說在第三五八頁。

\tc{0:44:36}
\speaker{導讀人：亦凡}
% cue 1069 2676.980 2680.460
\hypertarget{w06:cue:1069}{}他就是英國跟美國的那個對照。

\tc{0:44:40}
\speaker{老師}
% cue 1070 2680.460 2682.610
\hypertarget{w06:cue:1070}{}看一下第三五八頁的第一大段，
% cue 1071 2682.610 2683.830
\hypertarget{w06:cue:1071}{}小註 36，
% cue 1072 2683.830 2686.660
\hypertarget{w06:cue:1072}{}In other countries' judicial review
% cue 1073 2686.660 2689.530
\hypertarget{w06:cue:1073}{}Take a place with regard to bill of rights
% cue 1074 2689.530 2695.520
\hypertarget{w06:cue:1074}{}That is not the specifically designated part of the structural constitution
% cue 1075 2695.520 2696.240
\hypertarget{w06:cue:1075}{}就這一句，
% cue 1076 2696.240 2698.590
\hypertarget{w06:cue:1076}{}有聽懂嗎？

\tc{0:44:58}
% cue 1077 2698.590 2700.870
\hypertarget{w06:cue:1077}{}我們剛才說分類吧，
% cue 1078 2700.870 2702.110
\hypertarget{w06:cue:1078}{}你要 A 跟 B，
% cue 1079 2702.110 2703.430
\hypertarget{w06:cue:1079}{}你現在是第二大類，
% cue 1080 2703.430 2706.090
\hypertarget{w06:cue:1080}{}第二大類是看你講的人權，
% cue 1081 2706.090 2707.270
\hypertarget{w06:cue:1081}{}它放的位置，
% cue 1082 2707.270 2710.120
\hypertarget{w06:cue:1082}{}人權在憲法體系中的位置，
% cue 1083 2710.120 2712.080
\hypertarget{w06:cue:1083}{}人權是放進憲法法典，
% cue 1084 2712.080 2714.400
\hypertarget{w06:cue:1084}{}還是人權是放進其他的地方，
% cue 1085 2714.400 2716.210
\hypertarget{w06:cue:1085}{}但你不在憲法法典，
% cue 1086 2716.210 2718.840
\hypertarget{w06:cue:1086}{}聽懂嗎，
% cue 1087 2718.840 2722.330
\hypertarget{w06:cue:1087}{}有聽懂嗎？

\tc{0:45:22}
% cue 1088 2722.330 2726.570
\hypertarget{w06:cue:1088}{}注意看他這一句，
% cue 1089 2726.570 2731.510
\hypertarget{w06:cue:1089}{}1358 頁就是剛才講小註 36 的下面句，
% cue 1090 2731.510 2734.040
\hypertarget{w06:cue:1090}{}Weak Judicial Review 答答答答，
% cue 1091 2734.040 2735.700
\hypertarget{w06:cue:1091}{}Is of this kind
% cue 1092 2735.700 2736.580
\hypertarget{w06:cue:1092}{}英國就是這樣，
% cue 1093 2736.580 2740.420
\hypertarget{w06:cue:1093}{}Because most cases of strong judicial review
% cue 1094 2740.420 2744.310
\hypertarget{w06:cue:1094}{}Are associated with constitutional review
% cue 1095 2744.310 2746.900
\hypertarget{w06:cue:1095}{}I should focus on these cases
% cue 1096 2746.900 2748.540
\hypertarget{w06:cue:1096}{}它主要是要討論這個，
% cue 1097 2748.540 2754.930
\hypertarget{w06:cue:1097}{}主要是要討論立法者通過的法律，
% cue 1098 2754.930 2759.040
\hypertarget{w06:cue:1098}{}抵觸了放在憲法之中的人權條款，
% cue 1099 2759.040 2759.900
\hypertarget{w06:cue:1099}{}這種情況，
% cue 1100 2759.900 2761.910
\hypertarget{w06:cue:1100}{}好 以上是第二個分類。

\section{第三種區分：事前審查與事後審查}

\tc{0:46:01}
% cue 1101 2761.910 2763.150
\hypertarget{w06:cue:1101}{}第三個分類，
% cue 1102 2763.150 2765.500
\hypertarget{w06:cue:1102}{}1358 頁最底下，
% cue 1103 2765.500 2767.020
\hypertarget{w06:cue:1103}{}第三個分類，
% cue 1104 2767.020 2768.680
\hypertarget{w06:cue:1104}{}我已經講過很多遍，
% cue 1105 2768.680 2771.800
\hypertarget{w06:cue:1105}{}你讀書要找到脈絡，
% cue 1106 2771.800 2774.840
\hypertarget{w06:cue:1106}{}不然你就沒有辦法接著做報告，
% cue 1107 2774.840 2776.820
\hypertarget{w06:cue:1107}{}它現在不是有脈絡嗎，
% cue 1108 2776.820 2778.300
\hypertarget{w06:cue:1108}{}它的脈絡你要把它找出來，
% cue 1109 2778.300 2779.820
\hypertarget{w06:cue:1109}{}那第三個分類，
% cue 1110 2779.820 2781.200
\hypertarget{w06:cue:1110}{}請講。

\tc{0:46:21}
\speaker{導讀人：亦凡}
% cue 1111 2781.200 2785.120
\hypertarget{w06:cue:1111}{}第三個分類就是事前審查跟事後審查。

\tc{0:46:25}
\speaker{老師}
% cue 1112 2785.120 2785.480
\hypertarget{w06:cue:1112}{}好。

\tc{0:46:25}
% cue 1113 2785.480 2788.200
\hypertarget{w06:cue:1113}{}那你現在看到的第一個是事前還是事後，
% cue 1114 2788.200 2790.400
\hypertarget{w06:cue:1114}{}posteriori review\footnote{拉丁文 a posteriori 意為「由後」，ex ante 意為「事前」。刊本 1358 至 1359 頁的第三種區分，一方是美國式的 a posteriori review，在具體訴訟中進行，有時在法律制定很久之後；另一方是 ex ante review，由專設的憲法法院在法案制定的最後階段作抽象評估。Waldron 說明本文集中討論前者。}
% cue 1115 2790.400 2791.440
\hypertarget{w06:cue:1115}{}事後，
% cue 1116 2791.440 2794.820
\hypertarget{w06:cue:1116}{}那前面有一個是 ex ante review，
% cue 1117 2794.820 2795.420
\hypertarget{w06:cue:1117}{}事前，
% cue 1118 2795.420 2797.620
\hypertarget{w06:cue:1118}{}現在我們大概全世界，
% cue 1119 2797.620 2801.040
\hypertarget{w06:cue:1119}{}違憲審查國家都是事後，
% cue 1120 2801.040 2803.550
\hypertarget{w06:cue:1120}{}就是法律生效以後去做審查，
% cue 1121 2803.550 2805.440
\hypertarget{w06:cue:1121}{}還沒有生效以前的審查，
% cue 1122 2805.440 2806.340
\hypertarget{w06:cue:1122}{}大概只有法國。

\tc{0:46:46}
% cue 1123 2806.340 2808.670
\hypertarget{w06:cue:1123}{}那在這個地方，
% cue 1124 2808.670 2810.370
\hypertarget{w06:cue:1124}{}你會非常懷疑，
% cue 1125 2810.370 2813.090
\hypertarget{w06:cue:1125}{}在法國從事違憲審查的那個機關，
% cue 1126 2813.090 2814.950
\hypertarget{w06:cue:1126}{}老師這個文章裡面有寫\footnote{湯德宗〈論違憲審查制度的改進：由「多元多軌」到「一元單軌」的改進方案〉，收於《違憲審查與動態平衡：權力分立新論 卷二》（增訂三版，元照，2005），刊本頁 175 起。該文壹部分以集權式與分權式、抽象與具體、事前與事後、諮詢性與拘束性四組區分比較各國制度；課綱所列教科書為該書增訂五版（2026）。}，
% cue 1127 2814.950 2816.210
\hypertarget{w06:cue:1127}{}書裡面有寫，
% cue 1128 2816.210 2817.980
\hypertarget{w06:cue:1128}{}這個叫做憲法委\footnote{指法國憲法委員會（Conseil constitutionnel），規定於第五共和憲法第七章。第 56 條規定委員九人，任期九年不得連任，每三年改任三分之一，由總統、國民議會議長、參議院議長各任命三人，卸任總統為當然終身委員。第 24 條規定國會由國民議會與參議院組成。\url{https://www.conseil-constitutionnel.fr/le-bloc-de-constitutionnalite/texte-integral-de-la-constitution-du-4-octobre-1958-en-vigueur}。}，
% cue 1129 2817.980 2820.340
\hypertarget{w06:cue:1129}{}他試院嗎，
% cue 1130 2820.340 2823.990
\hypertarget{w06:cue:1130}{}還沒有生效的法律，
% cue 1131 2823.990 2824.830
\hypertarget{w06:cue:1131}{}你可以審查，
% cue 1132 2824.830 2826.770
\hypertarget{w06:cue:1132}{}那我認為你是法院嗎，
% cue 1133 2826.770 2829.170
\hypertarget{w06:cue:1133}{}是在適用法律，
% cue 1134 2829.170 2831.180
\hypertarget{w06:cue:1134}{}還是你是立法者，
% cue 1135 2831.180 2834.500
\hypertarget{w06:cue:1135}{}你是兩院之外的第三院，
% cue 1136 2834.500 2835.500
\hypertarget{w06:cue:1136}{}那兩院，
% cue 1137 2835.500 2836.980
\hypertarget{w06:cue:1137}{}因為法國也是兩院，
% cue 1138 2836.980 2838.220
\hypertarget{w06:cue:1138}{}如果臺灣只有一院，
% cue 1139 2838.220 2839.540
\hypertarget{w06:cue:1139}{}那就變成第二院，
% cue 1140 2839.540 2841.740
\hypertarget{w06:cue:1140}{}大法官變成第二院，
% cue 1141 2841.740 2842.780
\hypertarget{w06:cue:1141}{}你是事前審查，
% cue 1142 2842.780 2844.060
\hypertarget{w06:cue:1142}{}立法院說了不算，
% cue 1143 2844.060 2845.240
\hypertarget{w06:cue:1143}{}你還可以再說一次，
% cue 1144 2845.240 2846.140
\hypertarget{w06:cue:1144}{}他通過了，
% cue 1145 2846.140 2847.040
\hypertarget{w06:cue:1145}{}你可以不讓他通過，
% cue 1146 2847.040 2848.750
\hypertarget{w06:cue:1146}{}我學高幫微信給你通過，
% cue 1147 2848.750 2850.020
\hypertarget{w06:cue:1147}{}事前審查，
% cue 1148 2850.020 2851.300
\hypertarget{w06:cue:1148}{}還沒有公布，
% cue 1149 2851.300 2852.080
\hypertarget{w06:cue:1149}{}還沒有生效，
% cue 1150 2852.080 2853.200
\hypertarget{w06:cue:1150}{}我就先審查。

\tc{0:47:33}
% cue 1151 2853.200 2855.770
\hypertarget{w06:cue:1151}{}這種情況下，
% cue 1152 2855.770 2857.020
\hypertarget{w06:cue:1152}{}最大的問題是，
% cue 1153 2857.020 2858.630
\hypertarget{w06:cue:1153}{}你要不要認定，
% cue 1154 2858.630 2859.810
\hypertarget{w06:cue:1154}{}那個審查機構，
% cue 1155 2859.810 2860.450
\hypertarget{w06:cue:1155}{}那個事前，
% cue 1156 2860.450 2862.670
\hypertarget{w06:cue:1156}{}那個事前違憲審查機關，
% cue 1157 2862.670 2863.750
\hypertarget{w06:cue:1157}{}是法院，
% cue 1158 2863.750 2866.390
\hypertarget{w06:cue:1158}{}它算 judicial review model，
% cue 1159 2866.390 2867.930
\hypertarget{w06:cue:1159}{}這樣的意思。

\section{第四種區分：普通法院審查與憲法法院審查}

\tc{0:47:47}
% cue 1160 2867.930 2868.830
\hypertarget{w06:cue:1160}{}第四個，
% cue 1161 2868.830 2872.150
\hypertarget{w06:cue:1161}{}第四個。

\tc{0:47:52}
\speaker{導讀人：亦凡}
% cue 1162 2872.150 2880.300
\hypertarget{w06:cue:1162}{}它就是看違憲審查是由普通法院\footnote{刊本 1359 頁的第四種區分：「Judicial review can be carried out by ordinary courts (as in the Massachusetts case we began with) or it can be carried out by a specialized constitutional court.」所稱麻州案件為導論開頭的 Goodridge v. Department of Public Health（2003）。Waldron 說明此區分與後文的論點相關：一般法院的法官受審判紀律牽引，就權利直接說理的能力常被高估。}進行，
% cue 1163 2880.300 2884.140
\hypertarget{w06:cue:1163}{}還是集中式的違憲審查，
% cue 1164 2884.140 2885.420
\hypertarget{w06:cue:1164}{}還是像美國一樣。

\tc{0:48:05}
\speaker{老師}
% cue 1165 2885.420 2888.740
\hypertarget{w06:cue:1165}{}在這個地方你要照它的字面翻，
% cue 1166 2888.740 2890.900
\hypertarget{w06:cue:1166}{}然後你才去用老師的話說，
% cue 1167 2890.900 2893.520
\hypertarget{w06:cue:1167}{}我們是普通法院審查，
% cue 1168 2893.520 2897.600
\hypertarget{w06:cue:1168}{}還是一個專業特別的憲法法院審查，
% cue 1169 2897.600 2900.100
\hypertarget{w06:cue:1169}{}如果是後者，
% cue 1170 2900.100 2902.430
\hypertarget{w06:cue:1170}{}通常只有憲法法院有審查權，
% cue 1171 2902.430 2904.890
\hypertarget{w06:cue:1171}{}所以它是集權式的違憲審查\footnote{湯德宗〈論違憲審查制度的改進〉壹、一界定：「集權式審查」（centralized review）指一國之內僅有某特定法院（或特定機關）可為違憲審查，原型見於奧地利 1920 年 10 月 1 日頒布的憲法，現存典範為德國聯邦憲法法院；「分權式審查」（decentralized review）指各個法院皆可為違憲審查，原型見於美國憲法。該文註中說明亦有稱「集中審查」者（蘇永欽），其意或為避免「集權」二字在中文的負面意涵，並主張兩詞意涵不同。}，
% cue 1172 2904.890 2907.030
\hypertarget{w06:cue:1172}{}然後如果前面，
% cue 1173 2907.030 2909.750
\hypertarget{w06:cue:1173}{}Ordinary courts 一定是很多法院，
% cue 1174 2909.750 2914.160
\hypertarget{w06:cue:1174}{}任何法院只要它在個案中都可以審查，
% cue 1175 2914.160 2917.090
\hypertarget{w06:cue:1175}{}所以我說這叫分權式的違憲審查。

\tc{0:48:37}
% cue 1176 2917.090 2919.940
\hypertarget{w06:cue:1176}{}但是在臺灣因為很多原因，
% cue 1177 2919.940 2922.470
\hypertarget{w06:cue:1177}{}大家不喜歡用分權跟集權，
% cue 1178 2922.470 2924.590
\hypertarget{w06:cue:1178}{}所以大部分就叫分權式的審查，
% cue 1179 2924.590 2926.200
\hypertarget{w06:cue:1179}{}跟集中式的審查，
% cue 1180 2926.200 2930.110
\hypertarget{w06:cue:1180}{}但我認為這個名詞是有問題的，
% cue 1181 2930.110 2931.530
\hypertarget{w06:cue:1181}{}因為它會造成誤會，
% cue 1182 2931.530 2933.410
\hypertarget{w06:cue:1182}{}你說集權式的審查，
% cue 1183 2933.410 2935.130
\hypertarget{w06:cue:1183}{}你說很多法院都有權，
% cue 1184 2935.130 2936.980
\hypertarget{w06:cue:1184}{}只不過把它集中到這個法院，
% cue 1185 2936.980 2938.410
\hypertarget{w06:cue:1185}{}實際叫集中式，
% cue 1186 2938.410 2940.090
\hypertarget{w06:cue:1186}{}如果我講集權式，
% cue 1187 2940.090 2942.720
\hypertarget{w06:cue:1187}{}那只有這個法院有其他法院是沒有的，
% cue 1188 2942.720 2944.520
\hypertarget{w06:cue:1188}{}這不一樣的。

\tc{0:49:04}
% cue 1189 2944.520 2947.240
\hypertarget{w06:cue:1189}{}集權式的審查必然是集中審查，
% cue 1190 2947.240 2949.800
\hypertarget{w06:cue:1190}{}集中審查不一定是集權審查，
% cue 1191 2949.800 2951.020
\hypertarget{w06:cue:1191}{}三個法院都有權，
% cue 1192 2951.020 2953.320
\hypertarget{w06:cue:1192}{}我為了抗訊原因把它集中到這個法院，
% cue 1193 2953.320 2955.630
\hypertarget{w06:cue:1193}{}所以這個名詞我是不同意的，
% cue 1194 2955.630 2957.990
\hypertarget{w06:cue:1194}{}那我也不認為集權一定是壞事，
% cue 1195 2957.990 2960.030
\hypertarget{w06:cue:1195}{}中央集權一定是壞事，
% cue 1196 2960.030 2965.460
\hypertarget{w06:cue:1196}{}所以我不認為這個是有什麼必要在這裡做文章，
% cue 1197 2965.460 2968.680
\hypertarget{w06:cue:1197}{}總之這個是書裡面有講到的就是，
% cue 1198 2968.680 2971.220
\hypertarget{w06:cue:1198}{}文章叫什麼的分線還在修改，
% cue 1199 2971.220 2972.540
\hypertarget{w06:cue:1199}{}書還沒出來。

\tc{0:49:32}
% cue 1200 2972.540 2973.780
\hypertarget{w06:cue:1200}{}好了第二個部分，
% cue 1201 2973.780 2975.220
\hypertarget{w06:cue:1201}{}剛剛第一部分講完了，
% cue 1202 2975.220 2976.980
\hypertarget{w06:cue:1202}{}有關於違憲審查，
% cue 1203 2976.980 2979.290
\hypertarget{w06:cue:1203}{}我要攻擊的是違憲審查，
% cue 1204 2979.290 2980.330
\hypertarget{w06:cue:1204}{}什麼叫違憲審查，
% cue 1205 2980.330 2981.250
\hypertarget{w06:cue:1205}{}大概差不多講完了，
% cue 1206 2981.250 2982.760
\hypertarget{w06:cue:1206}{}它主要是針對的，
% cue 1207 2982.760 2985.440
\hypertarget{w06:cue:1207}{}strong 的或事後的，
% cue 1208 2985.440 2988.030
\hypertarget{w06:cue:1208}{}或是寫在線反法用的，
% cue 1209 2988.030 2990.170
\hypertarget{w06:cue:1209}{}所以也是 constitutional review 的一種，
% cue 1210 2990.170 2994.110
\hypertarget{w06:cue:1210}{}然後呢至於說是哪一個法院來審查，
% cue 1211 2994.110 2995.250
\hypertarget{w06:cue:1211}{}他有沒有在意嗎？

\tc{0:49:55}
\speaker{導讀人：亦凡}
% cue 1212 2995.250 2999.790
\hypertarget{w06:cue:1212}{}就是沒有講很多。

\tc{0:49:59}
\speaker{老師}
% cue 1213 2999.790 3002.150
\hypertarget{w06:cue:1213}{}這個部分我們看不出來，
% cue 1214 3002.150 3004.730
\hypertarget{w06:cue:1214}{}他好像不太在意，
% cue 1215 3004.730 3008.170
\hypertarget{w06:cue:1215}{}他不在意他要批評的是哪一個，
% cue 1216 3008.170 3009.240
\hypertarget{w06:cue:1216}{}換句話說，
% cue 1217 3009.240 3011.800
\hypertarget{w06:cue:1217}{}憲法法院也是他批評的對象，
% cue 1218 3011.800 3013.680
\hypertarget{w06:cue:1218}{}美國聯邦最高法院也是，
% cue 1219 3013.680 3016.730
\hypertarget{w06:cue:1219}{}反正只要是違憲審查就好，
% cue 1220 3016.730 3019.010
\hypertarget{w06:cue:1220}{}這個第四個他沒有特別在意。

\section{四項假設之設定目的}

\tc{0:50:19}
% cue 1221 3019.010 3020.720
\hypertarget{w06:cue:1221}{}好來講四個假設，
% cue 1222 3020.720 3026.020
\hypertarget{w06:cue:1222}{}現在進行到一天三百五十九夜。

\tc{0:50:26}
\speaker{導讀人：亦凡}
% cue 1223 3026.020 3028.830
\hypertarget{w06:cue:1223}{}這個假設就是，
% cue 1224 3028.830 3030.190
\hypertarget{w06:cue:1224}{}大家可以看那個，
% cue 1225 3030.190 3031.320
\hypertarget{w06:cue:1225}{}附錄，
% cue 1226 3031.320 3032.540
\hypertarget{w06:cue:1226}{}附錄整理的，
% cue 1227 3032.540 3040.350
\hypertarget{w06:cue:1227}{}如果直接講的話。

\tc{0:50:40}
\speaker{老師}
% cue 1228 3040.350 3042.150
\hypertarget{w06:cue:1228}{}先講一下他要說什麼，
% cue 1229 3042.150 3044.070
\hypertarget{w06:cue:1229}{}他為什麼需要假設，
% cue 1230 3044.070 3046.190
\hypertarget{w06:cue:1230}{}My fault to focus my argument
% cue 1231 3046.190 3049.340
\hypertarget{w06:cue:1231}{}and to distinguish the core case
% cue 1232 3049.340 3053.510
\hypertarget{w06:cue:1232}{}in which objection to judicial review
% cue 1233 3053.510 3056.080
\hypertarget{w06:cue:1233}{}is at its clearest
% cue 1234 3056.080 3058.980
\hypertarget{w06:cue:1234}{}from the core case
% cue 1235 3058.980 3060.340
\hypertarget{w06:cue:1235}{}in which judicial review
% cue 1236 3060.340 3061.920
\hypertarget{w06:cue:1236}{}might be deemed appropriate
% cue 1237 3061.920 3063.300
\hypertarget{w06:cue:1237}{}這兩個不一樣，
% cue 1238 3063.300 3068.780
\hypertarget{w06:cue:1238}{}As an anomalous provision to deal with a special pathologies
% cue 1239 3068.780 3071.590
\hypertarget{w06:cue:1239}{}I shall set up some assumptions。

\tc{0:51:11}
% cue 1240 3071.590 3074.980
\hypertarget{w06:cue:1240}{}它為什麼需要設這四個假設的，
% cue 1241 3074.980 3076.370
\hypertarget{w06:cue:1241}{}這句話你要看著，
% cue 1242 3076.370 3077.480
\hypertarget{w06:cue:1242}{}你如果看不懂，
% cue 1243 3077.480 3080.480
\hypertarget{w06:cue:1243}{}這一段就沒有辦法理解它為什麼要講一句，
% cue 1244 3080.480 3082.140
\hypertarget{w06:cue:1244}{}這才是你覺得要講的，
% cue 1245 3082.140 3083.000
\hypertarget{w06:cue:1245}{}它怎麼那麼奇怪？

\tc{0:51:23}
\speaker{導讀人：亦凡}
% cue 1246 3083.000 3087.240
\hypertarget{w06:cue:1246}{}是為了論證上面的區分，
% cue 1247 3087.240 3088.040
\hypertarget{w06:cue:1247}{}如果是。

\tc{0:51:28}
\speaker{老師}
% cue 1248 3088.040 3089.440
\hypertarget{w06:cue:1248}{}你先不要抽象，
% cue 1249 3089.440 3093.560
\hypertarget{w06:cue:1249}{}你不要去先講剛才那麼一句抽象的話，
% cue 1250 3093.560 3096.520
\hypertarget{w06:cue:1250}{}你就把他這句話講完你就懂了嘛，
% cue 1251 3096.520 3097.840
\hypertarget{w06:cue:1251}{}是嗎，
% cue 1252 3097.840 3100.870
\hypertarget{w06:cue:1252}{}他說我主要，
% cue 1253 3100.870 3103.970
\hypertarget{w06:cue:1253}{}因為我區別了 the core case，
% cue 1254 3103.970 3105.570
\hypertarget{w06:cue:1254}{}和 non core case，
% cue 1255 3105.570 3106.370
\hypertarget{w06:cue:1255}{}對吧，
% cue 1256 3106.370 3109.040
\hypertarget{w06:cue:1256}{}我要批評的是核心的，
% cue 1257 3109.040 3111.190
\hypertarget{w06:cue:1257}{}違憲審查，
% cue 1258 3111.190 3112.030
\hypertarget{w06:cue:1258}{}核心的，
% cue 1259 3112.030 3113.350
\hypertarget{w06:cue:1259}{}典型的，
% cue 1260 3113.350 3115.600
\hypertarget{w06:cue:1260}{}我覺得 the core case 很難翻的，
% cue 1261 3115.600 3117.500
\hypertarget{w06:cue:1261}{}你把它翻成是典型的違憲審查，
% cue 1262 3117.500 3119.820
\hypertarget{w06:cue:1262}{}與非典型的違憲審查，
% cue 1263 3119.820 3122.170
\hypertarget{w06:cue:1263}{}我批評的是典型違憲審查，
% cue 1264 3122.170 3124.540
\hypertarget{w06:cue:1264}{}那在這樣的情況之下呢，
% cue 1265 3124.540 3125.360
\hypertarget{w06:cue:1265}{}我一定要講，
% cue 1266 3125.360 3127.110
\hypertarget{w06:cue:1266}{}這個典型的違憲情況，
% cue 1267 3127.110 3132.480
\hypertarget{w06:cue:1267}{}是在什麼樣的假設下，
% cue 1268 3132.480 3133.390
\hypertarget{w06:cue:1268}{}它來運轉，
% cue 1269 3133.390 3136.080
\hypertarget{w06:cue:1269}{}所以我需要講一下這四個假設，
% cue 1270 3136.080 3138.800
\hypertarget{w06:cue:1270}{}如果這四個假設不存在，
% cue 1271 3138.800 3141.120
\hypertarget{w06:cue:1271}{}那就不是我講的典型的違憲審查，
% cue 1272 3141.120 3143.430
\hypertarget{w06:cue:1272}{}那這個時候有可能，
% cue 1273 3143.430 3145.510
\hypertarget{w06:cue:1273}{}違憲審查也許是好的，
% cue 1274 3145.510 3147.560
\hypertarget{w06:cue:1274}{}例外我允許，
% cue 1275 3147.560 3150.170
\hypertarget{w06:cue:1275}{}這樣講有聽懂嗎？

\tc{0:52:30}
% cue 1276 3150.170 3152.570
\hypertarget{w06:cue:1276}{}這樣子講完了，
% cue 1277 3152.570 3154.630
\hypertarget{w06:cue:1277}{}才你剛才加你剛才結論，
% cue 1278 3154.630 3156.310
\hypertarget{w06:cue:1278}{}這是為了它行文的需要，
% cue 1279 3156.310 3158.620
\hypertarget{w06:cue:1279}{}你不講在你這樣抽象的講，
% cue 1280 3158.620 3160.600
\hypertarget{w06:cue:1280}{}你同好非常德國式，
% cue 1281 3160.600 3162.240
\hypertarget{w06:cue:1281}{}那你這樣講，
% cue 1282 3162.240 3163.740
\hypertarget{w06:cue:1282}{}我們聽不懂你要講什麼，
% cue 1283 3163.740 3166.320
\hypertarget{w06:cue:1283}{}你勢必還是要再解釋一遍，
% cue 1284 3166.320 3168.190
\hypertarget{w06:cue:1284}{}為了行文的需要是什麼，
% cue 1285 3168.190 3170.370
\hypertarget{w06:cue:1285}{}所以你就照它的脈絡講完。

\tc{0:52:50}
% cue 1286 3170.370 3171.330
\hypertarget{w06:cue:1286}{}所以你可以說，
% cue 1287 3171.330 3173.920
\hypertarget{w06:cue:1287}{}可見它是為了行文的需要，
% cue 1288 3173.920 3175.550
\hypertarget{w06:cue:1288}{}行文的需要是作者接近你，
% cue 1289 3175.550 3177.340
\hypertarget{w06:cue:1289}{}每個人我要寫文章，
% cue 1290 3177.340 3177.920
\hypertarget{w06:cue:1290}{}我要怎麼寫，
% cue 1291 3177.920 3178.740
\hypertarget{w06:cue:1291}{}你能管得到嗎，
% cue 1292 3178.740 3179.840
\hypertarget{w06:cue:1292}{}那我要這麼寫，
% cue 1293 3179.840 3180.590
\hypertarget{w06:cue:1293}{}我有這個考慮，
% cue 1294 3180.590 3182.460
\hypertarget{w06:cue:1294}{}但是你要記得，
% cue 1295 3182.460 3184.220
\hypertarget{w06:cue:1295}{}那現在整個文章，
% cue 1296 3184.220 3186.260
\hypertarget{w06:cue:1296}{}是以這四個假設作為前提，
% cue 1297 3186.260 3187.370
\hypertarget{w06:cue:1297}{}批評的，
% cue 1298 3187.370 3189.160
\hypertarget{w06:cue:1298}{}這四個假設如果不存在，
% cue 1299 3189.160 3190.340
\hypertarget{w06:cue:1299}{}它不見得反對違憲，
% cue 1300 3190.340 3191.860
\hypertarget{w06:cue:1300}{}它不見得反對違憲，
% cue 1301 3191.860 3193.730
\hypertarget{w06:cue:1301}{}這樣子可以了嗎？

\section{假設 A：民主制度}

\tc{0:53:13}
% cue 1302 3193.730 3194.660
\hypertarget{w06:cue:1302}{}好 先講，
% cue 1303 3194.660 3198.580
\hypertarget{w06:cue:1303}{}第一個 1361 頁 嗎？

\tc{0:53:18}
\speaker{導讀人：亦凡}
% cue 1304 3198.580 3201.290
\hypertarget{w06:cue:1304}{}對差不多。

\tc{0:53:21}
\speaker{老師}
% cue 1305 3201.290 3202.850
\hypertarget{w06:cue:1305}{}好來 趕快就要 1361 頁，
% cue 1306 3202.850 3203.970
\hypertarget{w06:cue:1306}{}這樣有 ABCD 嗎，
% cue 1307 3203.970 3205.300
\hypertarget{w06:cue:1307}{}我已經講過了，
% cue 1308 3205.300 3206.520
\hypertarget{w06:cue:1308}{}他把脈絡抓出來，
% cue 1309 3206.520 3208.770
\hypertarget{w06:cue:1309}{}A 在 1361 頁，
% cue 1310 3208.770 3210.890
\hypertarget{w06:cue:1310}{}B 在 1363 頁，
% cue 1311 3210.890 3213.640
\hypertarget{w06:cue:1311}{}C 在 1364 頁，
% cue 1312 3213.640 3215.420
\hypertarget{w06:cue:1312}{}在 1366 頁，
% cue 1313 3215.420 3216.230
\hypertarget{w06:cue:1313}{}好了啦，
% cue 1314 3216.230 3217.390
\hypertarget{w06:cue:1314}{}頭腳完了，
% cue 1315 3217.390 3218.490
\hypertarget{w06:cue:1315}{}你看腳腳尾巴，
% cue 1316 3218.490 3219.550
\hypertarget{w06:cue:1316}{}那這就結束了，
% cue 1317 3219.550 3220.870
\hypertarget{w06:cue:1317}{}A 是什麼？

\tc{0:53:40}
\speaker{導讀人：亦凡}
% cue 1318 3220.870 3226.680
\hypertarget{w06:cue:1318}{}民主的機制。

\tc{0:53:46}
\speaker{老師}
% cue 1319 3226.680 3229.020
\hypertarget{w06:cue:1319}{}來 民主的機制是什麼意思？

\tc{0:53:49}
\speaker{導讀人：亦凡}
% cue 1320 3229.020 3232.870
\hypertarget{w06:cue:1320}{}民主就是有，
% cue 1321 3232.870 3233.940
\hypertarget{w06:cue:1321}{}嗯，
% cue 1322 3233.940 3235.570
\hypertarget{w06:cue:1322}{}假設。

\tc{0:53:55}
\speaker{老師}
% cue 1323 3235.570 3251.010
\hypertarget{w06:cue:1323}{}經由選舉，
% cue 1324 3251.010 3253.580
\hypertarget{w06:cue:1324}{}然後產生代表，
% cue 1325 3253.580 3254.900
\hypertarget{w06:cue:1325}{}然後還有人民，
% cue 1326 3254.900 3256.280
\hypertarget{w06:cue:1326}{}人民和人民的，
% cue 1327 3256.280 3258.220
\hypertarget{w06:cue:1327}{}經由選舉產生的人民代表，
% cue 1328 3258.220 3259.760
\hypertarget{w06:cue:1328}{}他們來制定法律，
% cue 1329 3259.760 3262.620
\hypertarget{w06:cue:1329}{}然後來制定公共政策。

\tc{0:54:22}
% cue 1330 3262.620 3269.070
\hypertarget{w06:cue:1330}{}那 this society has a broadly democratic political system，
% cue 1331 3269.070 3271.090
\hypertarget{w06:cue:1331}{}我們說這樣的社會呢，
% cue 1332 3271.090 3275.770
\hypertarget{w06:cue:1332}{}就有一個廣義的民主的政治的體系，
% cue 1333 3275.770 3279.410
\hypertarget{w06:cue:1333}{}with universal adult suffrage
% cue 1334 3279.410 3279.990
\hypertarget{w06:cue:1334}{}什麼意思？

\tc{0:54:39}
\speaker{導讀人：亦凡}
% cue 1335 3279.990 3284.820
\hypertarget{w06:cue:1335}{}就是有一個大。

\tc{0:54:44}
\speaker{老師}
% cue 1336 3284.820 3285.120
\hypertarget{w06:cue:1336}{}什麼意思，
% cue 1337 3285.120 3287.740
\hypertarget{w06:cue:1337}{}universal 是全面的，
% cue 1338 3287.740 3289.420
\hypertarget{w06:cue:1338}{}adult suffrage
% cue 1339 3289.420 3292.750
\hypertarget{w06:cue:1339}{}成年人普遍的都有投票權，
% cue 1340 3292.750 3299.180
\hypertarget{w06:cue:1340}{}and it has a representative legislature
% cue 1341 3299.180 3302.900
\hypertarget{w06:cue:1341}{}由人民選出的代表組成的國會，
% cue 1342 3302.900 3305.250
\hypertarget{w06:cue:1342}{}有代議式的國會，
% cue 1343 3305.250 3311.700
\hypertarget{w06:cue:1343}{}To which elections are held on a fair and regular basis
% cue 1344 3311.700 3316.020
\hypertarget{w06:cue:1344}{}這個 legislature 的代表是怎麼來的呢，
% cue 1345 3316.020 3318.040
\hypertarget{w06:cue:1345}{}立法委員是怎麼來的呢，
% cue 1346 3318.040 3322.190
\hypertarget{w06:cue:1346}{}他們是經由公平和經常性的選舉。

\tc{0:55:22}
% cue 1347 3322.190 3326.640
\hypertarget{w06:cue:1347}{}我現在講的這幾句話，
% cue 1348 3326.640 3330.890
\hypertarget{w06:cue:1348}{}他也許不像大陸法系人講那麼清楚，
% cue 1349 3330.890 3333.020
\hypertarget{w06:cue:1349}{}但是你可以反過來講，
% cue 1350 3333.020 3336.550
\hypertarget{w06:cue:1350}{}這是他對於所謂的民主機制的定義，
% cue 1351 3336.550 3340.650
\hypertarget{w06:cue:1351}{}我們的決策是透過民主的程序完成，
% cue 1352 3340.650 3343.510
\hypertarget{w06:cue:1352}{}而這些民主的機關，
% cue 1353 3343.510 3347.800
\hypertarget{w06:cue:1353}{}他們的代表是人民有普遍的選舉，
% cue 1354 3347.800 3348.980
\hypertarget{w06:cue:1354}{}每個人都有成年人，
% cue 1355 3348.980 3350.060
\hypertarget{w06:cue:1355}{}每個人都有選舉權，
% cue 1356 3350.060 3352.100
\hypertarget{w06:cue:1356}{}然後定期的選舉，
% cue 1357 3352.100 3354.510
\hypertarget{w06:cue:1357}{}然後公平的選舉所產生，
% cue 1358 3354.510 3357.660
\hypertarget{w06:cue:1358}{}這樣才叫做 Democratic Institutions。

\tc{0:55:57}
% cue 1359 3357.660 3360.160
\hypertarget{w06:cue:1359}{}然後除此之外呢，
% cue 1360 3360.160 3361.840
\hypertarget{w06:cue:1360}{}還會有牽涉到立法程序，
% cue 1361 3361.840 3363.480
\hypertarget{w06:cue:1361}{}下面你看在下一句，
% cue 1362 3363.480 3365.740
\hypertarget{w06:cue:1362}{}註 46 的那一句，
% cue 1363 3365.740 3369.940
\hypertarget{w06:cue:1363}{}The legislators deliberate and vote on public issues
% cue 1364 3369.940 3375.570
\hypertarget{w06:cue:1364}{}立法者就公共議題來進行辯論來進行投票，
% cue 1365 3375.570 3380.260
\hypertarget{w06:cue:1365}{}and the procedures for lawmaking are elaborate and responsible
% cue 1366 3380.260 3381.940
\hypertarget{w06:cue:1366}{}還有立法程序，
% cue 1367 3381.940 3384.050
\hypertarget{w06:cue:1367}{}這裡就會用大陸法系的話，
% cue 1368 3384.050 3387.500
\hypertarget{w06:cue:1368}{}為什麼我們要強調正當立法程序，
% cue 1369 3387.500 3389.930
\hypertarget{w06:cue:1369}{}那如果沒有正當立法程序，
% cue 1370 3389.930 3393.430
\hypertarget{w06:cue:1370}{}那就沒有經過適當的討論辯論和投票，
% cue 1371 3393.430 3397.460
\hypertarget{w06:cue:1371}{}and incorporate various safeguards。

\tc{0:56:37}
% cue 1372 3397.460 3399.460
\hypertarget{w06:cue:1372}{}而且還要注意一下這句，
% cue 1373 3399.460 3400.880
\hypertarget{w06:cue:1373}{}這句就一般人寫不出來了，
% cue 1374 3400.880 3404.130
\hypertarget{w06:cue:1374}{}他說你要說它是個民主的機制，
% cue 1375 3404.130 3411.090
\hypertarget{w06:cue:1375}{}這個民主的機制本身還要蘊含著各種各樣的防衛措施，
% cue 1376 3411.090 3414.100
\hypertarget{w06:cue:1376}{}such as bicameralism
% cue 1377 3414.100 3415.120
\hypertarget{w06:cue:1377}{}兩院制，
% cue 1378 3415.120 3416.730
\hypertarget{w06:cue:1378}{}有一院發瘋了，
% cue 1379 3416.730 3418.150
\hypertarget{w06:cue:1379}{}還有另外一個院可以 check，
% cue 1380 3418.150 3419.920
\hypertarget{w06:cue:1380}{}兩院是互相制衡，
% cue 1381 3419.920 3424.240
\hypertarget{w06:cue:1381}{}robust committee scrutiny
% cue 1382 3424.240 3426.580
\hypertarget{w06:cue:1382}{}在各個委員會當中，
% cue 1383 3426.580 3428.630
\hypertarget{w06:cue:1383}{}他們進行強有力的審查，
% cue 1384 3428.630 3430.720
\hypertarget{w06:cue:1384}{}委員會不能虛掉，
% cue 1385 3430.720 3432.270
\hypertarget{w06:cue:1385}{}還要有委員會，
% cue 1386 3432.270 3434.750
\hypertarget{w06:cue:1386}{}and multiple levels of consideration
% cue 1387 3434.750 3437.360
\hypertarget{w06:cue:1387}{}而且討論還要經過很多次，
% cue 1388 3437.360 3438.850
\hypertarget{w06:cue:1388}{}Debate and Voting
% cue 1389 3438.850 3440.760
\hypertarget{w06:cue:1389}{}所以都不是單一的表決，
% cue 1390 3440.760 3442.620
\hypertarget{w06:cue:1390}{}有一次表決就瘋了，
% cue 1391 3442.620 3446.180
\hypertarget{w06:cue:1391}{}不可能三四次四次表決都瘋了。

\tc{0:57:26}
% cue 1392 3446.180 3449.880
\hypertarget{w06:cue:1392}{}你就用我的白話講，
% cue 1393 3449.880 3451.580
\hypertarget{w06:cue:1393}{}要不要採用兩院制，
% cue 1394 3451.580 3454.400
\hypertarget{w06:cue:1394}{}那兩院制在審議的過程中，
% cue 1395 3454.400 3455.280
\hypertarget{w06:cue:1395}{}有沒有委員會，
% cue 1396 3455.280 3457.980
\hypertarget{w06:cue:1396}{}以及我們有沒有三讀程序\footnote{立法院職權行使法第 7 條規定：「立法院依憲法第六十三條規定所議決之議案，除法律案、預算案應經三讀會議決外，其餘均經二讀會議決之。」\url{https://law.moj.gov.tw/LawClass/LawSingle.aspx?pcode=A0020058&flno=7}。}，
% cue 1397 3457.980 3459.540
\hypertarget{w06:cue:1397}{}我們大陸法新發獎，
% cue 1398 3459.540 3462.180
\hypertarget{w06:cue:1398}{}這些都是賬的機制，
% cue 1399 3462.180 3465.230
\hypertarget{w06:cue:1399}{}不然你說我們有投票，
% cue 1400 3465.230 3467.290
\hypertarget{w06:cue:1400}{}這些人是我們選的，
% cue 1401 3467.290 3468.610
\hypertarget{w06:cue:1401}{}這樣就講完了嗎？

\tc{0:57:48}
% cue 1402 3468.610 3470.740
\hypertarget{w06:cue:1402}{}你看它的論證，
% cue 1403 3470.740 3473.750
\hypertarget{w06:cue:1403}{}你要在這種地方去學習看，
% cue 1404 3473.750 3479.470
\hypertarget{w06:cue:1404}{}美國人或者是 British 的這種系統，
% cue 1405 3479.470 3481.470
\hypertarget{w06:cue:1405}{}Anglo-Saxon 系統下面的人，
% cue 1406 3481.470 3483.670
\hypertarget{w06:cue:1406}{}他們想事情解事情的方法，
% cue 1407 3483.670 3485.680
\hypertarget{w06:cue:1407}{}它是用你最容易講的先講，
% cue 1408 3485.680 3487.550
\hypertarget{w06:cue:1408}{}然後再一層層加上去，
% cue 1409 3487.550 3489.010
\hypertarget{w06:cue:1409}{}再一層層加上去，
% cue 1410 3489.010 3490.560
\hypertarget{w06:cue:1410}{}這跟大陸法系不同，
% cue 1411 3490.560 3493.610
\hypertarget{w06:cue:1411}{}我認為民主程序由下列三變，
% cue 1412 3493.610 3496.320
\hypertarget{w06:cue:1412}{}按一二三這樣的行為，
% cue 1413 3496.320 3499.260
\hypertarget{w06:cue:1413}{}所以他們的思維不是這樣的思維，
% cue 1414 3499.260 3500.780
\hypertarget{w06:cue:1414}{}它是一層層加上去，
% cue 1415 3500.780 3502.900
\hypertarget{w06:cue:1415}{}加上去好像夯土一樣，
% cue 1416 3502.900 3503.740
\hypertarget{w06:cue:1416}{}夯了三遍，
% cue 1417 3503.740 3504.800
\hypertarget{w06:cue:1417}{}這個腹就扎死，
% cue 1418 3504.800 3506.370
\hypertarget{w06:cue:1418}{}我先講底座，
% cue 1419 3506.370 3507.840
\hypertarget{w06:cue:1419}{}然後我再加一層，
% cue 1420 3507.840 3508.840
\hypertarget{w06:cue:1420}{}我再加一層，
% cue 1421 3508.840 3510.280
\hypertarget{w06:cue:1421}{}夯實上去，
% cue 1422 3510.280 3512.490
\hypertarget{w06:cue:1422}{}這樣有聽懂嗎？

\tc{0:58:32}
% cue 1423 3512.490 3514.950
\hypertarget{w06:cue:1423}{}那再往下，
% cue 1424 3514.950 3516.650
\hypertarget{w06:cue:1424}{}Members of Legislature
% cue 1425 3516.650 3519.440
\hypertarget{w06:cue:1425}{}有關於代議士立法委員，
% cue 1426 3519.440 3521.820
\hypertarget{w06:cue:1426}{}他們是怎麼思考他們自己的，
% cue 1427 3521.820 3524.760
\hypertarget{w06:cue:1427}{}這句話在一般的人也不會寫，
% cue 1428 3524.760 3527.370
\hypertarget{w06:cue:1428}{}他說有時候把他們當成是，
% cue 1429 3527.370 3529.890
\hypertarget{w06:cue:1429}{}他們會去思考，
% cue 1430 3529.890 3532.570
\hypertarget{w06:cue:1430}{}思考什麼，
% cue 1431 3532.570 3536.120
\hypertarget{w06:cue:1431}{}當成是代議代表，
% cue 1432 3536.120 3538.130
\hypertarget{w06:cue:1432}{}代表什麼叫代表，
% cue 1433 3538.130 3541.560
\hypertarget{w06:cue:1433}{}他去思考他們選民的利益跟意見，
% cue 1434 3541.560 3545.410
\hypertarget{w06:cue:1434}{}他投票的時候投下去想，
% cue 1435 3545.410 3547.610
\hypertarget{w06:cue:1435}{}我是淡水區選出來的利益，
% cue 1436 3547.610 3549.020
\hypertarget{w06:cue:1436}{}我這樣投下去，
% cue 1437 3549.020 3550.440
\hypertarget{w06:cue:1437}{}我反對淡水該打球，
% cue 1438 3550.440 3552.420
\hypertarget{w06:cue:1438}{}大概我不能打球，
% cue 1439 3552.420 3554.100
\hypertarget{w06:cue:1439}{}還有時是這樣，
% cue 1440 3554.100 3556.200
\hypertarget{w06:cue:1440}{}然後你注意看在下一句，
% cue 1441 3556.200 3558.720
\hypertarget{w06:cue:1441}{}Sometimes 有時候，
% cue 1442 3558.720 3562.110
\hypertarget{w06:cue:1442}{}他是把整個社會來做考慮的，
% cue 1443 3562.110 3564.750
\hypertarget{w06:cue:1443}{}是不是我講的一層一層，
% cue 1444 3564.750 3565.900
\hypertarget{w06:cue:1444}{}湯上去的，
% cue 1445 3565.900 3568.220
\hypertarget{w06:cue:1445}{}這就是寫作方法。

\tc{0:59:28}
% cue 1446 3568.220 3571.320
\hypertarget{w06:cue:1446}{}這裡用我們大陸法系的話講，
% cue 1447 3571.320 3575.560
\hypertarget{w06:cue:1447}{}就是民意代表他考慮的利益，
% cue 1448 3575.560 3577.980
\hypertarget{w06:cue:1448}{}一方面是選區的利益，
% cue 1449 3577.980 3580.890
\hypertarget{w06:cue:1449}{}一方面是社會整體或國家整體的，
% cue 1450 3580.890 3582.850
\hypertarget{w06:cue:1450}{}用我們的話講呢，
% cue 1451 3582.850 3584.310
\hypertarget{w06:cue:1451}{}它是法定代表，
% cue 1452 3584.310 3587.200
\hypertarget{w06:cue:1452}{}它不是單純的委任代表\footnote{議員依選區選民指示行事者為委任代表。德國基本法第 38 條第 1 項第 2 句規定聯邦議會議員「Vertreter des ganzen Volkes, an Aufträge und Weisungen nicht gebunden und nur ihrem Gewissen unterworfen」（為全體人民的代表，不受委託與指示拘束，只服從自己的良知）。刊本 1361 頁 Waldron 寫議員有時以選區選民的利益與意見為重，有時以 virtual representation（虛擬代表）的方式考量整個社會的利益與意見。\url{https://www.gesetze-im-internet.de/gg/art_38.html}。}，
% cue 1453 3587.200 3590.320
\hypertarget{w06:cue:1453}{}它不是要照原選區選民的意思，
% cue 1454 3590.320 3592.580
\hypertarget{w06:cue:1454}{}一定要照原選區選民的意思，
% cue 1455 3592.580 3593.950
\hypertarget{w06:cue:1455}{}它不是，
% cue 1456 3593.950 3596.540
\hypertarget{w06:cue:1456}{}簡單的用我的話講是這樣子。

\tc{0:59:56}
% cue 1457 3596.540 3598.030
\hypertarget{w06:cue:1457}{}還有最後它一點，
% cue 1458 3598.030 3600.880
\hypertarget{w06:cue:1458}{}它說我也假設這裡面有政黨，
% cue 1459 3600.880 3602.460
\hypertarget{w06:cue:1459}{}還有政黨屬性，
% cue 1460 3602.460 3603.820
\hypertarget{w06:cue:1460}{}政黨歸屬的問題，
% cue 1461 3603.820 3605.410
\hypertarget{w06:cue:1461}{}party affiliation
% cue 1462 3605.410 3606.970
\hypertarget{w06:cue:1462}{}就是你隸屬於政黨，
% cue 1463 3606.970 3609.370
\hypertarget{w06:cue:1463}{}所以小鷹下個命令，
% cue 1464 3609.370 3611.450
\hypertarget{w06:cue:1464}{}綠色的都會集結起來，
% cue 1465 3611.450 3614.690
\hypertarget{w06:cue:1465}{}我是說立法委員就是政黨，
% cue 1466 3614.690 3619.410
\hypertarget{w06:cue:1466}{}他還要假定在這個民主機制裡面也有政黨，
% cue 1467 3619.410 3621.930
\hypertarget{w06:cue:1467}{}然後委員會有政黨立數，
% cue 1468 3621.930 3623.950
\hypertarget{w06:cue:1468}{}這些都是我們要考慮的，
% cue 1469 3623.950 3624.950
\hypertarget{w06:cue:1469}{}這是我的假定，
% cue 1470 3624.950 3627.170
\hypertarget{w06:cue:1470}{}我說這是一個民主機制。

\tc{1:00:27}
% cue 1471 3627.170 3628.640
\hypertarget{w06:cue:1471}{}講到這裡，
% cue 1472 3628.640 3631.810
\hypertarget{w06:cue:1472}{}我覺得一般人大體上應該是講完，
% cue 1473 3631.810 3637.960
\hypertarget{w06:cue:1473}{}但是好像作者比較還有強調一些其他的東西，
% cue 1474 3637.960 3640.820
\hypertarget{w06:cue:1474}{}比方說 1362 頁，
% cue 1475 3640.820 3645.670
\hypertarget{w06:cue:1475}{}361 頁的底下，
% cue 1476 3645.670 3646.490
\hypertarget{w06:cue:1476}{}他有一個，
% cue 1477 3646.490 3649.060
\hypertarget{w06:cue:1477}{}他有好幾個 value，
% cue 1478 3649.060 3669.420
\hypertarget{w06:cue:1478}{}這裡為止，它好像比它強調的是一種民主的文化，
% cue 1479 3669.420 3681.540
\hypertarget{w06:cue:1479}{}這種民主的文化是非常看重負責任的自治\footnote{刊本 1361 頁原文為「a culture of democracy, valuing responsible deliberation and political equality」（重視負責任的審議與政治平等的民主文化）；1362 頁接著說明第二項價值 political equality。}以及政治上的平等，
% cue 1480 3681.540 3687.460
\hypertarget{w06:cue:1480}{}然後呢，下面馬上就 1361 跳到 1362 轉過來，
% cue 1481 3687.460 3694.930
\hypertarget{w06:cue:1481}{}第二個這種價值就是政治上的平等，
% cue 1482 3694.930 3699.020
\hypertarget{w06:cue:1482}{}大概是，它要講的是這樣。

\tc{1:01:39}
% cue 1483 3699.020 3701.720
\hypertarget{w06:cue:1483}{}那下面他就會牽涉到別的東西，
% cue 1484 3701.720 3702.500
\hypertarget{w06:cue:1484}{}我們就先付了，
% cue 1485 3702.500 3706.420
\hypertarget{w06:cue:1485}{}以上第 361 到 362，
% cue 1486 3706.420 3707.600
\hypertarget{w06:cue:1486}{}第一種價值講完了，
% cue 1487 3707.600 3709.600
\hypertarget{w06:cue:1487}{}也就是說，
% cue 1488 3709.600 3711.910
\hypertarget{w06:cue:1488}{}他先從制度的最根本講解，
% cue 1489 3711.910 3713.430
\hypertarget{w06:cue:1489}{}我們的一層層價格，
% cue 1490 3713.430 3715.410
\hypertarget{w06:cue:1490}{}價格完你們到底挖到了興趣，
% cue 1491 3715.410 3716.950
\hypertarget{w06:cue:1491}{}我現在之所以講解，
% cue 1492 3716.950 3718.880
\hypertarget{w06:cue:1492}{}他們會運行價值，
% cue 1493 3718.880 3720.540
\hypertarget{w06:cue:1493}{}運行民主的價值，
% cue 1494 3720.540 3723.060
\hypertarget{w06:cue:1494}{}第一種是 deliberation\footnote{英文，指審慎的討論與權衡，政治理論上通譯「審議」。課綱第 9 週 Ferejohn 與 Pasquino 的 Constitutional Courts as Deliberative Institutions、第 12 週 Castro 的 Deliberative Democracy and the Counter-majoritarian Difficulty 均以此詞立題。} 的價值，
% cue 1495 3723.060 3725.360
\hypertarget{w06:cue:1495}{}第二種是 equality 的價值。

\tc{1:02:05}
% cue 1496 3725.360 3728.800
\hypertarget{w06:cue:1496}{}你認為如果民主沒有 deliberation，
% cue 1497 3728.800 3729.720
\hypertarget{w06:cue:1497}{}這叫什麼民主，
% cue 1498 3729.720 3732.330
\hypertarget{w06:cue:1498}{}是臺灣現在的問題，
% cue 1499 3732.330 3734.310
\hypertarget{w06:cue:1499}{}我們有民主的形式，
% cue 1500 3734.310 3736.400
\hypertarget{w06:cue:1500}{}我們有民主的內核，
% cue 1501 3736.400 3739.030
\hypertarget{w06:cue:1501}{}Equality 臺灣是有的，
% cue 1502 3739.030 3740.660
\hypertarget{w06:cue:1502}{}基本上一人一票，
% cue 1503 3740.660 3742.160
\hypertarget{w06:cue:1503}{}然後選舉的制度，
% cue 1504 3742.160 3743.710
\hypertarget{w06:cue:1504}{}兼有區域代表，
% cue 1505 3743.710 3744.950
\hypertarget{w06:cue:1505}{}也兼有比例代表\footnote{中華民國憲法增修條文第 4 條第 1 項規定，立法委員 113 人，其中自由地區直轄市、縣市 73 人，平地原住民及山地原住民各 3 人，全國不分區及僑居國外國民共 34 人；第 2 項規定全國不分區及僑居國外國民名額依政黨名單投票選舉，由獲得百分之五以上政黨選舉票的政黨依得票比率選出。\url{https://law.moj.gov.tw/LawClass/LawSingle.aspx?pcode=A0000002&flno=4}。}，
% cue 1506 3744.950 3745.980
\hypertarget{w06:cue:1506}{}所以我們的 Equality，
% cue 1507 3745.980 3746.940
\hypertarget{w06:cue:1507}{}應該還可以。

\tc{1:02:26}
% cue 1508 3746.940 3748.420
\hypertarget{w06:cue:1508}{}但是我們現在的問題，
% cue 1509 3748.420 3749.540
\hypertarget{w06:cue:1509}{}可能是 Deliberation，
% cue 1510 3749.540 3751.360
\hypertarget{w06:cue:1510}{}但是這個 Deliberation，
% cue 1511 3751.360 3752.060
\hypertarget{w06:cue:1511}{}得慢慢來，
% cue 1512 3752.060 3753.790
\hypertarget{w06:cue:1512}{}如果我們選舉的時候，
% cue 1513 3753.790 3755.030
\hypertarget{w06:cue:1513}{}大家都要順起來，
% cue 1514 3755.030 3756.370
\hypertarget{w06:cue:1514}{}羊油很舒服，
% cue 1515 3756.370 3757.730
\hypertarget{w06:cue:1515}{}不要去討論政策，
% cue 1516 3757.730 3759.180
\hypertarget{w06:cue:1516}{}那就沒有 Deliberation，
% cue 1517 3759.180 3761.660
\hypertarget{w06:cue:1517}{}我有講什麼不對？

\tc{1:02:41}
% cue 1518 3761.660 3764.380
\hypertarget{w06:cue:1518}{}那為什麼他要強調 Deliberation，
% cue 1519 3764.380 3766.900
\hypertarget{w06:cue:1519}{}他也承認民主有缺點。

\tc{1:02:46}
\speaker{導讀人：亦凡}
% cue 1520 3766.900 3768.010
\hypertarget{w06:cue:1520}{}不是藍藍。

\tc{1:02:48}
\speaker{老師}
% cue 1521 3768.010 3773.060
\hypertarget{w06:cue:1521}{}但是他也要告訴你終極的理想是要有 deliberation，
% cue 1522 3773.060 3774.880
\hypertarget{w06:cue:1522}{}要有思辨 要有討論，
% cue 1523 3774.880 3778.350
\hypertarget{w06:cue:1523}{}要對公共政治議題本身要進行討論，
% cue 1524 3778.350 3783.110
\hypertarget{w06:cue:1524}{}所以你可以說蔣茫安吉的這個什麼切劍機長，
% cue 1525 3783.110 3786.080
\hypertarget{w06:cue:1525}{}方便讓你去講他為什麼，
% cue 1526 3786.080 3789.780
\hypertarget{w06:cue:1526}{}但是這個不討論這個就沒有 deliberation，
% cue 1527 3789.780 3791.520
\hypertarget{w06:cue:1527}{}這是我要講的意思，
% cue 1528 3791.520 3795.570
\hypertarget{w06:cue:1528}{}好了 到這裡所謂的 democratic institution，
% cue 1529 3795.570 3798.250
\hypertarget{w06:cue:1529}{}他講完了 很複雜 沒有那麼簡單。

\section{假設 B：司法制度}

\tc{1:03:18}
% cue 1530 3798.250 3800.470
\hypertarget{w06:cue:1530}{}好 現在講第二個節目，
% cue 1531 3800.470 3804.190
\hypertarget{w06:cue:1531}{}他現在講，在這個前提之下，我來接近違憲審查，
% cue 1532 3804.190 3806.550
\hypertarget{w06:cue:1532}{}因為我們民主是正常運作，
% cue 1533 3806.550 3810.340
\hypertarget{w06:cue:1533}{}所以我當然覺得沒有理由交給違憲審查來做最終局的解釋，
% cue 1534 3810.340 3811.650
\hypertarget{w06:cue:1534}{}第二個假設，
% cue 1535 3811.650 3818.990
\hypertarget{w06:cue:1535}{}好，司法的見識或者司法的機關，
% cue 1536 3818.990 3819.210
\hypertarget{w06:cue:1536}{}來。

\tc{1:03:39}
\speaker{導讀人：亦凡}
% cue 1537 3819.210 3823.650
\hypertarget{w06:cue:1537}{}I assume that the society we are considering has courts
% cue 1538 3823.650 3828.550
\hypertarget{w06:cue:1538}{}that is a well-established and politically independent judiciary
% cue 1539 3828.550 3832.110
\hypertarget{w06:cue:1539}{}again in reasonably good working order
% cue 1540 3832.110 3837.780
\hypertarget{w06:cue:1540}{}set up to hear lawsuits, settle disputes, and uphold the rule of law。

\tc{1:04:26}
\speaker{老師}
% cue 1541 3866.360 3867.860
\hypertarget{w06:cue:1541}{}就是這個意思，
% cue 1542 3867.860 3871.790
\hypertarget{w06:cue:1542}{}然後這裡面比較有問題的是，
% cue 1543 3871.790 3875.530
\hypertarget{w06:cue:1543}{}司法和其他的民意機關有什麼不同，
% cue 1544 3875.530 3881.840
\hypertarget{w06:cue:1544}{}他說 the courts are mostly not elective or representative institution，
% cue 1545 3881.840 3884.110
\hypertarget{w06:cue:1545}{}絕大多數的情況下，
% cue 1546 3884.110 3886.810
\hypertarget{w06:cue:1546}{}法院不是選舉的也不是代議的機關，
% cue 1547 3886.810 3888.570
\hypertarget{w06:cue:1547}{}不是代表民意，
% cue 1548 3888.570 3889.970
\hypertarget{w06:cue:1548}{}不是民意代表的機關，
% cue 1549 3889.970 3891.670
\hypertarget{w06:cue:1549}{}他做他的事情，
% cue 1550 3891.670 3893.270
\hypertarget{w06:cue:1550}{}就算是某些州，
% cue 1551 3893.270 3895.290
\hypertarget{w06:cue:1551}{}可是在美國有少數有一些州，
% cue 1552 3895.290 3896.350
\hypertarget{w06:cue:1552}{}他是用選的，
% cue 1553 3896.350 3897.970
\hypertarget{w06:cue:1553}{}即便是這樣子，
% cue 1554 3897.970 3900.090
\hypertarget{w06:cue:1554}{}他們也沒有所謂的。

\tc{1:05:00}
% cue 1555 3900.090 3901.310
\hypertarget{w06:cue:1555}{}再往下看四行，
% cue 1556 3901.310 3903.430
\hypertarget{w06:cue:1556}{}Electoral Accountability
% cue 1557 3903.430 3907.060
\hypertarget{w06:cue:1557}{}沒有所謂的選舉負責禁，
% cue 1558 3907.060 3908.890
\hypertarget{w06:cue:1558}{}因為我們這邊的話講就是，
% cue 1559 3908.890 3910.470
\hypertarget{w06:cue:1559}{}負政治責任，
% cue 1560 3910.470 3912.670
\hypertarget{w06:cue:1560}{}所以你做了一個判決，
% cue 1561 3912.670 3913.390
\hypertarget{w06:cue:1561}{}我不滿意，
% cue 1562 3913.390 3914.290
\hypertarget{w06:cue:1562}{}所以我向你罷免，
% cue 1563 3914.290 3915.010
\hypertarget{w06:cue:1563}{}我幫你罷免，
% cue 1564 3915.010 3916.210
\hypertarget{w06:cue:1564}{}這個意思嗎，
% cue 1565 3916.210 3916.760
\hypertarget{w06:cue:1565}{}可以。

\tc{1:05:16}
% cue 1566 3916.760 3918.510
\hypertarget{w06:cue:1566}{}這個不像立法者，
% cue 1567 3918.510 3921.600
\hypertarget{w06:cue:1567}{}立法者我們剛剛去年經過了大罷免\footnote{2025 年 7 月 26 日與 8 月 23 日，第 11 屆立法委員罷免案分兩次投票，計 24 件與 7 件。中央選舉委員會 2025 年 8 月 1 日中選務字第 1143150602 號、8 月 29 日中選務字第 1143150700 號公告的投票結果均為否決。\url{https://web.cec.gov.tw/api/file/b35ac130-d35d-4a9e-a145-cee7b7c4b8d2.pdf}。}，
% cue 1568 3921.600 3925.190
\hypertarget{w06:cue:1568}{}我們現在不管這個偉大的柯總招，
% cue 1569 3925.190 3926.330
\hypertarget{w06:cue:1569}{}那是什麼想法，
% cue 1570 3926.330 3928.090
\hypertarget{w06:cue:1570}{}我們是講正常的情況，
% cue 1571 3928.090 3930.070
\hypertarget{w06:cue:1571}{}是說因為你做不好，
% cue 1572 3930.070 3931.070
\hypertarget{w06:cue:1572}{}所以我要罷免你，
% cue 1573 3931.070 3933.120
\hypertarget{w06:cue:1573}{}我要你負政治責任，
% cue 1574 3933.120 3936.320
\hypertarget{w06:cue:1574}{}我現在使用的負責手段就是罷免，
% cue 1575 3936.320 3939.870
\hypertarget{w06:cue:1575}{}這是 Political Accountability 的解釋，
% cue 1576 3939.870 3942.360
\hypertarget{w06:cue:1576}{}我們為什麼要罷免，
% cue 1577 3942.360 3945.110
\hypertarget{w06:cue:1577}{}就是要讓客與他政治責任，
% cue 1578 3945.110 3947.480
\hypertarget{w06:cue:1578}{}在政治上他必須向我負責，
% cue 1579 3947.480 3949.050
\hypertarget{w06:cue:1579}{}可是你很難想象，
% cue 1580 3949.050 3951.950
\hypertarget{w06:cue:1580}{}法官可以因為判決做不好被罷免，
% cue 1581 3951.950 3953.470
\hypertarget{w06:cue:1581}{}就算他是選的，
% cue 1582 3953.470 3955.130
\hypertarget{w06:cue:1582}{}他也不是負政治責任，
% cue 1583 3955.130 3956.060
\hypertarget{w06:cue:1583}{}這是他要講的，
% cue 1584 3956.060 3957.280
\hypertarget{w06:cue:1584}{}OK。

\tc{1:05:57}
% cue 1585 3957.280 3958.610
\hypertarget{w06:cue:1585}{}Many Defenders
% cue 1586 3958.610 3959.310
\hypertarget{w06:cue:1586}{}請繼續。

\tc{1:05:59}
\speaker{導讀人：亦凡}
% cue 1587 3959.310 3961.630
\hypertarget{w06:cue:1587}{}Many Defenders of Judicial Review。

\tc{1:06:01}
\speaker{老師}
% cue 1588 3961.630 3967.550
\hypertarget{w06:cue:1588}{}對，很多人都認為我們司法機關來從事違憲審查，
% cue 1589 3967.550 3969.310
\hypertarget{w06:cue:1589}{}因為司法機關是獨立的，很好。

\tc{1:06:09}
\speaker{導讀人：亦凡}
% cue 1590 3969.310 3975.280
\hypertarget{w06:cue:1590}{}because it means courts can deliberate on issues of principle
% cue 1591 3975.280 3982.820
\hypertarget{w06:cue:1591}{}undistracted by popular pressures and invulnerable to public anger。

\tc{1:06:22}
\speaker{老師}
% cue 1592 3982.820 3984.840
\hypertarget{w06:cue:1592}{}好，很好，這句話就是大家講的意思。

\tc{1:06:24}
% cue 1593 3984.840 3987.190
\hypertarget{w06:cue:1593}{}來，為什麼我們特別贊成，
% cue 1594 3987.190 3989.990
\hypertarget{w06:cue:1594}{}很多人特別贊成由司法來進行審查，
% cue 1595 3989.990 3993.290
\hypertarget{w06:cue:1595}{}因為法院呢？

\tc{1:06:33}
\speaker{導讀人：亦凡}
% cue 1596 3993.290 3995.700
\hypertarget{w06:cue:1596}{}不受民意幹擾。

\tc{1:06:35}
\speaker{老師}
% cue 1597 3995.700 3996.100
\hypertarget{w06:cue:1597}{}對，
% cue 1598 3996.100 3997.240
\hypertarget{w06:cue:1598}{}這是一個，
% cue 1599 3997.240 3999.060
\hypertarget{w06:cue:1599}{}沒有選舉壓力，
% cue 1600 3999.060 4000.800
\hypertarget{w06:cue:1600}{}然後還有一個句你沒講，
% cue 1601 4000.800 4003.240
\hypertarget{w06:cue:1601}{}deliberate of issues of the principle
% cue 1602 4003.240 4006.250
\hypertarget{w06:cue:1602}{}就原則問題來進行辯論，
% cue 1603 4006.250 4010.070
\hypertarget{w06:cue:1603}{}覺得法官是不是就原則問題進行辯論，
% cue 1604 4010.070 4014.620
\hypertarget{w06:cue:1604}{}這樣子可以了嗎？

\tc{1:06:54}
% cue 1605 4014.620 4020.220
\hypertarget{w06:cue:1605}{}going to assume that
% cue 1606 4020.220 4022.490
\hypertarget{w06:cue:1606}{}in the society we are considering
% cue 1607 4022.490 4025.870
\hypertarget{w06:cue:1607}{}courts are capable of performing the functions
% cue 1608 4025.870 4027.990
\hypertarget{w06:cue:1608}{}that will be assigned to them
% cue 1609 4027.990 4039.330
\hypertarget{w06:cue:1609}{}They could review legislation. The question is whether they should, and if so, whether their determination should be final and binding.
% cue 1610 4039.330 4053.110
\hypertarget{w06:cue:1610}{}他不反對審查，但是重點是，你可不可以審查完了，你就做成最終的決定，既有拘束力，而且還可以把法律宣告不效。
% cue 1611 4053.110 4060.700
\hypertarget{w06:cue:1611}{}這是他質疑的，我沒有說你不能審查，所以他不反對 weak，他反對 strong judicial review.
% cue 1612 4060.700 4064.460
\hypertarget{w06:cue:1612}{}然後小註 49 下一句，
% cue 1613 4064.460 4073.630
\hypertarget{w06:cue:1613}{}我們面對的這種法院，
% cue 1614 4073.630 4075.670
\hypertarget{w06:cue:1614}{}這種法院是一二三，
% cue 1615 4075.670 4083.650
\hypertarget{w06:cue:1615}{}這句話什麼意思？
% cue 1616 4083.650 4086.200
\hypertarget{w06:cue:1616}{}我們講的法院不是指，
% cue 1617 4086.200 4091.030
\hypertarget{w06:cue:1617}{}什麼叫做 act on their own motion，
% cue 1618 4091.030 4093.030
\hypertarget{w06:cue:1618}{}這句話不太不難懂，
% cue 1619 4093.030 4094.390
\hypertarget{w06:cue:1619}{}這句話是什麼意思？
% cue 1620 4094.390 4097.740
\hypertarget{w06:cue:1620}{}所有人都可以回答，
% cue 1621 4097.740 4102.430
\hypertarget{w06:cue:1621}{}這個法院不是他想要怎麼樣，
% cue 1622 4102.430 4103.290
\hypertarget{w06:cue:1622}{}他就可以怎麼樣，
% cue 1623 4103.290 4104.770
\hypertarget{w06:cue:1623}{}意思是說這個法院，
% cue 1624 4104.770 4106.710
\hypertarget{w06:cue:1624}{}是要應人民的申請，
% cue 1625 4106.710 4109.190
\hypertarget{w06:cue:1625}{}或者他是被動的，
% cue 1626 4109.190 4110.230
\hypertarget{w06:cue:1626}{}簡單講就是被動，
% cue 1627 4110.230 4110.970
\hypertarget{w06:cue:1627}{}他不是主動，
% cue 1628 4110.970 4113.910
\hypertarget{w06:cue:1628}{}on their own motion
% cue 1629 4113.910 4115.000
\hypertarget{w06:cue:1629}{}他不可以主動，
% cue 1630 4115.000 4117.220
\hypertarget{w06:cue:1630}{}你們覺得法院可以主動嗎？
% cue 1631 4117.220 4119.280
\hypertarget{w06:cue:1631}{}我路見不平拔刀相助，
% cue 1632 4119.280 4121.040
\hypertarget{w06:cue:1632}{}老子愛拔刀我就拔刀，
% cue 1633 4121.040 4122.440
\hypertarget{w06:cue:1633}{}沒有人來申請，
% cue 1634 4122.440 4123.980
\hypertarget{w06:cue:1634}{}也沒有人來告訴我可以拔刀，
% cue 1635 4123.980 4126.260
\hypertarget{w06:cue:1635}{}他不是指這種法院。

\tc{1:08:46}
% cue 1636 4126.260 4128.260
\hypertarget{w06:cue:1636}{}然後第二個就比較難講了，
% cue 1637 4128.260 4130.440
\hypertarget{w06:cue:1637}{}by abstract reference\footnote{刊本 1363 頁 Waldron 假設所討論的法院「do not act on their own motion or by abstract reference, but rather respond to particular claims brought by particular litigants」。abstract reference 的例子如加拿大最高法院法第 53 條：Governor in Council（總督會同內閣）得就憲法的解釋、聯邦或省法律的合憲性等重要法律或事實問題，提交最高法院審理並陳述意見。Waldron 在 1358 至 1359 頁第三種區分已提到憲法法院就法案所作的抽象評估，並說明本文集中討論事後審查。}
% cue 1638 4130.440 4132.870
\hypertarget{w06:cue:1638}{}這個我就不知道該怎麼回答，
% cue 1639 4132.870 4134.570
\hypertarget{w06:cue:1639}{}因為我剛才講過了，
% cue 1640 4134.570 4135.990
\hypertarget{w06:cue:1640}{}他這個文章最大的漏洞，
% cue 1641 4135.990 4138.350
\hypertarget{w06:cue:1641}{}是他忘記了大陸法系的法院，
% cue 1642 4138.350 4140.450
\hypertarget{w06:cue:1642}{}是有抽象違憲審查，
% cue 1643 4140.450 4144.450
\hypertarget{w06:cue:1643}{}我可以比方說某些國家的法院，
% cue 1644 4144.450 4149.390
\hypertarget{w06:cue:1644}{}他們甚至於可以因應政府的請求，
% cue 1645 4149.390 4152.090
\hypertarget{w06:cue:1645}{}就某個問題提供諮詢意見，
% cue 1646 4152.090 4153.980
\hypertarget{w06:cue:1646}{}我們如果要這樣做，
% cue 1647 4153.980 4154.780
\hypertarget{w06:cue:1647}{}會不會違憲的，
% cue 1648 4154.780 4156.020
\hypertarget{w06:cue:1648}{}你可以來問一問，
% cue 1649 4156.020 4158.540
\hypertarget{w06:cue:1649}{}這是 abstract reference，
% cue 1650 4158.540 4160.280
\hypertarget{w06:cue:1650}{}五月的個案，
% cue 1651 4160.280 4161.690
\hypertarget{w06:cue:1651}{}所以它是抽象的，
% cue 1652 4161.690 4164.360
\hypertarget{w06:cue:1652}{}Reference 就是送來給我那麼久，
% cue 1653 4164.360 4166.750
\hypertarget{w06:cue:1653}{}你有申請抽象申請，
% cue 1654 4166.750 4169.320
\hypertarget{w06:cue:1654}{}用白話講可能是比較相似的，
% cue 1655 4169.320 4171.730
\hypertarget{w06:cue:1655}{}好了不是指這個法院，
% cue 1656 4171.730 4173.430
\hypertarget{w06:cue:1656}{}那現在在這裡我就不懂，
% cue 1657 4173.430 4175.270
\hypertarget{w06:cue:1657}{}它也會包含德國的法院，
% cue 1658 4175.270 4177.530
\hypertarget{w06:cue:1658}{}德國的法院抽象審查，
% cue 1659 4177.530 4179.870
\hypertarget{w06:cue:1659}{}它們這邊都沒有講到抽象審查，
% cue 1660 4179.870 4180.950
\hypertarget{w06:cue:1660}{}它是不在了法院。

\tc{1:09:40}
% cue 1661 4180.950 4182.090
\hypertarget{w06:cue:1661}{}好第二點，

\tc{1:09:42}
\speaker{導讀人：亦凡}
% cue 1662 4182.090 4183.730
\hypertarget{w06:cue:1662}{}老師他這邊講是不是，
% cue 1663 4183.730 4186.650
\hypertarget{w06:cue:1663}{}不是做憲法審查的法院，

\tc{1:09:46}
\speaker{老師}
% cue 1664 4186.650 4188.880
\hypertarget{w06:cue:1664}{}你先不管，
% cue 1665 4188.880 4190.000
\hypertarget{w06:cue:1665}{}等一下再讓他講，
% cue 1666 4190.000 4191.280
\hypertarget{w06:cue:1666}{}第二點是什麼？
% cue 1667 4191.280 4193.360
\hypertarget{w06:cue:1667}{}他說我們現在處理的，
% cue 1668 4193.360 4194.580
\hypertarget{w06:cue:1668}{}我們現在講的法院，
% cue 1669 4194.580 4195.780
\hypertarget{w06:cue:1669}{}我所謂的法院，
% cue 1670 4195.780 4197.260
\hypertarget{w06:cue:1670}{}一，不是那種法院，

\tc{1:09:57}
\speaker{導讀人：亦凡}
% cue 1671 4197.260 4207.680
\hypertarget{w06:cue:1671}{}二，deal with issues in the context of binary adversarial presentations，

\tc{1:10:07}
\speaker{老師}
% cue 1672 4207.680 4208.760
\hypertarget{w06:cue:1672}{}好，這句話什麼意思？
% cue 1673 4208.760 4212.130
\hypertarget{w06:cue:1673}{}兩造的並且是爭訟式的，
% cue 1674 4212.130 4218.030
\hypertarget{w06:cue:1674}{}他們是在兩造爭訟式的環境裡面的，
% cue 1675 4218.030 4220.590
\hypertarget{w06:cue:1675}{}這是個典型的訴訟。

\tc{1:10:20}
% cue 1676 4220.590 4222.330
\hypertarget{w06:cue:1676}{}第三，

\tc{1:10:22}
\speaker{導讀人：亦凡}
% cue 1677 4222.330 4233.240
\hypertarget{w06:cue:1677}{}就是什麼意思？
% cue 1678 4233.240 4240.180
\hypertarget{w06:cue:1678}{}會有受到就是之前相關聯的案件，

\tc{1:10:40}
\speaker{老師}
% cue 1679 4240.180 4242.350
\hypertarget{w06:cue:1679}{}就是講先例，
% cue 1680 4242.350 4246.150
\hypertarget{w06:cue:1680}{}他會討論我們這個案子，
% cue 1681 4246.150 4248.480
\hypertarget{w06:cue:1681}{}跟我們已經的判決先例，
% cue 1682 4248.480 4250.590
\hypertarget{w06:cue:1682}{}是一樣的還是不一樣的。

\tc{1:10:50}
% cue 1683 4250.590 4253.650
\hypertarget{w06:cue:1683}{}我現在講這個人，
% cue 1684 4253.650 4255.030
\hypertarget{w06:cue:1684}{}因為他是美國人，
% cue 1685 4255.030 4257.150
\hypertarget{w06:cue:1685}{}他是講英美法系的，
% cue 1686 4257.150 4260.060
\hypertarget{w06:cue:1686}{}正在大陸法系這件事情，
% cue 1687 4260.060 4264.400
\hypertarget{w06:cue:1687}{}我至少我的理解沒有那麼 popular，
% cue 1688 4264.400 4265.820
\hypertarget{w06:cue:1688}{}我是討論的，
% cue 1689 4265.820 4268.890
\hypertarget{w06:cue:1689}{}我認為法官作為最高級的法院，
% cue 1690 4268.890 4271.610
\hypertarget{w06:cue:1690}{}而且是唯一的有權的釋憲法院機關\footnote{中華民國憲法第 78 條規定：「司法院解釋憲法，並有統一解釋法律及命令之權。」憲法訴訟法第 1 條第 1 項規定，司法院大法官組成憲法法庭，審理法規範憲法審查及裁判憲法審查、機關爭議、總統副總統彈劾、政黨違憲解散、地方自治保障與統一解釋法律及命令等案件。\url{https://law.moj.gov.tw/LawClass/LawSingle.aspx?pcode=A0000001&flno=78}。}，
% cue 1691 4271.610 4275.090
\hypertarget{w06:cue:1691}{}法官必須要受到他的判決協定，
% cue 1692 4275.090 4277.590
\hypertarget{w06:cue:1692}{}你不可以自己違反自己，
% cue 1693 4277.590 4281.030
\hypertarget{w06:cue:1693}{}即便說那個你是 300 年前的你還是你，
% cue 1694 4281.030 4281.960
\hypertarget{w06:cue:1694}{}你可以推翻，
% cue 1695 4281.960 4285.660
\hypertarget{w06:cue:1695}{}你沒有嚴格的判決協定拘束的原則，
% cue 1696 4285.660 4286.880
\hypertarget{w06:cue:1696}{}但是不是沒有，
% cue 1697 4286.880 4289.730
\hypertarget{w06:cue:1697}{}你不能說我不能推翻，
% cue 1698 4289.730 4291.810
\hypertarget{w06:cue:1698}{}但是你要自己知道，
% cue 1699 4291.810 4293.770
\hypertarget{w06:cue:1699}{}我現在講的已經苛責啦。

\tc{1:11:33}
% cue 1700 4293.770 4295.740
\hypertarget{w06:cue:1700}{}我們現在大法官哪裡有這種能力，
% cue 1701 4295.740 4297.740
\hypertarget{w06:cue:1701}{}知道自己知道他講的，
% cue 1702 4297.740 4299.740
\hypertarget{w06:cue:1702}{}跟以前的判決先例到底一樣還不一樣，
% cue 1703 4299.740 4302.090
\hypertarget{w06:cue:1703}{}我現在講是貼物業團，
% cue 1704 4302.090 4303.090
\hypertarget{w06:cue:1704}{}這是需要工作的，
% cue 1705 4303.090 4305.090
\hypertarget{w06:cue:1705}{}你想看我們現在將近有一千個判決，
% cue 1706 4305.090 4307.150
\hypertarget{w06:cue:1706}{}是不是每一判決，
% cue 1707 4307.150 4309.750
\hypertarget{w06:cue:1707}{}有一個助理就有去查，
% cue 1708 4309.750 4312.670
\hypertarget{w06:cue:1708}{}我們這麼解釋有沒有抵觸到監情講過的話，
% cue 1709 4312.670 4314.190
\hypertarget{w06:cue:1709}{}現在誰還記得呢，
% cue 1710 4314.190 4316.410
\hypertarget{w06:cue:1710}{}因為沒有一個人可以活那麼久，
% cue 1711 4316.410 4317.410
\hypertarget{w06:cue:1711}{}我已經離開了，
% cue 1712 4317.410 4320.080
\hypertarget{w06:cue:1712}{}我在的時候通通都是我在講的，
% cue 1713 4320.080 4321.360
\hypertarget{w06:cue:1713}{}那我現在走啊，
% cue 1714 4321.360 4324.360
\hypertarget{w06:cue:1714}{}那我怎麼知道誰會講呢，
% cue 1715 4324.360 4326.600
\hypertarget{w06:cue:1715}{}只會靠助理去查啦。

\tc{1:12:06}
% cue 1716 4326.600 4329.960
\hypertarget{w06:cue:1716}{}然後繼續它也假設認識有層級的，
% cue 1717 4329.960 4332.300
\hypertarget{w06:cue:1717}{}然後第三個，
% cue 1718 4332.300 4336.550
\hypertarget{w06:cue:1718}{}1364 它也假設法院裡面的法官，
% cue 1719 4336.550 4338.690
\hypertarget{w06:cue:1719}{}至少會有五個到九個，
% cue 1720 4338.690 4339.450
\hypertarget{w06:cue:1720}{}會有很多人，
% cue 1721 4339.450 4342.140
\hypertarget{w06:cue:1721}{}所以合議制的法院，
% cue 1722 4342.140 4343.360
\hypertarget{w06:cue:1722}{}不是單議法官，
% cue 1723 4343.360 4344.800
\hypertarget{w06:cue:1723}{}不是獨任制的法院，
% cue 1724 4344.800 4350.030
\hypertarget{w06:cue:1724}{}然後我們也會認識到所謂的法院，
% cue 1725 4350.030 4351.880
\hypertarget{w06:cue:1725}{}它在這個政治系統裡面，
% cue 1726 4351.880 4353.580
\hypertarget{w06:cue:1726}{}會享有一定的高，
% cue 1727 4353.580 4354.980
\hypertarget{w06:cue:1727}{}很高的崇高的地位，
% cue 1728 4354.980 4359.480
\hypertarget{w06:cue:1728}{}這樣可以了嗎？
% cue 1729 4359.480 4360.130
\hypertarget{w06:cue:1729}{}好，
% cue 1730 4360.130 4362.810
\hypertarget{w06:cue:1730}{}這是它講的第二個前提，
% cue 1731 4362.810 4365.040
\hypertarget{w06:cue:1731}{}有司法機關。

\section{權利保障之承諾}

\tc{1:12:45}
% cue 1732 4365.040 4367.100
\hypertarget{w06:cue:1732}{}好 第三個，

\tc{1:12:47}
\speaker{導讀人：亦凡}
% cue 1733 4367.100 4370.360
\hypertarget{w06:cue:1733}{}a commitment to right\footnote{英文，Waldron 第三項假設之標題作 A Commitment to Rights，通譯「對權利之承諾」。該項假設社會大多數成員對個人權利與少數者權利之理念有堅定的承諾（a strong commitment on the part of most members of the society we are contemplating to the idea of individual and minority rights），並肯認個人之若干利益與自由，不得僅因多數人便利而遭否定。見 Jeremy Waldron, The Core of the Case Against Judicial Review, 115 Yale L.J. 1346, 1364 (2006)。}

\tc{1:12:50}
\speaker{老師}
% cue 1734 4370.360 4372.140
\hypertarget{w06:cue:1734}{}這需要解釋一下，
% cue 1735 4372.140 4373.950
\hypertarget{w06:cue:1735}{}這很難翻，
% cue 1736 4373.950 4375.810
\hypertarget{w06:cue:1736}{}什麼叫 commitment to right，
% cue 1737 4375.810 4378.660
\hypertarget{w06:cue:1737}{}這個要翻成什麼，

\tc{1:12:58}
\speaker{導讀人：亦凡}
% cue 1738 4378.660 4381.350
\hypertarget{w06:cue:1738}{}我覺得承諾不太夠，
% cue 1739 4381.350 4382.070
\hypertarget{w06:cue:1739}{}就是有一種，
% cue 1740 4382.070 4386.870
\hypertarget{w06:cue:1740}{}大家都會花心思去關注那些，
% cue 1741 4386.870 4388.210
\hypertarget{w06:cue:1741}{}不管是多數還是少數，
% cue 1742 4388.210 4389.410
\hypertarget{w06:cue:1742}{}長期還是短期，
% cue 1743 4389.410 4390.830
\hypertarget{w06:cue:1743}{}或者是那種，
% cue 1744 4390.830 4395.320
\hypertarget{w06:cue:1744}{}就是比較有參與感的那種感覺，

\tc{1:13:15}
\speaker{老師}
% cue 1745 4395.320 4395.880
\hypertarget{w06:cue:1745}{}真的嗎？
% cue 1746 4395.880 4397.320
\hypertarget{w06:cue:1746}{}那你要翻成什麼，
% cue 1747 4397.320 4398.970
\hypertarget{w06:cue:1747}{}總是要翻一個嗎？

\tc{1:13:18}
\speaker{導讀人：亦凡}
% cue 1748 4398.970 4400.520
\hypertarget{w06:cue:1748}{}就翻成承諾吧，
% cue 1749 4400.520 4401.300
\hypertarget{w06:cue:1749}{}就是比較，

\tc{1:13:21}
\speaker{老師}
% cue 1750 4401.300 4401.800
\hypertarget{w06:cue:1750}{}承諾，
% cue 1751 4401.800 4403.620
\hypertarget{w06:cue:1751}{}承諾權利嗎？
% cue 1752 4403.620 4405.380
\hypertarget{w06:cue:1752}{}這聽起來有點奇怪，
% cue 1753 4405.380 4406.170
\hypertarget{w06:cue:1753}{}對不對？

\tc{1:13:26}
% cue 1754 4406.170 4408.240
\hypertarget{w06:cue:1754}{}我覺得你可不可以考慮，
% cue 1755 4408.240 4411.010
\hypertarget{w06:cue:1755}{}把它翻成是尊重權利，
% cue 1756 4411.010 4413.450
\hypertarget{w06:cue:1756}{}或者是敬重權利，
% cue 1757 4413.450 4416.940
\hypertarget{w06:cue:1757}{}commitment to right
% cue 1758 4416.940 4419.480
\hypertarget{w06:cue:1758}{}就是我們是誠心誠意的，
% cue 1759 4419.480 4420.480
\hypertarget{w06:cue:1759}{}要保護權利，
% cue 1760 4420.480 4423.440
\hypertarget{w06:cue:1760}{}這個意思，
% cue 1761 4423.440 4424.640
\hypertarget{w06:cue:1761}{}我的 commitment，
% cue 1762 4424.640 4425.960
\hypertarget{w06:cue:1762}{}就是我同意，
% cue 1763 4425.960 4427.000
\hypertarget{w06:cue:1763}{}我宣示，
% cue 1764 4427.000 4428.600
\hypertarget{w06:cue:1764}{}我要保護權利，
% cue 1765 4428.600 4432.320
\hypertarget{w06:cue:1765}{}或者是賬權利，
% cue 1766 4432.320 4435.950
\hypertarget{w06:cue:1766}{}我覺得尊重可能太多，
% cue 1767 4435.950 4437.250
\hypertarget{w06:cue:1767}{}也許是保障，
% cue 1768 4437.250 4439.950
\hypertarget{w06:cue:1768}{}賬應該比它接近中文，
% cue 1769 4439.950 4443.160
\hypertarget{w06:cue:1769}{}對權利是有保障的，
% cue 1770 4443.160 4444.360
\hypertarget{w06:cue:1770}{}尊重我們要保障，
% cue 1771 4444.360 4444.960
\hypertarget{w06:cue:1771}{}I believe。

\tc{1:14:04}
% cue 1772 4444.960 4448.280
\hypertarget{w06:cue:1772}{}下面這裡要講什麼，
% cue 1773 4448.280 4449.440
\hypertarget{w06:cue:1773}{}這裡比較難幹，
% cue 1774 4449.440 4450.300
\hypertarget{w06:cue:1774}{}你要講哪裡，

\tc{1:14:10}
\speaker{導讀人：亦凡}
% cue 1775 4450.300 4454.800
\hypertarget{w06:cue:1775}{}要不要跳過，
% cue 1776 4454.800 4456.060
\hypertarget{w06:cue:1776}{}跳過，

\tc{1:14:16}
\speaker{老師}
% cue 1777 4456.060 4457.760
\hypertarget{w06:cue:1777}{}那你總是要跟我們解釋一下，
% cue 1778 4457.760 4459.160
\hypertarget{w06:cue:1778}{}它什麼意思啊，
% cue 1779 4459.160 4460.980
\hypertarget{w06:cue:1779}{}對不對？
% cue 1780 4460.980 4462.430
\hypertarget{w06:cue:1780}{}這總是，
% cue 1781 4462.430 4464.970
\hypertarget{w06:cue:1781}{}這是它第三個假設，
% cue 1782 4464.970 4467.970
\hypertarget{w06:cue:1782}{}我們如果沒有這個假設，
% cue 1783 4467.970 4468.430
\hypertarget{w06:cue:1783}{}我就，
% cue 1784 4468.430 4472.650
\hypertarget{w06:cue:1784}{}我的反對違憲是它，
% cue 1785 4472.650 4473.690
\hypertarget{w06:cue:1785}{}它就不反了，
% cue 1786 4473.690 4475.810
\hypertarget{w06:cue:1786}{}它是在這個四個假設下，
% cue 1787 4475.810 4476.790
\hypertarget{w06:cue:1787}{}它就反了，
% cue 1788 4476.790 4482.860
\hypertarget{w06:cue:1788}{}那這個這個假設一定是我們大家都同意要去保障基本權，
% cue 1789 4482.860 4487.190
\hypertarget{w06:cue:1789}{}我們如果不同意保障基本權我就反對 RU 看什麼，
% cue 1790 4487.190 4491.650
\hypertarget{w06:cue:1790}{}只是我們同意保障基本權是怎麼個保障法，
% cue 1791 4491.650 4492.410
\hypertarget{w06:cue:1791}{}是這個意思，
% cue 1792 4492.410 4494.160
\hypertarget{w06:cue:1792}{}有沒有？

\tc{1:14:54}
% cue 1793 4494.160 4497.670
\hypertarget{w06:cue:1793}{}來我們看一下 1365 頁，
% cue 1794 4497.670 4500.450
\hypertarget{w06:cue:1794}{}to make this third assumption more concrete
% cue 1795 4500.450 4502.090
\hypertarget{w06:cue:1795}{}你為什麼沒看到這句，
% cue 1796 4502.090 4504.760
\hypertarget{w06:cue:1796}{}它要使得這個假設變得比較具體，
% cue 1797 4504.760 4521.610
\hypertarget{w06:cue:1797}{}We may assume also that the society cherishes rights to an extent that has led to the adoption of official written bill or declaration of rights of the familiar kind.
% cue 1798 4521.610 4529.150
\hypertarget{w06:cue:1798}{}I shall refer to this throughout as the Bill of Rights，人權清單 of the society concerned.
% cue 1799 4529.150 4538.530
\hypertarget{w06:cue:1799}{}他取了這些例子，這些例子都足以證明這個社會是看重並且保障人權。
% cue 1800 4538.530 4540.110
\hypertarget{w06:cue:1800}{}不然你定住幹什麼，
% cue 1801 4540.110 4542.420
\hypertarget{w06:cue:1801}{}好了，
% cue 1802 4542.420 4544.820
\hypertarget{w06:cue:1802}{}這樣可以了嗎？

\section{關於權利之分歧}

\tc{1:15:44}
% cue 1803 4544.820 4546.450
\hypertarget{w06:cue:1803}{}4
% cue 1804 4546.450 4548.870
\hypertarget{w06:cue:1804}{}4366
% cue 1805 4548.870 4549.890
\hypertarget{w06:cue:1805}{}動不快，

\tc{1:15:49}
\speaker{導讀人：亦凡}
% cue 1806 4549.890 4551.550
\hypertarget{w06:cue:1806}{}disagreement about rights\footnote{英文，「關於權利之分歧」，Waldron 第四項假設之標題。該項假設社會對權利之共識不免於一般政治問題上普遍存在之分歧，社會成員對於有哪些權利、權利之內容為何，存有相當程度之歧見（substantial dissensus as to what rights there are and what they amount to）。見 Jeremy Waldron, The Core of the Case Against Judicial Review, 115 Yale L.J. 1346, 1366–1367 (2006)。}

\tc{1:15:51}
\speaker{老師}
% cue 1807 4551.550 4554.140
\hypertarget{w06:cue:1807}{}這邊是最麻煩，
% cue 1808 4554.140 4555.980
\hypertarget{w06:cue:1808}{}最麻煩的是這個，
% cue 1809 4555.980 4558.220
\hypertarget{w06:cue:1809}{}一般人大概寫不到這裡，
% cue 1810 4558.220 4559.970
\hypertarget{w06:cue:1810}{}這是作者厲害的地方，
% cue 1811 4559.970 4561.520
\hypertarget{w06:cue:1811}{}你們寫寫看，
% cue 1812 4561.520 4563.400
\hypertarget{w06:cue:1812}{}你們會寫 3 跟 4 嗎？
% cue 1813 4563.400 4564.700
\hypertarget{w06:cue:1813}{}我覺得你寫到 1 跟 2，
% cue 1814 4564.700 4565.830
\hypertarget{w06:cue:1814}{}究竟不得了，
% cue 1815 4565.830 4567.550
\hypertarget{w06:cue:1815}{}讓你寫你不會寫 3 跟 4，
% cue 1816 4567.550 4569.350
\hypertarget{w06:cue:1816}{}寫到 3 也不會寫到 4，
% cue 1817 4569.350 4570.130
\hypertarget{w06:cue:1817}{}他也寫到 4，
% cue 1818 4570.130 4570.330
\hypertarget{w06:cue:1818}{}來，
% cue 1819 4570.330 4571.450
\hypertarget{w06:cue:1819}{}什麼意思？

\tc{1:16:11}
\speaker{導讀人：亦凡}
% cue 1820 4571.450 4574.520
\hypertarget{w06:cue:1820}{}my final and crucial

\tc{1:16:14}
\speaker{老師}
% cue 1821 4574.520 4575.620
\hypertarget{w06:cue:1821}{}這個翻譯一下，

\tc{1:16:15}
\speaker{導讀人：亦凡}
% cue 1822 4575.620 4578.440
\hypertarget{w06:cue:1822}{}關於權利的分歧，
% cue 1823 4578.440 4578.960
\hypertarget{w06:cue:1823}{}對，

\tc{1:16:18}
\speaker{老師}
% cue 1824 4578.960 4582.040
\hypertarget{w06:cue:1824}{}關於權利是有分歧的，
% cue 1825 4582.040 4584.590
\hypertarget{w06:cue:1825}{}或者是你想不同意也可以，
% cue 1826 4584.590 4589.910
\hypertarget{w06:cue:1826}{}關於權利的見解分歧，
% cue 1827 4589.910 4592.940
\hypertarget{w06:cue:1827}{}關於權利的認知分歧，
% cue 1828 4592.940 4593.880
\hypertarget{w06:cue:1828}{}還是什麼分歧，
% cue 1829 4593.880 4595.470
\hypertarget{w06:cue:1829}{}好來 怎麼個分歧，
% cue 1830 4595.470 4597.230
\hypertarget{w06:cue:1830}{}My final and crucial assumption
% cue 1831 4597.230 4598.390
\hypertarget{w06:cue:1831}{}is that

\tc{1:16:38}
\speaker{導讀人：亦凡}
% cue 1832 4598.390 4603.840
\hypertarget{w06:cue:1832}{}is that the consensus about rights

\tc{1:16:43}
\speaker{老師}
% cue 1833 4603.840 4605.300
\hypertarget{w06:cue:1833}{}我們對於權利的共識，

\tc{1:16:45}
\speaker{導讀人：亦凡}
% cue 1834 4605.300 4612.540
\hypertarget{w06:cue:1834}{}is not exempt from the incidence
% cue 1835 4612.540 4614.600
\hypertarget{w06:cue:1835}{}of general disagreement
% cue 1836 4614.600 4618.620
\hypertarget{w06:cue:1836}{}您的意思？

\tc{1:16:58}
\speaker{老師}
% cue 1837 4618.620 4621.380
\hypertarget{w06:cue:1837}{}我說我們對權利，
% cue 1838 4621.380 4623.330
\hypertarget{w06:cue:1838}{}我剛才講要保障人權，
% cue 1839 4623.330 4625.350
\hypertarget{w06:cue:1839}{}保障人權是我們的共識，
% cue 1840 4625.350 4627.650
\hypertarget{w06:cue:1840}{}我們對權利保障有共識，
% cue 1841 4627.650 4629.090
\hypertarget{w06:cue:1841}{}並不排除，
% cue 1842 4629.090 4634.000
\hypertarget{w06:cue:1842}{}the incidence of general disagreement
% cue 1843 4634.000 4636.610
\hypertarget{w06:cue:1843}{}不能說我們對權利保障，
% cue 1844 4636.610 4638.060
\hypertarget{w06:cue:1844}{}這件事情有共識，
% cue 1845 4638.060 4641.840
\hypertarget{w06:cue:1845}{}就是說我們對於所有重要的政治問題，
% cue 1846 4641.840 4645.830
\hypertarget{w06:cue:1846}{}不會出現一般性的分歧，
% cue 1847 4645.830 4649.920
\hypertarget{w06:cue:1847}{}有共識我們認為權利很重要，
% cue 1848 4649.920 4650.540
\hypertarget{w06:cue:1848}{}要保障，
% cue 1849 4650.540 4652.460
\hypertarget{w06:cue:1849}{}可是這個權利保障到哪裡，
% cue 1850 4652.460 4654.380
\hypertarget{w06:cue:1850}{}這個權利有沒有包含墮胎，
% cue 1851 4654.380 4657.190
\hypertarget{w06:cue:1851}{}這個總是大家可以討論嗎？
% cue 1852 4657.190 4658.350
\hypertarget{w06:cue:1852}{}所以我們有分歧，
% cue 1853 4658.350 4660.020
\hypertarget{w06:cue:1853}{}有共識，
% cue 1854 4660.020 4661.620
\hypertarget{w06:cue:1854}{}共識之中有分歧，
% cue 1855 4661.620 4662.620
\hypertarget{w06:cue:1855}{}因為我們中文話講，
% cue 1856 4662.620 4663.420
\hypertarget{w06:cue:1856}{}應該是這樣講，
% cue 1857 4663.420 4664.320
\hypertarget{w06:cue:1857}{}好了。

\tc{1:17:44}
% cue 1858 4664.320 4664.820
\hypertarget{w06:cue:1858}{}然後，
% cue 1859 4664.820 4668.240
\hypertarget{w06:cue:1859}{}I assume there is substantial
% cue 1860 4668.240 4670.320
\hypertarget{w06:cue:1860}{}dissensus as to
% cue 1861 4670.320 4671.500
\hypertarget{w06:cue:1861}{}哪兩件事，
% cue 1862 4671.500 4673.400
\hypertarget{w06:cue:1862}{}what rights there are
% cue 1863 4673.400 4675.260
\hypertarget{w06:cue:1863}{}他們是什麼權利，
% cue 1864 4675.260 4677.670
\hypertarget{w06:cue:1864}{}and what they amount to
% cue 1865 4677.670 4678.910
\hypertarget{w06:cue:1865}{}他們相當於什麼，
% cue 1866 4678.910 4681.660
\hypertarget{w06:cue:1866}{}所以言論自由，
% cue 1867 4681.660 4683.200
\hypertarget{w06:cue:1867}{}言論自由是什麼權，
% cue 1868 4683.200 4684.770
\hypertarget{w06:cue:1868}{}言論自由相當於什麼，
% cue 1869 4684.770 4686.650
\hypertarget{w06:cue:1869}{}這裡面有不同的層次。

\tc{1:18:06}
% cue 1870 4686.650 4692.150
\hypertarget{w06:cue:1870}{}好了 這裡還要講什麼，
% cue 1871 4692.150 4694.500
\hypertarget{w06:cue:1871}{}需要講什麼嗎？

\tc{1:18:14}
\speaker{導讀人：亦凡}
% cue 1872 4694.500 4697.260
\hypertarget{w06:cue:1872}{}後面要有更多，

\tc{1:18:17}
\speaker{老師}
% cue 1873 4697.260 4699.280
\hypertarget{w06:cue:1873}{}好 你就要照顧，
% cue 1874 4699.280 4701.500
\hypertarget{w06:cue:1874}{}我讓你照顧，
% cue 1875 4701.500 4703.470
\hypertarget{w06:cue:1875}{}他要總結一下嗎？
% cue 1876 4703.470 4705.950
\hypertarget{w06:cue:1876}{}不要了吧。

\section{論證之形式}

\tc{1:18:25}
% cue 1877 4705.950 4709.770
\hypertarget{w06:cue:1877}{}第 3 來，

\tc{1:18:29}
\speaker{導讀人：亦凡}
% cue 1878 4709.770 4719.570
\hypertarget{w06:cue:1878}{}我論證的這個形式，
% cue 1879 4719.570 4721.520
\hypertarget{w06:cue:1879}{}就是承接上面的四個假設，
% cue 1880 4721.520 4723.660
\hypertarget{w06:cue:1880}{}我就可以直接看到，
% cue 1881 4723.660 4725.660
\hypertarget{w06:cue:1881}{}1371

\tc{1:18:45}
\speaker{老師}
% cue 1882 4725.660 4727.660
\hypertarget{w06:cue:1882}{}1371 來，
% cue 1883 4727.660 4730.930
\hypertarget{w06:cue:1883}{}你跳哪裡啊？講啊，

\tc{1:18:50}
\speaker{導讀人：亦凡}
% cue 1884 4730.930 4738.910
\hypertarget{w06:cue:1884}{}我自己覺得比較重要的句子是，
% cue 1885 4738.910 4742.910
\hypertarget{w06:cue:1885}{}在什麼地方？
% cue 1886 4742.910 4745.480
\hypertarget{w06:cue:1886}{}在 1371，

\tc{1:19:05}
\speaker{老師}
% cue 1887 4745.480 4747.480
\hypertarget{w06:cue:1887}{}你要讓人家跟得到你啊？
% cue 1888 4747.480 4749.480
\hypertarget{w06:cue:1888}{}這這麼多是第幾段啊？

\tc{1:19:09}
\speaker{導讀人：亦凡}
% cue 1889 4749.480 4751.480
\hypertarget{w06:cue:1889}{}1371 的第一段，

\tc{1:19:11}
\speaker{老師}
% cue 1890 4751.480 4753.730
\hypertarget{w06:cue:1890}{}第一段？
% cue 1891 4753.730 4756.200
\hypertarget{w06:cue:1891}{}Let's assume 那裡嗎？

\tc{1:19:16}
\speaker{導讀人：亦凡}
% cue 1892 4756.200 4757.200
\hypertarget{w06:cue:1892}{}對，

\tc{1:19:17}
\speaker{老師}
% cue 1893 4757.200 4758.200
\hypertarget{w06:cue:1893}{}The question we face
% cue 1894 4758.200 4759.200
\hypertarget{w06:cue:1894}{}是這句嗎？
% cue 1895 4759.200 4760.200
\hypertarget{w06:cue:1895}{}對，
% cue 1896 4760.200 4763.120
\hypertarget{w06:cue:1896}{}這是這一段的最後一句嗎？
% cue 1897 4763.120 4764.240
\hypertarget{w06:cue:1897}{}對，
% cue 1898 4764.240 4765.640
\hypertarget{w06:cue:1898}{}你怎麼每次都不會講到，
% cue 1899 4765.640 4766.640
\hypertarget{w06:cue:1899}{}老師幫你去找。

\tc{1:19:26}
% cue 1900 4766.640 4769.880
\hypertarget{w06:cue:1900}{}我們面對的問題是，

\tc{1:19:29}
\speaker{導讀人：亦凡}
% cue 1901 4769.880 4775.020
\hypertarget{w06:cue:1901}{}whether that resolution of the legislature
% cue 1902 4775.020 4779.460
\hypertarget{w06:cue:1902}{}should be dispositive
% cue 1903 4779.460 4783.640
\hypertarget{w06:cue:1903}{}or whether there is reason to have it second guessed
% cue 1904 4783.640 4786.260
\hypertarget{w06:cue:1904}{}and perhaps overruled by the judiciary。

\tc{1:19:46}
\speaker{老師}
% cue 1905 4786.260 4788.080
\hypertarget{w06:cue:1905}{}好這句話要解釋一下，
% cue 1906 4788.080 4788.480
\hypertarget{w06:cue:1906}{}可以，
% cue 1907 4788.480 4791.470
\hypertarget{w06:cue:1907}{}但是這裡他只是還是，
% cue 1908 4791.470 4794.780
\hypertarget{w06:cue:1908}{}他還沒有講到他對於這個 form 的 argument，
% cue 1909 4794.780 4796.380
\hypertarget{w06:cue:1909}{}他還是在講問題，

\tc{1:19:56}
\speaker{導讀人：亦凡}
% cue 1910 4796.380 4799.010
\hypertarget{w06:cue:1910}{}它下面就是，

\tc{1:19:59}
\speaker{老師}
% cue 1911 4799.010 4799.810
\hypertarget{w06:cue:1911}{}這句話什麼意思？

\tc{1:19:59}
\speaker{導讀人：亦凡}
% cue 1912 4799.810 4800.410
\hypertarget{w06:cue:1912}{}先講完，

\tc{1:20:00}
\speaker{老師}
% cue 1913 4800.410 4801.510
\hypertarget{w06:cue:1913}{}沒有說。

\tc{1:20:01}
% cue 1914 4801.510 4803.810
\hypertarget{w06:cue:1914}{}這句話的意思是說，
% cue 1915 4803.810 4804.910
\hypertarget{w06:cue:1915}{}我們面對的問題，
% cue 1916 4804.910 4807.160
\hypertarget{w06:cue:1916}{}我們到底反不反對違憲審查，
% cue 1917 4807.160 4808.190
\hypertarget{w06:cue:1917}{}問題在於，
% cue 1918 4808.190 4810.270
\hypertarget{w06:cue:1918}{}立法者所做的決定，
% cue 1919 4810.270 4811.870
\hypertarget{w06:cue:1919}{}是不是就是最終的決定，
% cue 1920 4811.870 4813.270
\hypertarget{w06:cue:1920}{}還是不對，
% cue 1921 4813.270 4817.610
\hypertarget{w06:cue:1921}{}我們有理由讓他再被審查一次，
% cue 1922 4817.610 4820.190
\hypertarget{w06:cue:1922}{}而且可以對司法推翻，
% cue 1923 4820.190 4821.710
\hypertarget{w06:cue:1923}{}這就是我們講的問題，
% cue 1924 4821.710 4823.440
\hypertarget{w06:cue:1924}{}我們講的違憲審查，
% cue 1925 4823.440 4824.640
\hypertarget{w06:cue:1925}{}就因為司法可以推翻，
% cue 1926 4824.640 4828.800
\hypertarget{w06:cue:1926}{}司法可以做 second guess，再查一次，再審查一次，
% cue 1927 4828.800 4830.060
\hypertarget{w06:cue:1927}{}是這樣的問題。

\tc{1:20:30}
% cue 1928 4830.060 4833.410
\hypertarget{w06:cue:1928}{}但是還是沒有回到我剛才講的，
% cue 1929 4833.410 4836.310
\hypertarget{w06:cue:1929}{}那他講的 form of argument，
% cue 1930 4836.310 4840.220
\hypertarget{w06:cue:1930}{}我們論證的形式是怎樣的，
% cue 1931 4840.220 4842.330
\hypertarget{w06:cue:1931}{}所以這一段是多好辦的，
% cue 1932 4842.330 4846.130
\hypertarget{w06:cue:1932}{}這一段是不容易找的。

\tc{1:20:46}
\speaker{導讀人：亦凡}
% cue 1933 4846.130 4854.220
\hypertarget{w06:cue:1933}{}下面是以那個 Hobbes\footnote{Thomas Hobbes（1588–1679），英國政治哲學家，《利維坦》（Leviathan）於 1651 年出版。Waldron 於此句之註 64 引 Leviathan 第 123 頁（Richard Tuck 編，1996 年版），說明社會須先設定一套決策程序，以其運作平息原本促成該程序之爭議（settle, not reignite, the controversies）。見 Jeremy Waldron, The Core of the Case Against Judicial Review, 115 Yale L.J. 1346, 1371 (2006)。\url{https://plato.stanford.edu/entries/hobbes/}。}，

\tc{1:20:54}
\speaker{老師}
% cue 1934 4854.220 4859.510
\hypertarget{w06:cue:1934}{}Hobbes，再下一段，來，講吧，

\tc{1:20:59}
\speaker{導讀人：亦凡}
% cue 1935 4859.510 4861.790
\hypertarget{w06:cue:1935}{}The point is as old as Hobbes
% cue 1936 4861.790 4865.300
\hypertarget{w06:cue:1936}{}We must set up a decision procedure

\tc{1:21:05}
\speaker{老師}
% cue 1937 4865.300 4881.690
\hypertarget{w06:cue:1937}{}好，什麼意思？
% cue 1938 4881.690 4883.830
\hypertarget{w06:cue:1938}{}用你的白話講，
% cue 1939 4883.830 4886.170
\hypertarget{w06:cue:1939}{}我們現在面臨的，

\tc{1:21:26}
\speaker{導讀人：亦凡}
% cue 1940 4886.170 4890.430
\hypertarget{w06:cue:1940}{}在解決這個問題，
% cue 1941 4890.430 4891.990
\hypertarget{w06:cue:1941}{}他這邊強調的是，
% cue 1942 4891.990 4894.930
\hypertarget{w06:cue:1942}{}那個程序上面的。

\tc{1:21:34}
\speaker{老師}
% cue 1943 4894.930 4897.250
\hypertarget{w06:cue:1943}{}你不要讓你的句子，
% cue 1944 4897.250 4898.370
\hypertarget{w06:cue:1944}{}講得那麼模糊，
% cue 1945 4898.370 4900.230
\hypertarget{w06:cue:1945}{}我們就把他的意思翻出來，
% cue 1946 4900.230 4910.540
\hypertarget{w06:cue:1946}{}我們要設定一個決策程序，
% cue 1947 4910.540 4915.050
\hypertarget{w06:cue:1947}{}Whose operation will settle the controversy
% cue 1948 4915.050 4921.040
\hypertarget{w06:cue:1948}{}我們要設定一個解決爭議的決策程序，
% cue 1949 4921.040 4924.420
\hypertarget{w06:cue:1949}{}而不是重啟爭議的決策程序，
% cue 1950 4924.420 4926.530
\hypertarget{w06:cue:1950}{}聽懂嗎？
% cue 1951 4926.530 4928.530
\hypertarget{w06:cue:1951}{}這是他的任務，
% cue 1952 4928.530 4929.520
\hypertarget{w06:cue:1952}{}然後呢，
% cue 1953 4929.520 4933.900
\hypertarget{w06:cue:1953}{}Whose existence called for a decision procedure
% cue 1954 4933.900 4941.410
\hypertarget{w06:cue:1954}{}故事是什麼東西，
% cue 1955 4941.410 4944.840
\hypertarget{w06:cue:1955}{}是程序嗎？

\tc{1:22:24}
% cue 1956 4944.840 4950.190
\hypertarget{w06:cue:1956}{}其他同學看一下這個句子，
% cue 1957 4950.190 4951.540
\hypertarget{w06:cue:1957}{}你的問題就是，
% cue 1958 4951.540 4953.340
\hypertarget{w06:cue:1958}{}我不能說你，
% cue 1959 4953.340 4954.740
\hypertarget{w06:cue:1959}{}普遍的問題，
% cue 1960 4954.740 4956.290
\hypertarget{w06:cue:1960}{}都是這樣子，
% cue 1961 4956.290 4957.760
\hypertarget{w06:cue:1961}{}碰到了，
% cue 1962 4957.760 4959.800
\hypertarget{w06:cue:1962}{}你現在這個問題，
% cue 1963 4959.800 4960.640
\hypertarget{w06:cue:1963}{}不是只有英文，
% cue 1964 4960.640 4961.760
\hypertarget{w06:cue:1964}{}因為德文比英文，
% cue 1965 4961.760 4965.310
\hypertarget{w06:cue:1965}{}用更多的關係代名詞，
% cue 1966 4965.310 4967.230
\hypertarget{w06:cue:1966}{}帶一大堆東西，
% cue 1967 4967.230 4969.150
\hypertarget{w06:cue:1967}{}所以你要找到它帶的是什麼，
% cue 1968 4969.150 4970.320
\hypertarget{w06:cue:1968}{}你找不到了，
% cue 1969 4970.320 4971.500
\hypertarget{w06:cue:1969}{}這個句子就解不通了，
% cue 1970 4971.500 4972.830
\hypertarget{w06:cue:1970}{}德文一個句子裡面，
% cue 1971 4972.830 4974.290
\hypertarget{w06:cue:1971}{}可以有八個關係代名詞，
% cue 1972 4974.290 4976.720
\hypertarget{w06:cue:1972}{}那我們這裡才一個兩個，
% cue 1973 4976.720 4977.500
\hypertarget{w06:cue:1973}{}這裡有兩個，
% cue 1974 4977.500 4978.590
\hypertarget{w06:cue:1974}{}第一個是故事，
% cue 1975 4978.590 4979.630
\hypertarget{w06:cue:1975}{}第二個也是故事，

\tc{1:22:59}
\speaker{導讀人：亦凡}
% cue 1976 4979.630 4982.850
\hypertarget{w06:cue:1976}{}第二個故事是 controversies，

\tc{1:23:02}
\speaker{老師}
% cue 1977 4982.850 4983.450
\hypertarget{w06:cue:1977}{}controversies

\tc{1:23:03}
\speaker{導讀人：亦凡}
% cue 1978 4983.450 4984.170
\hypertarget{w06:cue:1978}{}OK。

\tc{1:23:04}
\speaker{老師}
% cue 1979 4984.170 4988.350
\hypertarget{w06:cue:1979}{}那就是由那個爭議的存在，
% cue 1980 4988.350 4994.010
\hypertarget{w06:cue:1980}{}要求我們首先要建立一個決策的程序，
% cue 1981 4994.010 4997.910
\hypertarget{w06:cue:1981}{}因為爭議要解決爭議，
% cue 1982 4997.910 5001.410
\hypertarget{w06:cue:1982}{}所以我們要先講我們怎麼解決這個爭議的程序，
% cue 1983 5001.410 5002.610
\hypertarget{w06:cue:1983}{}這個意思，
% cue 1984 5002.610 5003.800
\hypertarget{w06:cue:1984}{}然後呢，
% cue 1985 5003.800 5005.800
\hypertarget{w06:cue:1985}{}我們要設定這個程序，
% cue 1986 5005.800 5007.810
\hypertarget{w06:cue:1986}{}這個爭議的程序，
% cue 1987 5007.810 5009.750
\hypertarget{w06:cue:1987}{}是能夠解決這個爭議，
% cue 1988 5009.750 5011.770
\hypertarget{w06:cue:1988}{}而就是重燃這個爭議，
% cue 1989 5011.770 5015.950
\hypertarget{w06:cue:1989}{}聽懂。

\tc{1:23:36}
% cue 1990 5016.120 5017.630
\hypertarget{w06:cue:1990}{}迫步這句話是什麼，
% cue 1991 5017.630 5021.740
\hypertarget{w06:cue:1991}{}用你的白話講給我們聽，
% cue 1992 5021.740 5023.360
\hypertarget{w06:cue:1992}{}現在我這句都翻完了，
% cue 1993 5023.360 5024.180
\hypertarget{w06:cue:1993}{}白話是什麼意思？
% cue 1994 5024.180 5028.600
\hypertarget{w06:cue:1994}{}講註 64 前面這句白話是什麼意思？

\tc{1:23:48}
\speaker{導讀人：亦凡}
% cue 1995 5028.600 5034.060
\hypertarget{w06:cue:1995}{}我覺得跟後面他講那個程序面是呼應的。

\tc{1:23:54}
\speaker{老師}
% cue 1996 5034.060 5039.060
\hypertarget{w06:cue:1996}{}意思就是說，
% cue 1997 5039.060 5044.370
\hypertarget{w06:cue:1997}{}什麼問題都會兼有程序跟，
% cue 1998 5044.370 5047.270
\hypertarget{w06:cue:1998}{}我不用說實體，
% cue 1999 5047.270 5048.550
\hypertarget{w06:cue:1999}{}我應該說結果，
% cue 2000 5048.550 5051.970
\hypertarget{w06:cue:2000}{}我們現在是不是想要解決爭議，
% cue 2001 5051.970 5053.710
\hypertarget{w06:cue:2001}{}想要解決爭議，
% cue 2002 5053.710 5057.610
\hypertarget{w06:cue:2002}{}我們就首先要創造一個解決爭議的程序，
% cue 2003 5057.610 5060.340
\hypertarget{w06:cue:2003}{}然後這個解決爭議的程序，
% cue 2004 5060.340 5061.240
\hypertarget{w06:cue:2004}{}做完了以後，
% cue 2005 5061.240 5063.100
\hypertarget{w06:cue:2005}{}我們希望要解決爭議，
% cue 2006 5063.100 5064.410
\hypertarget{w06:cue:2006}{}而不是重燃爭議，
% cue 2007 5064.410 5065.400
\hypertarget{w06:cue:2007}{}這是結果。

\tc{1:24:25}
% cue 2008 5065.400 5070.220
\hypertarget{w06:cue:2008}{}看一遍這個句子，
% cue 2009 5070.220 5072.320
\hypertarget{w06:cue:2009}{}所有人都看，
% cue 2010 5072.320 5073.710
\hypertarget{w06:cue:2010}{}我知道這個對你們難，
% cue 2011 5073.710 5076.080
\hypertarget{w06:cue:2011}{}但是這就是老師存在的價值，
% cue 2012 5076.080 5077.460
\hypertarget{w06:cue:2012}{}如果這麼簡單，
% cue 2013 5077.460 5078.430
\hypertarget{w06:cue:2013}{}就不要老師了，
% cue 2014 5078.430 5080.370
\hypertarget{w06:cue:2014}{}那一般的老師是不會講這個，
% cue 2015 5080.370 5081.940
\hypertarget{w06:cue:2015}{}因為講這個我也很累，
% cue 2016 5081.940 5083.870
\hypertarget{w06:cue:2016}{}而且我也不一定會對，
% cue 2017 5083.870 5085.930
\hypertarget{w06:cue:2017}{}但是這個就是你要了解，
% cue 2018 5085.930 5087.050
\hypertarget{w06:cue:2018}{}這個句子的難處。

\tc{1:24:47}
% cue 2019 5087.050 5089.070
\hypertarget{w06:cue:2019}{}我們剛才前面讀到那個序言，
% cue 2020 5089.070 5089.610
\hypertarget{w06:cue:2020}{}他有講，
% cue 2021 5089.610 5092.570
\hypertarget{w06:cue:2021}{}他的論證形式是兼具，
% cue 2022 5092.570 5097.460
\hypertarget{w06:cue:2022}{}既有結果有關的理由，
% cue 2023 5097.460 5100.250
\hypertarget{w06:cue:2023}{}也有程序有關的理由，
% cue 2024 5100.250 5102.140
\hypertarget{w06:cue:2024}{}那為什麼會是這樣呢，
% cue 2025 5102.140 5103.890
\hypertarget{w06:cue:2025}{}他用霍布斯對於講，
% cue 2026 5103.890 5106.570
\hypertarget{w06:cue:2026}{}霍布斯說我們因為有爭議，
% cue 2027 5106.570 5108.900
\hypertarget{w06:cue:2027}{}所以我們首先要去，
% cue 2028 5108.900 5112.270
\hypertarget{w06:cue:2028}{}一個解決爭議的程序，
% cue 2029 5112.270 5115.490
\hypertarget{w06:cue:2029}{}其次我們經過的，
% cue 2030 5115.490 5117.190
\hypertarget{w06:cue:2030}{}這個解決爭議程序的操作，
% cue 2031 5117.190 5119.780
\hypertarget{w06:cue:2031}{}我們希望它能解決爭議，
% cue 2032 5119.780 5122.030
\hypertarget{w06:cue:2032}{}而不是再度點燃爭議，
% cue 2033 5122.030 5124.730
\hypertarget{w06:cue:2033}{}後面那句是結果，
% cue 2034 5124.730 5126.430
\hypertarget{w06:cue:2034}{}前面那句是程序，
% cue 2035 5126.430 5130.080
\hypertarget{w06:cue:2035}{}所以既有結果，
% cue 2036 5130.080 5134.460
\hypertarget{w06:cue:2036}{}outcome-related,rational,reasons
% cue 2037 5134.460 5137.470
\hypertarget{w06:cue:2037}{}也有 process-related reasons。

\tc{1:25:37}
% cue 2038 5137.470 5140.130
\hypertarget{w06:cue:2038}{}好啦到這裡就好了，
% cue 2039 5140.130 5141.790
\hypertarget{w06:cue:2039}{}還需要再講了嗎？
% cue 2040 5141.790 5142.890
\hypertarget{w06:cue:2040}{}應該我不管啦，
% cue 2041 5142.890 5144.350
\hypertarget{w06:cue:2041}{}應該可以不必了吧，
% cue 2042 5144.350 5148.080
\hypertarget{w06:cue:2042}{}後面有很多，
% cue 2043 5148.080 5150.720
\hypertarget{w06:cue:2043}{}它還有很多那個介紹。

\tc{1:25:50}
% cue 2044 5150.720 5155.900
\hypertarget{w06:cue:2044}{}來我們看一下那個 1357，
% cue 2045 5155.900 5158.140
\hypertarget{w06:cue:2045}{}1375 說是後，
% cue 2046 5158.140 5160.570
\hypertarget{w06:cue:2046}{}What about the normative process
% cue 2047 5160.570 5163.260
\hypertarget{w06:cue:2047}{}about the process-related reasons
% cue 2048 5163.260 5167.330
\hypertarget{w06:cue:2048}{}剛才講的都是有關於 outcome reasons，
% cue 2049 5167.330 5168.370
\hypertarget{w06:cue:2049}{}然後再講一遍，
% cue 2050 5168.370 5169.670
\hypertarget{w06:cue:2050}{}第三七三頁，
% cue 2051 5169.670 5171.940
\hypertarget{w06:cue:2051}{}一段第一大段，
% cue 2052 5171.940 5172.780
\hypertarget{w06:cue:2052}{}有沒有？
% cue 2053 5172.780 5188.840
\hypertarget{w06:cue:2053}{}這叫做，
% cue 2054 5188.840 5194.150
\hypertarget{w06:cue:2054}{}這種這種這一種，
% cue 2055 5194.150 5197.030
\hypertarget{w06:cue:2055}{}這種頭腦，
% cue 2056 5197.030 5200.310
\hypertarget{w06:cue:2056}{}我們通常在中文的著作裡面是看的，
% cue 2057 5200.310 5202.050
\hypertarget{w06:cue:2057}{}用了腦袋沒有東西，
% cue 2058 5202.050 5205.020
\hypertarget{w06:cue:2058}{}連終翻譯用了，
% cue 2059 5205.020 5207.290
\hypertarget{w06:cue:2059}{}更不要說想到，
% cue 2060 5207.290 5208.050
\hypertarget{w06:cue:2060}{}還沒有寫下，
% cue 2061 5208.050 5210.110
\hypertarget{w06:cue:2061}{}現在只講第一半哦，
% cue 2062 5210.110 5211.170
\hypertarget{w06:cue:2062}{}還沒講完咯，
% cue 2063 5211.170 5213.040
\hypertarget{w06:cue:2063}{}是程序，
% cue 2064 5213.040 5215.220
\hypertarget{w06:cue:2064}{}而是結果有關的理由，
% cue 2065 5215.220 5218.200
\hypertarget{w06:cue:2065}{}結果有關的理由是指這個東西，
% cue 2066 5218.200 5220.390
\hypertarget{w06:cue:2066}{}比方說好的，
% cue 2067 5220.390 5221.750
\hypertarget{w06:cue:2067}{}公正的，
% cue 2068 5221.750 5223.570
\hypertarget{w06:cue:2068}{}正確的。

\tc{1:27:03}
% cue 2069 5223.570 5226.100
\hypertarget{w06:cue:2069}{}然後往下，

\tc{1:27:06}
\speaker{導讀人：亦凡}
% cue 2070 5226.100 5227.950
\hypertarget{w06:cue:2070}{}程序在前，

\tc{1:27:07}
\speaker{老師}
% cue 2071 5227.950 5229.310
\hypertarget{w06:cue:2071}{}1375 頁的，
% cue 2072 5229.310 5232.700
\hypertarget{w06:cue:2072}{}什麼叫做 process related reasons，
% cue 2073 5232.700 5233.360
\hypertarget{w06:cue:2073}{}是在這裡，
% cue 2074 5233.360 5235.280
\hypertarget{w06:cue:2074}{}聽懂嗎？

\tc{1:27:15}
% cue 2075 5235.280 5237.600
\hypertarget{w06:cue:2075}{}你要說什麼，
% cue 2076 5237.600 5238.000
\hypertarget{w06:cue:2076}{}請講，

\tc{1:27:18}
\speaker{導讀人：亦凡}
% cue 2077 5238.000 5241.430
\hypertarget{w06:cue:2077}{}程序的那一段在 1372 的最後面，

\tc{1:27:21}
\speaker{老師}
% cue 2078 5241.430 5243.360
\hypertarget{w06:cue:2078}{}1372 的最後面，
% cue 2079 5243.360 5243.680
\hypertarget{w06:cue:2079}{}好，
% cue 2080 5243.680 5245.890
\hypertarget{w06:cue:2080}{}1372 最後面它有講對，
% cue 2081 5245.890 5247.370
\hypertarget{w06:cue:2081}{}它現在這個區別，
% cue 2082 5247.370 5249.500
\hypertarget{w06:cue:2082}{}主要是在 1372 區分的，
% cue 2083 5249.500 5250.510
\hypertarget{w06:cue:2083}{}對，
% cue 2084 5250.510 5252.800
\hypertarget{w06:cue:2084}{}那 1372 講的是什麼？
% cue 2085 5252.800 5254.960
\hypertarget{w06:cue:2085}{}你如果不講這個，
% cue 2086 5254.960 5256.900
\hypertarget{w06:cue:2086}{}你如果不講 1375，
% cue 2087 5256.900 5258.120
\hypertarget{w06:cue:2087}{}也可以講 1372，
% cue 2088 5258.120 5259.260
\hypertarget{w06:cue:2088}{}但我猜想，
% cue 2089 5259.260 5261.410
\hypertarget{w06:cue:2089}{}應該是 1375 的詳細，
% cue 2090 5261.410 5263.590
\hypertarget{w06:cue:2090}{}對吧，
% cue 2091 5263.590 5264.070
\hypertarget{w06:cue:2091}{}對，

\tc{1:27:44}
\speaker{導讀人：亦凡}
% cue 2092 5264.070 5266.820
\hypertarget{w06:cue:2092}{}但是如果是後面的話，
% cue 2093 5266.820 5268.080
\hypertarget{w06:cue:2093}{}它會更展開講，

\tc{1:27:48}
\speaker{老師}
% cue 2094 5268.080 5268.520
\hypertarget{w06:cue:2094}{}好，
% cue 2095 5268.520 5270.180
\hypertarget{w06:cue:2095}{}那你先講 1372，
% cue 2096 5270.180 5274.810
\hypertarget{w06:cue:2096}{}它們兩個是相關的，
% cue 2097 5274.810 5275.370
\hypertarget{w06:cue:2097}{}沒錯，
% cue 2098 5275.370 5276.670
\hypertarget{w06:cue:2098}{}對。

\tc{1:27:56}
% cue 2099 5276.670 5277.470
\hypertarget{w06:cue:2099}{}然後呢，

\tc{1:27:57}
\speaker{導讀人：亦凡}
% cue 2100 5277.470 5278.310
\hypertarget{w06:cue:2100}{}1372

\tc{1:27:58}
\speaker{老師}
% cue 2101 5278.310 5279.770
\hypertarget{w06:cue:2101}{}你要講 1372 是什麼？

\tc{1:27:59}
\speaker{導讀人：亦凡}
% cue 2102 5279.770 5282.840
\hypertarget{w06:cue:2102}{}就是它就是簡略的，
% cue 2103 5282.840 5285.240
\hypertarget{w06:cue:2103}{}我們來講一下它後面，

\tc{1:28:05}
\speaker{老師}
% cue 2104 5285.240 5285.730
\hypertarget{w06:cue:2104}{}是 林遊戲講了，
% cue 2105 5285.730 5287.450
\hypertarget{w06:cue:2105}{}reasons for

\tc{1:28:07}
\speaker{導讀人：亦凡}
% cue 2106 5287.450 5292.470
\hypertarget{w06:cue:2106}{}reasons for insisting that some person make
% cue 2107 5292.470 5296.130
\hypertarget{w06:cue:2107}{}or participate in making a given decision
% cue 2108 5296.130 5300.290
\hypertarget{w06:cue:2108}{}that stands independently of considerations
% cue 2109 5300.290 5302.290
\hypertarget{w06:cue:2109}{}about the appropriate outcome。

\tc{1:28:22}
\speaker{老師}
% cue 2110 5302.290 5304.570
\hypertarget{w06:cue:2110}{}我們不管結論如何，
% cue 2111 5304.570 5309.340
\hypertarget{w06:cue:2111}{}我們要求的是某人他應該做成決定，
% cue 2112 5309.340 5313.160
\hypertarget{w06:cue:2112}{}或者某人應該參與做成決定，
% cue 2113 5313.160 5314.690
\hypertarget{w06:cue:2113}{}這就是程序。

\tc{1:28:34}
% cue 2114 5314.690 5319.320
\hypertarget{w06:cue:2114}{}比方說我在座的沒有經過你同意，
% cue 2115 5319.320 5322.670
\hypertarget{w06:cue:2115}{}說我們下禮拜不讓你來上課，
% cue 2116 5322.670 5324.180
\hypertarget{w06:cue:2116}{}你會很生氣啊，
% cue 2117 5324.180 5327.100
\hypertarget{w06:cue:2117}{}這個結果當然是會你不高興的，
% cue 2118 5327.100 5328.140
\hypertarget{w06:cue:2118}{}如果對你有利的結果，
% cue 2119 5328.140 5330.320
\hypertarget{w06:cue:2119}{}你可能不太在意說我有沒有在意，
% cue 2120 5330.320 5332.690
\hypertarget{w06:cue:2120}{}但是如果這對你，
% cue 2121 5332.690 5333.770
\hypertarget{w06:cue:2121}{}尤其是對你不利，
% cue 2122 5333.770 5334.770
\hypertarget{w06:cue:2122}{}你會馬上問一句，
% cue 2123 5334.770 5336.520
\hypertarget{w06:cue:2123}{}為什麼沒有通知我，
% cue 2124 5336.520 5337.880
\hypertarget{w06:cue:2124}{}為什麼沒有得我同意，
% cue 2125 5337.880 5339.320
\hypertarget{w06:cue:2125}{}為什麼沒有參與決定，
% cue 2126 5339.320 5340.780
\hypertarget{w06:cue:2126}{}為什麼把我排除在外，
% cue 2127 5340.780 5342.020
\hypertarget{w06:cue:2127}{}這就是這個意思。

\tc{1:29:02}
% cue 2128 5342.020 5345.200
\hypertarget{w06:cue:2128}{}那剛才講的跟結果有關的是說，
% cue 2129 5345.200 5347.980
\hypertarget{w06:cue:2129}{}這個決定做出來是正確的嗎？
% cue 2130 5347.980 5349.200
\hypertarget{w06:cue:2130}{}是好的嗎？
% cue 2131 5349.200 5350.280
\hypertarget{w06:cue:2131}{}是對的嗎？
% cue 2132 5350.280 5351.540
\hypertarget{w06:cue:2132}{}是這方面的事情，
% cue 2133 5351.540 5352.530
\hypertarget{w06:cue:2133}{}好了 講完了，
% cue 2134 5352.530 5353.300
\hypertarget{w06:cue:2134}{}不講了。

\section{結果相關理由}

\tc{1:29:13}
% cue 2135 5353.300 5353.480
\hypertarget{w06:cue:2135}{}來，
% cue 2136 5353.480 5356.160
\hypertarget{w06:cue:2136}{}201376 年第四大塊，
% cue 2137 5356.160 5360.380
\hypertarget{w06:cue:2137}{}Outcome related reasons
% cue 2138 5360.380 5361.890
\hypertarget{w06:cue:2138}{}來了，
% cue 2139 5361.890 5372.350
\hypertarget{w06:cue:2139}{}這裡有很多東西，
% cue 2140 5372.350 5374.480
\hypertarget{w06:cue:2140}{}真是蠻多的，
% cue 2141 5374.480 5376.960
\hypertarget{w06:cue:2141}{}Outcome
% cue 2142 5376.960 5379.080
\hypertarget{w06:cue:2142}{}跟結果有關的理由，
% cue 2143 5379.080 5380.970
\hypertarget{w06:cue:2143}{}According to Raz\footnote{Joseph Raz（1939–2022），法律實證主義代表學者，任教於牛津大學、哥倫比亞大學法學院與倫敦國王學院，2022 年 5 月 2 日於倫敦去世。Waldron 所引為 Raz 1998 年刊於 American Journal of Jurisprudence 第 43 卷之 Disagreement in Politics，此句引文在第 45 頁。見 Jeremy Waldron, The Core of the Case Against Judicial Review, 115 Yale L.J. 1346, 1376 (2006)。\url{https://www.kcl.ac.uk/news/professor-joseph-raz-1939-2022}。}
% cue 2144 5380.970 5381.450
\hypertarget{w06:cue:2144}{}來，

\tc{1:29:41}
\speaker{導讀人：亦凡}
% cue 2145 5381.450 5382.970
\hypertarget{w06:cue:2145}{}According to Raz
% cue 2146 5382.970 5386.720
\hypertarget{w06:cue:2146}{}A natural way to proceed
% cue 2147 5386.720 5388.080
\hypertarget{w06:cue:2147}{}is to assume that
% cue 2148 5388.080 5390.420
\hypertarget{w06:cue:2148}{}the enforcement of fundamental rights
% cue 2149 5390.420 5391.960
\hypertarget{w06:cue:2149}{}should be entrusted
% cue 2150 5391.960 5395.160
\hypertarget{w06:cue:2150}{}to whichever political decision procedure is
% cue 2151 5395.160 5401.070
\hypertarget{w06:cue:2151}{}in the circumstances of the time and place
% cue 2152 5401.070 5403.700
\hypertarget{w06:cue:2152}{}most likely to enforce them well
% cue 2153 5403.700 5409.270
\hypertarget{w06:cue:2153}{}with the fewest adverse side effects。

\tc{1:30:09}
\speaker{老師}
% cue 2154 5409.270 5410.850
\hypertarget{w06:cue:2154}{}好，講這句就可以，
% cue 2155 5410.850 5414.200
\hypertarget{w06:cue:2155}{}根據這個 Raz 的見解，
% cue 2156 5414.200 5417.440
\hypertarget{w06:cue:2156}{}一個非常自然的方式，
% cue 2157 5417.440 5419.080
\hypertarget{w06:cue:2157}{}就是去假定，
% cue 2158 5419.080 5425.670
\hypertarget{w06:cue:2158}{}假定基本權的執行，
% cue 2159 5425.670 5428.080
\hypertarget{w06:cue:2159}{}誰來執行基本權，
% cue 2160 5428.080 5435.910
\hypertarget{w06:cue:2160}{}應該要把它委託給哪一個政治決策程序，
% cue 2161 5435.910 5440.580
\hypertarget{w06:cue:2161}{}在特定的時間跟地點下，
% cue 2162 5440.580 5445.670
\hypertarget{w06:cue:2162}{}哪一個政治決策程序，
% cue 2163 5445.670 5450.550
\hypertarget{w06:cue:2163}{}最可能把它們執行好，
% cue 2164 5450.550 5453.670
\hypertarget{w06:cue:2164}{}會產生最少的副作用，
% cue 2165 5453.670 5455.940
\hypertarget{w06:cue:2165}{}那你就把它委託給他。

\tc{1:30:55}
% cue 2166 5455.940 5457.940
\hypertarget{w06:cue:2166}{}這句話的意思說，
% cue 2167 5457.940 5465.250
\hypertarget{w06:cue:2167}{}如果立法者他就能夠最好的執行基本權，
% cue 2168 5465.250 5467.170
\hypertarget{w06:cue:2168}{}那你就委託給他，
% cue 2169 5467.170 5468.760
\hypertarget{w06:cue:2169}{}那反過來，
% cue 2170 5468.760 5473.350
\hypertarget{w06:cue:2170}{}如果是司法者最能執行，
% cue 2171 5473.350 5475.190
\hypertarget{w06:cue:2171}{}落實基本權，
% cue 2172 5475.190 5477.990
\hypertarget{w06:cue:2172}{}而且最不產生副作用，
% cue 2173 5477.990 5479.630
\hypertarget{w06:cue:2173}{}那你就委託給司法者，
% cue 2174 5479.630 5481.370
\hypertarget{w06:cue:2174}{}所以他的標準是中性，
% cue 2175 5481.370 5484.250
\hypertarget{w06:cue:2175}{}我們是講結果，
% cue 2176 5484.250 5486.620
\hypertarget{w06:cue:2176}{}誰最能執行，
% cue 2177 5486.620 5488.640
\hypertarget{w06:cue:2177}{}是不是這樣子，
% cue 2178 5488.640 5492.160
\hypertarget{w06:cue:2178}{}most likely to enforce them well
% cue 2179 5492.160 5495.410
\hypertarget{w06:cue:2179}{}誰最有辦法落實基本權，
% cue 2180 5495.410 5497.050
\hypertarget{w06:cue:2180}{}那我們就委託給他，
% cue 2181 5497.050 5498.830
\hypertarget{w06:cue:2181}{}如果是法院證的落實，
% cue 2182 5498.830 5500.070
\hypertarget{w06:cue:2182}{}那為什麼不讓法院省，
% cue 2183 5500.070 5501.990
\hypertarget{w06:cue:2183}{}那如果是立法者，
% cue 2184 5501.990 5502.970
\hypertarget{w06:cue:2184}{}那你為什麼不讓法院，
% cue 2185 5502.970 5504.010
\hypertarget{w06:cue:2185}{}我就聽不懂。

\tc{1:31:44}
% cue 2186 5504.010 5505.780
\hypertarget{w06:cue:2186}{}然後往下，
% cue 2187 5505.780 5508.650
\hypertarget{w06:cue:2187}{}I guess the discussion at this point
% cue 2188 5508.650 5511.850
\hypertarget{w06:cue:2188}{}ought to be continuous with the broader debate
% cue 2189 5511.850 5514.590
\hypertarget{w06:cue:2189}{}about the institutional competence\footnote{英文，「機關能力」，指各機關依其組織與程序適於決定何種問題。Waldron 說此一論辯由法律過程學派（legal process school）開啟，註 78 引 Henry M. Hart, Jr. 與 Albert M. Sacks 合著 The Legal Process 第 640–647 頁（William N. Eskridge, Jr. 與 Philip P. Frickey 編，1994 年版）。見 Jeremy Waldron, The Core of the Case Against Judicial Review, 115 Yale L.J. 1346, 1376 (2006)。}
% cue 2190 5514.590 5516.450
\hypertarget{w06:cue:2190}{}這個常常看到，
% cue 2191 5516.450 5517.700
\hypertarget{w06:cue:2191}{}什麼意思？
% cue 2192 5517.700 5520.540
\hypertarget{w06:cue:2192}{}法院有這個能力嗎？
% cue 2193 5520.540 5523.980
\hypertarget{w06:cue:2193}{}國會有這個能力嗎？
% cue 2194 5523.980 5526.170
\hypertarget{w06:cue:2194}{}這個是不是牽涉到，
% cue 2195 5526.170 5530.030
\hypertarget{w06:cue:2195}{}法院法系德國人喜歡講的，
% cue 2196 5530.030 5531.970
\hypertarget{w06:cue:2196}{}功能最適理論\footnote{中華民國憲法法庭 111 年憲判字第 9 號判決（2022 年 6 月 24 日）理由第 20 段闡釋權力分立原則，謂將不同性質之權力「分別歸屬於組織及程序等功能最適之不同機關行使，以發揮功效；原則上並禁止同一機關行使兩種以上權力，以避免權力集中」。\url{https://cons.judicial.gov.tw/docdata.aspx?fid=38&id=339757}。}，
% cue 2197 5531.970 5534.850
\hypertarget{w06:cue:2197}{}這是不是機關功能的問題，
% cue 2198 5534.850 5538.550
\hypertarget{w06:cue:2198}{}我想不然你們聽不懂，
% cue 2199 5538.550 5540.500
\hypertarget{w06:cue:2199}{}competence
% cue 2200 5540.500 5543.100
\hypertarget{w06:cue:2200}{}哪一個機關比較有能力，
% cue 2201 5543.100 5546.210
\hypertarget{w06:cue:2201}{}那這裡是 institutional competence，
% cue 2202 5546.210 5547.750
\hypertarget{w06:cue:2202}{}是機關能力的問題。

\tc{1:32:27}
% cue 2203 5547.750 5549.920
\hypertarget{w06:cue:2203}{}那現在講到你要把它交給法院，
% cue 2204 5549.920 5552.140
\hypertarget{w06:cue:2204}{}那我們一般是不是討論到，
% cue 2205 5552.140 5554.040
\hypertarget{w06:cue:2205}{}法院有能力做違憲審查，
% cue 2206 5554.040 5556.710
\hypertarget{w06:cue:2206}{}法院有能力落實基本權嗎？
% cue 2207 5556.710 5559.430
\hypertarget{w06:cue:2207}{}大家有聽懂嗎？

\tc{1:32:39}
% cue 2208 5559.430 5560.610
\hypertarget{w06:cue:2208}{}OK 好繼續，
% cue 2209 5560.610 5561.690
\hypertarget{w06:cue:2209}{}Courts

\tc{1:32:41}
\speaker{導讀人：亦凡}
% cue 2210 5561.690 5566.750
\hypertarget{w06:cue:2210}{}Courts are good at deciding some issues and not others。

\tc{1:32:46}
\speaker{老師}
% cue 2211 5566.750 5569.390
\hypertarget{w06:cue:2211}{}好但是這個作者的意見不一樣，
% cue 2212 5569.390 5573.810
\hypertarget{w06:cue:2212}{}每一個機關它的功能都是有限的，
% cue 2213 5573.810 5576.580
\hypertarget{w06:cue:2213}{}你們去看老師的《釋憲學思心路\footnote{湯德宗著，上、中、下冊，二版，2026 年，課綱所列四部指定教科書之一。}》的見解，
% cue 2214 5576.580 5578.030
\hypertarget{w06:cue:2214}{}基本上也是這樣的，
% cue 2215 5578.030 5580.630
\hypertarget{w06:cue:2215}{}我從來不認為司法院是萬能，
% cue 2216 5580.630 5583.700
\hypertarget{w06:cue:2216}{}那所以我們要考慮，
% cue 2217 5583.700 5585.540
\hypertarget{w06:cue:2217}{}每一個個案是我們要考慮，
% cue 2218 5585.540 5591.340
\hypertarget{w06:cue:2218}{}那這裡他講說法院在解決某些問題很棒，
% cue 2219 5591.340 5594.560
\hypertarget{w06:cue:2219}{}但是不是解決所有的問題都很棒，
% cue 2220 5594.560 5596.680
\hypertarget{w06:cue:2220}{}對某些問題他毫無能為力，
% cue 2221 5596.680 5597.200
\hypertarget{w06:cue:2221}{}OK。

\tc{1:33:17}
% cue 2222 5597.200 5599.670
\hypertarget{w06:cue:2222}{}Technically 技術上講，
% cue 2223 5599.670 5601.830
\hypertarget{w06:cue:2223}{}We use the term right
% cue 2224 5601.830 5604.030
\hypertarget{w06:cue:2224}{}我們用權利這個術語，
% cue 2225 5604.030 5605.160
\hypertarget{w06:cue:2225}{}to denote
% cue 2226 5605.160 5607.610
\hypertarget{w06:cue:2226}{}去指涉，
% cue 2227 5607.610 5609.200
\hypertarget{w06:cue:2227}{}指涉什麼，
% cue 2228 5609.200 5614.070
\hypertarget{w06:cue:2228}{}用權利這個字來指涉什麼，

\tc{1:33:34}
\speaker{導讀人：亦凡}
% cue 2229 5614.070 5617.870
\hypertarget{w06:cue:2229}{}指涉法院通常會決定的問題，

\tc{1:33:37}
\speaker{老師}
% cue 2230 5617.870 5623.290
\hypertarget{w06:cue:2230}{}可以吧，
% cue 2231 5623.290 5624.340
\hypertarget{w06:cue:2231}{}你的意思？

\tc{1:33:44}
% cue 2232 5624.340 5625.320
\hypertarget{w06:cue:2232}{}你的問題是，
% cue 2233 5625.320 5629.550
\hypertarget{w06:cue:2233}{}你的這個困難是出在 characteristically decide，
% cue 2234 5629.550 5631.190
\hypertarget{w06:cue:2234}{}你把它翻成通常決定，
% cue 2235 5631.190 5634.860
\hypertarget{w06:cue:2235}{}因為，
% cue 2236 5634.860 5636.440
\hypertarget{w06:cue:2236}{}因為繼續繼續，
% cue 2237 5636.440 5637.280
\hypertarget{w06:cue:2237}{}你不要看我，
% cue 2238 5637.280 5638.220
\hypertarget{w06:cue:2238}{}慢慢往下看，
% cue 2239 5638.220 5639.460
\hypertarget{w06:cue:2239}{}繼續因為，
% cue 2240 5639.460 5640.040
\hypertarget{w06:cue:2240}{}because

\tc{1:34:00}
\speaker{導讀人：亦凡}
% cue 2241 5640.040 5647.350
\hypertarget{w06:cue:2241}{}because a plaintiff has to state a claim of right to be heard in a court at all

\tc{1:34:07}
\speaker{老師}
% cue 2242 5647.350 5651.210
\hypertarget{w06:cue:2242}{}原告他認為他有一個權利的主張，
% cue 2243 5651.210 5654.180
\hypertarget{w06:cue:2243}{}應該在法院裡面被廳審，
% cue 2244 5654.180 5657.890
\hypertarget{w06:cue:2244}{}But, but，繼續，

\tc{1:34:17}
\speaker{導讀人：亦凡}
% cue 2245 5657.890 5673.180
\hypertarget{w06:cue:2245}{}But as Lon Fuller\footnote{Lon L. Fuller（1902–1978），美國法理學者，1940 年至 1978 年任教於哈佛法學院，1948 年至 1972 年任 Carter 一般法理學講座教授。Waldron 所引為 Fuller, The Forms and Limits of Adjudication, 92 Harv. L. Rev. 353 (1978) 第 368–370 頁，據以說明部分權利主張具兩造二元之性質，部分則具多面向性質（multifaceted），向來被認為不宜由司法結構決定。見 Jeremy Waldron, The Core of the Case Against Judicial Review, 115 Yale L.J. 1346, 1376 (2006)。\url{https://hls.harvard.edu/historical-faculty/lon-luvois-fuller-1902-1978}。} observed, it does not follow that courts are therefore the appropriate forum for dealing with claims of rights in the less technical sense under consideration here。

\tc{1:34:33}
\speaker{老師}
% cue 2246 5673.180 5678.880
\hypertarget{w06:cue:2246}{}但是 Lon Fuller 好像不是當然這麼認為，
% cue 2247 5678.880 5681.540
\hypertarget{w06:cue:2247}{}你可能會說權利，
% cue 2248 5681.540 5683.360
\hypertarget{w06:cue:2248}{}權利就是指的呢，
% cue 2249 5683.360 5686.810
\hypertarget{w06:cue:2249}{}原告有權利提出他的主張，
% cue 2250 5686.810 5690.020
\hypertarget{w06:cue:2250}{}然後讓法院來聽審，
% cue 2251 5690.020 5692.260
\hypertarget{w06:cue:2251}{}這就是我們為什麼講的權利，
% cue 2252 5692.260 5695.190
\hypertarget{w06:cue:2252}{}但是如果照了觀察呢，
% cue 2253 5695.190 5697.730
\hypertarget{w06:cue:2253}{}他說也許是這樣子，
% cue 2254 5697.730 5699.980
\hypertarget{w06:cue:2254}{}但是我們不能說，
% cue 2255 5699.980 5709.530
\hypertarget{w06:cue:2255}{}那法院就一定是有關於權利主張適當場所，
% cue 2256 5709.530 5711.730
\hypertarget{w06:cue:2256}{}換句話講你要主張權利，
% cue 2257 5711.730 5715.070
\hypertarget{w06:cue:2257}{}除了法院之外也有可能別的地方，
% cue 2258 5715.070 5717.740
\hypertarget{w06:cue:2258}{}不是說只有紀念。

\tc{1:35:17}
% cue 2259 5717.740 5724.200
\hypertarget{w06:cue:2259}{}Some claims of right have the character of the sort of binary issue
% cue 2260 5724.200 5726.160
\hypertarget{w06:cue:2260}{}聽得懂嗎？
% cue 2261 5726.160 5730.410
\hypertarget{w06:cue:2261}{}有一些權利的主張是具有雙面性，
% cue 2262 5730.410 5733.240
\hypertarget{w06:cue:2262}{}就是兩造當事人，
% cue 2263 5733.240 5738.030
\hypertarget{w06:cue:2263}{}有一些是具有 multifaceted，
% cue 2264 5738.030 5741.260
\hypertarget{w06:cue:2264}{}多面性多方的，
% cue 2265 5741.260 5744.410
\hypertarget{w06:cue:2265}{}前面兩造當事人，
% cue 2266 5744.410 5746.260
\hypertarget{w06:cue:2266}{}那你可以說 A 跟 B，
% cue 2267 5746.260 5747.920
\hypertarget{w06:cue:2267}{}誰有權利誰贏誰輸，
% cue 2268 5747.920 5749.540
\hypertarget{w06:cue:2268}{}這是債權人債務人，
% cue 2269 5749.540 5751.230
\hypertarget{w06:cue:2269}{}哪一個是債權人是債務人，
% cue 2270 5751.230 5753.000
\hypertarget{w06:cue:2270}{}但是你多方的時候呢，
% cue 2271 5753.000 5756.240
\hypertarget{w06:cue:2271}{}這個時候法院就不是很有能力。

\tc{1:35:56}
% cue 2272 5756.240 5758.670
\hypertarget{w06:cue:2272}{}那你先去想想看，
% cue 2273 5758.670 5760.690
\hypertarget{w06:cue:2273}{}我們講到基本人權的時候，
% cue 2274 5760.690 5763.270
\hypertarget{w06:cue:2274}{}他真正的只是原告跟被告兩方，
% cue 2275 5763.270 5766.440
\hypertarget{w06:cue:2275}{}那如果是多方的時候，
% cue 2276 5766.440 5768.350
\hypertarget{w06:cue:2276}{}法院你還覺得你很有能力嗎？
% cue 2277 5768.350 5770.910
\hypertarget{w06:cue:2277}{}在這裡面他已經暗示了好幾遍，
% cue 2278 5770.910 5773.380
\hypertarget{w06:cue:2278}{}法院能是比較適合，
% cue 2279 5773.380 5774.680
\hypertarget{w06:cue:2279}{}我覺得我不能說，
% cue 2280 5774.680 5776.010
\hypertarget{w06:cue:2280}{}我沒有看到明文，
% cue 2281 5776.010 5777.920
\hypertarget{w06:cue:2281}{}但是我感覺他已經在暗示你，
% cue 2282 5777.920 5779.120
\hypertarget{w06:cue:2282}{}法院說的是，
% cue 2283 5779.120 5781.120
\hypertarget{w06:cue:2283}{}法院適合做的事情是，
% cue 2284 5781.120 5782.510
\hypertarget{w06:cue:2284}{}Binary。

\tc{1:36:22}
% cue 2285 5782.510 5785.700
\hypertarget{w06:cue:2285}{}然後呢，上面剛才前面還有一個字，
% cue 2286 5785.700 5792.640
\hypertarget{w06:cue:2286}{}Advisor,Adversarial,Adversarial，兩造爭訟式，
% cue 2287 5792.640 5795.120
\hypertarget{w06:cue:2287}{}剛才那個字你念錯，我有糾正你在哪裡？
% cue 2288 5795.120 5800.640
\hypertarget{w06:cue:2288}{}Adversarial，剛才你念錯，我有糾正你在哪裡？
% cue 2289 5800.640 5809.170
\hypertarget{w06:cue:2289}{}這道嗎，我們知道太多了，
% cue 2290 5809.170 5813.300
\hypertarget{w06:cue:2290}{}好了，我找到了，今天請我示範，
% cue 2291 5813.300 5819.440
\hypertarget{w06:cue:2291}{}第 363 頁，告訴第三方，
% cue 2292 5819.440 5823.930
\hypertarget{w06:cue:2292}{}這是不是一再出現，
% cue 2293 5823.930 5827.740
\hypertarget{w06:cue:2293}{}Binary,Adversarial,Presentation
% cue 2294 5827.740 5830.490
\hypertarget{w06:cue:2294}{}法院是不是都是這樣子，
% cue 2295 5830.490 5832.070
\hypertarget{w06:cue:2295}{}兩造爭訟式，
% cue 2296 5832.070 5834.330
\hypertarget{w06:cue:2296}{}兩造又是爭訟的，
% cue 2297 5834.330 5835.750
\hypertarget{w06:cue:2297}{}這種比較適合法院，
% cue 2298 5835.750 5837.500
\hypertarget{w06:cue:2298}{}那如果是多面向的，
% cue 2299 5837.500 5838.330
\hypertarget{w06:cue:2299}{}是多造的，
% cue 2300 5838.330 5840.230
\hypertarget{w06:cue:2300}{}那你叫法院怎麼辦，
% cue 2301 5840.230 5842.590
\hypertarget{w06:cue:2301}{}你們自己去想想看。

\tc{1:37:22}
% cue 2302 5842.590 5843.790
\hypertarget{w06:cue:2302}{}我們在臺灣講，
% cue 2303 5843.790 5846.230
\hypertarget{w06:cue:2303}{}大陸法系的話會說，
% cue 2304 5846.230 5847.750
\hypertarget{w06:cue:2304}{}哎呀功能最適合理論，
% cue 2305 5847.750 5849.990
\hypertarget{w06:cue:2305}{}這個功能最適合哪個機構，
% cue 2306 5849.990 5852.870
\hypertarget{w06:cue:2306}{}可是我們常常很少去討論，
% cue 2307 5852.870 5856.380
\hypertarget{w06:cue:2307}{}那你的功能最適合定義是什麼？
% cue 2308 5856.380 5857.710
\hypertarget{w06:cue:2308}{}法院適合做什麼，
% cue 2309 5857.710 5859.190
\hypertarget{w06:cue:2309}{}立法者適合做什麼，
% cue 2310 5859.190 5860.830
\hypertarget{w06:cue:2310}{}其他部門適合做什麼，
% cue 2311 5860.830 5863.650
\hypertarget{w06:cue:2311}{}講話有聽懂嗎？

\tc{1:37:43}
% cue 2312 5863.650 5871.700
\hypertarget{w06:cue:2312}{}我不知道我沒怎麼講，我的意思是，我沒有否定誰講對誰說的錯，
% cue 2313 5871.700 5877.430
\hypertarget{w06:cue:2313}{}我只是覺得，你們應該要在比較憲法的這些課程活動裡面，
% cue 2314 5877.430 5884.560
\hypertarget{w06:cue:2314}{}去把不同地方學到的東西，有意識的，有意識的，我想到它們把它串組起來，
% cue 2315 5884.560 5889.450
\hypertarget{w06:cue:2315}{}然後你們將來用了之後，會比較容易融會貫通。

\tc{1:38:09}
% cue 2316 5889.450 5892.220
\hypertarget{w06:cue:2316}{}好，來，It's tempting，下一段，

\tc{1:38:12}
\speaker{導讀人：亦凡}
% cue 2317 5892.220 5901.950
\hypertarget{w06:cue:2317}{}It is tempting to associate outcome-related reasons with a case for judicial review

\tc{1:38:21}
\speaker{老師}
% cue 2318 5901.950 5903.510
\hypertarget{w06:cue:2318}{}好像有這樣的傾向，
% cue 2319 5903.510 5907.590
\hypertarget{w06:cue:2319}{}我們會把跟結果相關的理由，
% cue 2320 5907.590 5910.400
\hypertarget{w06:cue:2320}{}把它跟違憲審查在一起，
% cue 2321 5910.400 5911.920
\hypertarget{w06:cue:2321}{}然後呢，
% cue 2322 5911.920 5913.980
\hypertarget{w06:cue:2322}{}說跟程序相關的呢，
% cue 2323 5913.980 5915.580
\hypertarget{w06:cue:2323}{}我們就反對違憲審查，
% cue 2324 5915.580 5916.660
\hypertarget{w06:cue:2324}{}他說呢，
% cue 2325 5916.660 5917.860
\hypertarget{w06:cue:2325}{}This is a mistake
% cue 2326 5917.860 5918.560
\hypertarget{w06:cue:2326}{}不對，

\tc{1:38:38}
\speaker{導讀人：亦凡}
% cue 2327 5918.560 5922.590
\hypertarget{w06:cue:2327}{}Oh, this is a mistake
% cue 2328 5922.590 5928.590
\hypertarget{w06:cue:2328}{}It is true that many of the more important process-related reasons
% cue 2329 5928.590 5934.630
\hypertarget{w06:cue:2329}{}are participatory and therefore favor
% cue 2330 5934.630 5939.330
\hypertarget{w06:cue:2330}{}and therefore favor elective or representative institutions。

\tc{1:38:59}
\speaker{老師}
% cue 2331 5939.330 5942.720
\hypertarget{w06:cue:2331}{}因為要講因為程序上的理由，
% cue 2332 5942.720 5946.150
\hypertarget{w06:cue:2332}{}通常我們強調的是我要參加我要參加，
% cue 2333 5946.150 5951.610
\hypertarget{w06:cue:2333}{}那你會把它跟所謂的選舉和代議機構連在一起，
% cue 2334 5951.610 5958.280
\hypertarget{w06:cue:2334}{}所以你通常會以司法因為只限於當事人，
% cue 2335 5958.280 5960.730
\hypertarget{w06:cue:2335}{}所以他沒有辦法讓大家廣泛的參與，
% cue 2336 5960.730 5963.360
\hypertarget{w06:cue:2336}{}會用這個理由去反對違憲審查，
% cue 2337 5963.360 5965.890
\hypertarget{w06:cue:2337}{}但是從結果的觀點來看，
% cue 2338 5965.890 5969.550
\hypertarget{w06:cue:2338}{}你會認為結果有關的應該是支持司法審查，
% cue 2339 5969.550 5972.000
\hypertarget{w06:cue:2339}{}說這個是錯的。

\tc{1:39:32}
% cue 2340 5972.000 5972.580
\hypertarget{w06:cue:2340}{}but

\tc{1:39:32}
\speaker{導讀人：亦凡}
% cue 2341 5972.580 5983.230
\hypertarget{w06:cue:2341}{}but it does not follow that all or most outcome related reasons argue the other way

\tc{1:39:43}
\speaker{老師}
% cue 2342 5983.230 5984.450
\hypertarget{w06:cue:2342}{}你不能反過來講，
% cue 2343 5984.450 5985.680
\hypertarget{w06:cue:2343}{}再加一句，
% cue 2344 5985.680 5989.650
\hypertarget{w06:cue:2344}{}outcome related reason cut in both directions
% cue 2345 5989.650 5992.240
\hypertarget{w06:cue:2345}{}它兩個方向它都有，
% cue 2346 5992.240 5996.830
\hypertarget{w06:cue:2346}{}結果相關的理由而言，
% cue 2347 5996.830 6000.220
\hypertarget{w06:cue:2347}{}它有時支持違憲審查，
% cue 2348 6000.220 6002.910
\hypertarget{w06:cue:2348}{}有時反對違憲審查，
% cue 2349 6002.910 6004.920
\hypertarget{w06:cue:2349}{}不能假設所，
% cue 2350 6004.920 6006.600
\hypertarget{w06:cue:2350}{}這個一定是支持的，
% cue 2351 6006.600 6007.740
\hypertarget{w06:cue:2351}{}這個一定是反對的，
% cue 2352 6007.740 6008.660
\hypertarget{w06:cue:2352}{}不是這樣的。

\tc{1:40:08}
% cue 2353 6008.660 6009.820
\hypertarget{w06:cue:2353}{}下面他舉了很多，
% cue 2354 6009.820 6011.460
\hypertarget{w06:cue:2354}{}Raz 講的一些話，
% cue 2355 6011.460 6013.080
\hypertarget{w06:cue:2355}{}講的一些什麼理由，
% cue 2356 6013.080 6014.200
\hypertarget{w06:cue:2356}{}我們來看一下，
% cue 2357 6014.200 6015.080
\hypertarget{w06:cue:2357}{}你喜歡講案子，
% cue 2358 6015.080 6016.240
\hypertarget{w06:cue:2358}{}我們講案子 1377，
% cue 2359 6016.240 6017.660
\hypertarget{w06:cue:2359}{}這裡有好幾個案子，
% cue 2360 6017.660 6020.130
\hypertarget{w06:cue:2360}{}是嗎？

\tc{1:40:20}
\speaker{導讀人：亦凡}
% cue 2361 6020.130 6025.480
\hypertarget{w06:cue:2361}{}就是他那個腳註，
% cue 2362 6025.480 6027.130
\hypertarget{w06:cue:2362}{}81 82

\tc{1:40:27}
\speaker{老師}
% cue 2363 6027.130 6029.560
\hypertarget{w06:cue:2363}{}可能 82 就可以，

\tc{1:40:29}
\speaker{導讀人：亦凡}
% cue 2364 6029.560 6030.300
\hypertarget{w06:cue:2364}{}82

\tc{1:40:30}
\speaker{老師}
% cue 2365 6030.300 6032.190
\hypertarget{w06:cue:2365}{}The same
% cue 2366 6032.190 6035.370
\hypertarget{w06:cue:2366}{}Sectarian pressure
% cue 2367 6035.370 6037.150
\hypertarget{w06:cue:2367}{}often explain
% cue 2368 6037.150 6041.800
\hypertarget{w06:cue:2368}{}假如你認為，
% cue 2369 6041.800 6045.170
\hypertarget{w06:cue:2369}{}立法者他不能免俗，
% cue 2370 6045.170 6047.380
\hypertarget{w06:cue:2370}{}會受到很多社會的壓力，
% cue 2371 6047.380 6048.870
\hypertarget{w06:cue:2371}{}那你認為，
% cue 2372 6048.870 6051.280
\hypertarget{w06:cue:2372}{}立法者就可以免俗嗎？

\tc{1:40:51}
% cue 2373 6051.280 6054.040
\hypertarget{w06:cue:2373}{}那剛才下面他講了，
% cue 2374 6054.040 6054.630
\hypertarget{w06:cue:2374}{}Korematsu\footnote{Korematsu v. United States, 323 U.S. 214，聯邦最高法院 1944 年 12 月 18 日判決，Black 大法官主筆，維持一名日裔美國公民違反 Civilian Exclusion Order No. 34 之有罪判決；該命令要求日裔人士於 1942 年 5 月 9 日後離開加州 San Leandro 所在之軍事區。聯邦最高法院 2018 年 6 月 26 日 Trump v. Hawaii 判決之多數意見稱 Korematsu 案「was gravely wrong the day it was decided」。\url{https://www.law.cornell.edu/supremecourt/text/323/214}。}
% cue 2375 6054.630 6056.520
\hypertarget{w06:cue:2375}{}這個是有名的案子，
% cue 2376 6056.520 6058.190
\hypertarget{w06:cue:2376}{}在小註 82 裡面，
% cue 2377 6058.190 6058.710
\hypertarget{w06:cue:2377}{}說什麼，
% cue 2378 6058.710 6060.770
\hypertarget{w06:cue:2378}{}Korematsu vs. US
% cue 2379 6060.770 6063.170
\hypertarget{w06:cue:2379}{}說什麼，
% cue 2380 6063.170 6065.250
\hypertarget{w06:cue:2380}{}講掛鼓，

\tc{1:41:05}
\speaker{導讀人：亦凡}
% cue 2381 6065.250 6067.080
\hypertarget{w06:cue:2381}{}就是，

\tc{1:41:07}
\speaker{老師}
% cue 2382 6067.080 6069.840
\hypertarget{w06:cue:2382}{}拒絕，

\tc{1:41:09}
\speaker{導讀人：亦凡}
% cue 2383 6069.840 6081.220
\hypertarget{w06:cue:2383}{}Refusing to protect citizens of Japanese

\tc{1:41:21}
\speaker{老師}
% cue 2384 6081.220 6082.360
\hypertarget{w06:cue:2384}{}Japanese
% cue 2385 6082.360 6088.890
\hypertarget{w06:cue:2385}{}日本人的後代。

\tc{1:41:28}
% cue 2386 6088.890 6094.470
\hypertarget{w06:cue:2386}{}你們知道，
% cue 2387 6094.470 6096.760
\hypertarget{w06:cue:2387}{}二次世界大戰剛開始，
% cue 2388 6096.760 6098.900
\hypertarget{w06:cue:2388}{}尤其是珍珠港被偷襲了以後，
% cue 2389 6098.900 6101.710
\hypertarget{w06:cue:2389}{}美國人對於日裔的美國人都非常懷疑，
% cue 2390 6101.710 6104.470
\hypertarget{w06:cue:2390}{}因為這些日本人效忠於他的電話，
% cue 2391 6104.470 6106.490
\hypertarget{w06:cue:2391}{}就把他們通通抓去關起來，
% cue 2392 6106.490 6109.140
\hypertarget{w06:cue:2392}{}是另外的美國式的集中營，
% cue 2393 6109.140 6113.080
\hypertarget{w06:cue:2393}{}這個事情一直到好久好久好久以後，
% cue 2394 6113.080 6114.100
\hypertarget{w06:cue:2394}{}才有黑場。

\tc{1:41:54}
% cue 2395 6114.100 6122.950
\hypertarget{w06:cue:2395}{}那你剛剛不是說司法是不會受到壓力的嗎？
% cue 2396 6122.950 6126.210
\hypertarget{w06:cue:2396}{}司法是可以不受到世俗壓力的嗎？
% cue 2397 6126.210 6128.620
\hypertarget{w06:cue:2397}{}只有立法者因為有選人的壓力，
% cue 2398 6128.620 6132.460
\hypertarget{w06:cue:2398}{}所以立法者會做一些優俗或者免俗的決定嗎？
% cue 2399 6132.460 6133.420
\hypertarget{w06:cue:2399}{}不會啊，
% cue 2400 6133.420 6134.660
\hypertarget{w06:cue:2400}{}那你跟我解釋一下，
% cue 2401 6134.660 6136.930
\hypertarget{w06:cue:2401}{}Korematsu怎麼把這些人都關掉那麼久，
% cue 2402 6136.930 6140.340
\hypertarget{w06:cue:2402}{}關到二次大戰截後我們把它放掉，
% cue 2403 6140.340 6141.910
\hypertarget{w06:cue:2403}{}聽懂沒有聽懂，
% cue 2404 6141.910 6144.810
\hypertarget{w06:cue:2404}{}美國有拍一些片子講他們。

\tc{1:42:24}
% cue 2405 6144.810 6146.490
\hypertarget{w06:cue:2405}{}然後 Schenck\footnote{Schenck v. United States, 249 U.S. 47，聯邦最高法院 1919 年 3 月 3 日判決，Holmes 大法官主筆，提出「明顯而立即的危險」（clear and present danger）標準；判決理由謂「The most stringent protection of free speech would not protect a man in falsely shouting fire in a theatre and causing a panic.」Waldron 於註 82 概括其意旨為：第一次世界大戰期間批評徵兵，有如在擁擠戲院中呼喊失火。見 Jeremy Waldron, The Core of the Case Against Judicial Review, 115 Yale L.J. 1346, 1377 (2006)。\url{https://www.law.cornell.edu/supremecourt/text/249/47}。} 繼續，
% cue 2406 6146.490 6149.770
\hypertarget{w06:cue:2406}{}第二名案子 Fox or Schenck versus Lewis，

\tc{1:42:29}
\speaker{導讀人：亦凡}
% cue 2407 6149.770 6152.960
\hypertarget{w06:cue:2407}{}holding that criticising

\tc{1:42:32}
\speaker{老師}
% cue 2408 6152.960 6153.900
\hypertarget{w06:cue:2408}{}conscription

\tc{1:42:33}
\speaker{導讀人：亦凡}
% cue 2409 6153.900 6155.010
\hypertarget{w06:cue:2409}{}conscription

\tc{1:42:35}
\speaker{老師}
% cue 2410 6155.010 6156.380
\hypertarget{w06:cue:2410}{}什麼叫 conscription，
% cue 2411 6156.380 6157.500
\hypertarget{w06:cue:2411}{}徵兵制，
% cue 2412 6157.500 6159.970
\hypertarget{w06:cue:2412}{}繼續，

\tc{1:42:39}
\speaker{導讀人：亦凡}
% cue 2413 6159.970 6161.610
\hypertarget{w06:cue:2413}{}during the first world war
% cue 2414 6161.610 6166.000
\hypertarget{w06:cue:2414}{}was like shouting a fire in a crowded theater。

\tc{1:42:46}
\speaker{老師}
% cue 2415 6166.000 6168.890
\hypertarget{w06:cue:2415}{}這個 這一句話是我覺得寫的最好，
% cue 2416 6168.890 6171.810
\hypertarget{w06:cue:2416}{}什麼意思 有沒有人可以解釋一下，
% cue 2417 6171.810 6174.500
\hypertarget{w06:cue:2417}{}這個人可以 這個人答得出來，
% cue 2418 6174.500 6176.540
\hypertarget{w06:cue:2418}{}表示你的憲法知識 非常好，
% cue 2419 6176.540 6179.390
\hypertarget{w06:cue:2419}{}知識 非常好，
% cue 2420 6179.390 6185.700
\hypertarget{w06:cue:2420}{}這個 美國有一個很有名的大法官叫 Holmes\footnote{Oliver Wendell Holmes, Jr.（1841–1935），1902 年 12 月 8 日至 1932 年 1 月 12 日任聯邦最高法院大法官。\url{https://www.fjc.gov/history/judges/holmes-oliver-wendell-jr}。}，
% cue 2421 6185.700 6190.790
\hypertarget{w06:cue:2421}{}他說言論自由保障的人有發言的權力，
% cue 2422 6190.790 6192.490
\hypertarget{w06:cue:2422}{}但是不是條件，
% cue 2423 6192.490 6199.390
\hypertarget{w06:cue:2423}{}請問人可以在戲院就是昏暗的情況下，
% cue 2424 6199.390 6202.270
\hypertarget{w06:cue:2424}{}謊稱失火嗎？
% cue 2425 6202.270 6203.450
\hypertarget{w06:cue:2425}{}可以這樣嗎？
% cue 2426 6203.450 6205.760
\hypertarget{w06:cue:2426}{}所以當然不可以啊。

\tc{1:43:25}
% cue 2427 6205.760 6207.800
\hypertarget{w06:cue:2427}{}所以呢，
% cue 2428 6207.800 6211.810
\hypertarget{w06:cue:2428}{}所以你在一次世界大戰期間，
% cue 2429 6211.810 6213.970
\hypertarget{w06:cue:2429}{}你可以批評徵兵嗎？
% cue 2430 6213.970 6216.920
\hypertarget{w06:cue:2430}{}一次世界大戰批評徵兵，
% cue 2431 6216.920 6220.750
\hypertarget{w06:cue:2431}{}就相當於一個人在黑暗的戲院當中，
% cue 2432 6220.750 6221.890
\hypertarget{w06:cue:2432}{}謊稱失火一樣，
% cue 2433 6221.890 6224.110
\hypertarget{w06:cue:2433}{}這不是言論自由保障的範圍，
% cue 2434 6224.110 6226.320
\hypertarget{w06:cue:2434}{}聽起來是不是很不小？

\tc{1:43:46}
% cue 2435 6226.320 6229.920
\hypertarget{w06:cue:2435}{}誒，這好像事情也會發生啊，
% cue 2436 6229.920 6234.330
\hypertarget{w06:cue:2436}{}最近一個照片啊，掐了那個總統的什麼就把他趕出去，
% cue 2437 6234.330 6239.850
\hypertarget{w06:cue:2437}{}那他也沒有言論自由啊，聽得懂嗎？
% cue 2438 6239.850 6242.720
\hypertarget{w06:cue:2438}{}言論自由包含。

\tc{1:44:02}
% cue 2439 6242.720 6247.420
\hypertarget{w06:cue:2439}{}Dred Scott\footnote{Dred Scott v. Sandford, 60 U.S. 393，聯邦最高法院 1857 年 3 月 6 日判決，Taney 首席大法官主筆，認定祖先以奴隸身分被輸入美國之非洲裔人「are not included, and were not intended to be included, under the word 'citizens' in the Constitution」，不得主張憲法賦予美國公民之權利。\url{https://www.archives.gov/milestone-documents/dred-scott-v-sandford}。}，來，
% cue 2440 6247.420 6256.040
\hypertarget{w06:cue:2440}{}這兩個人一起的，來繼續，
% cue 2441 6256.040 6257.890
\hypertarget{w06:cue:2441}{}說什麼，說什麼，

\tc{1:44:17}
\speaker{導讀人：亦凡}
% cue 2442 6257.890 6263.610
\hypertarget{w06:cue:2442}{}簡單來講就是認為他們不是美國人啦，

\tc{1:44:23}
\speaker{老師}
% cue 2443 6263.610 6269.010
\hypertarget{w06:cue:2443}{}他們不受到這個憲法上人的保護，
% cue 2444 6269.010 6277.610
\hypertarget{w06:cue:2444}{}抓了他們送回給他們的歷史財產，
% cue 2445 6277.610 6282.670
\hypertarget{w06:cue:2445}{}這幾案子都是美國憲法上的嚴重點。

\tc{1:44:42}
% cue 2446 6282.670 6283.300
\hypertarget{w06:cue:2446}{}好啦，
% cue 2447 6283.300 6285.880
\hypertarget{w06:cue:2447}{}他舉這個例子意思已經講完了，
% cue 2448 6285.880 6291.080
\hypertarget{w06:cue:2448}{}你們不是說從結果觀點來講，
% cue 2449 6291.080 6292.310
\hypertarget{w06:cue:2449}{}並比較好嗎？
% cue 2450 6292.310 6295.200
\hypertarget{w06:cue:2450}{}因為必法則不能免俗啊，
% cue 2451 6295.200 6296.300
\hypertarget{w06:cue:2451}{}有壓力啊，
% cue 2452 6296.300 6298.160
\hypertarget{w06:cue:2452}{}或向社會投降啊，
% cue 2453 6298.160 6301.720
\hypertarget{w06:cue:2453}{}那你現在回答我這幾個最高法院的案子，
% cue 2454 6301.720 6302.860
\hypertarget{w06:cue:2454}{}是不是也有壓力，
% cue 2455 6302.860 6304.510
\hypertarget{w06:cue:2455}{}是不是下次會跟我講，
% cue 2456 6304.510 6306.540
\hypertarget{w06:cue:2456}{}為什麼要做這樣的局勢。

\tc{1:45:06}
% cue 2457 6306.540 6308.630
\hypertarget{w06:cue:2457}{}好，
% cue 2458 6308.630 6312.300
\hypertarget{w06:cue:2458}{}第三七八頁，
% cue 2459 6312.300 6315.560
\hypertarget{w06:cue:2459}{}有沒有要講的，

\tc{1:45:15}
\speaker{導讀人：亦凡}
% cue 2460 6315.560 6317.240
\hypertarget{w06:cue:2460}{}我後面更多，

\tc{1:45:17}
\speaker{老師}
% cue 2461 6317.240 6318.000
\hypertarget{w06:cue:2461}{}好，
% cue 2462 6318.000 6318.940
\hypertarget{w06:cue:2462}{}隨便你隨便你，
% cue 2463 6318.940 6320.670
\hypertarget{w06:cue:2463}{}下面多你講來，
% cue 2464 6320.670 6322.270
\hypertarget{w06:cue:2464}{}那為什麼有一二三四，
% cue 2465 6322.270 6326.840
\hypertarget{w06:cue:2465}{}三七九頁講什麼，
% cue 2466 6326.840 6329.430
\hypertarget{w06:cue:2466}{}你如果不再多回答我的問題，
% cue 2467 6329.430 6331.250
\hypertarget{w06:cue:2467}{}大概都跳完了。

\tc{1:45:31}
% cue 2468 6331.250 6333.800
\hypertarget{w06:cue:2468}{}一三七九，

\tc{1:45:33}
\speaker{導讀人：亦凡}
% cue 2469 6333.800 6336.640
\hypertarget{w06:cue:2469}{}三，
% cue 2470 6336.640 6339.290
\hypertarget{w06:cue:2470}{}謝謝他就是，

\tc{1:45:39}
\speaker{老師}
% cue 2471 6339.290 6341.040
\hypertarget{w06:cue:2471}{}他說有三個，

\tc{1:45:41}
\speaker{導讀人：亦凡}
% cue 2472 6341.040 6342.120
\hypertarget{w06:cue:2472}{}有三個優點，

\tc{1:45:42}
\speaker{老師}
% cue 2473 6342.120 6343.330
\hypertarget{w06:cue:2473}{}優點，
% cue 2474 6343.330 6344.940
\hypertarget{w06:cue:2474}{}有的時候呢，

\tc{1:45:44}
\speaker{導讀人：亦凡}
% cue 2475 6344.940 6346.770
\hypertarget{w06:cue:2475}{}有的時候，

\tc{1:45:46}
\speaker{老師}
% cue 2476 6346.770 6348.040
\hypertarget{w06:cue:2476}{}claimed for courts
% cue 2477 6348.040 6349.310
\hypertarget{w06:cue:2477}{}有人認為，
% cue 2478 6349.310 6352.530
\hypertarget{w06:cue:2478}{}有法院來進行違憲審查，
% cue 2479 6352.530 6354.500
\hypertarget{w06:cue:2479}{}會具有下面三個優點，
% cue 2480 6354.500 6355.420
\hypertarget{w06:cue:2480}{}第一。

\tc{1:45:56}
\speaker{導讀人：亦凡}
% cue 2481 6356.650 6367.020
\hypertarget{w06:cue:2481}{}A. That issues of rights are presented to courts in the context of specific cases。

\tc{1:46:07}
\speaker{老師}
% cue 2482 6367.020 6372.670
\hypertarget{w06:cue:2482}{}B. 權利的爭點送到法院的呈現，
% cue 2483 6372.670 6374.290
\hypertarget{w06:cue:2483}{}呈現到法院的時候，
% cue 2484 6374.290 6377.910
\hypertarget{w06:cue:2484}{}是在一個具體個案的情境下送進去。

\tc{1:46:17}
% cue 2485 6377.910 6380.310
\hypertarget{w06:cue:2485}{}所以下面的 1393 的 A，
% cue 2486 6380.310 6383.960
\hypertarget{w06:cue:2486}{}叫做 Orientation to particular cases，
% cue 2487 6383.960 6387.070
\hypertarget{w06:cue:2487}{}我把它翻成叫做個案導向，
% cue 2488 6387.070 6388.510
\hypertarget{w06:cue:2488}{}導向於個案，
% cue 2489 6388.510 6390.600
\hypertarget{w06:cue:2489}{}聽得懂嗎？
% cue 2490 6390.600 6393.720
\hypertarget{w06:cue:2490}{}原告告進來是不是有個個案？

\tc{1:46:33}
% cue 2491 6393.720 6395.820
\hypertarget{w06:cue:2491}{}臺灣我已經講過了，
% cue 2492 6395.820 6397.240
\hypertarget{w06:cue:2492}{}因為我們是抽象違憲審查，
% cue 2493 6397.240 6398.590
\hypertarget{w06:cue:2493}{}一個機關來申請，
% cue 2494 6398.590 6399.810
\hypertarget{w06:cue:2494}{}他也是有個案，
% cue 2495 6399.810 6401.210
\hypertarget{w06:cue:2495}{}但是他不限於這個個案，
% cue 2496 6401.210 6404.050
\hypertarget{w06:cue:2496}{}他問的也不是這個法條在個案的適用，
% cue 2497 6404.050 6406.610
\hypertarget{w06:cue:2497}{}他要的是你把這個法條整個宣告違憲，
% cue 2498 6406.610 6407.870
\hypertarget{w06:cue:2498}{}這是另外一件事，
% cue 2499 6407.870 6411.360
\hypertarget{w06:cue:2499}{}但是在美國你一定是一個個案才能進來，
% cue 2500 6411.360 6414.400
\hypertarget{w06:cue:2500}{}否則就不是我司法權的範圍，
% cue 2501 6414.400 6416.190
\hypertarget{w06:cue:2501}{}所以這當然是個案導向。

\tc{1:46:56}
% cue 2502 6416.190 6419.050
\hypertarget{w06:cue:2502}{}那這個他認為是對法院主張，
% cue 2503 6419.050 6420.990
\hypertarget{w06:cue:2503}{}我們應該給法院違憲審查，
% cue 2504 6420.990 6422.960
\hypertarget{w06:cue:2504}{}因為我們從結果來看，
% cue 2505 6422.960 6436.250
\hypertarget{w06:cue:2505}{}現在我們是講結果嘛，這個是好的，那為什麼結果是好的呢，看 1380 頁，它對這裡看你去看，有沒有，1380 頁 B 的前面，
% cue 2506 6436.250 6463.720
\hypertarget{w06:cue:2506}{}你們認為個案導向是比較適合由違憲審查機關、司法機關來進行的。

\tc{1:47:43}
% cue 2507 6463.720 6466.620
\hypertarget{w06:cue:2507}{}但是如果我們從權力的觀點來看，
% cue 2508 6466.620 6468.850
\hypertarget{w06:cue:2508}{}我們要知道這個權利的爭議，
% cue 2509 6468.850 6470.430
\hypertarget{w06:cue:2509}{}那這個爭議要怎麼消除，
% cue 2510 6470.430 6472.570
\hypertarget{w06:cue:2510}{}你不是只要考慮個案兩個人，
% cue 2511 6472.570 6474.090
\hypertarget{w06:cue:2511}{}要考慮到全面，
% cue 2512 6474.090 6474.810
\hypertarget{w06:cue:2512}{}所以他認為，
% cue 2513 6474.810 6476.100
\hypertarget{w06:cue:2513}{}你這個個案導向，
% cue 2514 6476.100 6477.060
\hypertarget{w06:cue:2514}{}其實不是優點。

\tc{1:47:57}
% cue 2515 6477.060 6477.650
\hypertarget{w06:cue:2515}{}好，
% cue 2516 6477.650 6478.810
\hypertarget{w06:cue:2516}{}B 第二個優點，
% cue 2517 6478.810 6481.440
\hypertarget{w06:cue:2517}{}對於違憲審查來講，
% cue 2518 6481.440 6482.580
\hypertarget{w06:cue:2518}{}你喜歡違憲審查，
% cue 2519 6482.580 6485.570
\hypertarget{w06:cue:2519}{}第二個結果上的理由。

\tc{1:48:05}
\speaker{導讀人：亦凡}
% cue 2520 6485.570 6491.660
\hypertarget{w06:cue:2520}{}人權清單導向。

\tc{1:48:11}
\speaker{老師}
% cue 2521 6491.660 6493.280
\hypertarget{w06:cue:2521}{}這句話什麼意思？
% cue 2522 6493.280 6494.960
\hypertarget{w06:cue:2522}{}用講的就好像是很複雜。

\tc{1:48:14}
\speaker{導讀人：亦凡}
% cue 2523 6494.960 6498.440
\hypertarget{w06:cue:2523}{}就是他覺得人權清單，
% cue 2524 6498.440 6501.100
\hypertarget{w06:cue:2524}{}是用書面去。

\tc{1:48:21}
\speaker{老師}
% cue 2525 6501.100 6504.780
\hypertarget{w06:cue:2525}{}你現在要講的是 1381 嗎？
% cue 2526 6504.780 6507.390
\hypertarget{w06:cue:2526}{}The written formulation of
% cue 2527 6507.390 6508.670
\hypertarget{w06:cue:2527}{}The help
% cue 2528 6508.670 6510.170
\hypertarget{w06:cue:2528}{}你想要講這個嗎？
% cue 2529 6510.170 6511.580
\hypertarget{w06:cue:2529}{}你為什麼要講書面？
% cue 2530 6511.580 6512.400
\hypertarget{w06:cue:2530}{}是這個意思嗎？
% cue 2531 6512.400 6512.820
\hypertarget{w06:cue:2531}{}對，
% cue 2532 6512.820 6514.320
\hypertarget{w06:cue:2532}{}對啊。

\tc{1:48:34}
% cue 2533 6514.320 6515.460
\hypertarget{w06:cue:2533}{}好來看一下，
% cue 2534 6515.460 6517.280
\hypertarget{w06:cue:2534}{}One reason for thinking
% cue 2535 6517.280 6518.520
\hypertarget{w06:cue:2535}{}It is a good idea
% cue 2536 6518.520 6520.100
\hypertarget{w06:cue:2536}{}為什麼 is that，
% cue 2537 6520.100 6522.500
\hypertarget{w06:cue:2537}{}The written formulation of
% cue 2538 6522.500 6524.940
\hypertarget{w06:cue:2538}{}The real right
% cue 2539 6524.940 6527.690
\hypertarget{w06:cue:2539}{}Can help focus on
% cue 2540 6527.690 6530.280
\hypertarget{w06:cue:2540}{}Abstract right issue
% cue 2541 6530.280 6531.720
\hypertarget{w06:cue:2541}{}把這幾個字圈起來，
% cue 2542 6531.720 6535.310
\hypertarget{w06:cue:2542}{}抽象的權利爭點，
% cue 2543 6535.310 6537.200
\hypertarget{w06:cue:2543}{}At stake
% cue 2544 6537.200 6546.690
\hypertarget{w06:cue:2544}{}就是一個被文字固定下來的權利法案。

\tc{1:49:06}
\speaker{導讀人：亦凡}
% cue 2545 6546.690 6548.930
\hypertarget{w06:cue:2545}{}它就是已經是一個抽象者。

\tc{1:49:08}
\speaker{老師}
% cue 2546 6548.930 6550.910
\hypertarget{w06:cue:2546}{}可以有助於爭議的人。

\tc{1:49:10}
\speaker{導讀人：亦凡}
% cue 2547 6550.910 6553.170
\hypertarget{w06:cue:2547}{}可以有助於爭議的人，
% cue 2548 6553.170 6557.490
\hypertarget{w06:cue:2548}{}關注到那些抽象的權利問題。

\tc{1:49:17}
\speaker{老師}
% cue 2549 6557.490 6559.500
\hypertarget{w06:cue:2549}{}對這個意思，
% cue 2550 6559.500 6562.650
\hypertarget{w06:cue:2550}{}我們不是要講這個權利有還是沒有，
% cue 2551 6562.650 6564.710
\hypertarget{w06:cue:2551}{}這個權利的範圍邊界等等，
% cue 2552 6564.710 6566.160
\hypertarget{w06:cue:2552}{}這就是抽象的問題，
% cue 2553 6566.160 6567.280
\hypertarget{w06:cue:2553}{}這不是個案的問題，
% cue 2554 6567.280 6568.590
\hypertarget{w06:cue:2554}{}不是個別的問題。

\tc{1:49:28}
% cue 2555 6568.590 6572.010
\hypertarget{w06:cue:2555}{}but there are powerful reasons on the other side
% cue 2556 6572.010 6574.540
\hypertarget{w06:cue:2556}{}但是另外一面也有情有念力，
% cue 2557 6574.540 6578.710
\hypertarget{w06:cue:2557}{}所以看起來好像也不是這樣子，
% cue 2558 6578.710 6582.390
\hypertarget{w06:cue:2558}{}那 1382 頁那裡面有講到，
% cue 2559 6582.390 6585.640
\hypertarget{w06:cue:2559}{}你把它放到人權清單，
% cue 2560 6585.640 6587.870
\hypertarget{w06:cue:2560}{}那這時候還會出現的問題是，
% cue 2561 6587.870 6589.410
\hypertarget{w06:cue:2561}{}人權清單裡面的權利，
% cue 2562 6589.410 6592.420
\hypertarget{w06:cue:2562}{}有可能 1382 頁第四行，
% cue 2563 6592.420 6596.830
\hypertarget{w06:cue:2563}{}是 positive rights\footnote{英文，積極權利，指要求國家作為或給付之權利，如社會經濟權；negative rights 為消極權利，指要求國家不加干預之自由權。Waldron 以社會對權利法案應否納入積極權利之分歧為例，說明權利法案之文本因其未列入者，亦可能改變法官對已列入之權利的理解，例如加重財產權與契約自由之分量。見 Jeremy Waldron, The Core of the Case Against Judicial Review, 115 Yale L.J. 1346, 1382 (2006)。} 和 negative rights 的問題，
% cue 2564 6596.830 6599.350
\hypertarget{w06:cue:2564}{}positive rights 你可以把它分成經濟權，
% cue 2565 6599.350 6601.340
\hypertarget{w06:cue:2565}{}也就是所謂的經濟社會權，
% cue 2566 6601.340 6603.760
\hypertarget{w06:cue:2566}{}要求國家要給我什麼東西，
% cue 2567 6603.760 6606.360
\hypertarget{w06:cue:2567}{}做什麼東西 拿什麼東西給我，
% cue 2568 6606.360 6609.980
\hypertarget{w06:cue:2568}{}他們自由權要求國家不要干涉，
% cue 2569 6609.980 6612.360
\hypertarget{w06:cue:2569}{}那這兩個權利，
% cue 2570 6612.360 6616.080
\hypertarget{w06:cue:2570}{}你認為當我們把他交給法院去解釋會，
% cue 2571 6616.080 6618.240
\hypertarget{w06:cue:2571}{}你認為支持違憲審查的時候，
% cue 2572 6618.240 6621.100
\hypertarget{w06:cue:2572}{}他的結果或者他的強度會是這樣的嗎？

\tc{1:50:21}
% cue 2573 6621.100 6622.570
\hypertarget{w06:cue:2573}{}好 第三個，
% cue 2574 6622.570 6623.650
\hypertarget{w06:cue:2574}{}13821
% cue 2575 6623.650 6625.180
\hypertarget{w06:cue:2575}{}支持違憲審查，
% cue 2576 6625.180 6627.000
\hypertarget{w06:cue:2576}{}因為法院呢。

\tc{1:50:27}
\speaker{導讀人：亦凡}
% cue 2577 6627.000 6628.520
\hypertarget{w06:cue:2577}{}和說明理由。

\tc{1:50:28}
\speaker{老師}
% cue 2578 6628.520 6629.900
\hypertarget{w06:cue:2578}{}Stating reasons
% cue 2579 6629.900 6631.270
\hypertarget{w06:cue:2579}{}然後呢，
% cue 2580 6631.270 6633.990
\hypertarget{w06:cue:2580}{}作者認為法院是有說明理由，
% cue 2581 6633.990 6634.510
\hypertarget{w06:cue:2581}{}不過，
% cue 2582 6634.510 6635.870
\hypertarget{w06:cue:2582}{}不過。

\tc{1:50:35}
\speaker{導讀人：亦凡}
% cue 2583 6635.870 6638.130
\hypertarget{w06:cue:2583}{}那個說明的品質，
% cue 2584 6638.130 6642.170
\hypertarget{w06:cue:2584}{}可能就是他覺得會有不同的，
% cue 2585 6642.170 6646.070
\hypertarget{w06:cue:2585}{}跟理想上有差距嗎？

\tc{1:50:46}
\speaker{老師}
% cue 2586 6646.070 6648.610
\hypertarget{w06:cue:2586}{}他的意思是說，
% cue 2587 6648.610 6651.370
\hypertarget{w06:cue:2587}{}這句話真的是比較難解釋的，
% cue 2588 6651.370 6653.010
\hypertarget{w06:cue:2588}{}你要找到句子來給我們講，
% cue 2589 6653.010 6654.450
\hypertarget{w06:cue:2589}{}你要想講哪裡。

\tc{1:50:54}
\speaker{導讀人：亦凡}
% cue 2590 6654.450 6662.080
\hypertarget{w06:cue:2590}{}基本上法院是有說明的理由。

\tc{1:51:02}
\speaker{老師}
% cue 2591 6662.080 6664.220
\hypertarget{w06:cue:2591}{}但是法院說明的理由，
% cue 2592 6664.220 6670.440
\hypertarget{w06:cue:2592}{}有助於解決我們大家對於權利的爭議或者分析嗎？

\tc{1:51:10}
\speaker{導讀人：亦凡}
% cue 2593 6670.440 6676.810
\hypertarget{w06:cue:2593}{}如果看那個 1383 的最後，
% cue 2594 6676.810 6692.450
\hypertarget{w06:cue:2594}{}但是法院會花很多的力去討論先例。

\tc{1:51:32}
\speaker{老師}
% cue 2595 6692.450 6696.300
\hypertarget{w06:cue:2595}{}然後這裡還會牽涉到解釋方法，
% cue 2596 6696.300 6700.570
\hypertarget{w06:cue:2596}{}這些枝枝節佔掉了大部分的篇幅，
% cue 2597 6700.570 6703.970
\hypertarget{w06:cue:2597}{}而真正這個權利的本質，
% cue 2598 6703.970 6708.140
\hypertarget{w06:cue:2598}{}或者是對關於這個權利的抽象的問題點，
% cue 2599 6708.140 6711.910
\hypertarget{w06:cue:2599}{}法院卻往有非常少的篇幅。

\tc{1:51:51}
% cue 2600 6711.910 6714.950
\hypertarget{w06:cue:2600}{}他舉的那個例子就是 Roe vs Wade\footnote{Roe v. Wade, 410 U.S. 113，聯邦最高法院 1973 年 1 月 22 日判決，認定德州刑法之墮胎罪規定違反第十四修正案正當程序條款所保障之隱私權；2022 年 Dobbs v. Jackson Women's Health Organization, 597 U.S. 215 判決將之推翻。Waldron 指出該案五十頁之意見中，討論生育權相對於隱私之道德重要性者僅有兩三段（第 153–155 頁）。見 Jeremy Waldron, The Core of the Case Against Judicial Review, 115 Yale L.J. 1346, 1383 (2006)。\url{https://www.law.cornell.edu/supremecourt/text/410/113}。}，
% cue 2601 6714.950 6717.340
\hypertarget{w06:cue:2601}{}你看這句話了嗎？
% cue 2602 6717.340 6718.180
\hypertarget{w06:cue:2602}{}看到了嗎？

\tc{1:51:58}
\speaker{導讀人：亦凡}
% cue 2603 6718.180 6719.920
\hypertarget{w06:cue:2603}{}他就比較 Roe vs Wade。

\tc{1:51:59}
\speaker{老師}
% cue 2604 6719.920 6722.840
\hypertarget{w06:cue:2604}{}他說在法院最高法院的判決裡面，
% cue 2605 6722.840 6727.200
\hypertarget{w06:cue:2605}{}這個一共有 50 頁的意見書裡，
% cue 2606 6727.200 6730.730
\hypertarget{w06:cue:2606}{}那只有幾個段落，
% cue 2607 6730.730 6737.200
\hypertarget{w06:cue:2607}{}討論到在生殖的權利的道德上的重要性，
% cue 2608 6737.200 6739.680
\hypertarget{w06:cue:2608}{}相對於隱私權，
% cue 2609 6739.680 6742.690
\hypertarget{w06:cue:2609}{}因為這個 Roe vs Wade 是墮胎嘛，
% cue 2610 6742.690 6748.750
\hypertarget{w06:cue:2610}{}那這是不是生殖的重要性 道德上的重要性，
% cue 2611 6748.750 6752.320
\hypertarget{w06:cue:2611}{}跟隱私 這哪個 哪個重哪個 哪個輕。

\tc{1:52:32}
% cue 2612 6752.320 6756.480
\hypertarget{w06:cue:2612}{}很像最近有一次八月老師在高雄的演講影片出來，
% cue 2613 6756.480 6760.370
\hypertarget{w06:cue:2613}{}就是在講 這個算是基本權的衝突，
% cue 2614 6760.370 6766.430
\hypertarget{w06:cue:2614}{}這裡有隱私權 然後有生命權 或者是生殖這樣的頻率，
% cue 2615 6766.430 6770.530
\hypertarget{w06:cue:2615}{}這個怎麼比 這個案子的重點其實是在這裡，
% cue 2616 6770.530 6773.950
\hypertarget{w06:cue:2616}{}但是法院只花了很少的篇幅講這個。

\tc{1:52:53}
% cue 2617 6773.950 6776.480
\hypertarget{w06:cue:2617}{}然後，
% cue 2618 6776.480 6784.490
\hypertarget{w06:cue:2618}{}結論也許你贊成，
% cue 2619 6784.490 6787.340
\hypertarget{w06:cue:2619}{}可是論證其實是很蒼白的，
% cue 2620 6787.340 6789.700
\hypertarget{w06:cue:2620}{}很有限的。

\tc{1:53:09}
% cue 2621 6789.700 6791.520
\hypertarget{w06:cue:2621}{}所以簡單的講就是，
% cue 2622 6791.520 6792.960
\hypertarget{w06:cue:2622}{}法院會給你理，
% cue 2623 6792.960 6794.810
\hypertarget{w06:cue:2623}{}但是法院給的理由，
% cue 2624 6794.810 6796.570
\hypertarget{w06:cue:2624}{}因為種種的考慮，
% cue 2625 6796.570 6798.310
\hypertarget{w06:cue:2625}{}包含法院認為，
% cue 2626 6798.310 6799.820
\hypertarget{w06:cue:2626}{}它的正當性不足，
% cue 2627 6799.820 6801.850
\hypertarget{w06:cue:2627}{}願意跟你公開的討論，
% cue 2628 6801.850 6803.070
\hypertarget{w06:cue:2628}{}這個問題的生死，
% cue 2629 6803.070 6805.630
\hypertarget{w06:cue:2629}{}它只能跟你講，
% cue 2630 6805.630 6806.990
\hypertarget{w06:cue:2630}{}我是哪裡得到的授權，
% cue 2631 6806.990 6809.070
\hypertarget{w06:cue:2631}{}我們以前的先例是怎麼樣，
% cue 2632 6809.070 6810.880
\hypertarget{w06:cue:2632}{}所以我們應該是這樣的解釋，
% cue 2633 6810.880 6814.650
\hypertarget{w06:cue:2633}{}他不能向地法則說，
% cue 2634 6814.650 6815.630
\hypertarget{w06:cue:2634}{}我就是反對，
% cue 2635 6815.630 6816.950
\hypertarget{w06:cue:2635}{}我為什麼什麼什麼反對，
% cue 2636 6816.950 6817.870
\hypertarget{w06:cue:2636}{}我就是贊成，
% cue 2637 6817.870 6819.230
\hypertarget{w06:cue:2637}{}我為什麼什麼什麼贊成，
% cue 2638 6819.230 6820.690
\hypertarget{w06:cue:2638}{}他不是這樣的，
% cue 2639 6820.690 6822.530
\hypertarget{w06:cue:2639}{}所以作者不認為，
% cue 2640 6822.530 6826.330
\hypertarget{w06:cue:2640}{}他會覺得交給法院有效。

\tc{1:53:46}
% cue 2641 6826.330 6828.650
\hypertarget{w06:cue:2641}{}你看哪一段 1386 頁，
% cue 2642 6828.650 6834.630
\hypertarget{w06:cue:2642}{}在回應那個 Robert，
% cue 2643 6834.630 6836.370
\hypertarget{w06:cue:2643}{}那就是 Ronald Dworkin\footnote{Ronald Dworkin（1931–2013），美國法律哲學家，1975 年起任教於紐約大學法學院。法院為「原則的論壇」（forum of principle）之說出自其 A Matter of Principle（1985），Waldron 註 108 引該書第 33、69–71 頁。見 Jeremy Waldron, The Core of the Case Against Judicial Review, 115 Yale L.J. 1346, 1385 (2006)。\url{https://law.nyu.edu/news/RONALD_DWORKIN_MEMORIAM}。}，
% cue 2644 6836.370 6838.000
\hypertarget{w06:cue:2644}{}他有一句名言，
% cue 2645 6838.000 6840.030
\hypertarget{w06:cue:2645}{}他說那法院呢，
% cue 2646 6840.030 6844.360
\hypertarget{w06:cue:2646}{}是最是一個 forum of principle，
% cue 2647 6844.360 6847.380
\hypertarget{w06:cue:2647}{}是一個講原則的地點，
% cue 2648 6847.380 6849.360
\hypertarget{w06:cue:2648}{}找到那句話嗎？
% cue 2649 6849.360 6855.120
\hypertarget{w06:cue:2649}{}找到那個 Dworkin 的 phrase 嗎？
% cue 2650 6855.120 6873.210
\hypertarget{w06:cue:2650}{}第三號我們說法官應該要做違憲審查，
% cue 2651 6873.210 6876.790
\hypertarget{w06:cue:2651}{}那有沒有墮胎權應該要回到法院去，
% cue 2652 6876.790 6885.410
\hypertarget{w06:cue:2652}{}因為法院應該被視為是講原則的地方，
% cue 2653 6885.410 6888.340
\hypertarget{w06:cue:2653}{}原則的場所，
% cue 2654 6888.340 6890.960
\hypertarget{w06:cue:2654}{}這是 Ronald Dworkin 的名言。

\tc{1:54:50}
% cue 2655 6890.960 6893.410
\hypertarget{w06:cue:2655}{}針對這句話，
% cue 2656 6893.410 6896.590
\hypertarget{w06:cue:2656}{}那現在他對比的是，
% cue 2657 6896.590 6900.960
\hypertarget{w06:cue:2657}{}一個是美國聯邦最高法院說墮胎怎樣怎樣，
% cue 2658 6900.960 6905.000
\hypertarget{w06:cue:2658}{}另外一個是沒有像美國這種違憲審查制度的英國，
% cue 2659 6905.000 6907.620
\hypertarget{w06:cue:2659}{}他們在國會裡面針對墮胎的辯論，
% cue 2660 6907.620 6910.110
\hypertarget{w06:cue:2660}{}你覺得誰的說理比較清楚？
% cue 2661 6910.110 6911.580
\hypertarget{w06:cue:2661}{}哪一個說理比較充分？
% cue 2662 6911.580 6915.650
\hypertarget{w06:cue:2662}{}講的是這個對比，
% cue 2663 6915.650 6917.500
\hypertarget{w06:cue:2663}{}有趣嗎？

\tc{1:55:17}
% cue 2664 6917.500 6920.990
\hypertarget{w06:cue:2664}{}然後最後是 1386 頁他所做的結論，
% cue 2665 6920.990 6922.710
\hypertarget{w06:cue:2665}{}有沒有，
% cue 2666 6922.710 6925.530
\hypertarget{w06:cue:2666}{}這個正常的數過來，
% cue 2667 6925.530 6927.070
\hypertarget{w06:cue:2667}{}12345678 左右，
% cue 2668 6927.070 6928.190
\hypertarget{w06:cue:2668}{}Despite
% cue 2669 6928.190 6930.830
\hypertarget{w06:cue:2669}{}儘管 Ronald Dworkin，
% cue 2670 6930.830 6933.770
\hypertarget{w06:cue:2670}{}他說 Forum of Principles 是法院，
% cue 2671 6933.770 6940.360
\hypertarget{w06:cue:2671}{}I think courts are expected to behave in the ways that I have criticized
% cue 2672 6940.360 6946.270
\hypertarget{w06:cue:2672}{}focusing on precedent text, the doctrine, and the other legalism
% cue 2673 6946.270 6949.290
\hypertarget{w06:cue:2673}{}法院的論證都是這種論證，
% cue 2674 6949.290 6955.340
\hypertarget{w06:cue:2674}{}法院很少針對問題核心抽象的來討論說，
% cue 2675 6955.340 6957.060
\hypertarget{w06:cue:2675}{}墮胎是好的還是壞的，
% cue 2676 6957.060 6958.820
\hypertarget{w06:cue:2676}{}我們贊成還是我們反對，
% cue 2677 6958.820 6963.020
\hypertarget{w06:cue:2677}{}因為法院會覺得我們沒有這個基礎講這個話，
% cue 2678 6963.020 6964.610
\hypertarget{w06:cue:2678}{}我們也不是民意代表，
% cue 2679 6964.610 6966.310
\hypertarget{w06:cue:2679}{}這種事情不是哲學家，
% cue 2680 6966.310 6967.470
\hypertarget{w06:cue:2680}{}我不是法官，
% cue 2681 6967.470 6969.780
\hypertarget{w06:cue:2681}{}能不能聽懂了，
% cue 2682 6969.780 6970.900
\hypertarget{w06:cue:2682}{}好了 這裡可以了。

\section{程序面理由（第伍部分）}

\tc{1:56:10}
% cue 2683 6970.900 6971.920
\hypertarget{w06:cue:2683}{}來 第五大段，
% cue 2684 6971.920 6973.440
\hypertarget{w06:cue:2684}{}Process related cases
% cue 2685 6973.440 6975.140
\hypertarget{w06:cue:2685}{}我要加速了，
% cue 2686 6975.140 6976.040
\hypertarget{w06:cue:2686}{}來 快，
% cue 2687 6976.040 6977.830
\hypertarget{w06:cue:2687}{}這裡是講什麼？

\tc{1:56:17}
\speaker{導讀人：亦凡}
% cue 2688 6977.830 6979.430
\hypertarget{w06:cue:2688}{}就是程序面。

\tc{1:56:19}
\speaker{老師}
% cue 2689 6979.430 6982.000
\hypertarget{w06:cue:2689}{}就是像剛剛前面講的那個，
% cue 2690 6982.000 6983.520
\hypertarget{w06:cue:2690}{}然後你可以跟我講，
% cue 2691 6983.520 6985.680
\hypertarget{w06:cue:2691}{}他為什麼認為程序面而言，
% cue 2692 6985.680 6987.560
\hypertarget{w06:cue:2692}{}也是比較適合由國會，
% cue 2693 6987.560 6989.810
\hypertarget{w06:cue:2693}{}來做成正義的決定，
% cue 2694 6989.810 6991.230
\hypertarget{w06:cue:2694}{}而不適合由法院，
% cue 2695 6991.230 6993.210
\hypertarget{w06:cue:2695}{}這就比較難講了。

\tc{1:56:33}
\speaker{導讀人：亦凡}
% cue 2696 6993.210 6995.350
\hypertarget{w06:cue:2696}{}他就設想一個例，
% cue 2697 6995.350 6999.210
\hypertarget{w06:cue:2697}{}一個 abstract algebra。

\tc{1:56:39}
\speaker{老師}
% cue 2698 6999.210 7005.430
\hypertarget{w06:cue:2698}{}你們大家找到這個字嗎？
% cue 2699 7005.430 7008.180
\hypertarget{w06:cue:2699}{}說一三八七頁嗎？

\tc{1:56:48}
\speaker{導讀人：亦凡}
% cue 2700 7008.180 7010.570
\hypertarget{w06:cue:2700}{}對對對，
% cue 2701 7010.570 7015.690
\hypertarget{w06:cue:2701}{}他就是舉一個抽象的例子，
% cue 2702 7015.690 7017.630
\hypertarget{w06:cue:2702}{}去做寫，
% cue 2703 7017.630 7020.630
\hypertarget{w06:cue:2703}{}為了寫作上的方式。

\tc{1:57:00}
\speaker{老師}
% cue 2704 7020.630 7022.960
\hypertarget{w06:cue:2704}{}沒關係他舉了什麼例子是重點。

\tc{1:57:02}
\speaker{導讀人：亦凡}
% cue 2705 7022.960 7023.660
\hypertarget{w06:cue:2705}{}他就是。

\tc{1:57:03}
\speaker{老師}
% cue 2706 7023.660 7026.000
\hypertarget{w06:cue:2706}{}你怎麼記得多名這樣講話這麼抽象。

\tc{1:57:06}
\speaker{導讀人：亦凡}
% cue 2707 7026.000 7030.780
\hypertarget{w06:cue:2707}{}他就是假設說有一群公民，
% cue 2708 7030.780 7034.350
\hypertarget{w06:cue:2708}{}然後有一群公民，
% cue 2709 7034.350 7037.000
\hypertarget{w06:cue:2709}{}然後組成一個集合，
% cue 2710 7037.000 7042.450
\hypertarget{w06:cue:2710}{}但是其中的有一部分的人是依據某些程序，
% cue 2711 7042.450 7049.340
\hypertarget{w06:cue:2711}{}然後他們有能夠做出這種有影響的 decision，
% cue 2712 7049.340 7056.440
\hypertarget{w06:cue:2712}{}但是如果是在這個決策團體之外的其他的公民，
% cue 2713 7056.440 7062.690
\hypertarget{w06:cue:2713}{}如果就是對先前的做出的。

\tc{1:57:42}
\speaker{老師}
% cue 2714 7062.690 7064.390
\hypertarget{w06:cue:2714}{}你現在講的這一段呢，
% cue 2715 7064.390 7066.180
\hypertarget{w06:cue:2715}{}你需要去處理一個問題，
% cue 2716 7066.180 7068.490
\hypertarget{w06:cue:2716}{}看有沒有講到所謂的，
% cue 2717 7068.490 7072.150
\hypertarget{w06:cue:2717}{}political legitimacy of decision procedure
% cue 2718 7072.150 7074.910
\hypertarget{w06:cue:2718}{}我們做一個決定，
% cue 2719 7074.910 7080.130
\hypertarget{w06:cue:2719}{}這個決定的程序裡面會牽涉到政治正當性的人，
% cue 2720 7080.130 7083.280
\hypertarget{w06:cue:2720}{}為什麼把我排除啊？
% cue 2721 7083.280 7086.490
\hypertarget{w06:cue:2721}{}我沒有排除我需要被拘束啊，
% cue 2722 7086.490 7089.590
\hypertarget{w06:cue:2722}{}那立法者是比較沒有問題，
% cue 2723 7089.590 7090.930
\hypertarget{w06:cue:2723}{}一人一票完去選，
% cue 2724 7090.930 7092.870
\hypertarget{w06:cue:2724}{}選出來的代表，
% cue 2725 7092.870 7095.630
\hypertarget{w06:cue:2725}{}然後很多個代表在裡面經過討論，
% cue 2726 7095.630 7097.860
\hypertarget{w06:cue:2726}{}所以你可以說他不代表你嗎？

\tc{1:58:17}
% cue 2727 7097.860 7100.800
\hypertarget{w06:cue:2727}{}那司法程序就不一樣啊，
% cue 2728 7100.800 7101.930
\hypertarget{w06:cue:2728}{}你不是原告你，
% cue 2729 7101.930 7103.230
\hypertarget{w06:cue:2729}{}不是被告你怎麼參加，
% cue 2730 7103.230 7105.460
\hypertarget{w06:cue:2730}{}怎麼參加？

\tc{1:58:25}
% cue 2731 7105.460 7109.570
\hypertarget{w06:cue:2731}{}所以你可以換個角度，
% cue 2732 7109.570 7113.850
\hypertarget{w06:cue:2732}{}所以大法官或者是抽象違憲審查，
% cue 2733 7113.850 7115.600
\hypertarget{w06:cue:2733}{}他也注意到這個問題，
% cue 2734 7115.600 7117.690
\hypertarget{w06:cue:2734}{}所以他要歡迎法庭之友\footnote{憲法訴訟法第 20 條第 1 項規定，當事人以外之人民、機關或團體，認其與憲法法庭審理之案件有關聯性，得聲請憲法法庭裁定許可，於所定期間內提出具參考價值之專業意見或資料；第 19 條第 1 項規定，憲法法庭審理案件認有必要時，得指定專家學者、機關或團體就相關問題提供專業意見或資料。\cite[第 19 條、第 20 條]{csfa}。}，
% cue 2735 7117.690 7119.760
\hypertarget{w06:cue:2735}{}請各界提供意見，
% cue 2736 7119.760 7123.560
\hypertarget{w06:cue:2736}{}大法官也會指定所謂的，
% cue 2737 7123.560 7129.190
\hypertarget{w06:cue:2737}{}專家學者陳述意見給我們聽，
% cue 2738 7129.190 7131.190
\hypertarget{w06:cue:2738}{}就是我們要知道，
% cue 2739 7131.190 7133.330
\hypertarget{w06:cue:2739}{}我們知道這件事影響很多人，
% cue 2740 7133.330 7134.600
\hypertarget{w06:cue:2740}{}如果只有原告被告，
% cue 2741 7134.600 7136.980
\hypertarget{w06:cue:2741}{}他的程序的觀點來看，
% cue 2742 7136.980 7138.770
\hypertarget{w06:cue:2742}{}他的正當性是不足的，
% cue 2743 7138.770 7140.110
\hypertarget{w06:cue:2743}{}所以讓他打開，
% cue 2744 7140.110 7143.420
\hypertarget{w06:cue:2744}{}我這樣講應該可以聽得懂吧，
% cue 2745 7143.420 7145.150
\hypertarget{w06:cue:2745}{}我們現在在不同的情景，
% cue 2746 7145.150 7148.110
\hypertarget{w06:cue:2746}{}我是用不同的話講給你們聽他講的事情。

\tc{1:59:08}
% cue 2747 7148.110 7148.630
\hypertarget{w06:cue:2747}{}好，
% cue 2748 7148.630 7150.390
\hypertarget{w06:cue:2748}{}那現在你要告訴我，
% cue 2749 7150.390 7152.160
\hypertarget{w06:cue:2749}{}這個地方有Scalia\footnote{Antonin Scalia（1936–2016），1986 年 9 月 26 日起任聯邦最高法院大法官，至 2016 年 2 月 13 日去世止。\url{https://www.fjc.gov/history/judges/scalia-antonin}。}，
% cue 2750 7152.160 7153.660
\hypertarget{w06:cue:2750}{}他引了很多Scalia的話，
% cue 2751 7153.660 7154.980
\hypertarget{w06:cue:2751}{}在一三九零頁，
% cue 2752 7154.980 7156.080
\hypertarget{w06:cue:2752}{}對不對，
% cue 2753 7156.080 7157.310
\hypertarget{w06:cue:2753}{}Scalia講什麼？
% cue 2754 7157.310 7159.560
\hypertarget{w06:cue:2754}{}有沒有？
% cue 2755 7159.560 7160.660
\hypertarget{w06:cue:2755}{}他有沒有講Scalia？
% cue 2756 7160.660 7163.060
\hypertarget{w06:cue:2756}{}Scalia在這件事情上還算坦白的。

\tc{1:59:23}
% cue 2757 7163.060 7166.060
\hypertarget{w06:cue:2757}{}Scalia的這個案子，
% cue 2758 7166.060 7168.420
\hypertarget{w06:cue:2758}{}是 Planned Parenthood\footnote{Planned Parenthood of Southeastern Pennsylvania v. Casey, 505 U.S. 833，聯邦最高法院 1992 年 6 月 29 日判決，審查賓州 1982 年墮胎管制法（Abortion Control Act）之五項規定，維持 Roe 案之核心見解。Scalia 大法官之不同意見在第 999–1000 頁，Waldron 於第 1390–1391 頁長段引之。見 Jeremy Waldron, The Core of the Case Against Judicial Review, 115 Yale L.J. 1346, 1390–1391 (2006)。\url{https://www.law.cornell.edu/supremecourt/text/505/833}。}，
% cue 2759 7168.420 7173.020
\hypertarget{w06:cue:2759}{}這個案子是牽涉到那個 Roe versus Wade，
% cue 2760 7173.020 7174.470
\hypertarget{w06:cue:2760}{}差一點被推翻，
% cue 2761 7174.470 7175.290
\hypertarget{w06:cue:2761}{}但是沒有推翻，
% cue 2762 7175.290 7176.900
\hypertarget{w06:cue:2762}{}基本上他們說，
% cue 2763 7176.900 7180.710
\hypertarget{w06:cue:2763}{}我們要所謂的 reasoned judgment\footnote{英文，「理性判斷」。Casey 案判決謂實體正當程序之裁判「may require this Court to exercise its reasoned judgment in determining the boundaries between the individual's liberty and the demands of organized society」，Scalia 大法官之不同意見即針對此一用語。\url{https://www.law.cornell.edu/supremecourt/text/505/833}。}，
% cue 2764 7180.710 7185.340
\hypertarget{w06:cue:2764}{}然後他認為這個 reasoned judgment 是非常荒謬的，
% cue 2765 7185.340 7188.450
\hypertarget{w06:cue:2765}{}我是說Scalia，
% cue 2766 7188.450 7191.440
\hypertarget{w06:cue:2766}{}Scalia認為看到了嗎？
% cue 2767 7191.440 7193.360
\hypertarget{w06:cue:2767}{}第三九零頁最底下一段，
% cue 2768 7193.360 7197.680
\hypertarget{w06:cue:2768}{}In reality, our process of constitutional adjudication
% cue 2769 7197.680 7201.600
\hypertarget{w06:cue:2769}{}我們憲法訴訟的程序的實質，
% cue 2770 7201.600 7206.450
\hypertarget{w06:cue:2770}{}真實，consists primarily of making value judgment，
% cue 2771 7206.450 7210.160
\hypertarget{w06:cue:2771}{}主要說穿的就是在做價值判斷，
% cue 2772 7210.160 7213.860
\hypertarget{w06:cue:2772}{}非民選的人，
% cue 2773 7213.860 7216.810
\hypertarget{w06:cue:2773}{}有什麼理由替你們大家做價值恩惡，
% cue 2774 7216.810 7220.450
\hypertarget{w06:cue:2774}{}那你們大家都沒有參與，
% cue 2775 7220.450 7223.260
\hypertarget{w06:cue:2775}{}覺得程序上這樣子說，
% cue 2776 7223.260 7225.170
\hypertarget{w06:cue:2776}{}我們有正當性嗎？
% cue 2777 7225.170 7227.530
\hypertarget{w06:cue:2777}{}你們聽懂嗎？

\tc{2:00:27}
% cue 2778 7227.530 7230.940
\hypertarget{w06:cue:2778}{}第三九一頁他又再次的強調了這個，
% cue 2779 7230.940 7232.480
\hypertarget{w06:cue:2779}{}有沒有？
% cue 2780 7232.480 7234.080
\hypertarget{w06:cue:2780}{}As the court says
% cue 2781 7234.080 7236.190
\hypertarget{w06:cue:2781}{}Undefined and unbounded
% cue 2782 7236.190 7241.090
\hypertarget{w06:cue:2782}{}Then the people should demonstrate to protest
% cue 2783 7241.090 7243.060
\hypertarget{w06:cue:2783}{}人民應該抗議，
% cue 2784 7243.060 7247.850
\hypertarget{w06:cue:2784}{}That we do not implement their values instead of ours
% cue 2785 7247.850 7253.460
\hypertarget{w06:cue:2785}{}我們現在說的這個意思，
% cue 2786 7253.460 7255.560
\hypertarget{w06:cue:2786}{}是我們大法官的價值判斷，
% cue 2787 7255.560 7257.510
\hypertarget{w06:cue:2787}{}還是人民的價值判斷，
% cue 2788 7257.510 7259.710
\hypertarget{w06:cue:2788}{}是我們說的，
% cue 2789 7259.710 7263.990
\hypertarget{w06:cue:2789}{}這樣有清楚嗎？

\tc{2:01:03}
% cue 2790 7263.990 7265.780
\hypertarget{w06:cue:2790}{}那後面它還有一個，
% cue 2791 7265.780 7267.020
\hypertarget{w06:cue:2791}{}12345
% cue 2792 7267.020 7269.510
\hypertarget{w06:cue:2792}{}它有五個什麼什麼東西，
% cue 2793 7269.510 7271.650
\hypertarget{w06:cue:2793}{}第三九三，
% cue 2794 7271.650 7293.200
\hypertarget{w06:cue:2794}{}第一點來講下，
% cue 2795 7293.200 7294.220
\hypertarget{w06:cue:2795}{}第三九三，
% cue 2796 7294.220 7296.080
\hypertarget{w06:cue:2796}{}贊成違憲審查，
% cue 2797 7296.080 7298.660
\hypertarget{w06:cue:2797}{}認為這個可以克服民主價值，
% cue 2798 7298.660 7299.870
\hypertarget{w06:cue:2798}{}來，講，
% cue 2799 7299.870 7304.140
\hypertarget{w06:cue:2799}{}找到了嗎？1393。

\tc{2:01:44}
\speaker{導讀人：亦凡}
% cue 2800 7304.140 7315.600
\hypertarget{w06:cue:2800}{}法官並沒有做他們自己的判斷。

\tc{2:01:55}
\speaker{老師}
% cue 2801 7315.600 7328.680
\hypertarget{w06:cue:2801}{}好，這樣子有聽懂嗎？
% cue 2802 7328.680 7333.890
\hypertarget{w06:cue:2802}{}法官說：沒有沒有 我不是在強加我的價值判斷於你的，
% cue 2803 7333.890 7339.540
\hypertarget{w06:cue:2803}{}我是在解釋放在憲法人權條款裡面，
% cue 2804 7339.540 7340.820
\hypertarget{w06:cue:2804}{}我在解釋。

\tc{2:02:20}
% cue 2805 7340.820 7344.120
\hypertarget{w06:cue:2805}{}這是我們常常看到的典型的例子，
% cue 2806 7344.120 7346.110
\hypertarget{w06:cue:2806}{}113 憲判八\footnote{中華民國憲法法庭 113 年憲判字第 8 號判決（死刑案），2024 年 9 月 20 日宣示，主筆黃昭元大法官。主文第一項謂刑法第 271 條第 1 項等規定以死刑為最重本刑部分，「僅得適用於個案犯罪情節屬最嚴重，且其刑事程序符合憲法最嚴密之正當法律程序要求之情形。於此範圍內，與憲法保障人民生命權之意旨尚屬無違。」\cite{ccj113-8}。}，
% cue 2807 7346.110 7348.060
\hypertarget{w06:cue:2807}{}我們憲法哪一條？
% cue 2808 7348.060 7353.710
\hypertarget{w06:cue:2808}{}廢死 所不可憂的思想，
% cue 2809 7353.710 7356.020
\hypertarget{w06:cue:2809}{}法官說：沒有 我講的都是憲法的意思，
% cue 2810 7356.020 7357.360
\hypertarget{w06:cue:2810}{}我沒有講我的意思，
% cue 2811 7357.360 7359.320
\hypertarget{w06:cue:2811}{}我現在講的聽得懂吧。

\tc{2:02:39}
% cue 2812 7359.320 7360.730
\hypertarget{w06:cue:2812}{}你們今天要一面，
% cue 2813 7360.730 7362.760
\hypertarget{w06:cue:2813}{}我都在盡量幫你們了嘛，
% cue 2814 7362.760 7365.470
\hypertarget{w06:cue:2814}{}我們這個學期 你要找到你可以寫的東西，
% cue 2815 7365.470 7368.340
\hypertarget{w06:cue:2815}{}不然你要怎麼交你的報告，
% cue 2816 7368.340 7371.620
\hypertarget{w06:cue:2816}{}我已經提醒了你要去看這些東西，
% cue 2817 7371.620 7373.820
\hypertarget{w06:cue:2817}{}因為它都是非常有爭議的案件，
% cue 2818 7373.820 7375.920
\hypertarget{w06:cue:2818}{}它的正當性在哪裡，
% cue 2819 7375.920 7377.720
\hypertarget{w06:cue:2819}{}你覺得它怎麼不正當怎麼正當，
% cue 2820 7377.720 7379.660
\hypertarget{w06:cue:2820}{}今天用了一大堆材料給你。

\tc{2:02:59}
% cue 2821 7379.660 7380.500
\hypertarget{w06:cue:2821}{}好第二個，
% cue 2822 7380.500 7380.900
\hypertarget{w06:cue:2822}{}Second
% cue 2823 7380.900 7382.740
\hypertarget{w06:cue:2823}{}第三九三對比一下。

\tc{2:03:02}
\speaker{導讀人：亦凡}
% cue 2824 7382.740 7384.980
\hypertarget{w06:cue:2824}{}Second and in my。

\tc{2:03:04}
\speaker{老師}
% cue 2825 7384.980 7386.980
\hypertarget{w06:cue:2825}{}Judges are simply enforcing。

\tc{2:03:06}
\speaker{導讀人：亦凡}
% cue 2826 7386.980 7392.520
\hypertarget{w06:cue:2826}{}Judges are simply enforcing the society's own pre-commitment to rights。

\tc{2:03:12}
\speaker{老師}
% cue 2827 7392.520 7392.760
\hypertarget{w06:cue:2827}{}對，
% cue 2828 7392.760 7394.920
\hypertarget{w06:cue:2828}{}法官只不過在執行，
% cue 2829 7394.920 7398.040
\hypertarget{w06:cue:2829}{}這個社會對於權利已經有的承諾，
% cue 2830 7398.040 7399.820
\hypertarget{w06:cue:2830}{}我們沒有說什麼。

\tc{2:03:19}
% cue 2831 7399.820 7403.420
\hypertarget{w06:cue:2831}{}好 第三 第三組是，
% cue 2832 7403.420 7421.450
\hypertarget{w06:cue:2832}{}好 你們覺得這樣講對不對，
% cue 2833 7421.450 7424.040
\hypertarget{w06:cue:2833}{}我們做了決定 我們有去梳理，
% cue 2834 7424.040 7425.840
\hypertarget{w06:cue:2834}{}你不服你小朋友去推翻嗎？

\tc{2:03:45}
% cue 2835 7425.840 7428.180
\hypertarget{w06:cue:2835}{}這裡有一個有趣的問題，
% cue 2836 7428.180 7430.420
\hypertarget{w06:cue:2836}{}美國要怎麼推翻？
% cue 2837 7430.420 7432.820
\hypertarget{w06:cue:2837}{}這也牽涉到老師的理論，
% cue 2838 7432.820 7435.380
\hypertarget{w06:cue:2838}{}為什麼我同意我贊成，
% cue 2839 7435.380 7438.190
\hypertarget{w06:cue:2839}{}而且我儘量擴大違憲審查的範圍，
% cue 2840 7438.190 7440.190
\hypertarget{w06:cue:2840}{}要建構一個無漏洞的違憲審查，
% cue 2841 7440.190 7442.650
\hypertarget{w06:cue:2841}{}但是我仍然覺得，
% cue 2842 7442.650 7445.230
\hypertarget{w06:cue:2842}{}它的效力應該只是，
% cue 2843 7445.230 7446.810
\hypertarget{w06:cue:2843}{}這個回合的最終，
% cue 2844 7446.810 7449.930
\hypertarget{w06:cue:2844}{}它不可以具有永久的最終去輸入，
% cue 2845 7449.930 7453.110
\hypertarget{w06:cue:2845}{}我也甚至於認為立法者，
% cue 2846 7453.110 7455.990
\hypertarget{w06:cue:2846}{}可以重新啟動下一輪的立法，
% cue 2847 7455.990 7458.240
\hypertarget{w06:cue:2847}{}跟你現在做的完全，
% cue 2848 7458.240 7459.830
\hypertarget{w06:cue:2848}{}內容完全相同，
% cue 2849 7459.830 7461.150
\hypertarget{w06:cue:2849}{}內容完全相反，
% cue 2850 7461.150 7462.280
\hypertarget{w06:cue:2850}{}你說這個違憲，
% cue 2851 7462.280 7463.540
\hypertarget{w06:cue:2851}{}違憲是一個可以再立法。

\tc{2:04:23}
% cue 2852 7463.540 7464.660
\hypertarget{w06:cue:2852}{}比方說你說，
% cue 2853 7464.660 7473.400
\hypertarget{w06:cue:2853}{}所以刑法當中那個焚燒國旗罪\footnote{刑法第 160 條第 1 項規定：「意圖侮辱中華民國，而公然損壞、除去或污辱中華民國之國徽、國旗者，處一年以下有期徒刑、拘役或九千元以下罰金。」\url{https://law.moj.gov.tw/LawClass/LawAll.aspx?pcode=C0000001}。}被我宣告違憲，
% cue 2854 7473.400 7476.460
\hypertarget{w06:cue:2854}{}那你可不可以再立一個同樣的規定？
% cue 2855 7476.460 7478.680
\hypertarget{w06:cue:2855}{}我會認為可以 如果你立得出來，
% cue 2856 7478.680 7480.960
\hypertarget{w06:cue:2856}{}然後如果立出來了 又有人來申請，
% cue 2857 7480.960 7483.090
\hypertarget{w06:cue:2857}{}那我不要再次的審查，
% cue 2858 7483.090 7486.940
\hypertarget{w06:cue:2858}{}我不贊成 405 號\footnote{司法院釋字第 405 號解釋（1996 年 6 月 7 日）解釋文謂，教育人員任用條例第 21 條第 2 項中關於「並得在各學校間調任」之規定，「使未經考試及格者與取得公務人員任用資格者之法律地位幾近相同，與憲法第八十五條、第七條及前開解釋意旨不符，應自本解釋公布之日起失其效力」；所稱前開解釋即釋字第 278 號。\url{https://cons.judicial.gov.tw/docdata.aspx?fid=100&id=310586}。}的見解，
% cue 2859 7486.940 7488.940
\hypertarget{w06:cue:2859}{}我見解就告訴你在 278 號\footnote{司法院釋字第 278 號解釋（1991 年 5 月 17 日）解釋文謂，教育人員任用條例第 21 條所定已遴用之學校現任職員「適用各該原有關法令」，「並不能使未經考試及格者取得與考試及格者相同之公務人員任用資格，因之，僅能繼續在原學校任職」。\url{https://cons.judicial.gov.tw/docdata.aspx?fid=100&id=310459}。}告訴你，
% cue 2860 7488.940 7491.870
\hypertarget{w06:cue:2860}{}怎麼再來問一次 那當然違憲退回去。

\tc{2:04:51}
% cue 2861 7491.870 7494.250
\hypertarget{w06:cue:2861}{}我不這麼認為，
% cue 2862 7494.250 7496.980
\hypertarget{w06:cue:2862}{}時過境遷 社會改變，
% cue 2863 7496.980 7499.980
\hypertarget{w06:cue:2863}{}剛才他不是講他是實行社會上大多數人的意思嗎？
% cue 2864 7499.980 7500.980
\hypertarget{w06:cue:2864}{}不是他的意思，
% cue 2865 7500.980 7503.170
\hypertarget{w06:cue:2865}{}那你怎麼可以堅持你過去的意志建立，
% cue 2866 7503.170 7505.170
\hypertarget{w06:cue:2866}{}你就說是不可以的，
% cue 2867 7505.170 7506.300
\hypertarget{w06:cue:2867}{}我認為這就要重新來，
% cue 2868 7506.300 7508.300
\hypertarget{w06:cue:2868}{}這是一個第三個看法。

\tc{2:05:08}
% cue 2869 7508.300 7509.300
\hypertarget{w06:cue:2869}{}第四個看法，
% cue 2870 7509.300 7517.720
\hypertarget{w06:cue:2870}{}這是我們臺灣人最喜歡講的，
% cue 2871 7517.720 7518.720
\hypertarget{w06:cue:2871}{}尤其那幾個大法官，
% cue 2872 7518.720 7521.720
\hypertarget{w06:cue:2872}{}誰說我們大法官沒有民主正當性，
% cue 2873 7521.720 7522.720
\hypertarget{w06:cue:2873}{}我們是總統提名\footnote{憲法增修條文第 5 條第 1 項規定，司法院設大法官十五人，並以其中一人為院長、一人為副院長，「由總統提名，經立法院同意任命之」。\url{https://law.moj.gov.tw/LawClass/LawAll.aspx?pcode=A0000002}。}的，
% cue 2874 7522.720 7523.720
\hypertarget{w06:cue:2874}{}立法委員同意的，
% cue 2875 7523.720 7524.720
\hypertarget{w06:cue:2875}{}總統是人民民選的，
% cue 2876 7524.720 7525.720
\hypertarget{w06:cue:2876}{}立法委員是民選的，
% cue 2877 7525.720 7528.720
\hypertarget{w06:cue:2877}{}所以我們有間接民意基礎。

\tc{2:05:28}
% cue 2878 7528.720 7530.160
\hypertarget{w06:cue:2878}{}記不記得這個話？
% cue 2879 7530.160 7532.560
\hypertarget{w06:cue:2879}{}在哪裡聽起來似曾相識，
% cue 2880 7532.560 7533.700
\hypertarget{w06:cue:2880}{}老師還批評了他，
% cue 2881 7533.700 7536.320
\hypertarget{w06:cue:2881}{}有興趣去看老師的影片，
% cue 2882 7536.320 7538.330
\hypertarget{w06:cue:2882}{}觀還是聽，
% cue 2883 7538.330 7540.800
\hypertarget{w06:cue:2883}{}廢死辯論，
% cue 2884 7540.800 7542.910
\hypertarget{w06:cue:2884}{}大法官廢死辯論案有感，
% cue 2885 7542.910 7545.430
\hypertarget{w06:cue:2885}{}我覺得那個是非常有問題，
% cue 2886 7545.430 7547.100
\hypertarget{w06:cue:2886}{}這就是這個想法，
% cue 2887 7547.100 7549.370
\hypertarget{w06:cue:2887}{}我講的每一件事情，
% cue 2888 7549.370 7550.750
\hypertarget{w06:cue:2888}{}不同的文本，
% cue 2889 7550.750 7552.570
\hypertarget{w06:cue:2889}{}你會找到他們的來源，
% cue 2890 7552.570 7554.910
\hypertarget{w06:cue:2890}{}作者也不同意，
% cue 2891 7554.910 7558.120
\hypertarget{w06:cue:2891}{}好了 第五個，
% cue 2892 7558.120 7561.610
\hypertarget{w06:cue:2892}{}不能因為是總統提名立法院同意，
% cue 2893 7561.610 7563.170
\hypertarget{w06:cue:2893}{}就說我們大法官代表民意，
% cue 2894 7563.170 7564.300
\hypertarget{w06:cue:2894}{}我們有民意基礎，
% cue 2895 7564.300 7565.560
\hypertarget{w06:cue:2895}{}所以我們可以做決定，
% cue 2896 7565.560 7566.380
\hypertarget{w06:cue:2896}{}不是這個意思。

\tc{2:06:06}
\speaker{導讀人：亦凡}
% cue 2897 7566.380 7567.100
\hypertarget{w06:cue:2897}{}第五個，
% cue 2898 7567.890 7571.500
\hypertarget{w06:cue:2898}{}最後，法律的法律的檢查證明，
% cue 2899 7571.500 7576.200
\hypertarget{w06:cue:2899}{}那個法律的法律可以製造成更加進行，
% cue 2900 7576.200 7579.810
\hypertarget{w06:cue:2900}{}民間的進行政策的進行政策。

\tc{2:06:19}
\speaker{老師}
% cue 2901 7579.810 7581.040
\hypertarget{w06:cue:2901}{}這句話看不懂，
% cue 2902 7581.040 7584.670
\hypertarget{w06:cue:2902}{}有時，民間的進行政策的進行政策，
% cue 2903 7584.670 7586.270
\hypertarget{w06:cue:2903}{}有時，民間的進行政策，
% cue 2904 7586.270 7588.280
\hypertarget{w06:cue:2904}{}有時，民間的進行政策。

\tc{2:06:28}
% cue 2905 7588.280 7591.070
\hypertarget{w06:cue:2905}{}那你現在是說，我作為訴訟當事人，
% cue 2906 7591.070 7592.150
\hypertarget{w06:cue:2906}{}我就可以參加，
% cue 2907 7592.150 7594.530
\hypertarget{w06:cue:2907}{}所以程序上面來講，你沒有錯，
% cue 2908 7594.530 7597.110
\hypertarget{w06:cue:2908}{}但是你不要忘記，那你再怎麼參加，
% cue 2909 7597.110 7601.970
\hypertarget{w06:cue:2909}{}除非你今天把憲法訴訟 搞成公民訴訟\footnote{環境影響評估法第 23 條第 8 項、第 9 項規定，開發單位違反該法而主管機關疏於執行時，受害人民或公益團體得以書面告知主管機關；主管機關於書面告知送達之日起六十日內仍未依法執行者，人民或公益團體得以該主管機關為被告，直接向行政法院提起訴訟，請求判令其執行。\url{https://law.moj.gov.tw/LawClass/LawSingle.aspx?pcode=O0090001&flno=23}。}，
% cue 2910 7601.970 7603.530
\hypertarget{w06:cue:2910}{}你有嗎？

\tc{2:06:43}
% cue 2911 7603.530 7605.540
\hypertarget{w06:cue:2911}{}我現在問你們很簡單的問題，
% cue 2912 7605.540 7609.600
\hypertarget{w06:cue:2912}{}臺灣人民可不可以提起憲法訴訟？
% cue 2913 7609.600 7611.520
\hypertarget{w06:cue:2913}{}人民可不可以申請大法官解釋？
% cue 2914 7611.520 7612.590
\hypertarget{w06:cue:2914}{}可以，
% cue 2915 7612.590 7615.390
\hypertarget{w06:cue:2915}{}那人民申請解釋\footnote{憲法訴訟法第 59 條第 1 項規定，人民於其憲法上所保障之權利遭受不法侵害，經依法定程序用盡審級救濟程序，對於所受不利確定終局裁判，或該裁判及其所適用之法規範，認有牴觸憲法者，得聲請憲法法庭為宣告違憲之判決。\cite[第 59 條]{csfa}。}是沒有前提要件，
% cue 2916 7615.390 7618.470
\hypertarget{w06:cue:2916}{}他必須受有對他福利的確定案件，
% cue 2917 7618.470 7619.710
\hypertarget{w06:cue:2917}{}所以他要嘛，
% cue 2918 7619.710 7622.540
\hypertarget{w06:cue:2918}{}他通常是被告或者他是原告，
% cue 2919 7622.540 7625.260
\hypertarget{w06:cue:2919}{}但總之他必須是個案當事人，
% cue 2920 7625.260 7626.900
\hypertarget{w06:cue:2920}{}我們並沒有承認，
% cue 2921 7626.900 7628.340
\hypertarget{w06:cue:2921}{}像環境法當中，
% cue 2922 7628.340 7630.100
\hypertarget{w06:cue:2922}{}任何人可以基於公益，
% cue 2923 7630.100 7631.570
\hypertarget{w06:cue:2923}{}來申請大法官解釋，
% cue 2924 7631.570 7632.730
\hypertarget{w06:cue:2924}{}所以你現在講說，
% cue 2925 7632.730 7634.530
\hypertarget{w06:cue:2925}{}資格可以變成瑞士格子，
% cue 2926 7634.530 7635.870
\hypertarget{w06:cue:2926}{}然後他就可以參加了，
% cue 2927 7635.870 7638.000
\hypertarget{w06:cue:2927}{}所以你不要跟我說你沒有機會參加，
% cue 2928 7638.000 7639.130
\hypertarget{w06:cue:2928}{}這個沒有正當性，
% cue 2929 7639.130 7640.510
\hypertarget{w06:cue:2929}{}這個還是有問題，
% cue 2930 7640.510 7642.480
\hypertarget{w06:cue:2930}{}這個好像有落差 嚴重落差。

\tc{2:07:22}
% cue 2931 7642.480 7644.350
\hypertarget{w06:cue:2931}{}現在相反的是機關，
% cue 2932 7644.350 7647.090
\hypertarget{w06:cue:2932}{}機關他就可以抽象的來申請，
% cue 2933 7647.090 7648.510
\hypertarget{w06:cue:2933}{}而我就是基於公益，
% cue 2934 7648.510 7649.890
\hypertarget{w06:cue:2934}{}那你要問 Hans Kelsen\footnote{Hans Kelsen（1881–1973），奧地利法學家，1920 年奧地利聯邦憲法之主要起草人，曾任奧地利憲法法院法官。其向德國國家法學者協會提出之報告〈Wesen und Entwicklung der Staatsgerichtsbarkeit〉（國事裁判權之本質與發展）認為，准許任何人聲請之 actio popularis（民眾訴訟）固為憲法最強之保障，但有恣意聲請與憲法法院負荷過重之虞，不宜採行。\cite[74]{kelsen1929}。}，
% cue 2935 7649.890 7651.370
\hypertarget{w06:cue:2935}{}Hans Kelsen 給你的理就上門，
% cue 2936 7651.370 7652.970
\hypertarget{w06:cue:2936}{}大家看到下層規範，
% cue 2937 7652.970 7654.330
\hypertarget{w06:cue:2937}{}不可以牴觸上層規範，
% cue 2938 7654.330 7657.690
\hypertarget{w06:cue:2938}{}大家是為了要維護這個法規範的體系，
% cue 2939 7657.690 7659.970
\hypertarget{w06:cue:2939}{}所以你當然有權利。

\tc{2:07:39}
% cue 2940 7659.970 7662.810
\hypertarget{w06:cue:2940}{}那我就不懂，Hans Kelsen，那你應該加一句啊，
% cue 2941 7662.810 7666.990
\hypertarget{w06:cue:2941}{}那人民是不是也有維護階層式規範，
% cue 2942 7666.990 7670.200
\hypertarget{w06:cue:2942}{}階層式規範構造的一般經歷，
% cue 2943 7670.200 7672.220
\hypertarget{w06:cue:2943}{}那為什麼人民不能申請，
% cue 2944 7672.220 7676.850
\hypertarget{w06:cue:2944}{}如果照 Hans Kelsen 的說法，機關可以，
% cue 2945 7676.850 7678.800
\hypertarget{w06:cue:2945}{}為什麼人民不能申請？
% cue 2946 7678.800 7681.810
\hypertarget{w06:cue:2946}{}我就聽不懂，但是我沒有看到哪一個國家，
% cue 2947 7681.810 7685.480
\hypertarget{w06:cue:2947}{}允許人民來提起一般的憲法訴訟，
% cue 2948 7685.480 7688.870
\hypertarget{w06:cue:2948}{}公民訴訟，我們一般是指特定的情況下，
% cue 2949 7688.870 7692.030
\hypertarget{w06:cue:2949}{}我們沒有看到憲法訴訟變成公民訴訟，
% cue 2950 7692.030 7693.920
\hypertarget{w06:cue:2950}{}好了 講完了。

\section{多數暴政（第陸部分）}

\tc{2:08:13}
% cue 2951 7693.920 7695.520
\hypertarget{w06:cue:2951}{}來 第四 第六點，
% cue 2952 7695.520 7697.630
\hypertarget{w06:cue:2952}{}Tyranny of Majority
% cue 2953 7697.630 7701.870
\hypertarget{w06:cue:2953}{}有關於這個所謂的，
% cue 2954 7701.870 7704.290
\hypertarget{w06:cue:2954}{}二十六分，
% cue 2955 7704.290 7705.850
\hypertarget{w06:cue:2955}{}所謂的多數暴力，
% cue 2956 7705.850 7707.450
\hypertarget{w06:cue:2956}{}所以我們必須要用，
% cue 2957 7707.450 7709.660
\hypertarget{w06:cue:2957}{}違憲審查來對抗多數暴力，
% cue 2958 7709.660 7710.780
\hypertarget{w06:cue:2958}{}作者講什麼？
% cue 2959 7710.780 7713.270
\hypertarget{w06:cue:2959}{}來 找到最精確的論。

\tc{2:08:33}
\speaker{導讀人：亦凡}
% cue 2960 7713.270 7716.260
\hypertarget{w06:cue:2960}{}I believe that。

\tc{2:08:36}
\speaker{老師}
% cue 2961 7716.260 7717.260
\hypertarget{w06:cue:2961}{}在哪裡。

\tc{2:08:37}
\speaker{導讀人：亦凡}
% cue 2962 7717.260 7720.260
\hypertarget{w06:cue:2962}{}在 1365 的最後面。

\tc{2:08:40}
\speaker{老師}
% cue 2963 7720.260 7721.260
\hypertarget{w06:cue:2963}{}1395 吧。

\tc{2:08:41}
\speaker{導讀人：亦凡}
% cue 2964 7721.260 7722.720
\hypertarget{w06:cue:2964}{}對 1395。

\tc{2:08:42}
\speaker{老師}
% cue 2965 7722.720 7723.880
\hypertarget{w06:cue:2965}{}好 來講一講，
% cue 2966 7723.880 7724.480
\hypertarget{w06:cue:2966}{}來哪裡？

\tc{2:08:44}
\speaker{導讀人：亦凡}
% cue 2967 7724.480 7725.380
\hypertarget{w06:cue:2967}{}走。

\tc{2:08:45}
\speaker{老師}
% cue 2968 7725.380 7726.300
\hypertarget{w06:cue:2968}{}什麼地方？

\tc{2:08:46}
\speaker{導讀人：亦凡}
% cue 2969 7726.300 7729.290
\hypertarget{w06:cue:2969}{}I believe that this common argument is seriously confused。

\tc{2:08:49}
\speaker{老師}
% cue 2970 7729.290 7730.490
\hypertarget{w06:cue:2970}{}你沒有找到嗎？
% cue 2971 7730.490 7733.290
\hypertarget{w06:cue:2971}{}我沒有看到"at least"的在哪裡了。
% cue 2972 7733.290 7734.570
\hypertarget{w06:cue:2972}{}你是"at least"的嗎？

\tc{2:08:54}
\speaker{導讀人：亦凡}
% cue 2973 7734.570 7735.580
\hypertarget{w06:cue:2973}{}I believe。

\tc{2:08:55}
\speaker{老師}
% cue 2974 7735.580 7737.370
\hypertarget{w06:cue:2974}{}I believe.
% cue 2975 7737.370 7740.800
\hypertarget{w06:cue:2975}{}第二段，下一段，最後一段，I believe。

\tc{2:09:00}
\speaker{導讀人：亦凡}
% cue 2976 7740.800 7744.820
\hypertarget{w06:cue:2976}{}I believe that this common argument is seriously confused。

\tc{2:09:04}
\speaker{老師}
% cue 2977 7744.820 7747.700
\hypertarget{w06:cue:2977}{}啊，嚴重的誤導，混淆。
% cue 2978 7747.700 7750.400
\hypertarget{w06:cue:2978}{}Let's grant，讓我們假定。

\tc{2:09:10}
\speaker{導讀人：亦凡}
% cue 2979 7750.400 7756.030
\hypertarget{w06:cue:2979}{}For now, that tyranny is what happens to someone
% cue 2980 7756.030 7757.730
\hypertarget{w06:cue:2980}{}when their rights are denied。

\tc{2:09:17}
\speaker{老師}
% cue 2981 7757.730 7759.970
\hypertarget{w06:cue:2981}{}假定是這樣子。

\tc{2:09:19}
\speaker{導讀人：亦凡}
% cue 2982 7759.970 7776.860
\hypertarget{w06:cue:2982}{}我現在是說他要怎麼解釋這個事情。

\tc{2:09:36}
\speaker{老師}
% cue 2983 7776.860 7781.800
\hypertarget{w06:cue:2983}{}我們一般不是認為這樣子嗎？
% cue 2984 7781.800 7784.840
\hypertarget{w06:cue:2984}{}一般認為為什麼違憲審查具有正當性，
% cue 2985 7784.840 7787.680
\hypertarget{w06:cue:2985}{}因為違憲審查可以用來對抗多數暴力，
% cue 2986 7787.680 7790.570
\hypertarget{w06:cue:2986}{}那他對這點他是怎麼解釋的，
% cue 2987 7790.570 7791.690
\hypertarget{w06:cue:2987}{}他反對喔。

\tc{2:09:51}
\speaker{導讀人：亦凡}
% cue 2988 7791.690 7792.630
\hypertarget{w06:cue:2988}{}他從根本，
% cue 2989 7792.630 7798.300
\hypertarget{w06:cue:2989}{}他從根本上去拆解多數暴力的這個情形。

\tc{2:09:58}
\speaker{老師}
% cue 2990 7798.300 7799.780
\hypertarget{w06:cue:2990}{}好 那你告訴我怎麼拆解，
% cue 2991 7799.780 7801.560
\hypertarget{w06:cue:2991}{}剛才這一段沒有講到。

\tc{2:10:01}
\speaker{導讀人：亦凡}
% cue 2992 7801.560 7803.130
\hypertarget{w06:cue:2992}{}他後面。

\tc{2:10:03}
\speaker{老師}
% cue 2993 7803.130 7805.040
\hypertarget{w06:cue:2993}{}就我叫你講啊，
% cue 2994 7805.040 7806.440
\hypertarget{w06:cue:2994}{}你現在也沒有時間了，
% cue 2995 7806.440 7808.080
\hypertarget{w06:cue:2995}{}我們現在已經 28 分了，
% cue 2996 7808.080 7808.580
\hypertarget{w06:cue:2996}{}對不對？
% cue 2997 7808.580 7810.140
\hypertarget{w06:cue:2997}{}也快要結束了，
% cue 2998 7810.140 7812.630
\hypertarget{w06:cue:2998}{}事實上還是有逼你把它講完，
% cue 2999 7812.630 7813.530
\hypertarget{w06:cue:2999}{}現在距離結束，
% cue 3000 7813.530 7814.650
\hypertarget{w06:cue:3000}{}只剩幾月了。

\tc{2:10:14}
\speaker{導讀人：亦凡}
% cue 3001 7814.650 7818.290
\hypertarget{w06:cue:3001}{}他從我可以看到那個 1397 頁，
% cue 3002 7818.290 7822.170
\hypertarget{w06:cue:3002}{}然後他就是區分那個決定上的多數少數，
% cue 3003 7822.170 7825.950
\hypertarget{w06:cue:3003}{}跟議題上的多數、少數。

\tc{2:10:25}
\speaker{老師}
% cue 3004 7825.950 7827.790
\hypertarget{w06:cue:3004}{}那你告訴我在什麼地方？

\tc{2:10:27}
\speaker{導讀人：亦凡}
% cue 3005 7827.790 7833.760
\hypertarget{w06:cue:3005}{}在 1397 的第二段，
% cue 3006 7833.760 7836.490
\hypertarget{w06:cue:3006}{}就是在 1397 的那一個段，
% cue 3007 7836.490 7837.870
\hypertarget{w06:cue:3007}{}Injustice。

\tc{2:10:37}
\speaker{老師}
% cue 3008 7837.870 7839.990
\hypertarget{w06:cue:3008}{}Injustice is what happened。

\tc{2:10:39}
\speaker{導讀人：亦凡}
% cue 3009 7839.990 7842.560
\hypertarget{w06:cue:3009}{}Injustice is what。

\tc{2:10:42}
\speaker{老師}
% cue 3010 7842.560 7843.920
\hypertarget{w06:cue:3010}{}那哪裡你要講哪裡，
% cue 3011 7843.920 7845.440
\hypertarget{w06:cue:3011}{}照你要去掉的地方講。

\tc{2:10:45}
\speaker{導讀人：亦凡}
% cue 3012 7845.440 7850.630
\hypertarget{w06:cue:3012}{}We need to distinguish at least in the first instance
% cue 3013 7850.630 7853.930
\hypertarget{w06:cue:3013}{}Between the decisional majority\footnote{英文，「決定上的多數」。Waldron 區分決定上的（decisional）多數與少數，即就某一決定投票之多數與少數，以及議題上的（topical）多數與少數，即其權利於該決定中受影響之群體；兩者成員一致之情形，例如白人議員投票支持白人特權，為其所稱多數暴政應特別關切者。見 Jeremy Waldron, The Core of the Case Against Judicial Review, 115 Yale L.J. 1346, 1397 (2006)。} and minority
% cue 3014 7853.930 7859.010
\hypertarget{w06:cue:3014}{}And what I shall call the topical majority and minority。

\tc{2:10:59}
\speaker{老師}
% cue 3015 7859.010 7860.770
\hypertarget{w06:cue:3015}{}好 大家跟他念這段開始吧，
% cue 3016 7860.770 7864.040
\hypertarget{w06:cue:3016}{}我們要區別 區別什麼。

\tc{2:11:04}
\speaker{導讀人：亦凡}
% cue 3017 7864.040 7868.850
\hypertarget{w06:cue:3017}{}區別決定上的多數少數跟。

\tc{2:11:08}
\speaker{老師}
% cue 3018 7868.850 7871.230
\hypertarget{w06:cue:3018}{}你覺得是決定上的 majority。

\tc{2:11:11}
\speaker{導讀人：亦凡}
% cue 3019 7871.230 7873.190
\hypertarget{w06:cue:3019}{}decisional。

\tc{2:11:13}
\speaker{老師}
% cue 3020 7873.190 7873.730
\hypertarget{w06:cue:3020}{}還有呢？

\tc{2:11:13}
\speaker{導讀人：亦凡}
% cue 3021 7873.730 7876.510
\hypertarget{w06:cue:3021}{}還有議題上的多數少數。

\tc{2:11:16}
\speaker{老師}
% cue 3022 7876.510 7878.200
\hypertarget{w06:cue:3022}{}好 然後呢。

\tc{2:11:18}
\speaker{導讀人：亦凡}
% cue 3023 7878.200 7881.370
\hypertarget{w06:cue:3023}{}就是他 大家可以看我整理的，
% cue 3024 7881.370 7883.990
\hypertarget{w06:cue:3024}{}那個表格在附錄的地方，
% cue 3025 7883.990 7889.780
\hypertarget{w06:cue:3025}{}就是有一個在第八十。

\tc{2:11:29}
\speaker{老師}
% cue 3026 7889.780 7891.330
\hypertarget{w06:cue:3026}{}那你現在只要做你的表格講的，
% cue 3027 7891.330 7893.050
\hypertarget{w06:cue:3027}{}因為其他我們姿態的話，
% cue 3028 7893.050 7893.630
\hypertarget{w06:cue:3028}{}好，
% cue 3029 7893.630 7894.290
\hypertarget{w06:cue:3029}{}八十二嗎？

\tc{2:11:34}
\speaker{導讀人：亦凡}
% cue 3030 7894.290 7895.410
\hypertarget{w06:cue:3030}{}八十六，
% cue 3031 7895.410 7895.990
\hypertarget{w06:cue:3031}{}八十六。

\tc{2:11:35}
\speaker{老師}
% cue 3032 7895.990 7899.820
\hypertarget{w06:cue:3032}{}你們看他這麼用心，
% cue 3033 7899.820 7901.950
\hypertarget{w06:cue:3033}{}還幫你們做了這麼多表格，
% cue 3034 7901.950 7904.060
\hypertarget{w06:cue:3034}{}嚇死人了，
% cue 3035 7904.060 7907.760
\hypertarget{w06:cue:3035}{}你花了多長時間了？
% cue 3036 7907.760 7912.230
\hypertarget{w06:cue:3036}{}繼續一起開始就再做這個，
% cue 3037 7912.230 7913.750
\hypertarget{w06:cue:3037}{}好來，
% cue 3038 7913.750 7915.080
\hypertarget{w06:cue:3038}{}八十六頁來，
% cue 3039 7915.080 7915.760
\hypertarget{w06:cue:3039}{}請講。

\tc{2:11:55}
\speaker{導讀人：亦凡}
% cue 3040 7915.760 7916.920
\hypertarget{w06:cue:3040}{}就是他。

\tc{2:11:56}
\speaker{老師}
% cue 3041 7916.920 7918.760
\hypertarget{w06:cue:3041}{}你用這個表解釋給我們聽，
% cue 3042 7918.760 7923.140
\hypertarget{w06:cue:3042}{}什麼叫 decisional and topical majorities？

\tc{2:12:03}
\speaker{導讀人：亦凡}
% cue 3043 7923.140 7924.920
\hypertarget{w06:cue:3043}{}區分兩個維度吧，
% cue 3044 7924.920 7927.480
\hypertarget{w06:cue:3044}{}就是議題上的少數，
% cue 3045 7927.480 7931.270
\hypertarget{w06:cue:3045}{}就是這個問題所涉及的權利。

\tc{2:12:11}
\speaker{老師}
% cue 3046 7931.270 7934.050
\hypertarget{w06:cue:3046}{}可是它的標題是多數與少數，
% cue 3047 7934.050 7938.040
\hypertarget{w06:cue:3047}{}那你為什麼只有少數沒有多數呢？

\tc{2:12:18}
\speaker{導讀人：亦凡}
% cue 3048 7938.040 7944.080
\hypertarget{w06:cue:3048}{}少數多數就是那個少數的對應的那個題。

\tc{2:12:24}
\speaker{老師}
% cue 3049 7944.080 7947.360
\hypertarget{w06:cue:3049}{}所以你是不是應該寫議題上的多數與少數？
% cue 3050 7947.360 7948.360
\hypertarget{w06:cue:3050}{}我不知道，
% cue 3051 7948.360 7950.060
\hypertarget{w06:cue:3051}{}好 您開始來繼續講，
% cue 3052 7950.160 7952.950
\hypertarget{w06:cue:3052}{}你現在講 86 業級，繼續說。

\tc{2:12:32}
\speaker{導讀人：亦凡}
% cue 3053 7952.950 7957.650
\hypertarget{w06:cue:3053}{}他的，他去，就是他去拆分，
% cue 3054 7957.650 7963.040
\hypertarget{w06:cue:3054}{}然後去定位出那個核心情形的那個位置，
% cue 3055 7963.040 7965.470
\hypertarget{w06:cue:3055}{}就是既是議題上的少數，
% cue 3056 7965.470 7969.670
\hypertarget{w06:cue:3056}{}又是，既是這個 topical 上的少數，
% cue 3057 7969.670 7971.890
\hypertarget{w06:cue:3057}{}又是 decisional 的少數，
% cue 3058 7971.890 7980.770
\hypertarget{w06:cue:3058}{}才是，就是所謂碰到這個 tyranny of the majority，
% cue 3059 7980.770 7984.820
\hypertarget{w06:cue:3059}{}所要討論的核心問題。

\tc{2:13:04}
\speaker{老師}
% cue 3060 7984.820 7989.400
\hypertarget{w06:cue:3060}{}第一個你這張表，
% cue 3061 7989.400 7992.640
\hypertarget{w06:cue:3061}{}我不知道大家決定的，
% cue 3062 7992.640 7995.280
\hypertarget{w06:cue:3062}{}大家現在看起來是怎樣的，
% cue 3063 7995.280 7996.380
\hypertarget{w06:cue:3063}{}我不曉得這個表，
% cue 3064 7996.380 7997.740
\hypertarget{w06:cue:3064}{}這個表是你畫的，
% cue 3065 7997.740 7998.140
\hypertarget{w06:cue:3065}{}對不對？
% cue 3066 7998.140 7999.810
\hypertarget{w06:cue:3066}{}是吧？
% cue 3067 7999.810 8001.310
\hypertarget{w06:cue:3067}{}那你要去想一想，
% cue 3068 8001.310 8001.870
\hypertarget{w06:cue:3068}{}你這個表，
% cue 3069 8001.870 8003.520
\hypertarget{w06:cue:3069}{}有沒有達成你的目的？
% cue 3070 8003.520 8005.950
\hypertarget{w06:cue:3070}{}你希望讓我們看得一目瞭然，
% cue 3071 8005.950 8007.870
\hypertarget{w06:cue:3071}{}我現在是看到四個框框，
% cue 3072 8007.870 8009.430
\hypertarget{w06:cue:3072}{}可是我沒有看懂，
% cue 3073 8009.430 8011.840
\hypertarget{w06:cue:3073}{}當然知道你把他，
% cue 3074 8011.840 8013.200
\hypertarget{w06:cue:3074}{}第三個左下角那個，
% cue 3075 8013.200 8014.370
\hypertarget{w06:cue:3075}{}你把他框起來，
% cue 3076 8014.370 8017.560
\hypertarget{w06:cue:3076}{}你要想辦法讓我們看得到的是，
% cue 3077 8017.560 8020.140
\hypertarget{w06:cue:3077}{}兩個重疊的那個部分，
% cue 3078 8020.140 8021.600
\hypertarget{w06:cue:3078}{}但是兩個重疊，
% cue 3079 8021.600 8023.990
\hypertarget{w06:cue:3079}{}既是決定上的少數，
% cue 3080 8023.990 8026.280
\hypertarget{w06:cue:3080}{}又是議題上的少數，
% cue 3081 8026.280 8029.190
\hypertarget{w06:cue:3081}{}才是我們這裡講的 tyranny 嗎？
% cue 3082 8029.190 8031.130
\hypertarget{w06:cue:3082}{}是呢，
% cue 3083 8031.130 8032.710
\hypertarget{w06:cue:3083}{}是擔心的嗎？
% cue 3084 8032.710 8036.010
\hypertarget{w06:cue:3084}{}才是這裡講的 tyranny of majority，
% cue 3085 8036.010 8037.340
\hypertarget{w06:cue:3085}{}所被欺負的，
% cue 3086 8037.340 8039.680
\hypertarget{w06:cue:3086}{}才是我們這裡需要關心的嗎？
% cue 3087 8039.680 8041.030
\hypertarget{w06:cue:3087}{}你的意思是這樣嗎？

\tc{2:14:01}
\speaker{導讀人：亦凡}
% cue 3088 8041.030 8042.710
\hypertarget{w06:cue:3088}{}對文章是這樣。

\tc{2:14:02}
\speaker{老師}
% cue 3089 8042.710 8043.230
\hypertarget{w06:cue:3089}{}好。

\tc{2:14:03}
% cue 3090 8043.230 8044.770
\hypertarget{w06:cue:3090}{}那現在問題是，
% cue 3091 8044.770 8045.550
\hypertarget{w06:cue:3091}{}這一點，
% cue 3092 8045.550 8047.230
\hypertarget{w06:cue:3092}{}那為什麼不適合，
% cue 3093 8047.230 8048.670
\hypertarget{w06:cue:3093}{}就這點而言，
% cue 3094 8048.670 8052.210
\hypertarget{w06:cue:3094}{}為什麼由法院來做違憲審查，
% cue 3095 8052.210 8055.180
\hypertarget{w06:cue:3095}{}還是不如由立法者來做決定，
% cue 3096 8055.180 8057.580
\hypertarget{w06:cue:3096}{}他有這樣講嗎？
% cue 3097 8057.580 8060.000
\hypertarget{w06:cue:3097}{}因為他反對啊。

\tc{2:14:20}
\speaker{導讀人：亦凡}
% cue 3098 8060.000 8063.410
\hypertarget{w06:cue:3098}{}他認為那些申請。

\tc{2:14:23}
\speaker{老師}
% cue 3099 8063.410 8064.530
\hypertarget{w06:cue:3099}{}你找地方給我，
% cue 3100 8064.530 8066.150
\hypertarget{w06:cue:3100}{}在哪裡講到這個？
% cue 3101 8066.150 8069.310
\hypertarget{w06:cue:3101}{}有嗎？
% cue 3102 8069.310 8071.960
\hypertarget{w06:cue:3102}{}他有沒有一個 sum up？

\tc{2:14:31}
\speaker{導讀人：亦凡}
% cue 3103 8071.960 8079.160
\hypertarget{w06:cue:3103}{}1397 的最下面到 1398 的那邊，
% cue 3104 8079.160 8083.410
\hypertarget{w06:cue:3104}{}1397 的最下面的那一段吧。

\tc{2:14:43}
\speaker{老師}
% cue 3105 8083.410 8084.430
\hypertarget{w06:cue:3105}{}那什麼意思你說？

\tc{2:14:44}
\speaker{導讀人：亦凡}
% cue 3106 8084.430 8085.630
\hypertarget{w06:cue:3106}{}就是。

\tc{2:14:45}
\speaker{老師}
% cue 3107 8085.630 8088.660
\hypertarget{w06:cue:3107}{}就是你現在講的 distinction，
% cue 3108 8088.660 8092.990
\hypertarget{w06:cue:3108}{}對不對？

\tc{2:14:52}
% cue 3109 8092.990 8094.830
\hypertarget{w06:cue:3109}{}你現在要講的是，
% cue 3110 8094.830 8097.110
\hypertarget{w06:cue:3110}{}你要讓你自己分的東西，
% cue 3111 8097.110 8098.490
\hypertarget{w06:cue:3111}{}能夠跟他的文本，
% cue 3112 8098.490 8099.890
\hypertarget{w06:cue:3112}{}因為你不是作者，
% cue 3113 8099.890 8101.430
\hypertarget{w06:cue:3113}{}他是作者，
% cue 3114 8101.430 8103.530
\hypertarget{w06:cue:3114}{}那你要把他的想法，
% cue 3115 8103.530 8105.920
\hypertarget{w06:cue:3115}{}講給我們的精確的表達，
% cue 3116 8105.920 8107.320
\hypertarget{w06:cue:3116}{}不然他才跑來跟我靠一，
% cue 3117 8107.320 8109.400
\hypertarget{w06:cue:3117}{}然後來接受講。

\tc{2:15:09}
\speaker{導讀人：亦凡}
% cue 3118 8109.400 8113.330
\hypertarget{w06:cue:3118}{}在 1397 的最後面一段，
% cue 3119 8113.330 8115.790
\hypertarget{w06:cue:3119}{}我認為他是想表示說，
% cue 3120 8115.790 8118.410
\hypertarget{w06:cue:3120}{}當我們在談論多數暴力的時候，
% cue 3121 8118.410 8122.890
\hypertarget{w06:cue:3121}{}我們要去證明自己所遭遇的，
% cue 3122 8122.890 8124.290
\hypertarget{w06:cue:3122}{}那個真的是，
% cue 3123 8124.290 8128.680
\hypertarget{w06:cue:3123}{}他這邊用這個表格所呈現出來的，
% cue 3124 8128.680 8130.340
\hypertarget{w06:cue:3124}{}這個既是議題上的少數，
% cue 3125 8130.340 8132.980
\hypertarget{w06:cue:3125}{}又是決定上的少數的這種情況，
% cue 3126 8132.980 8135.250
\hypertarget{w06:cue:3126}{}而他才，
% cue 3127 8135.250 8138.550
\hypertarget{w06:cue:3127}{}才需要司法還是才不需要司法。

\tc{2:15:38}
\speaker{老師}
% cue 3128 8138.550 8142.490
\hypertarget{w06:cue:3128}{}才需要法院還是才不需要法院，
% cue 3129 8142.490 8145.630
\hypertarget{w06:cue:3129}{}可能照你的這個論證，
% cue 3130 8145.630 8147.370
\hypertarget{w06:cue:3130}{}應該是才需要法院。

\tc{2:15:47}
\speaker{導讀人：亦凡}
% cue 3131 8147.370 8148.870
\hypertarget{w06:cue:3131}{}才需要法院。

\tc{2:15:48}
\speaker{老師}
% cue 3132 8148.870 8149.740
\hypertarget{w06:cue:3132}{}不是嗎？
% cue 3133 8149.740 8151.520
\hypertarget{w06:cue:3133}{}他投票又投不贏人家，
% cue 3134 8151.520 8154.030
\hypertarget{w06:cue:3134}{}立體上他又是孤立而隔絕的少數\footnote{語出 United States v. Carolene Products Co., 304 U.S. 144（1938 年 4 月 25 日判決）註 4：「prejudice against discrete and insular minorities may be a special condition, which tends seriously to curtail the operation of those political processes ordinarily to be relied upon to protect minorities」。Waldron 以此界定違憲審查論據較具說服力之非核心情形。見 Jeremy Waldron, The Core of the Case Against Judicial Review, 115 Yale L.J. 1346, 1403–1404 (2006)。\url{https://www.law.cornell.edu/supremecourt/text/304/144}。}，
% cue 3135 8154.030 8155.810
\hypertarget{w06:cue:3135}{}沒有人跟他一樣啊，
% cue 3136 8155.810 8158.170
\hypertarget{w06:cue:3136}{}這種人就是很少啊。

\tc{2:15:58}
% cue 3137 8158.170 8161.320
\hypertarget{w06:cue:3137}{}好像臺灣的原住民\footnote{內政部統計處 2026 年 2 月 21 日發布之內政統計通報：2025 年底戶籍登記為原住民身分者計 62 萬 9,456 人，占全國人口百分之 2.7。\url{https://www.moi.gov.tw/News_Content.aspx?n=2905&s=336847}。}是 2\%啊，
% cue 3138 8161.320 8162.710
\hypertarget{w06:cue:3138}{}也投敗投得贏嗎？
% cue 3139 8162.710 8166.190
\hypertarget{w06:cue:3139}{}不過原住民倒還好，
% cue 3140 8166.190 8167.410
\hypertarget{w06:cue:3140}{}因為我們的社會並沒有，
% cue 3141 8167.410 8169.830
\hypertarget{w06:cue:3141}{}讓他們變成孤立而隔絕的少數，
% cue 3142 8169.830 8172.740
\hypertarget{w06:cue:3142}{}孤立而隔絕少數不敢出貴的人才有，
% cue 3143 8172.740 8173.990
\hypertarget{w06:cue:3143}{}比較是這種，
% cue 3144 8173.990 8176.230
\hypertarget{w06:cue:3144}{}但是現在同時已經很開放了，
% cue 3145 8176.230 8178.730
\hypertarget{w06:cue:3145}{}他也不太像孤立而隔絕的少數，
% cue 3146 8178.730 8181.600
\hypertarget{w06:cue:3146}{}所以他也有可能去爭取選票變成多數，
% cue 3147 8181.600 8183.510
\hypertarget{w06:cue:3147}{}不是嗎？

\tc{2:16:23}
% cue 3148 8183.510 8186.990
\hypertarget{w06:cue:3148}{}然後你要講什麼，
% cue 3149 8186.990 8188.890
\hypertarget{w06:cue:3149}{}最主要是這個文章到這裡，
% cue 3150 8188.890 8190.550
\hypertarget{w06:cue:3150}{}沒有達成你要的任務，
% cue 3151 8190.550 8192.470
\hypertarget{w06:cue:3151}{}聽得懂聽不懂？

\tc{2:16:32}
% cue 3152 8192.470 8197.200
\hypertarget{w06:cue:3152}{}first 在 1398 有一個 first，
% cue 3153 8197.200 8198.740
\hypertarget{w06:cue:3153}{}然後有一個 second，
% cue 3154 8198.740 8200.010
\hypertarget{w06:cue:3154}{}second 在哪裡？
% cue 3155 8200.010 8201.390
\hypertarget{w06:cue:3155}{}1440 有一個 second，
% cue 3156 8201.390 8204.320
\hypertarget{w06:cue:3156}{}那你結論你又不結了，
% cue 3157 8204.320 8205.620
\hypertarget{w06:cue:3157}{}你要回答我，
% cue 3158 8205.620 8207.910
\hypertarget{w06:cue:3158}{}為什麼作者的觀點，
% cue 3159 8207.910 8212.160
\hypertarget{w06:cue:3159}{}認為即便是一般人認為，
% cue 3160 8212.160 8215.150
\hypertarget{w06:cue:3160}{}我們為了對抗多數暴力，
% cue 3161 8215.150 8216.460
\hypertarget{w06:cue:3161}{}所以我們需要違憲審查，
% cue 3162 8216.460 8218.080
\hypertarget{w06:cue:3162}{}作者還是認為，
% cue 3163 8218.080 8221.590
\hypertarget{w06:cue:3163}{}立法者仍然是比較有效的，
% cue 3164 8221.590 8223.840
\hypertarget{w06:cue:3164}{}是解決多數暴力的情況，
% cue 3165 8223.840 8227.010
\hypertarget{w06:cue:3165}{}或者多數暴力也不是像你們所想象的，
% cue 3166 8227.010 8229.120
\hypertarget{w06:cue:3166}{}那麼有意義還是怎麼樣，
% cue 3167 8229.120 8231.020
\hypertarget{w06:cue:3167}{}你覺得是怎麼樣？

\tc{2:17:11}
\speaker{導讀人：亦凡}
% cue 3168 8231.020 8234.780
\hypertarget{w06:cue:3168}{}就建立在他的前提的假設上面的。

\tc{2:17:14}
\speaker{老師}
% cue 3169 8234.780 8238.580
\hypertarget{w06:cue:3169}{}你講話都是非常抽象的，
% cue 3170 8238.580 8240.520
\hypertarget{w06:cue:3170}{}你現在這樣的大，
% cue 3171 8240.520 8242.280
\hypertarget{w06:cue:3171}{}你會覺得剛剛講那句話，
% cue 3172 8242.280 8243.620
\hypertarget{w06:cue:3172}{}放到什麼地方都可以說。

\tc{2:17:23}
\speaker{導讀人：亦凡}
% cue 3173 8243.620 8249.980
\hypertarget{w06:cue:3173}{}就是他就是認為有更好的，
% cue 3174 8249.980 8252.120
\hypertarget{w06:cue:3174}{}或者是更適切，
% cue 3175 8252.120 8253.400
\hypertarget{w06:cue:3175}{}在一個理想狀態當中，
% cue 3176 8253.400 8257.340
\hypertarget{w06:cue:3176}{}有更適切的程序跟機制。

\tc{2:17:37}
\speaker{老師}
% cue 3177 8257.340 8258.820
\hypertarget{w06:cue:3177}{}這個板一下，
% cue 3178 8258.820 8260.520
\hypertarget{w06:cue:3178}{}你下個禮拜找出來，
% cue 3179 8260.520 8261.320
\hypertarget{w06:cue:3179}{}我一直在告訴你，
% cue 3180 8261.320 8262.440
\hypertarget{w06:cue:3180}{}反正這個大家都帶著，
% cue 3181 8262.440 8263.640
\hypertarget{w06:cue:3181}{}你下個禮拜找出來，
% cue 3182 8263.640 8264.400
\hypertarget{w06:cue:3182}{}要不丟了。

\section{非核心情形與防衛式民主（第柒部分）}

\tc{2:17:44}
% cue 3183 8264.400 8266.260
\hypertarget{w06:cue:3183}{}現在 1401 頁講什麼，
% cue 3184 8266.260 8269.370
\hypertarget{w06:cue:3184}{}非核心的案子，
% cue 3185 8269.370 8270.840
\hypertarget{w06:cue:3185}{}非核心的案子，
% cue 3186 8270.840 8271.420
\hypertarget{w06:cue:3186}{}大家講什麼？

\tc{2:17:51}
\speaker{導讀人：亦凡}
% cue 3187 8271.420 8273.630
\hypertarget{w06:cue:3187}{}就是如果那些假設不成立的話。

\tc{2:17:53}
\speaker{老師}
% cue 3188 8273.630 8274.690
\hypertarget{w06:cue:3188}{}就前面四個假設。

\tc{2:17:54}
\speaker{導讀人：亦凡}
% cue 3189 8274.690 8275.850
\hypertarget{w06:cue:3189}{}如果有其中的。

\tc{2:17:55}
\speaker{老師}
% cue 3190 8275.850 8276.810
\hypertarget{w06:cue:3190}{}是 沒關係，
% cue 3191 8276.810 8277.450
\hypertarget{w06:cue:3191}{}就是不成立，
% cue 3192 8277.450 8278.350
\hypertarget{w06:cue:3192}{}然後呢，
% cue 3193 8278.350 8280.300
\hypertarget{w06:cue:3193}{}有一個不成立也是不成立，
% cue 3194 8280.300 8280.880
\hypertarget{w06:cue:3194}{}然後呢。

\tc{2:18:00}
\speaker{導讀人：亦凡}
% cue 3195 8280.880 8283.320
\hypertarget{w06:cue:3195}{}他因為違憲審查的這個制度，
% cue 3196 8283.320 8285.620
\hypertarget{w06:cue:3196}{}機制還是會有存在的必要性。

\tc{2:18:05}
\speaker{老師}
% cue 3197 8285.620 8286.220
\hypertarget{w06:cue:3197}{}為什麼？
% cue 3198 8286.220 8288.890
\hypertarget{w06:cue:3198}{}因為那些民主機制不能運作，
% cue 3199 8288.890 8292.130
\hypertarget{w06:cue:3199}{}那你在不靠司法而不意，那不可以打掉嗎？

\tc{2:18:12}
% cue 3200 8292.130 8295.470
\hypertarget{w06:cue:3200}{}那從這裡來看，他有沒有提到防衛式民主\footnote{德文 streitbare Demokratie 或 wehrhafte Demokratie，指憲法秩序設有對抗其敵對者之法律手段，概念可溯及 Karl Loewenstein 所稱 militant democracy。德國基本法第 18 條規定，濫用言論、集會、結社等基本權以對抗自由民主基本秩序者，喪失該基本權；第 21 條第 2 項規定，以損害或廢除自由民主基本秩序為目的之政黨違憲。憲法增修條文第 5 條第 5 項規定：「政黨之目的或其行為，危害中華民國之存在或自由民主之憲政秩序者為違憲。」\url{https://www.bpb.de/themen/rechtsextremismus/dossier-rechtsextremismus/41891/streitbare-demokratie/}。}？
% cue 3201 8295.470 8305.320
\hypertarget{w06:cue:3201}{}他有還是沒有？好像沒有，因為美國人不會講這個，
% cue 3202 8305.320 8308.830
\hypertarget{w06:cue:3202}{}但是我猜想，德國人一定會講，
% cue 3203 8308.830 8314.330
\hypertarget{w06:cue:3203}{}因為不是剛才講了嗎？民主的機制如果是正常運作的話，我當然反對，
% cue 3204 8314.330 8317.290
\hypertarget{w06:cue:3204}{}但是民主機制不正常運作，那我當然要靠司法，
% cue 3205 8317.290 8321.960
\hypertarget{w06:cue:3205}{}我猜想，這個也是一個防衛式民主非常原始的一個想法，
% cue 3206 8321.960 8324.960
\hypertarget{w06:cue:3206}{}他如果民主都正常，那幹嘛要防衛？
% cue 3207 8324.960 8327.580
\hypertarget{w06:cue:3207}{}自己防衛自己就好，
% cue 3208 8327.580 8328.500
\hypertarget{w06:cue:3208}{}自己不能防衛，
% cue 3209 8328.500 8329.740
\hypertarget{w06:cue:3209}{}所以要靠司法來防衛。

\section{結論與導讀人的最後陳述}

\tc{2:18:49}
% cue 3210 8329.740 8332.660
\hypertarget{w06:cue:3210}{}來 最後一個 conclusion 要講什麼，
% cue 3211 8332.660 8336.200
\hypertarget{w06:cue:3211}{}有什麼？
% cue 3212 8336.200 8337.140
\hypertarget{w06:cue:3212}{}有就有，
% cue 3213 8337.140 8337.800
\hypertarget{w06:cue:3213}{}沒有就沒有。

\tc{2:18:57}
\speaker{導讀人：亦凡}
% cue 3214 8337.800 8338.200
\hypertarget{w06:cue:3214}{}難怪，
% cue 3215 8338.200 8338.740
\hypertarget{w06:cue:3215}{}沒有。

\tc{2:18:58}
\speaker{老師}
% cue 3216 8338.740 8339.440
\hypertarget{w06:cue:3216}{}好。

\tc{2:18:59}
% cue 3217 8339.440 8342.210
\hypertarget{w06:cue:3217}{}那現在你自己給自己一個，
% cue 3218 8342.210 8343.250
\hypertarget{w06:cue:3218}{}最後的，
% cue 3219 8343.250 8345.450
\hypertarget{w06:cue:3219}{}老師一定給你一個最後陳述的機會，
% cue 3220 8345.450 8346.510
\hypertarget{w06:cue:3220}{}來 兩分鐘，
% cue 3221 8346.510 8348.570
\hypertarget{w06:cue:3221}{}你滿意嗎？
% cue 3222 8348.570 8349.470
\hypertarget{w06:cue:3222}{}你今天的報告，
% cue 3223 8349.470 8351.770
\hypertarget{w06:cue:3223}{}你們也覺得他做得很辛苦，
% cue 3224 8351.770 8352.730
\hypertarget{w06:cue:3224}{}他非常認真，
% cue 3225 8352.730 8354.820
\hypertarget{w06:cue:3225}{}然後你覺得今天，
% cue 3226 8354.820 8356.660
\hypertarget{w06:cue:3226}{}像這樣給你報告，
% cue 3227 8356.660 8358.150
\hypertarget{w06:cue:3227}{}你覺得滿意嗎？

\tc{2:19:18}
\speaker{導讀人：亦凡}
% cue 3228 8358.150 8359.950
\hypertarget{w06:cue:3228}{}不是，
% cue 3229 8359.950 8361.050
\hypertarget{w06:cue:3229}{}就不用讓他的，
% cue 3230 8361.050 8369.470
\hypertarget{w06:cue:3230}{}就他反對建立在一個太嚴格的假設上面了。

\tc{2:19:29}
\speaker{老師}
% cue 3231 8369.470 8370.090
\hypertarget{w06:cue:3231}{}太嚴格，
% cue 3232 8370.090 8371.450
\hypertarget{w06:cue:3232}{}那四個基礎嗎？
% cue 3233 8371.450 8372.790
\hypertarget{w06:cue:3233}{}你的意思是說什麼？

\tc{2:19:32}
\speaker{導讀人：亦凡}
% cue 3234 8372.790 8374.550
\hypertarget{w06:cue:3234}{}對因為他的那個四個基礎，
% cue 3235 8374.550 8378.900
\hypertarget{w06:cue:3235}{}都是很理想化的假設的情況下。

\tc{2:19:38}
\speaker{老師}
% cue 3236 8378.900 8381.490
\hypertarget{w06:cue:3236}{}他自己有承認這是很理想化嗎？
% cue 3237 8381.490 8383.840
\hypertarget{w06:cue:3237}{}他說他放得很快呢。

\tc{2:19:43}
\speaker{導讀人：亦凡}
% cue 3238 8383.840 8386.190
\hypertarget{w06:cue:3238}{}我自己感覺，
% cue 3239 8386.190 8388.040
\hypertarget{w06:cue:3239}{}我自己感覺就是。

\tc{2:19:48}
\speaker{老師}
% cue 3240 8388.040 8391.610
\hypertarget{w06:cue:3240}{}你應該這樣回答，
% cue 3241 8391.610 8396.040
\hypertarget{w06:cue:3241}{}他那四個假設放在我們一般的民主國家，
% cue 3242 8396.040 8398.390
\hypertarget{w06:cue:3242}{}有違憲審查制度的國家，
% cue 3243 8398.390 8399.960
\hypertarget{w06:cue:3243}{}是不是都是要講的這樣。

\tc{2:19:59}
\speaker{導讀人：亦凡}
% cue 3244 8399.960 8404.290
\hypertarget{w06:cue:3244}{}很難達到吧。

\tc{2:20:04}
\speaker{老師}
% cue 3245 8404.290 8405.490
\hypertarget{w06:cue:3245}{}是嗎？
% cue 3246 8405.490 8406.770
\hypertarget{w06:cue:3246}{}為什麼這麼覺得？

\tc{2:20:06}
% cue 3247 8406.770 8411.620
\hypertarget{w06:cue:3247}{}我現在講這個話，
% cue 3248 8411.620 8412.920
\hypertarget{w06:cue:3248}{}對你來講也不太公允，
% cue 3249 8412.920 8414.620
\hypertarget{w06:cue:3249}{}我只能提到這裡，
% cue 3250 8414.620 8416.000
\hypertarget{w06:cue:3250}{}大家回去想一想。

\tc{2:20:16}
% cue 3251 8416.000 8418.720
\hypertarget{w06:cue:3251}{}他作者是一直強調，
% cue 3252 8418.720 8422.660
\hypertarget{w06:cue:3252}{}民主本身是有缺陷的，
% cue 3253 8422.660 8425.050
\hypertarget{w06:cue:3253}{}但是我們期待這個缺陷，
% cue 3254 8425.050 8427.050
\hypertarget{w06:cue:3254}{}透過民主自己的程序來矯正，
% cue 3255 8427.050 8430.020
\hypertarget{w06:cue:3255}{}還是很有希望的，
% cue 3256 8430.020 8431.080
\hypertarget{w06:cue:3256}{}而且可以期待，
% cue 3257 8431.080 8435.320
\hypertarget{w06:cue:3257}{}比透過這個司法程序來矯正，
% cue 3258 8435.320 8437.020
\hypertarget{w06:cue:3258}{}他認為也是這樣的，
% cue 3259 8437.020 8441.460
\hypertarget{w06:cue:3259}{}那他也說有很多地方不完美的缺點，
% cue 3260 8441.460 8445.880
\hypertarget{w06:cue:3260}{}那他一直強調他好像不是希望得到，
% cue 3261 8445.880 8449.020
\hypertarget{w06:cue:3261}{}不是希望給人家像你得到了這樣的印象，
% cue 3262 8449.020 8450.870
\hypertarget{w06:cue:3262}{}他假設的很嚴厲，
% cue 3263 8450.870 8452.990
\hypertarget{w06:cue:3263}{}他一直是希望不要這樣的，
% cue 3264 8452.990 8455.960
\hypertarget{w06:cue:3264}{}但是你居然會得到一個跟他相反的假設，
% cue 3265 8455.960 8457.960
\hypertarget{w06:cue:3265}{}那表示他覺得說還蠻失敗的，
% cue 3266 8457.960 8459.320
\hypertarget{w06:cue:3266}{}他沒有說服你。

\tc{2:20:59}
% cue 3267 8459.320 8462.950
\hypertarget{w06:cue:3267}{}所以你的意思是說違憲審查還是很好的，
% cue 3268 8462.950 8463.790
\hypertarget{w06:cue:3268}{}我們不應該反對，
% cue 3269 8463.790 8467.440
\hypertarget{w06:cue:3269}{}因為他現在是反對的，
% cue 3270 8467.440 8470.430
\hypertarget{w06:cue:3270}{}然後他的反對是建立在，
% cue 3271 8470.430 8471.490
\hypertarget{w06:cue:3271}{}剛才那幾個假設，
% cue 3272 8471.490 8473.470
\hypertarget{w06:cue:3272}{}那你跟你現在的意思是說，
% cue 3273 8473.470 8474.890
\hypertarget{w06:cue:3273}{}那幾個假設他也可，
% cue 3274 8474.890 8476.070
\hypertarget{w06:cue:3274}{}我們更不可能做到，
% cue 3275 8476.070 8478.490
\hypertarget{w06:cue:3275}{}所以你現在是在烏托邦的反對，
% cue 3276 8478.490 8480.480
\hypertarget{w06:cue:3276}{}這意思嗎？

\tc{2:21:20}
\speaker{導讀人：亦凡}
% cue 3277 8480.480 8481.900
\hypertarget{w06:cue:3277}{}我不想說，
% cue 3278 8481.900 8485.070
\hypertarget{w06:cue:3278}{}其實我甚至不覺得，
% cue 3279 8485.070 8489.230
\hypertarget{w06:cue:3279}{}他的這篇文章有在反對違憲審查。

\tc{2:21:29}
\speaker{老師}
% cue 3280 8489.230 8491.350
\hypertarget{w06:cue:3280}{}這是因為他先前有告訴你，
% cue 3281 8491.350 8492.470
\hypertarget{w06:cue:3281}{}我剛剛講的兩個小子，
% cue 3282 8492.470 8494.260
\hypertarget{w06:cue:3282}{}他已經寫過很多文章，
% cue 3283 8494.260 8495.660
\hypertarget{w06:cue:3283}{}討論違憲審查，
% cue 3284 8495.660 8498.120
\hypertarget{w06:cue:3284}{}有很多地方哪些哪些哪些不好，
% cue 3285 8498.120 8502.720
\hypertarget{w06:cue:3285}{}他現在是把它通通都剝掉，
% cue 3286 8502.720 8504.120
\hypertarget{w06:cue:3286}{}我只要講最核心的，
% cue 3287 8504.120 8506.520
\hypertarget{w06:cue:3287}{}他是這個意思，
% cue 3288 8506.520 8508.990
\hypertarget{w06:cue:3288}{}他好像有個意思說，
% cue 3289 8508.990 8509.990
\hypertarget{w06:cue:3289}{}從此以後我不再發言，
% cue 3290 8509.990 8510.990
\hypertarget{w06:cue:3290}{}我講太多了，
% cue 3291 8510.990 8512.880
\hypertarget{w06:cue:3291}{}結果我們發現呢，
% cue 3292 8512.880 8514.880
\hypertarget{w06:cue:3292}{}他最後一次發言還是失敗，
% cue 3293 8514.880 8515.880
\hypertarget{w06:cue:3293}{}因為他沒有說服。

\section{下課後關於導讀資料與摘要的交代}

\tc{2:21:55}
% cue 3294 8515.880 8519.700
\hypertarget{w06:cue:3294}{}我們下課，
% cue 3295 8519.700 8521.110
\hypertarget{w06:cue:3295}{}下禮拜見，
% cue 3296 8521.110 8522.110
\hypertarget{w06:cue:3296}{}不會，
% cue 3297 8522.110 8523.050
\hypertarget{w06:cue:3297}{}大家加油，
% cue 3298 8523.050 8525.080
\hypertarget{w06:cue:3298}{}這個主題真的是不容易，
% cue 3299 8525.080 8527.050
\hypertarget{w06:cue:3299}{}我真的不知道有這麼猛，
% cue 3300 8527.050 8528.110
\hypertarget{w06:cue:3300}{}來，
% cue 3301 8528.110 8531.420
\hypertarget{w06:cue:3301}{}怎樣，
% cue 3302 8531.420 8532.510
\hypertarget{w06:cue:3302}{}謝謝，
% cue 3303 8532.510 8533.510
\hypertarget{w06:cue:3303}{}我還有下禮拜要。

\tc{2:22:13}
\speaker{學生}
% cue 3304 8533.510 8535.220
\hypertarget{w06:cue:3304}{}老師可以來。

\tc{2:22:15}
\speaker{老師}
% cue 3305 8535.220 8536.620
\hypertarget{w06:cue:3305}{}我不要我不來，
% cue 3306 8536.620 8538.420
\hypertarget{w06:cue:3306}{}你們都發現為什麼我要來，
% cue 3307 8538.420 8539.550
\hypertarget{w06:cue:3307}{}我當然不來，
% cue 3308 8539.550 8540.700
\hypertarget{w06:cue:3308}{}是不是，
% cue 3309 8540.700 8542.080
\hypertarget{w06:cue:3309}{}嘿嘿好。

\tc{2:22:22}
\speaker{導讀人：亦凡}
% cue 3310 8542.080 8549.200
\hypertarget{w06:cue:3310}{}那大家可以看一下那個，
% cue 3311 8549.200 8550.590
\hypertarget{w06:cue:3311}{}我上傳的那個資料夾，
% cue 3312 8550.590 8551.790
\hypertarget{w06:cue:3312}{}我有分析就是，
% cue 3313 8551.790 8554.290
\hypertarget{w06:cue:3313}{}老師以前做的導讀資料到裡面。

\tc{2:22:34}
\speaker{老師}
% cue 3314 8554.290 8555.490
\hypertarget{w06:cue:3314}{}大家可以看一下。

\tc{2:22:35}
\speaker{導讀人：亦凡}
% cue 3315 8555.490 8557.410
\hypertarget{w06:cue:3315}{}老師上學期做的。

\tc{2:22:37}
\speaker{老師}
% cue 3316 8557.410 8558.390
\hypertarget{w06:cue:3316}{}喔真的喔。

\tc{2:22:38}
\speaker{導讀人：亦凡}
% cue 3317 8558.390 8559.350
\hypertarget{w06:cue:3317}{}所以。

\tc{2:22:39}
\speaker{老師}
% cue 3318 8559.350 8560.850
\hypertarget{w06:cue:3318}{}這學期做的亂七八糟，
% cue 3319 8560.850 8561.760
\hypertarget{w06:cue:3319}{}沒力氣關的。

\tc{2:22:41}
\speaker{導讀人：亦凡}
% cue 3320 8561.760 8563.170
\hypertarget{w06:cue:3320}{}就那份導讀，
% cue 3321 8563.170 8565.170
\hypertarget{w06:cue:3321}{}就是老師可以看考驗老師做的樣子，
% cue 3322 8565.170 8566.370
\hypertarget{w06:cue:3322}{}大家要像一般這樣做，
% cue 3323 8566.370 8567.170
\hypertarget{w06:cue:3323}{}那就太累了。

\tc{2:22:47}
\speaker{老師}
% cue 3324 8567.170 8571.430
\hypertarget{w06:cue:3324}{}像亦凡這樣做的話，
% cue 3325 8571.430 8572.510
\hypertarget{w06:cue:3325}{}第一個要做很久，
% cue 3326 8572.510 8573.030
\hypertarget{w06:cue:3326}{}第二個，
% cue 3327 8573.030 8574.250
\hypertarget{w06:cue:3327}{}你就是要看那個，
% cue 3328 8574.250 8575.350
\hypertarget{w06:cue:3328}{}你就是要把它看，
% cue 3329 8575.350 8577.270
\hypertarget{w06:cue:3329}{}我想你這個要看完，
% cue 3330 8577.270 8578.230
\hypertarget{w06:cue:3330}{}不快要一小時？
% cue 3331 8578.230 8579.150
\hypertarget{w06:cue:3331}{}要不要？
% cue 3332 8579.150 8580.190
\hypertarget{w06:cue:3332}{}要吧。

\tc{2:23:00}
\speaker{導讀人：亦凡}
% cue 3333 8580.190 8581.490
\hypertarget{w06:cue:3333}{}不止一個小時。

\tc{2:23:01}
\speaker{老師}
% cue 3334 8581.490 8582.820
\hypertarget{w06:cue:3334}{}不止一個小時。

\tc{2:23:02}
\speaker{導讀人：亦凡}
% cue 3335 8582.820 8584.410
\hypertarget{w06:cue:3335}{}我沒有看。

\tc{2:23:04}
\speaker{老師}
% cue 3336 8584.410 8588.140
\hypertarget{w06:cue:3336}{}為什麼老師剛跟你那個討論，
% cue 3337 8588.140 8589.040
\hypertarget{w06:cue:3337}{}你去想一想看，
% cue 3338 8589.040 8589.940
\hypertarget{w06:cue:3338}{}今天有沒有錄音，
% cue 3339 8589.940 8591.710
\hypertarget{w06:cue:3339}{}根據我剛才問你，
% cue 3340 8591.710 8592.380
\hypertarget{w06:cue:3340}{}有沒有問題，
% cue 3341 8592.380 8594.120
\hypertarget{w06:cue:3341}{}把那個答案整出來，
% cue 3342 8594.120 8595.290
\hypertarget{w06:cue:3342}{}這就是要摘要，
% cue 3343 8595.290 8598.300
\hypertarget{w06:cue:3343}{}有聽懂嗎？

\tc{2:23:18}
% cue 3344 8598.300 8600.010
\hypertarget{w06:cue:3344}{}我已經講完了，
% cue 3345 8600.010 8601.290
\hypertarget{w06:cue:3345}{}我希望我講的話有聽懂，
% cue 3346 8601.290 8603.300
\hypertarget{w06:cue:3346}{}有人沒有聽懂嗎？
% cue 3347 8603.300 8604.960
\hypertarget{w06:cue:3347}{}你聽懂了沒有？
% cue 3348 8604.960 8606.210
\hypertarget{w06:cue:3348}{}聽懂了，
% cue 3349 8606.210 8607.900
\hypertarget{w06:cue:3349}{}老師就給你挑的重點牆，
% cue 3350 8607.900 8609.500
\hypertarget{w06:cue:3350}{}重點都是對著名字，
% cue 3351 8609.500 8611.480
\hypertarget{w06:cue:3351}{}然後裡面它有 123，
% cue 3352 8611.480 8612.610
\hypertarget{w06:cue:3352}{}我們就有 123，
% cue 3353 8612.610 8614.440
\hypertarget{w06:cue:3353}{}那你把那東西串起來，
% cue 3354 8614.440 8615.710
\hypertarget{w06:cue:3354}{}這個答案就出來了，
% cue 3355 8615.710 8617.020
\hypertarget{w06:cue:3355}{}這個摘要就做完了，
% cue 3356 8617.020 8618.750
\hypertarget{w06:cue:3356}{}你現在這是筆稿，
% cue 3357 8618.750 8620.030
\hypertarget{w06:cue:3357}{}你現在這是書稿，
% cue 3358 8620.030 8621.820
\hypertarget{w06:cue:3358}{}你現在把根據這個書稿，
% cue 3359 8621.820 8623.590
\hypertarget{w06:cue:3359}{}去回答老師剛才問你的問題，
% cue 3360 8623.590 8624.890
\hypertarget{w06:cue:3360}{}答案就有了。

\tc{2:23:44}
% cue 3361 8624.890 8626.440
\hypertarget{w06:cue:3361}{}那這樣子我猜想，
% cue 3362 8626.440 8627.680
\hypertarget{w06:cue:3362}{}應該三四頁，
% cue 3363 8627.680 8628.890
\hypertarget{w06:cue:3363}{}四五頁，
% cue 3364 8628.890 8630.330
\hypertarget{w06:cue:3364}{}應該摘要就做完了，
% cue 3365 8630.330 8633.760
\hypertarget{w06:cue:3365}{}這間同意，
% cue 3366 8633.760 8634.990
\hypertarget{w06:cue:3366}{}可是你要幫他。

\tc{2:23:54}
\speaker{導讀人：亦凡}
% cue 3367 8634.990 8636.960
\hypertarget{w06:cue:3367}{}三四頁太難了。

\tc{2:23:56}
\speaker{老師}
% cue 3368 8636.960 8638.100
\hypertarget{w06:cue:3368}{}太難了。

\tc{2:23:58}
\speaker{導讀人：亦凡}
% cue 3369 8638.100 8641.860
\hypertarget{w06:cue:3369}{}要把一個文章濃縮成三四頁太難了吧。

\tc{2:24:01}
\speaker{老師}
% cue 3370 8641.860 8642.640
\hypertarget{w06:cue:3370}{}太難了齁。

\tc{2:24:02}
\speaker{學生}
% cue 3371 8642.640 8643.340
\hypertarget{w06:cue:3371}{}對啊。

\tc{2:24:03}
\speaker{導讀人：亦凡}
% cue 3372 8643.340 8644.970
\hypertarget{w06:cue:3372}{}我們的功力可能還不夠。

\tc{2:24:04}
\speaker{老師}
% cue 3373 8644.970 8645.870
\hypertarget{w06:cue:3373}{}謝謝，
% cue 3374 8645.870 8647.910
\hypertarget{w06:cue:3374}{}那我們下下禮拜見，
% cue 3375 8647.910 8648.610
\hypertarget{w06:cue:3375}{}謝謝，
% cue 3376 8648.610 8650.050
\hypertarget{w06:cue:3376}{}拜拜。

\tc{2:24:10}
\speaker{導讀人：亦凡}
% cue 3377 8650.050 8653.370
\hypertarget{w06:cue:3377}{}大家可以參考一下就是老師做的一個導讀資料，
% cue 3378 8653.370 8655.810
\hypertarget{w06:cue:3378}{}就是分段分享然後寫一點摘要，
% cue 3379 8655.810 8657.570
\hypertarget{w06:cue:3379}{}然後把重點標出來，
% cue 3380 8657.570 8659.060
\hypertarget{w06:cue:3380}{}不要畫上點，
% cue 3381 8659.060 8661.480
\hypertarget{w06:cue:3381}{}我重點在附錄啊。

\tc{2:24:21}
\speaker{老師}
% cue 3382 8661.480 8662.390
\hypertarget{w06:cue:3382}{}呵呵呵。

\tc{2:24:22}
\speaker{學生}
% cue 3383 8662.390 8662.890
\hypertarget{w06:cue:3383}{}辛苦你了，
% cue 3384 8662.890 8664.700
\hypertarget{w06:cue:3384}{}我們慢慢回去慢慢看，
% cue 3385 8666.950 8670.940
\hypertarget{w06:cue:3385}{}還有手機我再打印點，
% cue 3386 8670.940 8671.780
\hypertarget{w06:cue:3386}{}你是說那個零，
% cue 3387 8671.780 8673.660
\hypertarget{w06:cue:3387}{}多號零的。

\tc{2:24:33}
\speaker{導讀人：亦凡}
% cue 3388 8673.660 8676.010
\hypertarget{w06:cue:3388}{}我就是直接放在那個工作屋裡面，
% cue 3389 8676.010 8676.490
\hypertarget{w06:cue:3389}{}這個，
% cue 3390 8676.490 8678.670
\hypertarget{w06:cue:3390}{}對的對的就是這個，
% cue 3391 8678.670 8679.490
\hypertarget{w06:cue:3391}{}就是這個。
